\documentclass[11pt,a4paper,openany,oneside]{book}
\usepackage[T1]{fontenc}
\usepackage{lmodern}
\usepackage{mathtools,amssymb,amsthm}
\usepackage{microtype}
\usepackage[left=29mm,right=29mm,top=27mm,bottom=25mm,headsep=7mm]{geometry}
\usepackage{xcolor}
\definecolor{inkgreen}{HTML}{213F39}
\usepackage[unicode,colorlinks,linkcolor=inkgreen,citecolor=inkgreen,urlcolor=inkgreen,bookmarksnumbered]{hyperref}
\usepackage{bookmark}
\usepackage{fancyhdr}
\setlength{\headheight}{14pt}
\pagestyle{fancy}
\fancyhf{}
\fancyhead[L]{\small\nouppercase{\leftmark}}
\fancyhead[R]{\small\thepage}
\renewcommand{\headrulewidth}{0.25pt}
\fancypagestyle{plain}{\fancyhf{}\fancyfoot[C]{\small\thepage}\renewcommand{\headrulewidth}{0pt}}
\setlength{\emergencystretch}{2em}
\allowdisplaybreaks[2]

\newtheoremstyle{resultstyle}
  {10pt plus 2pt}{7pt plus 2pt}{\normalfont}{}{\bfseries\color{inkgreen}}{.}{0.65em}{}
\theoremstyle{resultstyle}
\newtheorem{statement}{Result}[chapter]
\newenvironment{result}[2]{\clubpenalty=10000\begin{statement}[\textnormal{[#1]} #2]\leavevmode\label{res:#1}\ignorespaces}{\end{statement}}
\newtheorem{lemma}[statement]{Lemma}
\newtheorem{proposition}[statement]{Proposition}
\newtheorem{theorem}[statement]{Theorem}
\newtheorem{corollary}[statement]{Corollary}
\theoremstyle{definition}
\newtheorem{definition}[statement]{Definition}
\newtheorem{example}[statement]{Example}
\newtheorem{remark}[statement]{Remark}
\numberwithin{equation}{chapter}

\newcommand{\Gm}{\mathbb{G}_m}
\newcommand{\Ga}{\mathbb{G}_a}
\newcommand{\GL}{\mathrm{GL}}
\newcommand{\SL}{\mathrm{SL}}
\DeclareMathOperator{\Spec}{Spec}
\DeclareMathOperator{\Lie}{Lie}
\DeclareMathOperator{\Hom}{Hom}
\DeclareMathOperator{\Gal}{Gal}
\DeclareMathOperator{\Aut}{Aut}
\DeclareMathOperator{\Res}{Res}
\DeclareMathOperator{\id}{id}
\DeclareMathOperator{\rk}{rk}
\DeclareMathOperator{\Frob}{Frob}
\DeclareMathOperator{\Tr}{Tr}
\DeclareMathOperator{\End}{End}
\DeclareMathOperator{\Ker}{ker}
\DeclareMathOperator{\im}{im}
\newcommand{\OO}{\mathcal{O}}
\newcommand{\kk}{k}
\newcommand{\ks}{k_s}
\newcommand{\kb}{\overline{k}}
\newcommand{\ru}{R_u}
\newcommand{\der}{\mathrm{der}}
\newcommand{\ad}{\mathrm{ad}}
\newcommand{\ZZ}{\mathbb{Z}}
\newcommand{\RR}{\mathbb{R}}
\newcommand{\CC}{\mathbb{C}}
\newcommand{\NN}{\mathbb{N}}
\newcommand{\FF}{\mathbb{F}}
\newcommand{\QQ}{\mathbb{Q}}

\hypersetup{pdftitle={Linear Algebraic Groups},pdfauthor={MathGenome},pdfsubject={Group schemes, reductive groups, representations, and descent}}
\title{Linear Algebraic Groups}
\author{MathGenome}
\date{November 2026}

\begin{document}
\hypersetup{pageanchor=false}
\begin{titlepage}
\thispdfpagelabel{Cover}
\thispagestyle{empty}
\vspace*{0.17\textheight}
{\large\color{inkgreen}MathGenome\par}
\vspace{1.2em}
{\color{inkgreen}\noindent\rule{\linewidth}{0.7pt}\par}
\vspace{2.2em}
{\raggedright\Huge\bfseries Linear Algebraic Groups\par}
\vspace{1.5em}
{\Large Group Schemes, Reductive Structure,\\[0.25em] and Rational Representations\par}
\vfill
{\large November 2026\par}
\end{titlepage}
\frontmatter
\hypersetup{pageanchor=true}
\tableofcontents

\chapter*{Scope and conventions}
\addcontentsline{toc}{chapter}{Scope and conventions}

These notes develop affine group schemes over a field, the structure of reductive groups, rational representations, and descent. The central examples are matrix groups, split and nonsplit tori, flag varieties, and reductive forms. Standard external theorems are identified at the point of use, with their hypotheses and role in the argument. The companion source map locates all numbered results by stable identifier and by page.

We write $G(R)$ for the points on a commutative $k$-algebra $R$, and $G_{K}$ for extension of scalars to $K$. The fields $k_s$ and $\bar k$ denote a separable and an algebraic closure. Identities of schemes or sheaves can be tested on all $R$-points in the appropriate category. Equalities only on $k$-points or only on geometric points are stated at that level. A quotient $G/H$ uses right multiplication by $H$. Its geometric or rational points are discussed separately from its construction as a sheaf.

The characters and cocharacters of a split torus $T$ are $X^*(T)$ and $X_*(T)$, paired by composition. Whenever a positive root system is chosen, it is the one associated with a specified Borel $B^+$. $B^-$ denotes its opposite relative to $T$. Chapter~\ref{ch:representations} fixes the convention for association of line bundles to characters. Smoothness, geometric connectedness, reductivity, linear reductivity, and splitness are different properties and are specified independently.

A first reading proceeds through the six chapters in order. The algebra in Results~\ref{res:F01} and \ref{res:F02} is available after the Hopf construction and supplies descent for tori. Quotient representability in Result~\ref{res:D01} uses the linear stabilization from Chapter~1 and supplies the homogeneous spaces needed for Borel conjugacy. The big cell in Result~\ref{res:D03} is established using root subgroups. It is then available for the classification in Result~\ref{res:C08}.

The chapters give the field-level core. Three extensions lead beyond it. For reductive groups over general bases one studies the relative torus and parabolic functors and pinned classification \cite{conrad}. For nonsplit and nonperfect fields the next steps are the Tits--Selbach data and pseudo-reductive structure \cite[Sec.~12]{conradfeng}\cite{cgp}. For arithmetic groups the corresponding local theory uses buildings and parahoric models, and the global theory develops the general strong approximation theorem \cite{kalethaprasad,rapinchuk}. Chapter~6 proves the strong approximation criterion over global fields, starting with concrete cases. Chapter~5 derives the Steinberg tensor product formula from the Frobenius-kernel restriction theorem. The proof of that input uses distribution algebras in the modular-representation extension.

\mainmatter
\chapter{Affine group schemes, Hopf algebras, and Lie algebras}
\label{chap:affine-groups}

Unless a further hypothesis is stated, \(k\) is an arbitrary field and \(R\) is
a commutative \(k\)-algebra. An affine algebraic group is an affine group
scheme of finite type over \(k\). We write \(k[G]\) for its coordinate
algebra and \(G(R)\) for its \(R\)-points. Put \(k_s=k^{\mathrm{sep}}\) and
\(\bar k=k^{\mathrm{alg}}\). The term \emph{closed subgroup} includes a
possibly nonreduced scheme. The notation \(1\) for a group scheme means
\(\operatorname{Spec}k\) with its unique group structure.

\section{Functor of points and affine duality}

\begin{result}{A01}{The group law as a natural operation}
A group structure on a finite-type affine \(k\)-scheme \(G\) is equivalent
to a natural group structure on \(R\mapsto G(R)\). Morphisms of groups
correspond to natural transformations of group functors. Fiber products
and scalar extension of schemes admit their ordinary pointwise descriptions.
\end{result}

\begin{proof}
The group law, identity and inversion are morphisms \(m:G\times_kG\to G\),
\(e:\operatorname{Spec}k\to G\), and \(i:G\to G\). Associativity is the
equality \(m(m\times\mathrm{id})=m(\mathrm{id}\times m)\) on \(G^3\).
The two unit equalities hold on \(G\). Each of
\(m(\mathrm{id},i)\) and \(m(i,\mathrm{id})\) is the composite through
\(e\). Applying \(\operatorname{Hom}_k(\operatorname{Spec}R,-)\) gives
the group axioms. Conversely, write \(G=\operatorname{Spec}A\).
The natural multiplication is evaluated at the universal pair of
points in \(G(A\otimes_k A)^2\) to recover \(m\). The same argument with
the universal point in \(G(A)\) recovers \(i\), and evaluation at \(k\)
recovers \(e\). Naturality then recovers every operation by substitution.
The same universal-point argument verifies the identities of morphisms
and the assertion about homomorphisms. In particular,
\((G\times_H L)(R)=G(R)\times_{H(R)}L(R)\). For an extension \(K/k\)
and a \(K\)-algebra \(R\), \(\operatorname{Hom}_K(A\otimes_k K,R)\)
is naturally \(\operatorname{Hom}_k(A,R)\).
See \cite[Definitions 1.1--1.3 and \S1.h]{milne}.
\end{proof}

\begin{result}{A02}{The commutative Hopf algebra dictionary}
Finite-type affine \(k\)-groups and finitely generated commutative Hopf
\(k\)-algebras are contravariantly equivalent. The group is commutative
if and only if its comultiplication is cocommutative.
\end{result}

\begin{proof}
On \(A=k[G]\), pullback along \((m,e,i)\) gives respectively
\(\Delta:A\to A\otimes A\), \(\epsilon:A\to k\), and \(S:A\to A\).
The group equalities become
\[
 (\Delta\otimes1)\Delta=(1\otimes\Delta)\Delta,\qquad
 (\epsilon\otimes1)\Delta=(1\otimes\epsilon)\Delta=\mathrm{id}_A,
 \qquad
 \mu(S\otimes1)\Delta=\mu(1\otimes S)\Delta=\eta\epsilon,
\]
where \(\mu\) is algebra multiplication and \(\eta:k\to A\) the unit.
Conversely these equations imply the group equalities under
\(\operatorname{Spec}\). An algebra homomorphism in the opposite direction
commutes with \(\Delta,\epsilon,S\) exactly when the scheme morphism
preserves the group structure. The transposition of the factors of
\(G\times G\) pulls back to the flip on \(A\otimes A\), proving the last
claim. See \cite[Propositions 1.25, 3.1, 3.6 and Corollary 3.7]{milne}.
\end{proof}

\begin{result}{A03}{Additive, multiplicative, and matrix groups}
The affine groups \(\mathbb G_a,\mathbb G_m,\mathrm{GL}_n,\mathrm{SL}_n\)
represent addition, units, invertible matrices, and matrices of determinant
one, respectively, over every \(R\).
\end{result}

\begin{proof}
For \(\mathbb G_a\) use \(k[x]\) with
\(\Delta x=x\otimes1+1\otimes x\), \(\epsilon x=0\), \(Sx=-x\).
For \(\mathbb G_m\) use \(k[t,t^{-1}]\) with
\(\Delta t=t\otimes t\), \(\epsilon t=1\), \(St=t^{-1}\).
For \(\mathrm{GL}_n\), put \(X=(x_{ij})\) and invert \(d=\det X\) in
\(k[x_{ij}]\). The adjugate identity and multiplicativity of determinant
give the invertibility condition over arbitrary commutative rings. The
three maps on entries are
\[
 \Delta x_{ij}=\sum_{\ell=1}^n x_{i\ell}\otimes x_{\ell j},\qquad
 \epsilon x_{ij}=\delta_{ij},\qquad S(X)=d^{-1}\operatorname{adj}(X).
\]
The determinant is group-like, meaning \(\Delta d=d\otimes d\) and
\(\epsilon d=1\), and \(Sd=d^{-1}\). Thus \(d-1\) generates a Hopf ideal:
\(\Delta(d-1)=(d-1)\otimes d+1\otimes(d-1)\) and
\(S(d-1)=-d^{-1}(d-1)\). Its quotient is \(k[\mathrm{SL}_n]\).
Taking \(r\) tensor factors gives \(k[x_1,\ldots,x_r]\) for \(\mathbb G_a^r\)
and \(k[t_1^{\pm1},\ldots,t_r^{\pm1}]\) for \(\mathbb G_m^r\), with
the displayed rules on each generator. See \cite[\S\S2.a, 3.a]{milne}.
\end{proof}

\begin{result}{A04}{Closed subgroups, normality, and kernels}
A closed subgroup \(H\subset G\) corresponds to an ideal \(I\subset k[G]\)
satisfying
\(\Delta I\subset I\otimes k[G]+k[G]\otimes I\),
\(\epsilon I=0\), and \(S I\subset I\). A homomorphism \(f:G\to L\) has
normal kernel with defining ideal \(f^*(\ker\epsilon_L)\, k[G]\).
Intersections of closed subgroups are defined by sums of their ideals.
These descriptions commute with field extension.
\end{result}

\begin{proof}
A closed immersion is a surjective map on coordinate algebras. For
\(A\to A/I\), the kernel of the tensor-square quotient over the field is
\(I\otimes A+A\otimes I\), so the conditions are precisely those for
the Hopf operations to descend. For the kernel, form
\(G\times_L\operatorname{Spec}k\) using the identity of \(L\).
Conjugation of a kernel element stays in the kernel by functoriality of
\(f\). Equality of the factorization on all \(R\) is a scheme statement
by Result~\ref{res:A01}. More generally normality of \(H\) means that
\(G\times H\to G,\ (g,h)\mapsto ghg^{-1}\), factors through \(H\).
An intersection is the fiber product over \(G\), with coordinates
\(A/(I+J)\). Tensoring over \(k\) is exact, so it preserves all these
descriptions. A Hopf \emph{quotient} of \(k[G]\) gives a closed subgroup.
An included Hopf \emph{subalgebra} gives a map out of \(G\), as studied in
Result~\ref{res:A07}. See \cite[Propositions 1.41, 1.49, 3.9--3.15 and \S1.e]{milne}.
\end{proof}

\section{Representations and stabilizers}

\begin{result}{A05}{Rational representations as comodules}
Let \(A=k[G]\). A rational representation \(G\to\mathrm{GL}(V)\) on a
finite-dimensional \(k\)-space \(V\) is equivalent to a right \(A\)-coaction
\(\rho:V\to V\otimes A\). Equivariant maps and stable subspaces are
precisely comodule maps and subcomodules. This correspondence respects
tensor products, duals, and change of field. The invariant subspace is
the kernel of \(v\mapsto \rho(v)-v\otimes1\).
\end{result}

\begin{proof}
In a basis \(e_j\), write
\(\rho(e_j)=\sum_i e_i\otimes a_{ij}\), using the matrix coefficients of
the representation. They satisfy
\(\Delta a_{ij}=\sum_\ell a_{i\ell}\otimes a_{\ell j}\) and
\(\epsilon a_{ij}=\delta_{ij}\). Conversely these two identities are
the coaction axioms. After applying a point \(g:A\to R\), the matrix
\((g(a_{ij}))\) has inverse \((g(Sa_{ij}))\) by the antipode identity,
and the multiplication axiom follows from \(\Delta\).
The equivalence for maps and subspaces follows by testing at the
universal point \(A\to A\). Tensor products use
\(\rho(v\otimes w)=\sum v_{(0)}\otimes w_{(0)}
 \otimes v_{(1)}w_{(1)}\). The dual uses coefficients \(S(a_{ji})\).
The pairing is invariant since the two matrices multiply to the
identity. Formation of the kernel defining invariants commutes with
field extension since extension of scalars is flat. If \(G\) is smooth,
over \(\bar k\) its geometric points are schematically dense, so
vanishing of the matrix coefficients testing stability can be checked
on those points. See \cite[\S4.a--b, Proposition 4.31 and Corollary 4.34]{milne}.
\end{proof}

\begin{result}{A06}{Finite-dimensional faithful linearization}
Every vector in a comodule belongs to a finite-dimensional subcomodule.
Every affine algebraic group admits a closed immersion into
\(\mathrm{GL}(V)\) for some finite-dimensional \(V\).
\end{result}

\begin{proof}
For a comodule element \(v\), express
\(\rho(v)=\sum_{i=1}^s v_i\otimes a_i\) with the \(a_i\) linearly
independent. Coassociativity and linear functionals extracting each
\(a_i\) from the last factor show that
\(\rho(v_i)\in\langle v_1,\ldots,v_s\rangle\otimes A\).
The counit axiom gives \(v=\sum_i\epsilon(a_i)v_i\).
Sums of these subspaces give the directed-union assertion.
Apply it to \(\Delta:A\to A\otimes A\) and include a finite list of
algebra generators in one stable subspace \(V\). This is the action
\((g\cdot a)(x)=a(xg)\). If \(\Delta v_j=\sum_i v_i\otimes b_{ij}\)
for a basis of \(V\), then \(v_j=\sum_i\epsilon(v_i)b_{ij}\).
Thus the \(b_{ij}\) generate \(A\) as an algebra, so the coordinate
homomorphism from \(k[\mathrm{GL}(V)]\) is surjective.
See \cite[Proposition 4.7, Theorem 4.9 and Corollary 4.10]{milne}.
\end{proof}

\begin{result}{A08}{Stabilizers, centralizers, and normalizers}
If an affine algebraic group \(H\) acts on a separated finite-type
\(k\)-scheme \(X\), its fixed-point functor is a closed subscheme of \(X\).
For an action of \(G\) on such an \(X\) and \(x\in X(k)\), the stabilizer
of \(x\) is a closed subgroup of \(G\). For a closed subgroup \(H\subset G\)
the functors \(N_G(H)\) and \(C_G(H)\) are closed subgroup schemes.
These constructions commute with scalar extension.
\end{result}

\begin{proof}
To obtain the fixed-point equations, compare the action with projection
on \(H\times X\) by pulling back the closed diagonal of \(X\). On an
affine open \(U\subset X\) this equalizer is a closed subscheme of
\(H\times U\) defined by an ideal in \(k[H]\otimes_k k[U]\).
Expand the elements of the ideal in a \(k\)-basis of \(k[H]\) and
take all coefficient equations in \(k[U]\). They cut out exactly those
points \(u\) for which every point of \(H\) after any further base change
fixes \(u\), because the basis remains free after tensoring with any
algebra. The resulting closed subschemes glue on overlaps by their
functors of points. The stabilizer of \(x\) is the equalizer of the
orbit map \(G\to X\) and the constant map with value \(x\).

For the normalizer, pull back the equations of \(H\) along
\((g,h)\mapsto ghg^{-1}\) on \(G\times H\) and impose the vanishing of
all coefficients in \(k[H]\). Impose the same equations with \(g^{-1}\)
in place of \(g\). The two inclusions together give equality of
subgroup schemes. For the centralizer, impose equality
with the second projection as maps to \(G\). Closure under products
and inverse follows from the respective functorial conditions. In a
linear representation, stabilizing a subspace is equivalently the
vanishing of the corresponding off-block entries, with the same
condition for the inverse if working over general test rings.
All equations arise from base-compatible constructions. Testing only
on geometric points is justified when the scheme on which the
equations are tested is reduced over \(\bar k\). The resulting point
equalities need not determine a nonreduced subgroup.
See \cite[Propositions 1.79, 1.83--1.84, 1.92--1.93 and Theorem 1.105]{milne}.
\end{proof}

\section{Smoothness, images, and quotients}

\begin{result}{A09}{Smoothness and the characteristic-zero theorem}
For an affine algebraic group \(G\), smoothness over \(k\), geometric
reducedness, and smoothness at the identity are equivalent. In
characteristic zero every such \(G\) is smooth. Over a perfect field,
reducedness suffices.
\end{result}

\begin{proof}
Smoothness implies geometric reducedness. Conversely, a geometrically
reduced finite-type scheme is generically smooth after passage to
\(\bar k\). A smooth geometric point can be moved to the identity by
translation. Translations also move the identity to every geometric
point. The smooth locus is open, and if its complement over \(\bar k\)
were nonempty then it would have a closed geometric point. Smoothness
descends along field extension. Over a perfect field, reduction of a
finite-type scheme already implies geometric reduction.

Here is a local proof of the characteristic-zero conclusion. Replace
\(k\) by its algebraic closure and write \(\widehat A_e\) for the completed local algebra
at the identity, with maximal ideal \(\widehat{\mathfrak m}\). Choose
cotangent-space coordinates \(x_1,\ldots,x_n\). They induce a surjection
\(k[[X_1,\ldots,X_n]]\twoheadrightarrow\widehat A_e\): modulo each
successive power of the maximal ideal one can lift the next homogeneous
component as a polynomial in the chosen coordinates, and completeness
gives a lift in the power series algebra. Its kernel \(I\) is contained
in \((X)^2\). Multiplication gives a homomorphism to the completed local
algebra at \((e,e)\), which is \(k[[X,Y]]/(I(X)+I(Y))\).
The unit laws give \(x_i\mapsto X_i+Y_i+\) terms of order at least two.
If \(I\ne0\), let \(d\ge2\) be the smallest order of a nonzero relation
and \(h(X)\) one of its nonzero initial forms of degree \(d\).
Applying multiplication to this relation shows that the degree-\(d\)
polynomial \(h(X+Y)\) belongs to the span of degree-\(d\) initial forms
from the two separate ideals \(I(X)\), \(I(Y)\). No mixed form can occur
at that degree because both ideals start in degree \(d\). Restricting
to each factor yields \(h(X+Y)=h(X)+h(Y)\).
Substitution \(Y=X\) would give \((2^d-2)h(X)=0\), a contradiction.
Thus \(\widehat A_e\) is a power series algebra, the local ring at
\(e\) is regular, and over the perfect field it is smooth at \(e\).
The tangent dimension in this argument is the embedding dimension.
For a possibly singular point it is at least the local dimension.
Compare \cite[Propositions 1.26--1.28 and Theorem 3.23]{milne}
and \cite[p.~13 correction]{milneerrata}.
\end{proof}

\begin{result}{A07}{Scheme-theoretic image and faithful flatness}
An inclusion \(A\hookrightarrow B\) of finitely generated commutative
Hopf \(k\)-algebras makes \(B\) faithfully flat over \(A\). Consequently
a homomorphism \(f:G\to H\) factors uniquely through its
scheme-theoretic image as a faithfully flat morphism followed by a
closed immersion. If \(N=\ker f\), this image represents the fppf sheaf
quotient \(G/N\).
\end{result}

\begin{proof}
Put \(X=\operatorname{Spec}B\), \(Y=\operatorname{Spec}A\), and consider
the dominant map \(q:X\to Y\). Flatness and surjectivity descend along
field extension, so for the proof work over an algebraically closed field.

First let \(Y\) be reduced. Over the algebraically closed field,
Result~\ref{res:A09} makes it smooth, so its finitely many
irreducible components are open and closed. Indeed each local
ring is regular and hence a domain, so two components cannot meet.
On any component \(Y_i=\operatorname{Spec}A_i\), the restriction
comes from \(A_i\hookrightarrow B_i\). The ring \(A_i\) is a domain
and \(B_i\) is of finite type over it. Write
\(B_i=A_i[t_1,\ldots,t_n]/I\) and let \(K_i\) be the fraction
field of \(A_i\). The generic fiber is nonzero because the
inclusion survives localization.

Here is the freeness needed on an open set of each component.
Choose a monomial order on \(K_i[t_1,\ldots,t_n]\). The ideal
of leading monomials of \(I K_i[t_1,\ldots,t_n]\) is finitely
generated, since the polynomial ring over a field is noetherian.
Choose finitely many monic polynomials in the ideal whose leading
monomials generate this monomial ideal. After inverting a single
nonzero \(s_i\in A_i\), all their coefficients lie in
\((A_i)_{s_i}\) and the polynomials themselves belong to \(I_{s_i}\);
the latter requires only clearing the denominators in a finite
number of ideal expressions. Divide by the chosen monic polynomials
over \((A_i)_{s_i}\). Each step removes a leading monomial or places
it in the remainder, and terminates by the monomial order.
The remaining monomials span \((B_i)_{s_i}\). They are independent,
because a nonzero linear combination of them in \(I_{s_i}\)
would retain its leading monomial upon extension to \(K_i\).
Thus \((B_i)_{s_i}\) is free over \((A_i)_{s_i}\) with nonempty
basis. This shows \(q\) is faithfully flat above the union \(U\)
of these nonempty opens, which is dense in every component.

For \(y\in Y(k)\), the open \(U\cap yU\) is nonempty because both
sets are dense in every component. A closed point in it has the
form \(u=yv\) with \(u,v\in U(k)\). They lift along \(q\) by
faithful flatness over \(U\) and the existence of a \(k\)-point in
a nonempty finite-type fiber over the algebraically closed field.
Then \(y=uv^{-1}\) lifts as well. Every translate \(yU\) therefore
has faithfully flat inverse image, by translation by a chosen
lift. Those translates cover all closed points of \(Y\), hence
all of \(Y\). This proves the reduced-target case.

Next suppose \(A^+=\ker\epsilon_A\) is nilpotent. Let \(C=B/A^+B\), the
coordinate ring of \(\ker q\), and choose lifts \(e_j\in B\) of a
\(k\)-basis \(\bar e_j\) of \(C\). They define an \(A\)-linear surjection
\(u:A^{(J)}\to B\): the cokernel \(Q\) has \(Q=A^+Q\), hence \(Q=0\).
Identify \(B\otimes_A B\simeq B\otimes_k C\) using the isomorphism
\[
 X\times\ker q\longrightarrow X\times_Y X,\qquad
 (x,n)\longmapsto (x,xn),
\]
whose inverse on points is \((x,x')\mapsto(x,x^{-1}x')\).
After tensoring \(u\) on the left by \(B\) and reducing modulo \(A^+B\)
in that left factor, the induced \(C\)-linear map
\(C\otimes_k C\to C\otimes_k C\) is
\[
 c\otimes d\longmapsto \sum c\,d_{(1)}\otimes d_{(2)} ,
\]
with inverse \(c\otimes d\mapsto\sum c\,S(d_{(1)})\otimes d_{(2)}\).
Here the identifications send the chosen generators to \(\bar e_j\).
A map between free modules that is invertible modulo a nilpotent
ideal is invertible: lift the inverse on a basis, and invert the two
composition errors by their finite geometric series. This works for
arbitrary rank because their images are in successive powers of the
nilpotent ideal. Hence \(B\otimes_A u\) is an isomorphism.
The map from \(A^{(J)}\) to \(B^{(J)}\) is injective, so \(u\) is
injective as well. The rank of \(u\) is positive because \(C\) has
a counit. Thus \(B\) is faithfully flat over \(A\) in this case.

The remaining case is in characteristic \(p>0\), by
Result~\ref{res:A09}. For \(r\) sufficiently large,
\(A'=A^{p^r}\) is a reduced Hopf subalgebra: \(k\) is perfect and a
single power \(p^r\) kills all of the nilradical of the noetherian
algebra \(A\). Put \(Y'=\operatorname{Spec}A'\).
By the reduced case \(X\to Y'\) is faithfully flat. Its base change of
\(X\to Y\),
\[
 X\times_{Y'}X\longrightarrow Y\times_{Y'}X,
\]
is described over \(X\) as
\(\ker(X\to Y')\times_k X\to\ker(Y\to Y')\times_k X\), using quotients
of the two entries by the given lift in the second factor. Formation
of the map on coordinate algebras is \(A\hookrightarrow B\)
tensored over \(A'\) by \(B\), hence injective. Faithfulness of
extension from \(k\) to \(B\) now gives an inclusion on the two kernel
coordinate algebras before taking the products. Every augmentation
element of \(k[\ker(Y\to Y')]\) has \(p^r\)-th power zero.
The augmentation ideal is finitely generated, so it is nilpotent.
Apply the nilpotent case to the kernel map and base-change to \(X\).
This makes the displayed morphism faithfully flat. Faithful flatness
descends along \(X\to Y'\).

Finally \(k[H]\to k[G]\) splits through its Hopf algebra image, and the
factor just proved is faithfully flat. The image is uniquely determined
by the kernel of the coordinate homomorphism. Over any test scheme a
point of the image lifts fppf locally, and two lifts differ by right
multiplication by a point of \(N\). This is exactly the asserted
sheaf quotient. See \cite[Theorems 3.31 and 3.34, Definition 5.5 and
Propositions 5.7--5.10]{milne}. See also \cite{takeuchi}.
\end{proof}

For \(n>0\), the power map \([n]:\mathbb G_m\to\mathbb G_m\) is
faithfully flat: \(k[t,t^{-1}]\) is free on \(1,t,\ldots,t^{n-1}\)
over \(k[t^n,t^{-n}]\). Its map on \(k\)-points need not be surjective,
as witnessed by squaring over \(\mathbb R\). In characteristic \(p\),
the identity inclusion \(1\to\alpha_p\) of Result~\ref{res:A10} is
surjective on the underlying topological spaces but is not flat.
Indeed \(k[z]/(z^p)\to k\) is a nontrivial quotient of a local
Artin ring by its maximal ideal. A flat finite module over a local
ring is free, whereas \(k\) is not free over this ring.

\begin{result}{A10}{Relative Frobenius and infinitesimal kernels}
In characteristic \(p>0\), the relative Frobenius
\(F_{G/k}:G\to G^{(p)}\) is induced by
\(a\otimes c\mapsto c a^p\) on \(k[G]\otimes_{k,F_k}k\).
Its kernels on \(\mathbb G_a,\mathbb G_m\) are \(\alpha_p,\mu_p\)
with coordinate rings \(k[z]/(z^p)\) and \(k[t,t^{-1}]/(t^p-1)\).
Both have length \(p\), a single geometric point, and a
one-dimensional tangent space. Reducedness over a nonperfect field
does not imply geometric reducedness even for a group.
\end{result}

\begin{proof}
The displayed pullback is a Hopf map since \(p\)-th powers preserve
sums and products in this characteristic. Pulling back the unit ideals
in the two target rings gives the two stated quotients.
In the second case \(t^p-1=(t-1)^p\), so both unit local rings are
truncated polynomial rings. The dual-number points at the identity
are respectively \(z=a\varepsilon\) and \(t=1+a\varepsilon\), \(a\in k\).

For the last claim choose \(c\in k\setminus k^p\). The ideal
\((Y^p-cX^p)\) defines a closed subgroup of \(\mathbb G_a^2\), because
its generator is primitive for the additive coproduct.
In \(k(X)\) the equation has the single root \(c^{1/p}X\) over an
algebraic closure, not in \(k(X)\). Indeed a rational function whose
\(p\)-th power is a constant is itself constant. The corresponding
polynomial is irreducible, since a nontrivial purely inseparable
extension generated by a \(p\)-th root has degree \(p\).
Gauss's lemma gives that the coordinate ring over \(k\) is a domain.
Over \(k(c^{1/p})\) its defining relation is
\((Y-c^{1/p}X)^p=0\), with \(Y-c^{1/p}X\) a nonzero nilpotent class.
See \cite[Examples 1.27, 1.60, 12.28 and \S\S2.a, 2.e]{milne}.
\end{proof}

\section{Tangent vectors and Jordan components}

\begin{result}{A11}{Lie algebra and the adjoint representation}
Let \(\mathfrak m_e=\ker(k[G]\to k)\).
The group \(\ker(G(k[\varepsilon]/(\varepsilon^2))\to G(k))\)
is the additive group of
\((\mathfrak m_e/\mathfrak m_e^2)^\vee\). It has a natural Lie bracket
and commutes with extension of \(k\). The differential of conjugation
gives \(\operatorname{Ad}:G\to\mathrm{GL}(\operatorname{Lie}G)\), whose
differential is \(X\mapsto [X,-]\).
\end{result}

\begin{proof}
A length-two lift of the counit is \(\epsilon+\varepsilon D\) with
\(D(ab)=\epsilon(a)D(b)+D(a)\epsilon(b)\). The possible \(D\) are
precisely the linear forms on \(\mathfrak m_e/\mathfrak m_e^2\).
The coproduct gives addition of the forms on multiplying lifts,
because its linear terms are \(a\otimes1+1\otimes a\) for
\(a\in\mathfrak m_e\). Substitution \(\varepsilon\mapsto s\varepsilon\)
gives scalar multiplication by \(s\). The description by cotangent
spaces also gives field extension.

Work in \(k[s,t]/(s^2,t^2)\) and take the commutator of the \(s\)-lift
of \(X\) and the \(t\)-lift of \(Y\). Since it is the identity modulo
both \(s\) and \(t\), it has the form of a tangent vector times \(st\).
This defines \([X,Y]\) intrinsically. Via a closed immersion as in
Result~\ref{res:A06}, this commutator is
\(1+st(XY-YX)\). The tangent map of a closed immersion is injective,
so the matrix Jacobi identity and bilinearity descend to this bracket.
Naturality of the commutator makes every differential a Lie map.
Conjugation preserves the identity after arbitrary base change, hence
acts naturally on the relative tangent space there. The corresponding
matrix coefficients are algebraic in \(g\), either by differentiating
the conjugation pullback or by the identity
\(\operatorname{Ad}(g)(Y)=gYg^{-1}\) in a faithful representation.
Replacing \(g\) by \(1+sX\) proves \(d\operatorname{Ad}(X)(Y)=[X,Y]\).
See \cite[\S10.b, 10.6--10.24 and Propositions 10.28--10.29]{milne}.
\end{proof}

\begin{result}{A12}{Left exactness and the differential of a quotient}
The Lie functor preserves finite fiber products, in particular
\(\operatorname{Lie}(\ker f)=\ker(df)\). The matrix groups have Lie
algebras \(\mathfrak{gl}_n\) and \(\mathfrak{sl}_n\), with the
commutator bracket. The two one-dimensional groups \(\mathbb G_a\)
and \(\mathbb G_m\) have abelian Lie algebras. For a faithfully flat
homomorphism \(f:G\to H\) of smooth affine algebraic groups, the
kernel is smooth if and only if \(df\) is surjective.
\end{result}

\begin{proof}
Taking points over both \(k\) and the dual-number algebra preserves
fiber products. Taking the fibers over the identities yields the
stated Lie result and also the analogous formula for an intersection.
Every \(1+\varepsilon X\) is invertible, while
\(\det(1+\varepsilon X)=1+\varepsilon\,\operatorname{tr} X\) by the
permutation formula. This proves the matrix descriptions. In either
commutative one-dimensional group all commutators vanish.

For the smoothness claim, over an algebraic closure every geometric
fiber of \(f\) is a translate of the kernel. If the kernel is smooth,
flatness, finite presentation, and the smoothness of geometric fibers
show that \(f\) is smooth, so \(df\) is surjective. Conversely,
surjectivity of \(df\) at the identity of the smooth source,
with smooth target, gives smoothness at that point by the
differential criterion. The identity fiber is then smooth at its
identity, hence smooth by translation. In characteristic \(p\) the
map \(\mathbb G_a\to\mathbb G_a^{(p)}\) in Result~\ref{res:A10} is
faithfully flat since \(k[x]\) is free of rank \(p\) over \(k[x^p]\).
Its differential is zero and its kernel \(\alpha_p\) is not smooth.
See \cite[Proposition 1.63, 10.7--10.14 and Remark 10.24]{milne}.
\end{proof}

\begin{result}{A13}{Identity component and component group}
The identity component \(G^\circ\) is open and closed, normal, and
geometrically connected. There is a faithfully flat quotient
\(G\to\pi_0(G)\) onto a finite étale group scheme with kernel \(G^\circ\).
This is universal among homomorphisms to étale groups. It commutes
with field extension.
\end{result}

\begin{proof}
The finitely many connected components of \(G_{\bar k}\) are open and
closed. They form a finite group: multiplication sends the product
of two connected components into the component of the product of
any chosen points because products of connected schemes over an
algebraically closed field are connected. The same argument for
inversion gives the group inverse. The morphism to the discrete
component group has as kernel the component of the identity, so
normality includes the scheme structure.

Purely inseparable extensions induce universal homeomorphisms, so the
components over \(\bar k\) identify with those over \(k_s\), and
the finite component set has a continuous absolute Galois action by
group automorphisms. Each open-and-closed piece has an idempotent
in the affine coordinate ring over \(k_s\). All these idempotents descend
to a single finite Galois extension. Descent gives the matching finite
étale scheme and quotient over \(k\) (compare the affine descent in
Result~\ref{res:F02}). It is faithfully flat since this can be checked
over \(\bar k\), where each nonempty fiber is flat over a field.
Over \(\bar k\) a morphism to an étale group is constant on each
component and so factors uniquely. The factorization descends.
The original geometric pieces remain connected over larger
algebraically closed fields, so the construction commutes with
field extension. Over a nonperfect field, the reduced group of
Result~\ref{res:A10} becomes nonreduced after extension to
the algebraic closure. Thus reduction can depend on the field.
See \cite[Propositions 1.29--1.39, 1.52, 2.37 and Definition 2.38]{milne}
and \cite[p.~51 correction]{milneerrata}.
\end{proof}

\begin{result}{A14}{Rational Jordan--Chevalley decomposition}
Suppose \(k\) is perfect. Every \(g\in G(k)\) has unique commuting
factors \(g_s,g_u\in G(k)\) whose images in every finite-dimensional
rational representation are the semisimple and unipotent parts,
respectively. This decomposition is functorial.
\end{result}

\begin{proof}
For an invertible linear operator, over \(\bar k\) split the vector space
into its generalized eigenspaces. Its semisimple part acts as the
corresponding eigenvalue on each, and the quotient by this part
is unipotent. Both descend to a perfect field, for the construction is
invariant under the Galois action over a finite separable splitting
field. Uniqueness follows by restricting a commuting decomposition
to the generalized eigenspaces. Intertwining maps respect the
decomposition. For tensor products the pair of semisimple factors
gives a semisimple tensor product (diagonalize over \(\bar k\)),
and the pair of unipotent factors gives a unipotent one (triangularize).
The same factors commute, hence are the Jordan components of the
tensor product.

Apply this to right translations by \(g\) in all finite-dimensional
subcomodules of \(A=k[G]\). By Result~\ref{res:A06} they form a
directed union. Naturality under inclusions defines linear maps
\(T_s,T_u:A\to A\). Multiplication on \(A\) is equivariant. Restrict it
to the tensor product of any two finite-dimensional subcomodules.
Tensor compatibility and naturality give
\(T_\bullet(ab)=T_\bullet(a)T_\bullet(b)\). Constants are fixed.
Therefore \(\epsilon T_s,\epsilon T_u\) are \(k\)-points \(g_s,g_u\)
of \(G\). To see that their actions in a representation \(V\)
are the required ones, each functional \(\ell\in V^\vee\) gives an
equivariant matrix coefficient map \(V\to A\),
\(v\mapsto(g'\mapsto\ell(g'v))\). Its image is a finite-dimensional
subcomodule. Applying \(\epsilon T_\bullet\) and intertwining on that
subcomodule yields
\(\ell(g_\bullet v)=\ell((g|_V)_\bullet v)\).
A faithful representation detects equality of elements of \(G(k)\).
Thus these elements commute, multiply to \(g\), and are unique.
Applying the same argument to representations pulled back along
a homomorphism proves functoriality.

For a nonperfect field of characteristic \(p\), choose \(a\notin k^p\).
The matrix of multiplication by \(a^{1/p}\) on \(k(a^{1/p})\)
has characteristic polynomial \(X^p-a\) and is invertible. Its
semisimple part over \(\bar k\) is \(a^{1/p}\) times the identity,
which does not descend to \(k\).
See \cite[Theorems 9.2, 9.11, 9.18, Propositions 9.13--9.15 and
Example 9.16]{milne}.
\end{proof}

\section{Chevalley's stabilizer construction}

\begin{result}{A15}{Every closed subgroup stabilizes a line}
For every closed subgroup \(H\subset G\) there is a finite-dimensional
rational \(G\)-module and a line whose scheme-theoretic stabilizer is \(H\).
The assertion holds on points over every \(k\)-algebra, including
nonreduced algebras.
\end{result}

\begin{proof}
Let \(I\subset A=k[G]\) define \(H\). Choose a finite-dimensional
subcomodule \(V\subset A\) for right translation containing \(1\) and
generators of \(I\), and put \(W=V\cap I\). Flatness over the field
gives \(W_R=V_R\cap I_R\) inside \(A_R\).
Every \(h\in H(R)\) preserves \(I_R\) by right multiplication, and
thus preserves \(W_R\). Conversely if \(g\in G(R)\) preserves \(W_R\)
then the translate of each chosen generator lies in \(I_R\).
Evaluation at the identity shows that every generator vanishes at
\(g\). Hence \(g\) factors through \(H\).

Let \(d=\dim W\). If \(I=0\) the trivial module already works.
For \(0<d\le\dim V\) use the line \(\bigwedge^d W\) in
\(\bigwedge^d V\). To check the stabilizer over \(R\), choose a
complement to \(W\) over \(k\), giving a basis \(w_i\) of \(W_R\) and
the exterior generator \(w=w_1\wedge\cdots\wedge w_d\).
Then
\(\{v\in V_R:v\wedge w=0\}=W_R\), since the coefficient along each
complement basis vector is detected by a distinct exterior basis
element. An invertible linear map stabilizing \(Rw\) must act on it
by a unit, and the identity
\(g(v\wedge w)=g(v)\wedge g(w)\) therefore implies it stabilizes \(W_R\).
The converse follows by taking determinants on \(W_R\).
This proves the claim at the level of subgroup functors. It will be
used in Result~\ref{res:D01}.
See \cite[Theorem 4.27, Lemma 4.28 and Remark 4.30]{milne}.
\end{proof}

\chapter{Groups of multiplicative type and solvable groups}
\label{chap:solvable-groups}

The word \emph{split} will always refer to the displayed ground field.
For groups of multiplicative type it refers to a character group with
trivial Galois action. For a smooth connected unipotent group it refers
to a series with additive quotients. The two definitions are used
together when studying solvable groups.

\section{Diagonalizable groups, tori, and weights}

\begin{result}{B01}{Diagonalizable groups and their character groups}
For a finitely generated abelian group \(M\), put
\(D(M)=\operatorname{Spec}(k[M])\), with \(k\)-basis \(e^m\) and
Hopf structure
\[
 \Delta e^m=e^m\otimes e^m,\qquad \epsilon(e^m)=1,\qquad S(e^m)=e^{-m}.
\]
Then \(D(M)(R)=\operatorname{Hom}(M,R^\times)\). The functor \(D\)
is an antiequivalence onto the diagonalizable groups, whose inverse
is the group of characters to \(\mathbb G_m\). Closed subgroups and
quotient groups correspond respectively to quotients and subgroups of \(M\).
\end{result}

\begin{proof}
The identity \(e^{m+n}=e^me^n\) gives the formula for points.
If \(u=\sum_m a_me^m\) is group-like and \(\epsilon(u)=1\), comparing
the coefficients of \(e^m\otimes e^n\) shows that only one \(a_m\)
is nonzero and that it is \(1\). Thus Hopf maps between group
algebras are exactly the maps induced by abelian group homomorphisms
and characters are exactly the \(e^m\).

We also use the corresponding independence fact for any coalgebra
over a field: distinct group-like elements are linearly independent.
In a minimal dependence, express one group-like element as a
linear combination of an independent set. Applying the coproduct
and comparing tensor coefficients forces that combination to have
a unique nonzero term, with coefficient \(1\), contrary to
distinctness. In a Hopf quotient of \(k[M]\), the images of the
\(e^m\) are group-like, and identify exactly the cosets of
\(N=\{m:e^m\mapsto1\}\). The independence fact gives a basis indexed
by \(M/N\), proving the assertion about subgroups. For \(N\subset M\),
\(k[N]\hookrightarrow k[M]\) is faithfully flat directly:
choose one representative from each coset of \(N\) in \(M\).
The corresponding \(e^m\)'s are a nonempty basis over \(k[N]\).
The induced morphism is therefore an epimorphism of fppf sheaves
since it is also of finite presentation. Its kernel is \(D(M/N)\),
since setting all \(e^n=1\) gives \(k[M/N]\).
This realizes the quotient on the group side.
Conversely the characters in a
Hopf subalgebra of \(k[M]\) span it: if \(u\) is in it, applying
coordinate functionals to one factor of \(\Delta u\) shows that
each \(e^m\) in the support of \(u\) is in that subalgebra.
In particular \(D(\mathbb Z^r)=\mathbb G_m^r\) and
\(D(\mathbb Z/n\mathbb Z)=\mu_n\). To obtain the product decomposition,
present the finitely generated abelian group \(M\) as a quotient of
\(\mathbb Z^s\).
Its relation subgroup is finitely generated: projecting a subgroup onto
one coordinate gives an ideal of \(\mathbb Z\), and a lift of its generator
together with generators of the kernel suffices by induction on \(s\).
Use the resulting integer relation matrix. Among all nonzero entries
obtainable by invertible integer row and column operations, choose one
of least positive absolute value, and put it in the upper left corner
with positive sign. Division with remainder shows it divides every
entry in its row and column: otherwise subtraction of an integer
multiple would produce a smaller positive absolute value. Clear both
the row and column, and repeat on the remaining block. The cokernel of
the resulting diagonal matrix is a product of free and finite cyclic
factors. The tensor factorization of the group algebra of a direct
sum gives the asserted products of diagonalizable groups.
See \cite[Proposition 12.3, Lemma 12.4 and
Theorems 12.8--12.9]{milne}.
\end{proof}

\begin{result}{B02}{Multiplicative type and Galois descent of characters}
A group of multiplicative type is an affine group which becomes
diagonalizable over a separable extension of \(k\). Its characters over
\(k_s\) give an antiequivalence with finitely generated abelian groups
with continuous action of \(\Gamma_k=\operatorname{Gal}(k_s/k)\), where
the character group is discrete. Tori correspond to free abelian groups
of finite rank. The character group detects splitness: it has trivial
action precisely when the group is already diagonalizable over \(k\).
\end{result}

\begin{proof}
The coordinate algebras are finitely presented and the diagonalization
and its inverse use finitely many coefficients. Thus one can work over
a finite Galois splitting field \(K\subset k_s\).
For characters defined over \(K\), the descent action is the action
on the basis group-like elements together with the semilinear
action on coefficients. In particular it induces the asserted
action on the discrete group \(M\). Conversely, the action of
\(\Gamma_k\) on a finitely generated discrete group factors through a
finite quotient: take the intersection of the stabilizers of finitely
many generators, and take its open normal core. Let \(K/k\) be the
corresponding Galois extension. Give \(K[M]\) the action
\(\sigma(a e^m)=\sigma(a)e^{\sigma(m)}\).
All three Hopf maps respect it. The ordinary descent of semilinear
vector spaces, including tensor products, in Result~\ref{res:F01}
and its affine-algebra version in Result~\ref{res:F02} descend this
Hopf algebra to \(k\). The descended algebra is of finite type:
after choosing finitely many generators over \(K\), each is a
finite \(K\)-linear combination of descended elements, and the
subalgebra generated by all these elements becomes the full algebra
after faithful scalar extension. The same descent of maps applied
to Result~\ref{res:B01} gives the antiequivalence. In a larger finite
Galois splitting field the descended algebra is canonically the same,
by effectiveness and uniqueness of descent. For field extension
\(L/k\), choose compatible separable closures inside an algebraic
closure of \(L\). The character module over \(L\) is obtained by the
induced restriction of Galois action.
A torus is by definition a form of \(\mathbb G_m^r\) over a separable
extension, whence the free-lattice criterion. With trivial Galois
action the descended algebra is \(k[M]\).
See \cite[Definitions 12.14 and 12.17, Theorems 12.18 and 12.23]{milne}.
\end{proof}

\begin{result}{B03}{Smoothness and isogenies of multiplicative groups}
In characteristic \(p>0\), a group of multiplicative type is smooth
exactly when its geometric character group has no \(p\)-torsion.
In characteristic zero all such groups are smooth.
A smooth connected group of multiplicative type is a torus.
There is an fppf-exact sequence
\[
 1\longrightarrow \mu_n\longrightarrow\mathbb G_m
 \xrightarrow{[n]}\mathbb G_m\longrightarrow1
 \qquad(n>0).
\]
\end{result}

\begin{proof}
After passing to a separable splitting field use the product
decomposition of Result~\ref{res:B01}. A factor
\(k[t,t^{-1}]/(t^m-1)\) is geometrically reduced if \(m\) is prime
to the characteristic, by the derivative \(mt^{m-1}\).
If \(p\mid m\), the defining polynomial has repeated roots over
\(\bar k\). This proves the first two assertions by
Result~\ref{res:A09}. In the smooth case each nontrivial cyclic
torsion summand contributes a disconnected finite étale factor
over \(\bar k\). Geometric connectedness of a connected group
(Result~\ref{res:A13}) rules this out. Conversely the torus has
both properties. The exact sequence is the power-map example
following Result~\ref{res:A07}. Its kernel is nonreduced when
\(p\mid n\), although both terms \(\mathbb G_m\) are smooth.
The same lattice calculation identifies isogenies of tori with
inclusions of character lattices of finite index in the opposite
direction, and computes the kernel with the corresponding finite
cokernel. See \cite[Remark 12.5, Example 12.28 and
Theorems 12.30, 12.32--12.33]{milne}.
\end{proof}

\begin{result}{B04}{Affine Weil restriction and a norm-one torus}
For a finite extension \(K/k\) and an affine finite-type \(K\)-scheme
\(X\), \(\operatorname{Res}_{K/k}X\) represents
\(R\mapsto X(R\otimes_k K)\). The construction is compatible with base
change, interpreted as restriction of scalars for the finite free
algebra obtained after base change. For a separable quadratic field
extension, the kernel of the norm
\(\operatorname{Res}_{K/k}\mathbb G_m\to\mathbb G_m\) is a torus
whose geometric character lattice is \(\mathbb Z\) with the nontrivial
Galois automorphism acting by negation.
\end{result}

\begin{proof}
Write the coordinate algebra of \(X\) as \(K[x_1,\ldots,x_r]/I\).
Fix a \(k\)-basis \(b_1,\ldots,b_d\) of \(K\). Substitute
\(x_i=\sum_j X_{ij}b_j\) into finitely many generators of \(I\).
Take the coefficients in the \(b_j\)-basis as equations over \(k\).
For any \(R\), solving these equations is exactly specifying a
\(K\)-algebra map to \(R\otimes_k K\). This proves representability,
independence of basis, and the base-change assertion by the same
functorial description. Applied to a group it preserves multiplication
and inverses on points. For the torus of units, one can also write
the restriction as the open locus where multiplication by the element
of \(R\otimes_k K\) has invertible determinant. Invertibility of the
operator implies invertibility of the element by applying its inverse
to \(1\). This determinant defines the norm homomorphism.
If \(K/k\) is separable quadratic then \(K\otimes_k K\simeq K\times K\)
via the two embeddings, and the norm there is \((u,v)\mapsto uv\).
The kernel is the one-parameter subgroup \((u,u^{-1})\).
The two factors are exchanged by the nontrivial automorphism,
so this parameter and its character are inverted. For
\(\mathbb C/\mathbb R\), multiplication by \(a+ib\) has determinant
\(a^2+b^2\). The real norm-one group is the affine circle
\(a^2+b^2=1\), acting on the real plane by multiplication.
See \cite[\S2.i and Example 12.27]{milne}.
\end{proof}

\begin{result}{B05}{Weight spaces and exactness of invariants}
For a diagonalizable group \(D(M)\), a finite-dimensional rational
module decomposes canonically as the direct sum of its weight spaces
\(V_m=\{v:\rho(v)=v\otimes e^m\}\). Tensor products add weights and
duals negate them. For every group of multiplicative type, invariants
are exact on finite-dimensional rational modules.
\end{result}

\begin{proof}
Write \(\rho(v)=\sum_m p_m(v)\otimes e^m\) in the group-like basis,
with finite support for each \(v\). Coassociativity gives
\(p_n p_m=\delta_{n,m}p_m\) and the counit gives \(\sum_m p_m=\mathrm{id}\)
on each vector. Thus \(V_m=\operatorname{im}(p_m)\), and over a
finite-dimensional space only finitely many are nonzero. The product
and antipode rules of the group algebra give the two claims about
tensor products and duals. The zero-weight summand is the space of
invariants, so equivariant surjections of split modules remain
surjective on invariants. For a group of multiplicative type, use a
splitting field and the field-extension statement of
Result~\ref{res:A05}. Surjectivity of a linear map descends through a
faithful field extension. Hence invariants, already left exact as a
kernel, are exact. This argument uses no division by the order of
a torsion character. See \cite[Theorem 12.12, Definition 12.52
and Proposition 12.54]{milne}.
\end{proof}

\section{Unipotent groups and the imperfect-field boundary}

\begin{result}{B06}{Unipotence in representations and in coordinates}
Let \(U_n\) be the upper unitriangular group. For an affine algebraic
group \(U\), the following are equivalent:
\begin{enumerate}
\item every nonzero finite-dimensional rational representation has
a nonzero invariant vector.
\item each such representation admits a full invariant flag with
trivial action on one-dimensional successive quotients.
\item \(U\) embeds as a closed subgroup of some \(U_n\).
\item \(A=k[U]\) admits an exhaustive increasing filtration
\(C_0=k\subset C_1\subset\cdots\) such that
\(\Delta(C_r)\subset\sum_{i=0}^r C_i\otimes C_{r-i}\).
\end{enumerate}
We call a group with these properties \emph{unipotent}.
\end{result}

\begin{proof}
Given an invariant line, apply the first condition repeatedly to the
successive quotient representations. The converse follows from the
first nonzero flag term. Applying the flag statement to
Result~\ref{res:A06} gives an embedding in \(U_n\).
In the coordinate algebra of \(U_n\) give \(x_{ij}\), \(i<j\), weight
\(j-i\). The matrix coproduct is
\[
 \Delta x_{ij}=x_{ij}\otimes1+1\otimes x_{ij}
   +\sum_{i<\ell<j} x_{i\ell}\otimes x_{\ell j}.
\]
Each displayed tensor has total weight \(j-i\). Multiplication
proves the filtration identity on monomials. Pass to any Hopf
quotient by taking the images of the filtration.

For the converse let \(V_r=\{v:\rho(v)\in V\otimes C_r\}\).
A nonzero \(v\in V_0\) is fixed by the counit axiom. If \(V_r=0\),
the projected coaction \(V\to V\otimes(A/C_r)\) is injective.
Applying it twice gives an injection into
\(V\otimes(A/C_r)\otimes(A/C_r)\). For \(v\in V_{r+1}\),
the double coaction maps to zero by the filtration on \(\Delta\):
in each summand at least one index is at most \(r\).
Coassociativity then gives \(V_{r+1}=0\). Exhaustiveness gives
\(V\ne0\Rightarrow V_0\ne0\), proving the result.

For smooth \(U\) one may also test unipotence of every matrix in
\(U(\bar k)\) in a faithful representation. Over the algebraically
closed field, take a minimal nonzero invariant subspace \(W\) and let
\(A\) be the linear span on \(W\) of the acting group. This is a unital
subalgebra and \(W\) is a simple \(A\)-module. Each \(A\)-endomorphism
of \(W\) is scalar, since the eigenspace for any of its eigenvalues is
a nonzero \(A\)-submodule. For a basis \((w_1,\ldots,w_d)\), the
submodule \(A(w_1,\ldots,w_d)\) of \(W^d\) must be all of \(W^d\).
Otherwise choose a simple quotient of \(W^d\) killing that submodule.
The quotient is isomorphic to \(W\), as one coordinate copy of \(W\)
has nonzero image. The quotient map is then a nonzero linear
combination of the coordinates by the scalar endomorphism property.
It cannot kill the tuple of basis vectors. Thus \(A\) realizes every
prescribed image of the basis and is the full endomorphism algebra.
The span of \(g-1\) is a two-sided ideal therein, by multiplying on
either side by group elements and subtracting two differences of this
form. It consists of trace-zero matrices, since every group matrix is
unipotent, and is therefore zero. In fact multiplication on both sides
by matrix units would extract a nonzero entry of any nonzero element
of the ideal and produce a matrix unit of trace one. Thus the group
fixes \(W\) pointwise. Induction by quotienting out a fixed line gives
a unitriangular flag. Geometric points of a smooth group are
schematically dense, so the resulting inclusion into \(U_n\) holds
as a group scheme after extension to \(\bar k\), and the invariant
criterion descends by Result~\ref{res:A05}. Without smoothness
\(\mu_p(\bar k)\) is trivial but \(\mu_p\) is not unipotent, by
Results~\ref{res:B05} and \ref{res:B07}.
See \cite[Proposition 14.3, Theorem 14.5 and Corollaries 14.6, 14.11--14.12]{milne}.
\end{proof}

\begin{result}{B07}{Exact operations on unipotent groups}
Closed subgroups, faithfully flat quotients, and extensions of
unipotent groups are unipotent. Unipotence is preserved and detected
by extension of the ground field. A group that is both unipotent
and of multiplicative type is trivial. Every homomorphism between
these two kinds of groups, in either direction, is trivial.
\end{result}

\begin{proof}
For subgroups restrict an embedding into \(U_n\).
For a faithfully flat quotient, view a module for the quotient as
a module for the source. Invariance descends since the quotient
map is an epimorphism of fppf sheaves. In an extension with
unipotent kernel \(N\), a nonzero module \(V\) has \(V^N\ne0\).
Normality gives a rational action on \(V^N\), factoring through the
quotient by faithfully flat descent, so invariants for that quotient
provide an invariant vector for the full group. Ascent follows by
base-changing \(U_n\). Descent follows by scalar extension of
invariants as in Result~\ref{res:A05}. After passing to a splitting
field, each character of a group that is both unipotent and
diagonalizable must be trivial, by applying the invariant-vector
criterion to its one-dimensional weight spaces. A faithful
representation is then trivial. Groups of multiplicative type are
closed under subgroups and the faithfully flat images in question
by Result~\ref{res:B01} after scalar extension. Thus the
scheme-theoretic image of either homomorphism is of both types,
which proves triviality. In particular an intersection of a
unipotent and a multiplicative-type subgroup is trivial as a
scheme. See \cite[Corollaries 14.7, 14.9--14.10 and
Proposition 14.16, Corollaries 14.17--14.18]{milne}.
\end{proof}

\begin{result}{B08}{Split unipotent groups}
A unipotent group is split if it has a subnormal series with all
successive quotients isomorphic to \(\mathbb G_a\).
Every smooth connected unipotent group over a perfect field is split.
Over an arbitrary field a smooth connected unipotent group becomes
split over a finite purely inseparable extension, and if it becomes
split over a separable extension it was already split.
\end{result}

\begin{proof}
The filtration of \(U_n\) by diagonals above the main diagonal gives
a central series. More precisely, let \(U_n^{(i)}\) have vanishing
entries at heights \(1,\ldots,i-1\). Then multiplication of strictly
upper-triangular matrices adds heights, so
\([U_n^{(i)},U_n^{(j)}]\subset U_n^{(i+j)}\). Each height quotient
is an additive vector group. Its coordinate flag refines the series
to one with additive one-dimensional terms. Intersect with a
closed unipotent subgroup and take scheme-theoretic images in the
additive quotients to obtain a central series with factors that
are closed subgroup schemes of \(\mathbb G_a\). In particular all
unipotent groups are nilpotent. Here and below quotients by normal
closed subgroup schemes are the fppf quotients represented in
Result~\ref{res:D01}.

In characteristic zero a nontrivial closed subgroup of \(\mathbb G_a\)
is all of \(\mathbb G_a\). Indeed a proper subgroup is finite
and, by Result~\ref{res:A09}, étale, whereas the additive group of
an algebraically closed field of characteristic zero has no
nontrivial finite subgroup. The preceding central series therefore
gives splitness in this characteristic.

In positive characteristic, let \(k\) be perfect and \(U\ne1\) smooth
connected. In the lower central series, the subgroup generated by a
commutator image of smooth connected groups is smooth and connected:
over \(\bar k\), take the reduced closure of the group generated by
the image. Finite products of that image and its inverse are connected
and contain the identity, and the resulting reduced closed subgroup
descends over \(k\). The same argument is detailed for the derived
subgroup in Result~\ref{res:B10}. The last nontrivial lower central
term is therefore a positive-dimensional, smooth connected central
subgroup \(Z\). Via \(U_n\) it is killed by some power of \(p\):
if \(N\) is strictly upper triangular,
\((1+N)^{p^a}=1+N^{p^a}=1\) when \(p^a\ge n\).
For the smallest \(a\) such that \([p^a]\) kills \(Z\), the
scheme-theoretic image of \([p^{a-1}]\) on \(Z\) is nontrivial,
smooth and connected, and killed by \([p]\). Smoothness follows
from the inclusion of its coordinate ring in that of \(Z\) after
extension to \(\bar k\), and connectedness from the surjective
image map. It is commutative because \(Z\) is central.

For use over arbitrary fields, we prove that a smooth connected
commutative \(p\)-killed unipotent group \(H\) embeds in a vector
group. Use Result~\ref{res:B06} to embed it in \(U_n\), and let \(M\)
be the linear span of the matrix coordinates \(h-1\). More
precisely, \(M\) is cut out by those linear forms for which the
corresponding linear combination of coordinate functions vanishes
on \(H\). This construction commutes with field extension, so
over \(\bar k\) the space is spanned by matrices \(h-1\) for
geometric points \(h\). They commute and their \(p\)-th powers
vanish. Thus \(M\) is commutative and \(v^p=0\) for every \(v\in M\).
It is closed under multiplication because
\[
  (g-1)(h-1)=(gh-1)-(g-1)-(h-1).
\]
For a basis \((v_i)\) adapted to \(M\supset M^2\supset\cdots\), the
map
\[
 (t_i)\longmapsto
      \prod_i\left(\sum_{a=0}^{p-1}\frac{t_i^av_i^a}{a!}\right)
\]
is a homomorphism from a vector group to \(1+M\). For each factor
this follows from the binomial identity and \(v_i^p=0\), and the
factors commute. Modulo \(M^{j+1}\), the coordinates of degree
\(j\) enter linearly in the basis of \(M^j/M^{j+1}\) after the
coordinates of degrees less than \(j\) are fixed. Solving in
increasing degree gives a polynomial inverse. Hence \(H\) is a
closed subgroup of a vector group.

We also need the equations of such a subgroup over any field of
characteristic \(p\). Let \(I\) be a Hopf ideal of
\(A=k[x_1,\ldots,x_d]\), where each \(x_i\) is primitive, and fix
a graded monomial order. Start with an empty list of primitive
polynomials in \(I\). Choose a monic \(f\in I\) whose leading
monomial is least among those not divisible by the leading
monomials already on the list. Reduce its lower terms by the
polynomials on the list. In
\(\Delta f-f\otimes1-1\otimes f\) both entries of each remaining
tensor monomial are proper divisors of a monomial of \(f\) (its
constant term is zero). They are strictly smaller than the leading
monomial of \(f\), and none is divisible by a previous leading
monomial. By the choice of \(f\), these entries are outside the
initial ideal of \(I\), so the tensor monomials are independent in
\((A/I)\otimes(A/I)\). The Hopf-ideal property now forces
\(\Delta f=f\otimes1+1\otimes f\) already in \(A\otimes A\).
Primitive polynomials are precisely the sums
\(\sum_{i,j}c_{ij}x_i^{p^j}\): a monomial involving two variables
produces a mixed term of coefficient \(1\), and for one variable
the expansion of \((1+t)^m\) via the base-\(p\) digits of \(m\)
shows that its intermediate coefficients all vanish exactly for
\(m=p^j\). Terms from distinct monomials cannot cancel such a
coefficient. Thus the new leading monomial is a pure \(p\)-power
in a variable unused by the previous leaders. Indeed the new
minimum is greater than the previous minima, and a greater pure
power of the same variable would be divisible by its previous
leader. This process terminates in at most \(d\) steps. Division
by the chosen polynomials now proves that they generate \(I\).

For perfect \(k\) this also proves that \(H\) is a vector group.
Write its additive equations in a matrix over the skew polynomial
ring \(k[F]\), with \(Fa=a^p F\) and \(F(x)=x^p\). Both division
algorithms exist: one cancels leading coefficients by repeated
subtraction of a monomial multiple on the specified side, using
that Frobenius on \(k\) is an automorphism. Row and column additions
and interchanges change the equations by invertible operations
and the variables by automorphisms of the vector group. Place a
nonzero entry of least degree in the first position and divide
all entries of its column on the left and of its row on the right.
If a remainder is nonzero, use it as pivot, decreasing the degree.
When all entries are divisible on their respective sides, clear
the row and column and repeat on the remaining block. The resulting
diagonal equations exhibit \(H\) as a product of a vector group
and finite kernels of nonzero one-variable additive polynomials.
Upon base change a smooth geometrically connected such product
forces every finite factor to be trivial. This supplies the
central copy of \(\mathbb G_a\) in \(U\) from the \(p\)-killed
image constructed above. The quotient is smooth and connected:
its coordinate ring embeds in that of \(U\) after every base
extension, and its underlying space is a continuous image.
Induction on dimension gives a central split series.

Over an arbitrary field apply the perfect-field assertion to the
perfect closure. A finite series and the isomorphisms witnessing
its exact additive quotients descend to a finite level by finite
presentation and faithful field descent, hence to a finite
purely inseparable extension. Separable descent follows from
the canonical split subgroup construction after
Result~\ref{res:B09}. References for the perfect-field
classification and the additive-polynomial description are
\cite[Theorem 14.46, Propositions 14.49--14.54 and
14.61--14.66]{milne} and \cite[Proposition 5.2.10]{conradfeng}.
\end{proof}

\begin{result}{B09}{A smooth wound unipotent group}
Suppose \(\operatorname{char}k=p>0\) and \(c\notin k^p\).
The subgroup \(W\subset\mathbb G_a^2\) defined by
\(y^p=x-cx^p\) is smooth, geometrically connected and unipotent.
It becomes \(\mathbb G_a\) over \(k(c^{1/p})\), but every
\(k\)-morphism \(\mathbb A^1\to W\) is constant.
\end{result}

\begin{proof}
The relation \(y^p-x+cx^p=0\) is additive in each pair of variables,
so defines an additive subgroup. The partial derivative in \(x\)
is \(-1\), which gives smoothness. With \(d=c^{1/p}\) and
\(t=y+dx\), the equation becomes \(t^p=x\), with inverse
\((x,y)=(t^p,t-dt^p)\). These are additive maps. This proves the
isomorphism after scalar extension and geometric connectedness.
Unipotence also follows directly from Result~\ref{res:B07}.
A morphism from the affine line gives polynomials \(a(s),b(s)\)
with \(b^p=a-ca^p\). If \(a\) is nonconstant, comparison of highest
degrees gives \(\deg b=\deg a\), and the leading coefficients
\(\alpha,\beta\) satisfy \((\beta/\alpha)^p=-c\), impossible.
If \(a\) is constant so is \(b\). This is the wound property, and
in particular there is no subgroup isomorphic to \(\mathbb G_a\).
See \cite[14.69--14.75]{milne} and
\cite[Examples 5.2.7, 5.2.9 and Definition 5.2.8]{conradfeng}.
\end{proof}

We prove the canonical split-subgroup assertion used here. For a
smooth connected unipotent group \(U\) there is a unique smooth
connected split subgroup \(U_{\mathrm{spl}}\) containing every
split connected subgroup. It is normal, the quotient is wound,
and the construction commutes with separable field extensions.
In this setting woundness is equivalent to the absence of a
subgroup isomorphic to \(\mathbb G_a\).

Assume \(\operatorname{char}k=p>0\), and suppose there is a
nonconstant curve \(\phi:\mathbb A^1\to U\) after a separable field
extension. There is then such a curve over \(k_s\). Indeed, the
finitely many coefficients of its affine coordinate polynomials,
the relations between them and a nonzero nonconstant coefficient
spread to a finitely generated separably generated field over
\(k\). After clearing denominators in a separable presentation,
this field has a nonempty smooth affine model over \(k\) on
which these relations persist and the coefficient remains
nonzero. The model has a \(k_s\)-point by
Result~\ref{res:F04}, so evaluation gives the curve. Translate
so that \(\phi(0)=1\). If the curve is not central over \(k_s\),
smoothness gives an element \(g\in U(k_s)\) such that
\(g\phi(t)g^{-1}\phi(t)^{-1}\) is nonconstant. To see the needed
density, every nonempty open of \(U_{k_s}\) has a \(k_s\)-point
by Result~\ref{res:F04}, and a geometric vanishing equation on
these points descends coefficient by coefficient in a vector
space basis of its field of definition. The commutator has the
same value \(1\) at \(t=0\). Each such operation raises the
unitriangular height filtration used in Result~\ref{res:B08},
so after finitely many operations we find a nonconstant
central curve. Replace it by its last nonconstant \(p\)-power,
which is central and killed by \(p\).

Among smooth connected subgroups of \(Z(U_{k_s})[p]\) there is a
largest one \(F_s\). The schematic closure of the product of
two such subgroups is again such a subgroup, since the factors
are central. Its coordinate ring embeds in that of the product
even after geometric base change, so it is geometrically
reduced and connected. A subgroup of maximal dimension thus
contains every other. The central \(p\)-killed curve lands in
\(F_s\): the closures of the images of products of finitely
many copies of the curve and its inverse are geometrically
integral, contain the identity and stabilize by dimension,
giving a smooth connected subgroup. The largest subgroup is
Galois invariant by uniqueness. Its finitely many ideal
generators are defined after a finite Galois extension and
the ideal descends by Result~\ref{res:F02} to a smooth
connected central \(p\)-killed subgroup \(F\) over \(k\).

Use the vector-group embedding and additive equations of
Result~\ref{res:B08} for \(F\), and write the curve over a
finite Galois extension \(K/k\) in those coordinates with zero
constant terms. Group the powers in the coordinate polynomials
by their prime-to-\(p\) part. For any fixed \(e\) prime to \(p\),
replace each term \(c_{j,ep^n}t^{ep^n}\) in that group by
\(c_{j,ep^n}t^{p^n}\). This produces a homomorphism
\(\psi:\mathbb G_{a,K}\to F_K\), since each additive equation
preserves the prime-to-\(p\) part of exponents and the
substitution \(t\mapsto t^e\) is injective on \(K[t]\).
Some choice of \(e\) produces a nonzero \(\psi\).
For \(a\in K\), the sum of its Galois translates
\[
 t\longmapsto
 \sum_{\sigma\in\operatorname{Gal}(K/k)}
       \sigma(\psi)\bigl(\sigma(a)t\bigr)
\]
descends to a homomorphism \(\mathbb G_a\to F\). We can make it
nonzero, since for a nonzero coefficient \(c\) of \(\psi\) in
degree \(p^n\), its corresponding coefficient is
\(\operatorname{Tr}_{K/k}(ca^{p^n})\). The \(k\)-linear
span of \(K^{p^n}\) is a finite-dimensional subalgebra of
\(K\), hence a field. The extension over that span is purely
inseparable and separable, so the span is \(K\). Nondegeneracy
of trace for finite separable extensions gives the required
\(a\). The kernel of this nonzero homomorphism is finite,
defined by a monic nonzero additive polynomial by the Hopf-ideal
argument above. That polynomial induces a faithfully flat map
\(\mathbb G_a\to\mathbb G_a\) (the domain coordinate ring is
free with the usual monomial basis over the target). Hence
Result~\ref{res:A07} identifies the image of the original
homomorphism with a central copy of \(\mathbb G_a\).
This proves both the woundness criterion and its persistence
under separable extensions.

For the canonical subgroup, put \(U_{\mathrm{spl}}=1\) if \(U\)
is wound. Otherwise, take the central copy of \(\mathbb G_a\)
just constructed and induct on its smooth connected quotient.
Pulling back the quotient's normal split subgroup gives a
normal split subgroup of \(U\) with wound quotient. Every
homomorphism from a split group to a wound group is trivial:
apply the no-curve condition to the bottom additive subgroup
of the subnormal series, factor through its fppf quotient, and
induct on the series. Thus the construction contains all split
subgroups and is unique. After separable extension its
split/wound properties persist, which proves base change. In
characteristic zero Result~\ref{res:B08} already gives splitness.
A purely inseparable extension can enlarge the split subgroup
as the calculation for \(W\) demonstrates. Compare
\cite[Appendix B, Proposition B.3.4]{cgp} and
\cite[the paragraph after Proposition 5.2.10]{conradfeng}.

\section{Triangularization and the unipotent radical}

\begin{result}{B10}{Solvability and triangularizability}
Define \(D(G)\) as the smallest normal closed subgroup through which
the commutator map \(G\times G\to G\) factors. The group \(G\) is
solvable when its scheme-theoretic derived series terminates at \(1\).
Write \(T_n\) for the upper triangular invertible group. The following
conditions on an affine algebraic group are equivalent:
\begin{enumerate}
\item every simple finite-dimensional rational representation is
one-dimensional.
\item every finite-dimensional rational representation admits a full invariant flag.
\item \(G\) embeds as a closed subgroup of some \(T_n\).
\item \(G\) has a normal unipotent subgroup \(U\) such that \(G/U\)
is diagonalizable over \(k\).
\end{enumerate}
In particular \(T_n\) is split solvable.
\end{result}

\begin{proof}
The definition of \(D(G)\) is meaningful because arbitrary
intersections of closed subgroups are closed subgroups in this
affine setting, using sums of defining ideals. It is the kernel
needed to make the fppf quotient commutative. For a smooth connected
\(G\) over \(\bar k\), \(D(G)\) is smooth connected and its geometric
points form the ordinary derived subgroup. To see this, the set
\(C\) of commutators is a constructible irreducible subset containing
\(1\) and stable under inverse: it is the image on geometric points
of \(G^2\), and a smooth connected algebraic group over an
algebraically closed field is irreducible. The closures of the
products \(C^r\) are nested irreducible closed subsets.
Their dimensions eventually stop increasing. At the first stable
step, multiplication on either side by an element of \(C\) preserves
that closure, as it sends the dense \(C^r\) into \(C^{r+1}\).
All subsequent closures then coincide. The reduced closure of the
generated subgroup is consequently a closed connected subgroup,
normal by functoriality of commutators and scheme density.
One of the constructible sets \(C^r\) contains a dense open in that
closure. As translates of a dense open intersect, every geometric
point of the closure is a quotient of elements in the generated
subgroup. The reduced closure equals \(D(G)\): it contains the
scheme image of the commutator map by reducedness of \(G^2\), and
it is contained in any closed subgroup containing that image.
This argument also verifies smoothness and connectedness.

For the equivalences, choose a simple submodule of a nonzero module
and iterate on its quotient to obtain a flag. A flag on a faithful
representation gives the closed embedding in \(T_n\).
The diagonal map \(T_n\to\mathbb G_m^n\) has kernel \(U_n\).
For a closed subgroup \(G\subset T_n\) take the scheme image of its
diagonal map and the corresponding kernel \(G\cap U_n\).
Results~\ref{res:A07}, \ref{res:B01} and \ref{res:B07} give the
required quotient and normal kernel. Finally \(V^U\ne0\) in every
nonzero representation under the fourth condition. It is stable
under \(G\) and the quotient acts on it by weight spaces, giving
a one-dimensional \(G\)-submodule. This gives the first condition.
For \(T_n\), the central series of \(U_n\) from
Result~\ref{res:B08}, with character-compatible coordinate flags,
and a coordinate flag on the diagonal torus give a subnormal series
with factors \(\mathbb G_a\) or \(\mathbb G_m\). This is the definition
of \emph{split solvable}. See \cite[Corollary 6.19, Propositions
6.20, 6.30, Definition 6.33, Propositions 16.2--16.4]{milne}.
\end{proof}

\begin{result}{B11}{The Lie--Kolchin theorem}
Suppose \(k\) is algebraically closed and \(G\) is smooth, connected
and solvable. Every nonzero finite-dimensional rational representation
of \(G\) has an invariant line and hence a full invariant flag.
\end{result}

\begin{proof}
It suffices to consider a simple \(G\)-module \(V\).
Induct on \(\dim G\). If \(G=1\) the conclusion is immediate.
The smooth connected derived subgroup \(N=D(G)\) is proper if
\(G\ne1\) by solvability, hence has lower dimension. Induction gives
some nonzero \(N\)-eigenspace \(V_\chi\).
The nonzero \(N\)-eigenspaces for distinct characters are linearly
independent, by independence of group-like elements as in
Result~\ref{res:B01} applied to their coactions. In particular there
are only finitely many. Normality of \(N\) makes \(G(k)\) permute
these spaces. The orbit of \(V_\chi\) in its Grassmannian has
only finitely many geometric points and is connected, since it is
the image of a connected group. Therefore \(G(k)\) fixes
\(V_\chi\). By smoothness and Result~\ref{res:A05}, it fixes the
subspace scheme-theoretically. Simplicity gives \(V_\chi=V\).
The determinant character of \(G\) is trivial on \(N\).
Thus the scalar action of \(N\) satisfies \(\chi^{\dim V}=1\).
A morphism from a smooth connected group to \(\mu_{\dim V}\) is
trivial: it factors through the reduced part of this finite group
over \(k\) and lands in its identity component. Hence \(N\) acts
trivially.

The matrices of \(G(k)\) acting on \(V\) now commute. A family of
commuting endomorphisms over an algebraically closed field has a
common eigenvector: if some member is nonscalar, its proper
nonzero eigenspace is preserved by every member and one argues
by induction on the dimension. If all are scalar every nonzero
vector is suitable. Geometric point density again gives an invariant
line for the group scheme. Taking quotients gives the full flag.
See \cite[Theorem 4.25, Theorem 16.30 and
Corollaries 16.31--16.34]{milne}.
\end{proof}

\begin{result}{B12}{The unipotent radical of a solvable group}
If \(G\) is a smooth connected solvable group over a perfect field,
it has a smooth connected normal unipotent subgroup \(R_u(G)\)
containing every unipotent subgroup of \(G\). The quotient is a torus.
The construction commutes with extension of the field.
\end{result}

\begin{proof}
Over \(\bar k\), choose a triangular faithful representation by
Result~\ref{res:B11}. Let \(U\) be the kernel of the diagonal map
and \(D\) its scheme image. The diagonal group is split and
Result~\ref{res:B01} identifies \(D\) as diagonalizable.
By Result~\ref{res:A07} the map onto \(D\) is faithfully flat.
Its coordinate ring is a subring of the reduced coordinate ring
of \(G\), so it is smooth by Result~\ref{res:A09}. It is connected
as the surjective image of \(G\), and a smooth connected
diagonalizable group is a torus. The kernel \(U\) is a closed
subgroup of the upper unitriangular group and is unipotent by
Result~\ref{res:B06}.

To see that this scheme kernel is smooth and connected, enlarge
the algebraically closed field to an algebraic closure \(L\)
of the function field of \(D\). The generic point, viewed as
\(d\in D(L)\), has the property that its powers are schematically
dense in \(D_L\). Indeed the values \(\chi(d)\) of distinct
characters are distinct invertible elements of \(L\). If a
Laurent polynomial with \(m\) distinct characters vanishes at
\(d^0,\ldots,d^{m-1}\), its coefficients are zero by the
Vandermonde determinant on these distinct values.
Choose \(g\in G(L)\) lifting \(d\). Its semisimple part \(g_s\)
belongs to \(G(L)\) and maps to \(d\), by functoriality of
Result~\ref{res:A14}. The reduced closure \(M\) of its powers
is a closed subgroup of a diagonal torus after diagonalizing
\(g_s\) in the faithful representation. Its map to \(D_L\) is
faithfully flat, since its image contains all the powers of \(d\).
Its scheme kernel is trivial, as it is both a subgroup of \(M\)
and of the unipotent \(U_L\) (Result~\ref{res:B07}).
The map \(M\to D_L\) is therefore an isomorphism by the
quotient assertion of Result~\ref{res:A07}. Multiplication
using its inverse section identifies the underlying scheme
of \(G_L\) with \(U_L\times D_L\). Projection onto \(U_L\)
proves connectedness, and the inclusion of its coordinate
ring into the reduced ring \(L[G]\) proves smoothness over
the algebraically closed \(L\). Descent gives smoothness and
geometric connectedness of \(U\) over the original algebraic
closure.

Every unipotent subgroup maps trivially to \(D\) by
Result~\ref{res:B07}, hence lies in \(U\), regardless of
smoothness of the subgroup. This characterizes \(U\)
independently of the triangular representation. Over the
perfect field \(k\), the subgroup equations descend to a finite
separable extension. In a Galois enlargement the characterization
makes this subgroup invariant, so it descends along
Result~\ref{res:F02}. The diagonal quotient descends with it,
as the kernel quotient is characterized as an fppf sheaf.
It is a torus geometrically. The same quotient then
characterizes the base-changed subgroup over any extension
because unipotent groups have no homomorphism to a torus.
The construction of a section once the kernel is smooth
and connected is given in Result~\ref{res:B13}.

For comparison, over an arbitrary field a maximal-dimensional
smooth connected normal unipotent \(k\)-subgroup is still the
largest one. For two such subgroups use the closure of the
product morphism, whose coordinate ideal is its kernel on
functions and thus commutes with extension of the field.
The geometric product-closure argument of
Result~\ref{res:C01} verifies smoothness, connectedness,
normality and unipotence of this closure. The resulting
largest \(k\)-subgroup need not equal the geometric unipotent
radical after base change, as the Weil-restriction example
following that result shows.
Compare \cite[Theorems 16.26 and 16.33]{milne}.
\end{proof}

\begin{result}{B13}{Splitting and conjugacy of maximal tori}
Over an algebraically closed field a smooth connected solvable group is
\(R_u(G)\rtimes T\) for a maximal torus, and any two maximal tori
are conjugate by an element of \(R_u(G)(k)\).
More generally a split solvable group over any field admits such a
decomposition with a split torus, and its maximal split tori are
conjugate by \(R_u(G)(k)\).
Over a perfect field, a split torus acting on a smooth connected
unipotent group admits a stable central series of that group with
additive one-dimensional factors.
\end{result}

\begin{proof}
We first establish the last assertion. Let \(T\) act on \(U\).
The group \(U\rtimes T\) is triangularizable by
Result~\ref{res:B10}. In one faithful triangular embedding its
unipotent group \(U\) meets the height filtration of \(U_n\).
Within each height quotient refine by a flag of coordinate spaces.
Conjugation by \(T_n\) acts on a fixed height quotient through the
diagonal entries, so each flag is stable. The successive images
of \(U\) in these factors embed in one-dimensional additive groups.
The series is central because commuting with an element of \(U_n\)
raises the height. Take the reduced identity component of each
intersection. Over the perfect field these are smooth connected
and remain stable under \(T\). The commutator maps from their
products with \(U\) factor through the corresponding next reduced
identity component by reducedness and connectedness of the domain.
The resulting quotients have dimension at most one, since before
taking identity components their successive images were closed
in \(\mathbb G_a\) and the removal of other components does not
change dimension. A nontrivial smooth connected one-dimensional
unipotent quotient over a perfect field is \(\mathbb G_a\) by
Result~\ref{res:B08}. Delete duplicate terms to obtain the claimed
central series. This argument is the matrix-filtration construction
of \cite[Theorem 16.21 and Corollaries 16.22--16.24]{milne}.

We give the splitting and conjugacy argument for an extension
\(1\to U\to G\to T\to1\) with \(T\) split and \(U\) smooth connected
over a perfect field. In the same triangular embedding of \(G\)
(as allowed by Result~\ref{res:B10}), the intersection series just
constructed is \(G\)-stable since conjugation by \(G\) on the height
quotients factors through its diagonal map. Thus there is such a
series even before a group section of \(G\to T\) has been chosen.
The conjugation action on each one-dimensional factor is linear:
after passage to \(\bar k\), automorphisms of \(\mathbb G_a\) as an
algebraic group have degree one (their polynomial inverse forces
this), hence the higher coefficients vanish on geometric points.
The acting torus is smooth, so this gives a character action
scheme-theoretically. The action factors through \(T\) by centrality.

For a linear \(T\)-module \(W\) we use the following averaging
argument on regular cochains. The linear functional
\(k[T]=k[M]\to k\) taking the coefficient of \(e^0\) is invariant
under translations, even translations by points in algebras,
since \(\Delta(e^m)=e^m\otimes e^m\).
A homogeneous regular cochain of degree \(n\) is an equivariant
regular map \(T^{n+1}\to W\) for simultaneous left translations
and the given \(T\)-action on \(W\). The differential alternately
omits one argument. Sent to the average after insertion of a
first variable, such a cochain remains equivariant by translation
invariance. These averaging maps \(s\) satisfy
\(ds+sd=\mathrm{id}\) in positive degree: in the alternating sum
all terms with an inserted variable and an omitted old variable
cancel, and the remaining term is constant in the averaged
variable. Consequently regular \(1\)- and \(2\)-cocycles are
coboundaries. One passes to the usual inhomogeneous formulas
by expressing a tuple using its first entry and the successive
ratios of its entries.

At each step in lifting a section along the central series, its
inverse image is an extension of \(T\) by an additive group. The
underlying additive torsor over the affine scheme \(T\) has a
regular section: an additive torsor on an affine scheme is
trivial by faithfully flat descent of a quasi-coherent module
(the additive Čech complex of a faithfully flat affine cover
is exact in positive degrees, as becomes split after extending
scalars to the cover).
The defect of multiplicativity of a regular section \(j\) is
\(j(t)j(t')j(tt')^{-1}\in W\), a regular \(2\)-cocycle by
associativity. Correct \(j\) by a regular \(1\)-cochain whose
coboundary cancels that defect. This gives a group section, and
iteration gives a group section of \(G\to T\).
Once one group section exists, comparing it with another gives
a regular crossed homomorphism into \(U\). Induction on the same
series reduces first to conjugacy modulo the initial central
\(\mathbb G_a\), lifting the required conjugator on \(k\)-points
by the triviality of an additive torsor over \(\operatorname{Spec}k\).
The remaining difference is a regular \(1\)-cocycle into that
central additive group and hence is the coboundary associated
with an element of it. This proves conjugacy by \(U(k)\).

This proves splitting and conjugacy over an algebraically closed
field, using Result~\ref{res:B12}. We treat split solvable groups
over arbitrary \(k\) separately. A faithful representation is
triangularizable: the first, split-series part of the fixed-point
proof in Result~\ref{res:B14} applies to the full flag variety with
its rational point. The kernel \(U\) of the diagonal map is smooth
geometrically connected by the geometric triangular-kernel
argument of Result~\ref{res:B12}; its quotient \(D\) is a split torus.
In fact \(U\) is split. Intersect it with the terms \(G_i\) of the
given split solvable series. Their intersections \(U_i\) are the
smooth geometrically connected kernels of the diagonal maps of
\(G_i\), by the same argument. Each \(U_i/U_{i+1}\) is the
image in the one-dimensional quotient \(G_i/G_{i+1}\).
An image in \(\Gm\) is trivial by Result~\ref{res:B07}.
An image in \(\Ga\) is smooth and connected, since \(U_i\)
is geometrically reduced and connected. It is consequently
trivial or all of \(\Ga\). Deleting trivial successive terms
gives the split series.

Take a geometrically maximal \(k\)-torus by the existence argument
in Result~\ref{res:C02}, which uses only the geometric, algebraically
closed splitting and conjugacy just proved above, rather than this
arbitrary-field part. Its map to \(D\) is an isomorphism: the kernel
is trivial by Result~\ref{res:B07} and its geometric dimension is
\(\dim D\) by the geometric splitting above. Thus we get a split
torus section and \(G=U\rtimes T\). Any split torus in \(G\)
extends to one of these: take a maximal \(k\)-torus in its smooth
connected centralizer by Result~\ref{res:C02}. That torus contains
the given split torus by centrality, and is geometrically maximal
in \(G\), since a geometric maximal torus extending the given
torus belongs to the centralizer. This identifies maximal split
tori with complements of \(U\).

For their rational conjugacy we give a factorization of \(U\).
Use a faithful triangular embedding \(G\subset\GL(V)\) and
choose a basis in which the same invariant flag is respected
and the split \(T\) acts diagonally. This is possible by
the weight decomposition of each flag term.
Choose a cocharacter \(\lambda\) distinguishing the finitely
many distinct weights of \(T\) on \(V\).
The ambient unitriangular matrices factor uniquely into
the negative limit subgroup, the centralizer of \(\lambda\),
and the positive limit subgroup, in that order.
In fact solve the equation one superdiagonal at a time.
At a fixed height the entries of the product are the sums
of the entries of that height in each factor, plus terms
from smaller heights already known. Each position belongs
to exactly one factor according as its weight is negative,
zero, or positive. This proves existence and uniqueness,
with polynomial inverse, on every test algebra.

The factors of an element of \(U\) in this ambient product
belong to \(G\). Take a line \(kv\) in a representation
of \(\GL(V)\) with scheme stabilizer \(G\), by
Result~\ref{res:A15}. Move the middle factor \(z\) to the
right by normalizing the positive limit group, and write the
element as \(u_-u'_+z\), taking \(v\) to \(cv\) with \(c\)
a unit. Conjugation by \(\lambda(t)\) preserves this scalar
on the line. Thus on the punctured affine line
\[
 (\lambda(t)u'_+\lambda(t)^{-1})\,zv
   =c\,(\lambda(t)u_-^{-1}\lambda(t)^{-1})v .
\]
The left expression is polynomial at \(0\) and the right
is polynomial in \(t^{-1}\) at \(\infty\). A Laurent
polynomial with both properties is constant over any test
algebra. Its two endpoint values give \(zv=cv\);
substitution and \(t=1\) give \(u'_+v=v=u_-^{-1}v\).
Hence all factors belong to \(G\) as claimed, and being
unitriangular they belong to \(U\).
The central factor is \(Z_U(T)\) since the choice of
\(\lambda\) detects every nonzero weight difference.
Reordering the factors using that central factor's
normalization of the two limit groups gives a scheme
product \(U_-\times U_+\times Z_U(T)\simeq U\).

Over \(\bar k\), any other complement is \(uTu^{-1}\)
as already proved. Its conjugators in \(U\) form a right
coset of \(Z_U(T)\): a point of \(U\) normalizing \(T\)
centralizes it, since the quotient \(D\) restricts
isomorphically to \(T\). The product just proved gives a
\(k\)-section of the homogeneous quotient \(U/Z_U(T)\),
with image \(U_-U_+\).
The conjugator coset for a \(k\)-complement is a rational
point of this quotient: as a sheaf it consists of local
conjugating elements differing by \(Z_U(T)\), and nonempty
geometrically; equivalently the equations defining the
coset descend along a finite field extension introducing a
geometric conjugator. Therefore its representative in the
section belongs to \(U(k)\). This proves rational
conjugacy without a perfection hypothesis.
Compare \cite[Proposition 16.54]{milne} and
\cite[Appendix A.2]{conradfeng}.
\end{proof}

\begin{result}{B14}{The Borel fixed-point theorem}
Suppose a split solvable \(k\)-group \(G\) acts on a proper
finite-type \(k\)-scheme \(X\) with \(X(k)\ne\varnothing\).
Then the fixed-point subscheme has a \(k\)-point. In particular,
a smooth connected solvable group over an algebraically closed
field has a fixed point on every nonempty proper \(k\)-scheme
on which it acts.
\end{result}

\begin{proof}
First let \(G=\mathbb G_a\) or \(\mathbb G_m\).
For \(x\in X(k)\), extend the orbit map \(t\mapsto t\cdot x\)
from \(G\subset\mathbb P^1\) to \(\mathbb P^1\).
At each missing point the local ring of the nonsingular curve
is a discrete valuation ring, so the valuative criterion for
properness gives a unique extension, which is defined on
a neighborhood because \(X\) is of finite type.
Uniqueness along the overlap gives the desired global morphism.
The translation action of \(\mathbb G_a\) on \(\mathbb P^1\)
fixes the point \(\infty\). The multiplication action of
\(\mathbb G_m\) fixes \(0\) and \(\infty\).
The extended map is equivariant: the two maps on
\(G\times\mathbb P^1\) agree on \(G\times G\), which is dense in
the integral source, and the equalizer is closed by separation.
Thus the image of a boundary fixed point is a \(k\)-rational
scheme-theoretic fixed point.

For the general group choose the subnormal series with factors
\(\mathbb G_a\) or \(\mathbb G_m\). Work from the smallest
nontrivial term upwards. At each step the fixed-point scheme
for the previous term is closed by Result~\ref{res:A08} and
therefore proper. It is stable under the next term by normality.
Its action descends to the next quotient since the kernel acts
trivially and the quotient map is an epimorphism of fppf
sheaves. The first case now preserves nonemptiness of the
\(k\)-rational fixed locus along every step. Over an algebraically
closed field, Results~\ref{res:B08}, \ref{res:B12} and
\ref{res:B13} give the split solvable series, and a nonempty
finite-type scheme has a rational point. This proves the
stated special case. See \cite[Theorem 16.51 and Corollary 16.53]{milne}.
\end{proof}

\begin{result}{B15}{Boundary examples for triangularization}
The real norm-one torus has a two-dimensional rational
representation without an invariant real line. Over \(\mathbb C\),
the finite constant solvable group \(S_3\) has a simple
two-dimensional rational representation. In characteristic \(p\),
vanishing on geometric points alone does not detect an
infinitesimal group.
\end{result}

\begin{proof}
The action in Result~\ref{res:B04} includes multiplication
by \(i\), a quarter-turn on \(\mathbb R^2\).
It has no real eigenline. For \(S_3\), take the invariant plane
\(x_1+x_2+x_3=0\) in the permutation representation.
A \(3\)-cycle has distinct eigenvalues \(\zeta,\zeta^{-1}\)
there, for a primitive cube root of unity \(\zeta\).
A transposition exchanges their two eigenlines, since it conjugates
the cycle to its inverse. A line invariant under \(S_3\) would
be one of those eigenlines and would also be invariant under
the transposition. No such line exists. The group is solvable
since its subgroup \(A_3\) and the quotient \(S_3/A_3\) are
abelian, but its constant group scheme is disconnected.

Finally \(\mu_p(\bar k)\) and \(\alpha_p(\bar k)\) are both trivial,
while each algebra of Result~\ref{res:A10} has length \(p\).
For example the character \(\mu_p\to\mathbb G_m\) is
nontrivial on the dual-number point \(1+\varepsilon\).
These checks isolate the ground-field, connectedness, and
smoothness hypotheses in Results~\ref{res:B06} and \ref{res:B11}.
Compare \cite[Example 12.27 and Theorem 16.30]{milne}.
\end{proof}
\chapter{Reductive groups and root data}
\label{ch:reductive}

We fix a field \(k\). A torus is called split if it is isomorphic over \(k\) to a product of copies of \(\Gm\). A reductive group is called split if it contains a split torus that is a maximal torus after extension to an algebraic closure. Write \(k_s\) for a separable closure and \(\bar k\) for an algebraic closure. A statement about a group scheme will be read on all test algebras, unless a field of points is explicitly indicated.

\section{Reductivity and maximal tori}

\begin{result}{C01}{The geometric radical and reductivity}
Let \(G\) be a smooth geometrically connected affine group of finite type over \(k\). Over \(\bar k\), let \(R_u(G_{\bar k})\) be its largest smooth connected normal unipotent subgroup and \(R(G_{\bar k})\) its largest smooth connected normal solvable subgroup. We call \(G\) \emph{reductive} if \(R_u(G_{\bar k})=1\), and \emph{semisimple} if \(R(G_{\bar k})=1\). Both conditions are invariant under arbitrary extensions of \(k\). Over a perfect field the two geometric radicals are the base changes of the corresponding smooth connected normal radicals over \(k\).
\end{result}

\begin{proof}
The product of two smooth connected normal subgroups over an
algebraically closed field has a reduced scheme-theoretic closure \(H\)
that is again a smooth connected normal subgroup. Indeed, the
closure is irreducible because the source of multiplication is
irreducible. At field-valued points the product of the subgroups is
a normal subgroup of \(G(k)\). Taking closures under inversion,
multiplication and conjugation preserves these inclusions, using
density of the product of dense sets in the product of varieties.
The inclusions hold as morphisms of schemes because the source
varieties are reduced. Smoothness follows from
Result~\ref{res:A09}.

If both factors are solvable then their product of point groups is
solvable: the first factor is normal and the quotient is a quotient
of the other. Taking the reduced closure of each term of the
derived series of the point group gives a series in which each
commutator map lands in the next term. This follows from the same
product density argument. The series ends at the identity.
Equivalently, solvability of such a group with a given bound on
derived length is the identity of the corresponding derived word
(iterating the commutator on two disjoint sets of variables).
This equivalence follows by expanding commutators of products
using conjugates of commutators of the individual factors.
If both factors \(U_1,U_2\) are unipotent, let \(V\ne0\) be a
rational representation of \(H\). The subspace \(V^{U_1}\) is
nonzero by Result~\ref{res:B06} and \(H\)-stable by normality.
It contains a nonzero \(U_2\)-fixed vector. The products act
trivially on that vector and hence so does \(H\), since they
are schematically dense. Result~\ref{res:B06} proves unipotence
of \(H\). For each of the two classes choose a subgroup of
maximal dimension, starting with the trivial subgroup. Its product
closure with any other eligible subgroup is in the same class and
contains the chosen subgroup. The closure is irreducible and has
the same dimension as the chosen subgroup, hence coincides with it
as a reduced subgroup. This proves existence and uniqueness of the
largest subgroups.

We record the descent detail used for field extension. For an
extension \(F\subset L\) and a family of \(F\)-automorphisms of
\(L\) with common fixed field \(F\), any semilinearly stable subspace
of \(L\otimes_F F[G]\) descends. Choose an increasing sequence of
finite-dimensional subspaces of \(F[G]\) with union \(F[G]\).
In each, choose an \(F\)-basis and intersect its scalar extension
with the stable subspace. The unique reduced row echelon matrix
describing that intersection has entries fixed by every
automorphism, since applying an automorphism preserves its echelon
form and its row space. Its rows therefore give a basis over \(F\).
Taking the union gives the descended subspace. If the subspace is a
Hopf ideal defining a normal subgroup, the ideal, Hopf, and
normality conditions descend as inclusions of vector spaces,
detected by faithful flatness. Smoothness and geometric
connectedness descend under field extension. A derived-word identity
descends by equality of morphisms, and unipotence descends by
Result~\ref{res:B07}.

For an extension \(F\subset L\) of algebraically closed fields,
the automorphisms of \(L\) over \(F\) have fixed field \(F\).
An element \(a\notin F\) is transcendental over \(F\), can be
included in a transcendence basis, and can be moved by replacing
that coordinate by \(a+1\). This automorphism of the rational
function field extends to its algebraic closure \(L\). The two
largest subgroups over \(L\) have stable Hopf ideals by uniqueness
and hence descend as above. The descents qualify in the respective
classes, so they lie in the radicals over \(F\). Conversely the base
changes of the radicals over \(F\) qualify over \(L\). This proves
equality, and embedding an algebraic closure of \(k\) into one of
any extension proves the two geometric invariance assertions.
Finally, when \(k\) is perfect the fixed field of
\(\operatorname{Gal}(\bar k/k)\) is \(k\). Applying the same descent
to the unique geometric radicals constructs both subgroups over
\(k\). Any eligible \(k\)-subgroup becomes a subgroup of the
corresponding radical after base change, so its inclusion descends.
See \cite[Chapter 19]{milne} for the geometric radical construction.
\end{proof}

The condition in Result~\ref{res:C01} is geometric. To see why, let \(K/k\) be a nontrivial finite purely inseparable extension. Put \(G=\Res_{K/k}\Gm\), the affine Weil restriction from Result~\ref{res:B04}. This group is smooth: in the affine space underlying \(K\), it is the open set on which the determinant of multiplication is invertible. Over \(\bar k\), the algebra \(K\otimes_k\bar k\) is local with residue field \(\bar k\) and a nonzero nilpotent maximal ideal \(J\). There is a split exact sequence on units over \(\bar k\), with kernel \(1+J\) and quotient \(\Gm\). The kernel has additive successive quotients \(J^i/J^{i+1}\) via \(1+x\mapsto x\), so it is a nontrivial smooth connected unipotent group. The description as a product on the underlying schemes also proves geometric connectedness of \(G\). Over \(k\), any smooth connected unipotent subgroup \(U\subset G\) would, by adjunction, give a character \(U_K\to\Gm\). This is trivial, since a smooth connected unipotent group has no characters after algebraic closure. The inclusion \(U\to G\) must therefore be trivial. This is the basic pseudo-reductive example \cite[Example 1.1.3]{cgp}. Linear reductivity instead concerns complete reducibility of rational representations, the subject of Chapter~\ref{ch:representations}.

\begin{result}{C02}{Maximal tori and the Weyl group scheme}
An affine connected smooth algebraic group over \(k\) has a maximal \(k\)-torus that remains maximal over \(\bar k\). Geometric maximal tori are conjugate. If \(G\) is reductive and \(T\subset G\) is a maximal torus, then the scheme-theoretic centralizer is \(Z_G(T)=T\), and
\[
  W(G,T)=N_G(T)/T
\]
is finite \'etale. For split \(G\), its split maximal tori are conjugate by \(G(k)\), \(W(G,T)\) for such a torus is constant, and each Weyl element has a representative in \(N_G(T)(k)\).
\end{result}

\begin{proof}
Work first over an algebraically closed field. Geometric conjugacy
of maximal tori was proved in the geometric part of
Result~\ref{res:F09}: each is contained in a Borel subgroup, the
Borels are conjugate, and the tori inside one Borel are conjugate
by Result~\ref{res:B13}. We will use the homogeneous space \(G/B\)
for a Borel \(B\). It is proper: choose \(B\) as the reduced
identity component of a stabilizer \(H\) of a point of the
closed flag orbit used in that argument. Then \(G/H\) is the
closed orbit by Result~\ref{res:D01}. The quotient \(H/B\)
is zero-dimensional and separated of finite type, hence
finite over the field. After pulling back along \(G\to G/H\),
the map \(G/B\to G/H\) has the form of the projection with
fibre \(H/B\). Thus \(G/B\to G/H\) is finite by faithfully flat
descent. By conjugacy this argument applies to every Borel.

We first give some geometric centralizer details. On a smooth affine
variety with a split torus action, at a torus-fixed point the
maximal ideal and its square are graded. Choose homogeneous lifts
of a basis of its cotangent space. In the completed local ring
these lifts are formal power series coordinates: their maps on
associated graded rings are isomorphisms, since the point is
smooth. The torus fixed locus is defined by the nonzero-weight
functions. In the completion every nonzero-weight monomial has
a variable of nonzero weight, so this ideal is generated by those
variables, by checking the coefficient in every homogeneous
degree. Thus the fixed locus is smooth and its tangent space
is the zero-weight subspace. The same argument applies on a
torus-stable affine chart at any fixed point when a smooth variety
is given an equivariant embedding in projective space: invert an
eigenfunctional nonzero on the fixed line.

For a cocharacter \(\lambda\), in a faithful \(\lambda\)-weight
basis the functors \(P_G(\lambda)\), \(U_G(\lambda)\), and
\(Z_G(\lambda)\) are given by block triangular matrices,
unipotent block triangular matrices, and block diagonal
matrices, respectively, intersected with \(G\). This follows by
multiplying a block by its power of the parameter and taking
the limit at zero. The limit defines a homomorphism
\(P_G(\lambda)\to Z_G(\lambda)\); it is the identity on
\(Z_G(\lambda)\) and its kernel is \(U_G(\lambda)\).
It lands in \(G\) by schematic density of the punctured line.
As above, choose homogeneous formal coordinates at the identity
for the \(\mathbb G_m\)-action. The completed ideal for
\(P_G(\lambda)\) is obtained by killing the coordinates of
negative weight (weights here mean the exponents in
\(f(\lambda(t)g\lambda(t)^{-1})\)).
Indeed the limit exists exactly when the regular functions of
negative weight vanish. All negative-weight monomials contain a
negative-weight variable. The additional condition that the
limit be the identity kills the coordinates of weight zero;
after the two sets are killed, all remaining nonconstant monomials
have strictly positive weight. This proves smoothness at the
identity for the three groups, hence everywhere by translations
over the algebraic closure. Their tangent spaces are the
corresponding weight spaces in \(\Lie(G)\). Multiplication gives
an open immersion
\[
 U_G(-\lambda)\times Z_G(\lambda)\times U_G(\lambda)
 \longrightarrow G.
\]
Indeed combine the last two factors into \(P_G(\lambda)\) by
the split limit map. The product with \(U_G(-\lambda)\) is a
monomorphism on all algebras, because its intersection with
\(P_G(\lambda)\) is trivial by the block equations.
Its differential is the weight-space direct sum at the identity,
so the map is \'etale there. At every geometric point of the
source it is \'etale by equivariance under left
\(U_G(-\lambda)\) and right \(P_G(\lambda)\), which act
transitively on the factors. An \'etale monomorphism is an
open immersion. The source is therefore geometrically connected
when \(G\) is, since the image is open in an irreducible variety.
This proves geometric connectedness of the three factors.
Taking the centralizers of the successive one-dimensional factors
of a split torus proves smoothness and connectedness of \(Z_G(S)\)
for every geometric torus \(S\). The corresponding smoothness
and open-product assertions over a ground field follow by descent.
Compare \cite[Sections 13.a and 13.d]{milne}.

A smooth connected affine group with no nontrivial geometric
torus is unipotent. Over the algebraic closure, each Borel
of it is unipotent by Results~\ref{res:B12} and~\ref{res:B13}.
Apply Result~\ref{res:A15} to a Borel to get a line whose
scheme stabilizer is that Borel. As a unipotent group has no
characters, it fixes a spanning vector as well. The orbit
map to the affine space of the representation therefore factors
through the proper \(G/B\). A morphism from the proper connected
reduced \(G/B\) to an affine space is constant: a regular function
has closed proper image in the affine line, hence has finite
image, and connectedness makes it constant. The vector stabilizer
is thus \(G\), proving \(G=B\) and the assertion.

If \(T\) is maximal, \(C=Z_G(T)\) is smooth and connected, and
\(C/T\) is smooth and connected by Result~\ref{res:D01}.
If it contained a positive-dimensional torus, its preimage would
be a smooth connected extension of a torus by a torus, hence
solvable. The solvable decomposition by Results~\ref{res:B12}
and~\ref{res:B13} would give a larger torus in \(C\): the
unipotent radical maps trivially to the quotient torus and meets
the kernel trivially, so is trivial. This contradicts maximality.
Thus \(C/T\) is unipotent by the preceding argument, so \(C\)
is solvable and lies in one Borel containing \(T\). The normalizer
of \(T\) acts transitively on those Borels: conjugate two Borels,
then conjugate the image of \(T\) back to \(T\) inside the target
Borel. Consequently \(C\), which the normalizer preserves, lies
in every Borel containing \(T\).

For a split torus the group-like elements of \(R[T]\), for any
\(k\)-algebra \(R\), are sums of distinct characters with
pairwise orthogonal idempotent coefficients whose sum is one.
This follows by comparing coefficients in
\(\Delta f=f\otimes f\) and \(\epsilon f=1\).
In particular its automorphisms are given locally on open and
closed subsets of \(\operatorname{Spec} R\) by lattice
automorphisms. A faithful representation gives a finite set of
weights of \(T\) which generate its character lattice
(Result~\ref{res:B01}); a normalizer element permutes them,
since it gives an isomorphism of twisted and untwisted
representations. Its possible lattice automorphisms therefore
form a finite set, and over an algebraically closed field the
fibres over lattice automorphisms with a geometric point
are cosets of \(C\). It follows that the \(T\)-fixed subset
of \(G/B\) for \(B\supset T\) is finite on geometric points:
every such point \(gB\) satisfies \(g^{-1}Tg\subset B\),
and conjugacy within \(B\) gives a representative in \(N_G(T)\).
Each \(C\)-coset in that normalizer gives one such point
because \(C\subset B\). This action of \(N_G(T)\) on the fixed
points is transitive.

We next prove the precise intersection assertion needed for
reductivity of torus centralizers. Let \(J\) be the reduced
identity component of the intersection of the unipotent radicals
of the Borels containing \(T\). It is a smooth connected
unipotent subgroup. It suffices to show that \(J\) fixes \(G/B\),
since then it is contained in every Borel. Its conjugates all
have that property. The reduced identity component of the
intersection of all Borels is a smooth connected normal
solvable group, equal to the solvable radical: every smooth
connected normal solvable group is in a Borel by maximal
dimension and then in all by conjugacy. The unipotent radical
of this solvable radical is characteristic, so is normal
in \(G\), and contains \(J\) by Result~\ref{res:B12}. This
places \(J\) in \(R_u(G)\).

To prove the fixed-point assertion, embed \(X=G/B\) in
\(\mathbb P(V)\) by the line of Result~\ref{res:A15},
the orbit construction of Result~\ref{res:D01}, and properness.
Replace \(V\) by the span of the orbit lines. If \(T=1\) then
\(G\) is unipotent by the rank-zero argument and the intersection
assertion is immediate. Otherwise choose a cocharacter separating
the finitely many distinct weights on \(V\). Let \(V_{\min}\)
be the minimum-weight space. On the open of \(X\) where the
minimum component is nonzero, taking the limit along this
cocharacter is the map given by projection to
\(\mathbb P(V_{\min})\). This open is nonempty because the orbit
lines span \(V\). Its image is contained in the finite set
\(X^T\), so irreducibility makes it a single point \([v]\).
Moreover the projections of the orbit vectors span \(V_{\min}\);
it follows that \(V_{\min}=kv\).
Thus the open is an affine chart of \(X\). It is precisely
the set of \(x\) such that \([v]\in\overline{T x}\):
a nonzero minimum component gives the limit, and a zero
component puts the entire orbit in the invariant hyperplane.

Let \(\ell\) be the linear functional on the minimum-weight
line extended by zero on the others, with \(\ell(v)=1\).
The orbit of \([\ell]\) in \(\mathbb P(V^\vee)\) is closed.
In fact every nonzero functional is nonzero on some translate
of \(v\) because those vectors span \(V\); so its orbit
meets the chart of functionals nonzero on \(v\).
The inverse cocharacter contracts that chart to \([\ell]\).
Hence every orbit closure contains \([\ell]\), and a closed
orbit of minimal dimension in the closure of \(G[\ell]\)
must be \(G[\ell]\), proving the claim. The stabilizer
of \([\ell]\) contains \(T\) and a Borel, the latter since
the closed projective orbit has a Borel-fixed point by
Result~\ref{res:B14}. The reduced identity component of the
stabilizer contains both the Borel and \(T\). Conjugacy of
maximal tori within this connected group puts the Borel in it
in a position containing \(T\). It follows that \(J\) preserves
\([\ell]\) and therefore the affine chart cut out by \(\ell\ne0\).
The normalizer translates of the argument give a \(J\)-stable
affine chart at every \(T\)-fixed point, since \(J\) is normalized
by \(N_G(T)\). They cover \(X\): every \(T\)-orbit closure has
a \(T\)-fixed point by Result~\ref{res:B14}, and membership in
one of the charts is precisely the orbit-closure condition
proved above.

A smooth unipotent group has closed orbits on an affine variety.
To verify this, if an orbit is not closed, let \(A\) be its
closure with reduced structure. The boundary has a nonzero
vanishing ideal in \(k[A]\), stable under the unipotent group.
A nonzero function of the ideal lies in a finite-dimensional
stable subspace of the ideal by coalgebra local finiteness
as in Result~\ref{res:A06}. The invariant-vector property of
Result~\ref{res:B06} gives a nonzero invariant in that subspace.
It is constant on the dense orbit, hence constant on \(A\),
and vanishing on the nonempty boundary forces it to be zero.
This contradiction proves closedness. To apply it, the closure
of any \(J\)-orbit in \(X\) contains a \(J\)-fixed point by
Result~\ref{res:B14}. Choose a \(J\)-stable affine chart just
constructed containing that point. Its closed invariant
complement cannot contain any point of the orbit, for it would
then contain the whole orbit closure. Thus the orbit is contained
in the chart, where it is closed and contains the fixed point.
It is a singleton, proving the required fixed-point assertion
and \(J\subset R_u(G)\). The reverse inclusion follows since a
normal smooth connected unipotent group lies in every Borel
and maps trivially to each torus quotient by
Result~\ref{res:B07}. This is the unipotent-intersection argument
of \cite[Sections 17.g--17.h]{milne}.

For a subtorus \(S\) of \(T\), put \(C_S=Z_G(S)\).
For \(B\supset T\), \(C_S\cap B=Z_B(S)\) is smooth connected
solvable by the centralizer argument. Its quotient orbit
in \(G/B\) is open in the \(S\)-fixed locus. Indeed the
differential of \(G\to G/B\) at any point above an \(S\)-fixed
coset identifies, after translation, the quotient of Lie
algebras by the conjugate of \(\Lie(B)\). Taking \(S\)-invariants
is exact by the weight decomposition. Therefore the orbit map
from the smooth \(C_S\) has surjective differential to the
fixed-locus tangent space at that point. The fixed locus is
smooth by the chart argument above, so the map is smooth near
that point and its image contains an open neighbourhood.
The same argument at every fixed point shows that the
\(C_S\)-orbits are open, and each is also closed since its
complement is a union of open orbits. The orbit of \(eB\)
is thus proper. By Result~\ref{res:D01} this orbit is
\(C_S/(C_S\cap B)\): the orbit map is smooth surjective onto
it and the scheme stabilizer is \(C_S\cap B\).
A Borel of \(C_S\) fixes a point of this proper orbit,
hence is contained in a conjugate of \(C_S\cap B\),
which is smooth connected solvable. Maximality proves
that \(C_S\cap B\) is a Borel of \(C_S\).
The radical \(R_u(C_S)\) consequently lies in every
\(B\supset T\), and maps trivially to their torus quotients.
The intersection assertion gives \(R_u(C_S)\subset R_u(G)\).
If \(G\) is reductive this proves reductivity of \(C_S\).
For \(S=T\), the already solvable \(Z_G(T)\) is then a
torus by Result~\ref{res:B12}, and equals \(T\) by maximality.
The equality is scheme-theoretic and descends over the
ground field.

Now \(N_G(T)\) acts on the character lattice with kernel \(T\).
The finite constant weight-permutation target over a
splitting field contains the image as a union of components:
each nonempty geometric fibre is a translate of \(T\),
and the other fibres are empty. Thus the image is a
finite constant group over that field and the map onto it
is fppf, by faithful flatness and finite presentation on the
translated torus fibres. The finite group descends to a finite
\'etale group over \(k\) by Result~\ref{res:F02}, and represents
the quotient sheaf \(N_G(T)/T\). If \(T\) is split, the same
finite constant target and the same nonempty components
are already defined over \(k\), so the quotient is constant.
A fibre is a torsor of a split torus and has a \(k\)-point by
Hilbert 90 (Result~\ref{res:F04} applied to each factor).

Lastly suppose \(T,T'\) are split geometric maximal tori
over \(k\). A geometric conjugacy exists, and its
transporter scheme is geometrically a coset of \(N_G(T)\),
equivalently a finite union of cosets of the smooth \(T\).
It is therefore smooth and has a point \(g\) over a separable
closure by Result~\ref{res:F04}. The induced isomorphism
\(T\to T'\) is a map of split character lattices and is defined
over \(k\). Thus \(g^{-1}\sigma(g)\) centralizes \(T\)
for every Galois element \(\sigma\), so belongs to \(T(k_s)\).
Hilbert 90 writes this cocycle as \(a^{-1}\sigma(a)\)
with \(a\in T(k_s)\). The element \(ga^{-1}\) is rational
and conjugates \(T\) onto \(T'\).

Over finite fields the existence of a rational geometric
maximal torus follows from Result~\ref{res:F09}. For infinite
fields we give the existence argument, including the case in
which a nontrivial torus character has zero differential.

We use two matrix facts about \(\mathfrak g=\Lie(G)\).
For \(X\in\mathfrak g\), its additive Jordan components over
\(\bar k\) belong to \(\mathfrak g_{\bar k}\). In characteristic
zero use the invariant derivation \(D_X\) on \(k[G]\) obtained
by differentiating right translations. Its additive Jordan
components are compatible on the finite-dimensional
subrepresentations in Result~\ref{res:A06}. On a tensor
product the derivation is the sum of the two operators, and
the sums of the respective semisimple and nilpotent parts
give its Jordan decomposition. Multiplication intertwines
these operators, so each component is a derivation. Its
evaluation at the identity is a tangent vector. The matrix
coefficient argument of Result~\ref{res:A14}, now for
derivations, identifies its action in any representation
with the corresponding additive Jordan component.
This construction is rational when \(k\) is perfect, since
Jordan components of a matrix then commute with Galois
transport. In characteristic \(p>0\), \(D_X^p\) is again
an invariant derivation by the Leibniz rule, since the
intermediate binomial coefficients vanish; its matrix in
a representation is \(X^p\), and this operation is
preserved by differentials of group homomorphisms.
For a faithful matrix representation of dimension \(d\),
\(X^{p^r}\) is semisimple over \(\bar k\) when \(p^r\ge d\).
Over \(\bar k\), the span of the \(p\)-powers of the
semisimple part is a space of diagonal matrices on which
the \(p\)-operation is injective and thus bijective
semilinearly. The powers of order at least \(p^r\) therefore
span the semisimple part as well. This shows it belongs to
\(\mathfrak g_{\bar k}\), and the difference gives the
nilpotent part.

The centralizer \(Z_G(X)\) of a matrix-semisimple
\(X\in\mathfrak g(k)\) is smooth with tangent algebra
\(\ker(\operatorname{ad} X)\).
To verify smoothness geometrically, diagonalize \(X\) in a
faithful representation. The derivation from conjugation
by \(1+\epsilon X\) acts diagonally on the ambient matrix
coordinate functions, with eigenvalues the differences of
the diagonal entries of \(X\); it is zero on the determinant
and its inverse. It preserves the ideal of \(G\), since
\(1+\epsilon X\) is an infinitesimal \(G\)-point. The ideal
is therefore a direct sum of eigenspaces: any individual
polynomial has finitely many eigencomponents, which can be
recovered by polynomials in the derivation. The ideal of the
centralizer in \(k[G_{\bar k}]\) is generated by the nonzero
eigenspaces. In the ambient group these generate exactly the
ideal of the off-block entries between distinct eigenspaces
of \(X\), and the assertion descends through the homogeneous
quotient. At the identity lift a cotangent basis by
eigenfunctions. Smoothness of \(G\) makes these lifts formal
power series coordinates. The centralizer ideal in this
completion kills precisely the coordinates of nonzero
eigenvalue: an eigenmonomial of nonzero eigenvalue has such
a factor, since the eigenvalue of a product is the sum.
The completed local ring of the centralizer is thus a power
series ring on the remaining coordinates. Translation gives
smoothness at every geometric point. The tangent formula
follows by commuting \(X\) with \(1+\epsilon Y\).
In particular \([X,\mathfrak g]=0\) implies \(Z_G(X)=G\)
for connected \(G\), by dimension.

We first show that a smooth connected affine group over an
infinite field with a positive-dimensional geometric torus
has a nontrivial \(k\)-torus. Induct on its dimension.
If a geometric semisimple Lie element acts nontrivially
by the adjoint Lie operation, the characteristic polynomial
of \(\operatorname{ad} X\) on \(\mathfrak g\) differs from
a power of the variable on a nonempty open subset of the
affine Lie algebra. In fact the Jordan splitting above
gives the commuting semisimple operator
\(\operatorname{ad}(X_s)\) and nilpotent operator
\(\operatorname{ad}(X_n)\): these properties hold in the
matrix algebra and hence on the invariant subspace
\(\mathfrak g_{\bar k}\). The open condition is polynomial
over \(k\), and has a \(k\)-point because \(k\) is infinite.
In characteristic zero take its rational semisimple part.
In positive characteristic take \(X^{p^r}\) with \(p^r\ge d\);
the nonzero eigenvalues of \(\operatorname{ad}(X_s)\) remain
nonzero after this power, as
\(\operatorname{ad}(X^{p^r})=(\operatorname{ad} X)^{p^r}\).
This gives a nonzero semisimple rational vector \(Y\)
with proper smooth centralizer, whose smooth identity
component has smaller dimension. The component's Lie
algebra contains \(Y\); a smooth unipotent group has only
nilpotent elements in the Lie algebra in every representation,
by the upper unitriangular flag of Result~\ref{res:B06}.
Thus the centralizer component is geometrically nonunipotent,
and has a nontrivial geometric torus by the rank-zero
argument above. Induction gives a nontrivial \(k\)-torus.

It remains to reduce to that case or to the case of a
geometric central maximal torus, when all geometric
semisimple Lie elements are central. Let \(T\) be a
geometric maximal torus. If \(T\) is central, then \(G_{\bar k}\)
is solvable by the above calculation of \(Z_G(T)/T\).
Over \(\bar k\) write \(G=T\times U\) by
Results~\ref{res:B12} and~\ref{res:B13}; centrality makes the
semidirect product direct. In characteristic zero \(T\)
is unique by geometric conjugacy and descends by
Result~\ref{res:F02}. In positive characteristic choose
\(p^r\) large enough to kill every element of the upper
unitriangular realization of \(U\). The schematic image
of the morphism \(g\mapsto g^{p^r}\) then becomes exactly
\(T\) over \(\bar k\), since the power map on a torus is
faithfully flat. The coordinate ideal of the image is the
kernel of a ring map and commutes with extension of the
field. It therefore defines the required \(k\)-torus,
without assuming \(k\) perfect.

If \(T\) is not central, it acts by some nonzero weight on
\(\mathfrak g_{\bar k}\), because \(Z_G(T)\) is smooth
with the invariant tangent space. All these weights have
zero differential under the assumption on semisimple
elements. This is impossible in characteristic zero;
in characteristic \(p\) it means that the weights are
divisible by \(p\) in \(X^*(T)\).
Moreover the geometric semisimple elements are exactly
\(\mathfrak t=\Lie(T)\). Each such element is centralized
by \(G_{\bar k}\) by smoothness of its centralizer, so it
lies in \(\Lie(Z_G(T))\). Here \(Z_G(T)=T\times U_0\)
by its solvability, Results~\ref{res:B12}, \ref{res:B13},
and centrality of \(T\) in it. The tangent vectors of \(U_0\)
are nilpotent, commute with \(\mathfrak t\), and have zero
intersection with it. This proves the assertion by uniqueness
of the matrix Jordan decomposition.

The space \(\mathfrak t\) descends as a vector subspace:
it is the geometric span of all values of the polynomial
map \(X\mapsto X^{p^r}\) for \(p^r\ge d\).
Indeed every value is semisimple and on the diagonalized
\(\mathfrak t\) these powers are a semilinear bijection
over \(\bar k\). Any linear functional over \(\bar k\)
vanishing on the values for \(X\in\mathfrak g(k)\)
vanishes on every geometric value by polynomial density
of \(k\)-points of that affine space. Thus the span of the
rational values gives the descent.

We also descend the first Frobenius kernel \(E=T[F]\).
In the faithful representation, the character blocks of \(E\)
are exactly the common eigenblocks of \(\mathfrak t\):
for \(\mu_p^r\) its characters are determined by their
differentials, which are linear combinations over
\(\mathbf F_p\) in a diagonal basis. The commutative
\(k\)-subalgebra of matrices generated by the descended
\(\mathfrak t\) becomes the algebra of scalar operators
on these eigenblocks over \(\bar k\). Indeed the diagonal
eigenvalues of its elements separate the blocks, and
polynomial interpolation gives the individual projections.
The algebra is finite \'etale, since it is geometrically
a product of fields, so the blocks are defined over a
finite separable splitting field. In the split torus of
scalar operators on those blocks, \(E\) is determined by
its tangent space. Explicitly it is inside the product
of the \(\mu_p\)'s, and the kernel on characters modulo \(p\)
consists of the characters whose differentials vanish
on \(\mathfrak t\). These character equations are defined
over the splitting field and are Galois-stable by this
characterization. Thus \(E\) descends by Result~\ref{res:F02},
including its inclusion into \(G\).
It is central: conjugation over \(\bar k\) preserves the
same \(\mathfrak t\), hence the same \(E\), and it acts
trivially on its tangent space, so trivially on its
characters and on \(E\). This initially proves triviality
for conjugation by geometric points of \(G\), as morphisms
of \(E\). It proves scheme centrality as well because the
coefficients in a basis of the finite algebra \(k[E]\)
are regular functions on geometrically reduced \(G\),
so equality on geometric points forces their equality.

Pass to the affine quotient \(q:G\to G_1=G/E\)
from Result~\ref{res:D01}. It is finite by local
trivialization and descent, since \(E\) is finite.
The quotient is geometrically reduced, as its coordinate
algebra embeds in the geometrically reduced \(k[G]\);
it is then smooth by Result~\ref{res:A09}, and connected
by surjectivity. Every torus in \(G_1\) is the image of
a torus in \(G\), over its own field of definition:
if the torus is \(S_1\), put \(P=q^{-1}(S_1)\).
The commutator of \(P\) factors through
\(S_1\times S_1\to E\), since \(E\) is central. The
target is finite infinitesimal and the source
geometrically reduced, so the commutator is trivial.
Then the power map \([p]:P\to P\) is a homomorphism
killing \(E\), and factors through \(S_1\).
Its scheme image is a torus by the character description
of quotients of tori, and projects onto \(S_1\) by
\([p]:S_1\to S_1\). This proves the lifting assertion.
In particular \(T_1=q(T)\) is geometrically maximal.

The differential of \(q\) has kernel \(\Lie(E)=\mathfrak t\)
by Result~\ref{res:A12}. Its image consists of nilpotent
matrices in any faithful representation of \(G_1\), because
the semisimple part of each input is killed and the
\(p\)-operation is respected. Hence its intersection with
\(\Lie(T_1)\), consisting of semisimple matrices, is zero,
and dimensions give a direct sum decomposition of
\(\Lie(G_1)\) into these subspaces. The nonzero weights of
\(T_1\) on that algebra pull back exactly to the nonzero
weights of \(T\), by the equivariant identification of
\(\mathfrak g_{\bar k}/\mathfrak t\) with the differential
image. The injection on character lattices identifies
\(X^*(T_1)\) with \(pX^*(T)\). Thus for each nonzero
weight the exponent of its divisibility by \(p\) decreases
by one. Repeat these central quotients if necessary.
Since a nonzero vector in a finite-rank free abelian
group is divisible by only finitely many powers of \(p\),
the process gives a noncentral semisimple vector
after finitely many steps. All quotients have the
original dimension, and the smaller-dimensional
centralizer argument already proved gives a nontrivial
rational torus in the last quotient. Lift it through
the successive quotients as above. This completes the
dimension induction.

Finally choose a \(k\)-torus maximal by dimension.
A larger geometric torus containing it would give a
nontrivial geometric torus in \(Z_G(T)/T\).
This smooth connected affine quotient would then have
a nontrivial \(k\)-torus by the preceding argument.
Its inverse image in \(Z_G(T)\) is a smooth connected
extension of tori and is geometrically a torus:
it is solvable, and its geometric unipotent radical maps
trivially to the quotient and meets the kernel trivially
by Result~\ref{res:B07}. It would be a larger \(k\)-torus,
contradicting maximality. This proves geometric maximality.
Compare \cite[Section 13.3]{springer}
and \cite[Section 17.j]{milne}.
\end{proof}

The two kinds of maximality in this discussion are distinct. In \(\SL_2/\RR\) the diagonal torus is split, whereas the matrices
\[
 \begin{pmatrix}a&-b\\ b&a\end{pmatrix},\qquad a^2+b^2=1,
\]
form a nonsplit torus. Its complex eigencharacters are interchanged by conjugation. Both are maximal tori of dimension one, and rational conjugation cannot change whether a torus is split.

\begin{result}{C03}{The central torus and the derived group}
For a connected reductive \(G\), its scheme-theoretic derived group \(G^{\mathrm{der}}\) is connected semisimple and \(G/G^{\mathrm{der}}\) is a torus. Let \(Z_{\mathrm{tor}}\) be the maximal torus of the scheme centre. It commutes with field extension, and multiplication is a central isogeny
\[
 Z_{\mathrm{tor}}\times G^{\mathrm{der}}\longrightarrow G.
\]
Its kernel is the copy of \(Z_{\mathrm{tor}}\cap G^{\mathrm{der}}\) given by \(z\mapsto (z,z^{-1})\). A reductive group is semisimple exactly when its scheme centre is finite.
\end{result}

\begin{proof}
First construct the derived subgroup over an algebraically
closed field. For \(n\ge1\), take the reduced closure of the image
of the map \(G^{2n}\to G\) that multiplies \(n\) commutators.
These closures are nested irreducible subsets containing the identity,
and eventually two consecutive ones agree by the dimension bound.
If \(X_n=X_{n+1}\), multiplying its dense generating set by
one more commutator, then taking closures, shows that it contains
all subsequent products. The resulting \(X_n\) is stable under
multiplication, inverse, and conjugation, as can be checked on
the dense commutator products first. The source schemes for these
operations are reduced, so the operations restrict to the reduced
closure. This is a smooth connected normal subgroup by
Result~\ref{res:A09}. Its ideal is the kernel on coordinate rings
of the indicated product map. A field extension preserves that
kernel by flatness, hence the construction and its universal
property descend over \(k\): a normal closed subgroup gives a
commutative quotient exactly if it contains the commutator morphism,
and in that event it contains the closure of every finite product
of commutators. Denote the subgroup so constructed by \(G^{\mathrm{der}}\).

We check the remaining assertions geometrically first. The centre
is contained in a maximal torus by the centralizer equality of
Result~\ref{res:C02}. Its largest subtorus \(Z_0\) is characterized
on the character group as the quotient by its torsion subgroup:
if a diagonalizable group is \(D(M)\), maps to it from tori
factor through \(D(M/M_{\mathrm{tors}})\)
(Result~\ref{res:B01}). Over a separable field splitting a maximal
\(k\)-torus this is a characteristic construction on the centre.
The centre commutes with field extension: its equations
are the coefficients in the other variable of \(zg=gz\),
tested in the coordinate algebra of \(G\) as in
Result~\ref{res:A08}. The character construction consequently
descends \(Z_0\) to \(Z_{\mathrm{tor}}\) and commutes with every
field extension. This includes the case in which the centre has
infinitesimal multiplicative factors.

For the finite-intersection assertion, in a faithful representation
each \(Z_0\)-weight space \(V_\chi\) is \(G\)-stable because
\(Z_0\) is central. The determinant character on \(V_\chi\)
is trivial on \(G^{\mathrm{der}}\), so its restriction to
\(Z_0\cap G^{\mathrm{der}}\) is \(\chi^{\dim V_\chi}=1\).
The weights \(\chi\) generate \(X^*(Z_0)\) because the
representation restricted to \(Z_0\) is a closed immersion
and characters describe its kernel (Result~\ref{res:B01}).
The subgroup of the character lattice generated by
\((\dim V_\chi)\chi\) has finite index. The intersection,
which lies in the group defined by these determinant equations,
is therefore finite as a scheme.

For surjectivity, fix a geometric maximal torus \(T\).
Multiplication of \(T\) with all the root groups, in any
fixed order, has bijective differential at the identity by
Result~\ref{res:C04}. The source and target are smooth, so the
morphism is \'etale at that point and its image contains an open
subset of \(G\). In particular these factors generate \(G\)
as a closed group. Each root group belongs to \(G^{\mathrm{der}}\):
for \(t\in T(\bar k)\) with \(\alpha(t)\ne1\) the morphism
\[
 [t,x_\alpha(a)]=x_\alpha\big((\alpha(t)-1)a\big)
\]
covers \(U_\alpha\). Likewise the image of \(\alpha^\vee\)
belongs to \(G^{\mathrm{der}}\): its construction in
Result~\ref{res:C04} presents it as the product of the
rank-one root words \(w(a)w(1)^{-1}\).
The cocharacter space of \(Z_0\)
is the common kernel of the roots inside \(Y_\QQ=X_*(T)\otimes\QQ\).
Indeed a subtorus in that kernel centralizes \(T\) and all
\(U_\alpha\), and hence centralizes the schematically dense image
of their product, also after base change by a test algebra.
Conversely a central subtorus has trivial root characters on it.

The finite Weyl group of Result~\ref{res:C02} acts on \(Y_\QQ\).
Average a rational positive definite inner product over it.
The reflection for \(\alpha\) is orthogonal, with fixed
hyperplane \(\ker\alpha\) and negative eigenspace
\(\QQ\alpha^\vee\). Hence the span of the coroots is the
orthogonal complement of the common kernel.
Thus \(Z_0\) together with the images of the coroots generates
\(T\): the map of split tori from their product has full rank
on cocharacters and is faithfully flat by the character
description in Result~\ref{res:B01}. The scheme image of
\(Z_0\times G^{\mathrm{der}}\) is therefore all of \(G\),
since it contains \(T\) and every root group.
Result~\ref{res:A07} makes multiplication faithfully flat.
On every test algebra its kernel consists precisely of
\((z,z^{-1})\) for \(z\in Z_0\cap G^{\mathrm{der}}\).
These conclusions descend by faithful flatness, and the quotient
\(G/G^{\mathrm{der}}\) is the torus
\(Z_{\mathrm{tor}}/(Z_{\mathrm{tor}}\cap G^{\mathrm{der}})\),
whose character lattice is the finite-index sublattice of
characters trivial on the intersection.

Finally the geometric unipotent radical of \(G^{\mathrm{der}}\)
is normalized by \(G_{\bar k}\), by uniqueness in
Result~\ref{res:C01}. Normality first holds on geometric points
and then for the conjugation morphism from a reduced source.
Reductivity of \(G\) forces that radical to vanish. The smooth
connected solvable radical of \(G^{\mathrm{der}}\) is consequently
a torus by Result~\ref{res:B12}, and it too is normal in \(G_{\bar k}\).
A normal torus in a connected smooth group is central: the action
by conjugation is locally constant on its character lattice by
the group-like coefficient calculation in the proof of
Result~\ref{res:C02}, and equals the identity at the unit.
This torus therefore lies in the finite intersection just computed
and is trivial. The same argument for a reductive \(G\) itself
shows that its geometric solvable radical is exactly \(Z_0\).
Since a group of multiplicative type with zero-dimensional
maximal torus is finite, this proves the centre criterion.
For the central-isogeny description compare
\cite[Theorem 3.2.2]{conradfeng} and \cite[Chapter 19]{milne}.
\end{proof}

For \(\GL_n\) this map is \((z,A)\mapsto zA\), with \(A\in\SL_n\) and \(z\) scalar. A matrix \(M\) lifts after adjoining an \(n\)-th root of \(\det(M)\), a finite faithfully flat base extension when working over a test algebra, and then \(A=z^{-1}M\). Its kernel is \(\mu_n\). In particular, the kernel can be non-\'etale in positive characteristic.

\section{Absolute roots and coroots}

\begin{result}{C04}{Root groups from rank one}
Let \((G,T)\) be a split connected reductive pair. Put \(X=X^*(T)\), \(Y=X_*(T)\), and define \(\langle\chi,\lambda\rangle\in\ZZ\) by \(\chi(\lambda(t))=t^{\langle\chi,\lambda\rangle}\). The absolute roots \(\Phi\subset X\) are the nonzero weights of \(T\) on \(\Lie(G)\). The zero-weight space is \(\Lie(T)\), and each root space is one-dimensional. For \(\alpha\in\Phi\) there is a unique root subgroup \(U_\alpha\simeq\Ga\) normalized by \(T\) with that weight, a coroot \(\alpha^\vee\in Y\), and a representative \(n_\alpha\in N_G(T)(k)\) acting on \(X\) by
\[
 s_\alpha(x)=x-\langle x,\alpha^\vee\rangle\alpha,\qquad
 \langle\alpha,\alpha^\vee\rangle=2.
\]
In particular \(\dim G=\dim T+|\Phi|\).
\end{result}

\begin{proof}
The zero-weight space is the tangent space of \(Z_G(T)\),
by the torus fixed-locus calculation in Result~\ref{res:C02},
so it is \(\Lie(T)\). Fix \(\alpha\ne0\) and let \(S\) be the
largest subtorus of \(\ker\alpha\). The subgroup \(H=Z_G(S)\)
is smooth connected reductive by the geometric centralizer
argument in that result. It has the same maximal torus \(T\).
The weights in its Lie algebra are precisely those which
vanish on \(S\), hence those on the rational line through
\(\alpha\) together with the zero weights. Choose
\(\lambda:\Gm\to T\) with \(\langle\alpha,\lambda\rangle>0\).
Write \(U_\pm=U_H(\pm\lambda)\) for the two limit subgroups.
The dynamic calculation of Result~\ref{res:C02} identifies
\(P_H(\pm\lambda)=T\ltimes U_\pm\): in fact \(\lambda\)
and \(S\) generate \(T\), and \(Z_H(\lambda)=Z_H(T)=T\).
The subgroup \(P_H(\lambda)\) is proper, since otherwise
\(H\) would be solvable and hence a torus by reductivity
and Result~\ref{res:B12}, contrary to the \(\alpha\)-weight.
The opposite subgroup is proper as well by the same
argument. Their product contains an open subset by the
dynamic calculation.

Over \(\bar k\), the two \(P_H(\pm\lambda)\) are contained
in distinct Borels containing \(T\): a common Borel would
contain that open product and equal \(H\). The normalizer
action on the lattice of cocharacters of \(T\) fixes
\(X_*(S)\) pointwise because \(S\) is central in \(H\).
It acts on the one-dimensional rational quotient with
finite image, hence through \(\{1,-1\}\).
Its kernel on this quotient fixes the subspace and is a
shear with square-zero off-diagonal part; a finite-order
shear over \(\QQ\) is the identity. So \(W(H,T)\) has at
most two members. Its transitivity on Borels containing
\(T\), proved in Result~\ref{res:C02}, gives exactly two
Borels \(B^+,B^-\) containing \(T\), with the respective
limit subgroups. For \(X=H/B^+\), the fixed points of \(T\)
are indexed by \(N_H(T)(\bar k)/T(\bar k)\).
Indeed every fixed coset has a normalizer representative
by conjugacy of maximal tori in \(B^+\);
\(N_H(T)(\bar k)\cap B^+(\bar k)=T(\bar k)\) because
conjugation in \(B^+\) acts trivially on the torus quotient,
which restricts isomorphically to \(T\).
Thus \(X\) has exactly two geometric fixed points.

The action of \(T\) on this proper homogeneous space factors
through \(T/S\), since \(S\) is central and contained in \(B^+\).
To see that \(\dim X=1\), use its equivariant projective
embedding as the orbit of a line, as in Result~\ref{res:C02}.
The limiting points for \(\lambda\) and \(-\lambda\) of a
nonfixed point are distinct: in projective coordinates they
are the lowest- and highest-weight nonzero components, whose
weights differ precisely when the line is not fixed.
Here being \(\lambda\)-fixed is equivalent to being \(T\)-fixed,
since \(\lambda\) surjects onto \(T/S\). By the extremal-weight
chart argument of Result~\ref{res:C02}, there are affine charts
defined by the nonvanishing of the lowest and highest
components, and their limit points are the two fixed points. Here
all orbit lines lie in a single \(S\)-weight space: the possible
spaces have disjoint projectivizations and the orbit is irreducible.
Since the orbit lines span the chosen representation, distinct
\(T\)-weights therefore have distinct \(\lambda\)-weights.
These charts are defined by linear functionals because each of
the two extremal weight spaces on the span of the orbit is a line.
For any nonfixed point the two limits must be these extremal
points in weight order, hence it belongs to both charts. The
complement of either chart is thus the other fixed point.
A nonempty hyperplane section of the integral projective
\(X\) has pure codimension one, by the principal ideal
dimension theorem in its affine charts. This proves the
dimension assertion. The same argument with \(B^-\) gives
\(\dim(H/B^-)=1\).

The curve \(X\) is smooth and proper by Result~\ref{res:D01}.
A nonfixed point has an open orbit of the one-dimensional
torus \(T/S\), with finite stabilizer. That orbit is a
one-dimensional torus by the character quotient of
Result~\ref{res:B01}. Hence \(X\) has rational function field
over \(\bar k\). The smooth proper curve with this function
field is \(\mathbb P^1\): the isomorphism on dense opens extends
both ways using the valuation rings at the missing points and
properness, and the extensions are inverse by separatedness.

The projective-line automorphism functor is
\(\operatorname{PGL}_2\). One way to see the scheme
assertion is to use the images of \(0,\infty,1\): an ordered
triple of pairwise disjoint sections determines a projective
linear transformation fppf locally by choosing vectors for the
first two sections and scaling them so that their sum is on the
third line. The scalar ambiguity descends to
\(\operatorname{PGL}_2\). An automorphism fixing the three
sections preserves both standard affine charts. On the chart
with coordinate \(z\) it sends \(z\) to \(z\,u(z)\) with
\(u\) a unit, since it fixes the zero divisor as a closed
subscheme. On the opposite chart the formula is
\(z^{-1}v(z^{-1})\) with \(v\) a unit. Agreement on the
overlap gives \(u(z)^{-1}=v(z^{-1})\), so both are constant
by the intersection of the two polynomial rings in the
Laurent ring over any test algebra. Fixing \(1\) makes
the constant one. This proves the assertion on all algebras.

Now use the action of \(H\) on \(X_{\bar k}\) to map
\(H_{\bar k}\) into \(\operatorname{PGL}_2\).
Its kernel \(K\) is contained in all Borels of \(H\).
Therefore \(K_{\mathrm{red}}^\circ\) is contained in the
solvable radical, which is a central torus, by
Results~\ref{res:C01}, \ref{res:B12}, and the torus
rigidity in Result~\ref{res:C02}. In particular the smooth
unipotent groups \(U_+\) and \(U_-\) have zero-dimensional
kernel in this action. Their images in
\(\operatorname{PGL}_2\) fix the appropriate distinct fixed
points of \(B^+\) and \(B^-\), respectively:
a representative of the nonidentity Weyl element conjugates
\(B^+\) onto \(B^-\), since the two Borels are distinct
and the normalizer acts transitively.
An image of a unipotent group in a point stabilizer lies in its
one-dimensional unipotent radical, by its torus quotient.
Thus both \(U_\pm\) have dimension at most one. Both are nontrivial
by properness of both dynamic subgroups and the tangent weight
sum, so they have dimension one, and dimension equality
makes \(B^\pm=P_H(\pm\lambda)\) over the algebraic closure.
These Borels now have the indicated \(k\)-models.

Over \(k\) the open dynamic product \(U_-\times B^+\subset H\)
now gives an open chart \(U_-\subset X\) by its quotient.
Each \(U_\pm\) is a split additive group over \(k\):
its algebra is positively graded by the contracting
cocharacter, with degree-zero part \(k\) because the action
contracts to the identity. A homogeneous generator of the
one-dimensional cotangent space generates the algebra by
induction on positive degree. The algebra is a polynomial
algebra on that generator by smooth connected dimension one.
The generator is primitive, since a term with two positive
factors would have weight at least twice its degree. Thus
the chart is \(\mathbb A^1_k\) and the same rational-curve
argument gives \(X\simeq\mathbb P^1_k\).
The two fixed points are rational: one is \(eB^+\), and the
other is its translate by a rational representative of the
Weyl element of \(H\) from Result~\ref{res:C02}.

Choose projective coordinates in which \(B^+\) fixes
\(\infty\) and \(B^-\) fixes \(0\).
Under \(H\to\operatorname{PGL}_2\) the maps from \(U_\pm\)
to the upper and lower additive subgroups are isomorphisms.
They are surjective by the positive-dimensional image
calculation above. Their kernels are trivial as schemes:
the action kernel lies in \(B^+\cap B^-=T\) by the dynamic
block equations, and \(T\cap U_\pm=1\).
Use Result~\ref{res:A07} to conclude isomorphy.
Thus there are unique parametrizations \(x_+,x_-\)
mapping to the standard upper and lower transvections.
Since the two images have inverse diagonal characters, their
weights are \(\alpha,-\alpha\), with \(\alpha\) the positive
one by choice of \(\lambda\). This also proves that only
these two nonzero weights appear on the line through \(\alpha\),
each with multiplicity one. Any normalized subgroup isomorphic
to \(\Ga\) with that tangent weight has the scalar action
by that character: automorphisms of \(\Ga\) on geometric
points have degree one by comparing polynomial degrees
with their inverses, and the smooth acting torus makes the
coefficient identity schematic. It then lies in \(Z_G(S)\)
and contracts along the appropriate signed cocharacter,
so equals the corresponding \(U_\pm\) by dimension.

Set \(w(a)=x_+(a)x_-(-a^{-1})x_+(a)\) for an invertible
parameter \(a\), \(n_\alpha=w(1)\), and
\(\alpha^\vee(a)=w(a)w(1)^{-1}\).
In \(\operatorname{PGL}_2\) the matrix of \(w(a)\) is
the class of
\(\begin{psmallmatrix}0&a\\-a^{-1}&0\end{psmallmatrix}\).
It exchanges \(0,\infty\), and
\(w(a)w(1)^{-1}\) is diagonal.
The inverse image of the diagonal torus in \(H\) is
\(B^+\cap B^-=T\), so \(n_\alpha\) normalizes
\(T\) and \(\alpha^\vee\) is a morphism into \(T\) taking
\(1\) to the identity. A pointed morphism
\(\Gm\to T\) for a split torus is a cocharacter, since
the units of the Laurent polynomial ring over \(k\) are
scalar monomials. The projected diagonal matrix gives
\(\alpha(\alpha^\vee(a))=a^2\).
For \(t\in T\) the two character actions give
\(t w(a)t^{-1}=w(\alpha(t)a)\).
Putting \(a=1\) and rearranging gives
\[
 n_\alpha t n_\alpha^{-1}
   =t\,\alpha^\vee(\alpha(t))^{-1}.
\]
This induces precisely \(s_\alpha\) on the character lattice.
The dimension assertion follows from smoothness and all the
weight-space dimensions. The argument uses the original
characters, so a nontrivial character is still detected when
its differential is zero in the ground characteristic.
Compare \cite[Chapters 20--21]{milne}
and \cite[Sections 2.1--2.2]{conradfeng}.
\end{proof}

\begin{result}{C05}{The reduced root datum and the scheme centre}
The tuple \((X,\Phi,Y,\Phi^\vee)\), with the perfect integer character pairing, is a reduced root datum. Its Weyl group generated by the \(s_\alpha\) is naturally \(W(G,T)\). The scheme centre is
\[
 Z(G)=\bigcap_{\alpha\in\Phi}\ker(\alpha:T\to\Gm)
      =D\bigl(X/\ZZ\Phi\bigr),
\]
where \(D(M)\) denotes the diagonalizable group attached to the finitely generated abelian group \(M\).
\end{result}

\begin{proof}
The character and cocharacter lattices of the split torus are
dual. The reflection in Result~\ref{res:C04} acts by a
normalizer element, hence preserves the weights. It preserves
the attachment of coroots: a normalizer conjugates the rank-one
centralizer and exchanges the corresponding dynamic root groups.
In that centralizer the nontrivial Weyl action is unique, and its
formula determines the coroot from the specified root.
This gives both reflection axioms. The rank-one calculation shows
that the roots are reduced. The coroots are reduced as well:
two proportional coroots give reflections with the same negative
line. Average a positive definite inner product on the rational
cocharacter space over the finite Weyl group. Their fixed hyperplanes
are then both the orthogonal complement of that line, so their
roots are proportional and hence equal up to sign. The reflection
formula now gives equal or opposite coroots. These are the
reduced root-datum axioms.

Here are the geometric Borel and chamber details. Work over the
algebraic closure first. If \(B\) contains \(T\) and a root line
belongs to \(\Lie(B)\), then \(B\) contains that root group:
in \(Z_B(S)\), for \(S\) the subtorus used for that root in
Result~\ref{res:C04}, the relevant contracting subgroup has
the same one-dimensional tangent space as \(U_\alpha\).
By the smoothness and connectedness of the limit construction
in Result~\ref{res:C02}, this subgroup equals \(U_\alpha\).
Both opposite root groups cannot belong to \(B\), because
with \(T\) they fill an open subset of \(Z_G(S)\) by
the rank-one dynamic product. That connected reductive
centralizer is not solvable because it has nonzero root weights.
It follows that
\(\dim(B)\le\dim T+|\Phi|/2\).
If \(\lambda\) is regular on \(\Phi\), the centralizer
\(Z_G(\lambda)\) is smooth connected with Lie algebra
\(\Lie(T)\), so equals \(T\).
The dynamic subgroup \(P_G(\lambda)=T\ltimes U_G(\lambda)\)
is smooth connected solvable of dimension
\(\dim T+|\Phi|/2\).
It is contained in a Borel, hence is itself a Borel by
dimension. Its nonzero weights are exactly the roots positive
on \(\lambda\). Every Borel containing \(T\) arises this way,
since the normalizer of \(T\) acts transitively on these Borels
by Result~\ref{res:C02}. Distinct positive systems have distinct
tangent weights. The construction is by \(k\)-subgroups when
\(T\) is split, since every nonempty open rational chamber
contains a lattice cocharacter. Thus the positive systems
also classify Borels containing \(T\) over \(k\).

To see that the normalizer action on these Borels has no
stabilizer modulo \(T\), prove that a Borel is self-normalizing.
Continue to work geometrically and induct on the dimension
of the reductive group, with dimension zero trivial.
Let \(N=N_G(B)\). Its Lie algebra is contained in
\(\Lie(B)\) because a first-order normalizer has class
fixed by \(T\) in \(\Lie(G)/\Lie(B)\), and that
quotient has no weight-zero space. Indeed conjugate a
torus point by the first-order element and reduce its
displacement modulo \(\Lie(B)\). Since \(N\) contains the
smooth \(B\), equality of local and tangent dimensions
makes \(N\) smooth at the identity, hence smooth everywhere.
It suffices to prove the point equality. Given \(g\in N(\bar k)\),
adjust it by an element of \(B(\bar k)\) so it normalizes \(T\),
using Result~\ref{res:B13}. Consider \(t\mapsto gtg^{-1}t^{-1}\)
on \(T\). If its image is a proper subtorus, its kernel has a
positive-dimensional subtorus \(S\). If \(Z_G(S)\ne G\),
apply induction to \(Z_G(S)\) and its Borel \(Z_G(S)\cap B\),
both justified in Result~\ref{res:C02}. The element \(g\)
normalizes that Borel, so belongs to \(B\). If \(Z_G(S)=G\),
use \(G/S\) instead. It is smooth connected reductive:
the inverse image of its geometric unipotent radical would be
a smooth connected normal solvable extension of that radical
by \(S\), hence a torus by Results~\ref{res:C01}
and~\ref{res:B12}. A torus has no nontrivial unipotent
quotient by Result~\ref{res:B07}. Also \(B/S\) is Borel,
since the inverse image of any smooth connected solvable
overgroup is again smooth connected solvable.
Induction in this quotient puts \(g\) in \(B\).

If the image of \(t\mapsto gtg^{-1}t^{-1}\) is all of \(T\),
choose a line whose scheme stabilizer is \(N\), by
Result~\ref{res:A15}. The character through which \(N\)
acts on it is trivial on all those commutators, hence on \(T\).
It is also trivial on \(R_u(B)\). Thus \(B\) fixes a
spanning vector. Its orbit map to the affine representation
space factors through the proper \(G/B\), and must be
constant, by the proper-to-affine argument in
Result~\ref{res:C02}. Therefore \(G=N\), and the Lie dimension
calculation gives \(G=B\). This proves \(N_G(B)=B\)
in all cases. Inside \(B\) normalizing \(T\) amounts
to centralizing \(T\), because the torus quotient of \(B\)
restricts isomorphically to \(T\).
It follows that the Weyl action on the set of Borels containing
\(T\), or on the chambers, is free as well as transitive.
The reflections in roots generate a transitive subgroup:
a path between points in two chambers can be perturbed to
meet the hyperplanes one at a time, avoiding their pairwise
intersections. Each crossing goes between two chambers
interchanged by the reflection in that hyperplane, since the
reflection preserves the whole arrangement. Thus their group
is the full Weyl group. The scheme identification follows
since the two Weyl groups are constant in the split case.
Compare \cite[Sections 17.e and 21.d]{milne}.

For clarity, the simple roots of a chamber are the positive
roots that cannot be written as sums of two positive roots.
Every positive root is a sum of these with nonnegative integer
coefficients, by induction on its positive integer pairing with
a regular integral cocharacter in the chamber.
Two distinct simple roots \(a,b\) have nonpositive inner
product for a Weyl-invariant positive definite inner product
on the character space. If the inner product were positive,
the two Cartan pairings would be strictly positive integers
with product \(4(a,b)^2/((a,a)(b,b))<4\). One of the integers
would then be \(1\), and reflection would give \(a-b\) or
\(b-a\) as a root. One of these two opposite roots is positive,
decomposing a simple root, a contradiction.
The simple roots are therefore linearly independent.
Indeed a nontrivial dependence with both signs would equate
positive linear combinations on disjoint sets of simples.
The squared norm of their common sum is their mutual inner
product, which is nonpositive, whereas the sum is nonzero
by pairing with the chamber cocharacter. A relation with
only one sign already contradicts that pairing.
Thus all roots have integer coordinates in the simple basis,
with coefficients of one sign. In particular the intersection
of \(\Phi\) with the rational span of any simple subset is
its intersection with the integral span. The chamber is cut
out by the positive inequalities on the simple basis.
Independence makes each of these inequalities a wall, so
the simple roots are precisely the inward root normals to
the walls used in the gallery argument.

For use below, the positive root subgroups multiply
isomorphically, in any chosen order, onto \(U_G(\lambda)\).
The product map is equivariant for the contracting
\(\lambda\)-action and its differential at the identity is
the direct sum of the positive root lines
(Results~\ref{res:C02} and~\ref{res:C04}).
The coordinate algebras of its source and target are graded
in positive degrees away from the constants, because the
action extends to contraction to the identity at zero.
The map on cotangent spaces is an isomorphism.
Thus its coordinate homomorphism is surjective: choose
homogeneous preimages of the cotangent generators, and eliminate
successively the decomposable remainder of each degree,
whose products have factors of strictly smaller degree.
Its kernel is zero because these are integral algebras
of the same dimension after extending to the algebraic closure,
and a proper closed subset of an irreducible variety has
smaller dimension. This proves the product parametrization.
The same proof for a cocharacter on a wall uses just the
strictly positive weights: the tangent space of its contracting
group has no zero-weight summand, and the coordinate grading is
again strictly positive away from constants.
For a regular cocharacter the argument also shows that the ordered negative root product,
then \(T\), then the positive root product is an open
immersion into \(G\) by Result~\ref{res:C02}.
Compare \cite[Theorem 5.1.13]{conrad}.

A central element belongs to \(T=Z_G(T)\), scheme-theoretically.
Conjugation on \(U_\alpha\simeq\Ga\) is multiplication by
\(\alpha(t)\), so the equation for centralizing it is
\(\alpha(t)=1\), by testing the additive parameter \(1\).
Conversely an element of \(T\) satisfying these equations
centralizes every root group and \(T\), hence the product
open just constructed. This open is schematically dense
and remains so after every base change from \(k\), by flatness
of all \(k\)-algebras. The two morphisms to the separated
\(G\) therefore agree everywhere. Finally the character
description of Result~\ref{res:B01} identifies the simultaneous
equations with \(D(X/\ZZ\Phi)\). This proves the scheme
centre formula even for nonreduced test algebras.
\end{proof}

Thus \(p\)-torsion in the quotient lattice can produce a non-smooth centre in characteristic \(p\). For instance, the root lattice of \(\SL_2\) has index two in its character lattice, giving \(\mu_2\) even in characteristic two. Neither the rank of \(T\) alone nor the abstract set of roots records the full lattice data.

\section{Borel subgroups, pinnings, and classification}

\begin{result}{C06}{Borel pairs, positivity, and pinnings}
Over \(\bar k\), the maximal smooth connected solvable subgroups of \(G_{\bar k}\) are conjugate. They are called its Borel subgroups. The pairs \((B,T)\) consisting of a Borel subgroup and a maximal torus in it are conjugate as well. For a smooth closed \(P\subset G\), the variety \(G/P\) is proper exactly when \(P_{\bar k}\) contains a Borel subgroup. If \(T\) is a split maximal torus of \(G\), Borel subgroups containing \(T\) correspond to positive systems in \(\Phi(G,T)\). The Weyl group acts simply transitively on these systems. For \(B\supset T\), one has \(B=T\ltimes R_u(B)\), and an ordered product of the positive root groups parametrizes \(R_u(B)\).
\end{result}

\begin{proof}
Geometric conjugacy of Borels and maximal tori in each Borel
was proved in the geometric portion of Result~\ref{res:F09}.
The properness of the Borel quotient was deduced therefrom
in Result~\ref{res:C02}. If a smooth closed \(P_{\bar k}\)
contains \(B\), the map \(G_{\bar k}/B\to G_{\bar k}/P_{\bar k}\)
is surjective and the source is proper. The target is separated
and of finite type by Result~\ref{res:D01}. It is universally
closed over the field: after any base change, the image
of a closed subset of \(G_{\bar k}/P_{\bar k}\) under the
projection to the base is the image of its closed inverse
image in \(G_{\bar k}/B\). Thus it is proper. Conversely,
if \(G_{\bar k}/P_{\bar k}\) is proper, a \(B\)-fixed point
exists by Result~\ref{res:B14}. Its stabilizer is a conjugate
of \(P_{\bar k}\), giving the required containment.
Properness descends from the algebraic closure.

The correspondence between Borels through \(T\) and
positive systems, including simple transitivity of \(W\),
was established in the chamber argument of
Result~\ref{res:C05}. There each Borel was shown to
equal \(P_G(\lambda)=T\ltimes U_G(\lambda)\) for a
regular cocharacter. Its unipotent radical is \(U_G(\lambda)\):
it is smooth connected normal unipotent with torus quotient,
which characterizes it by Result~\ref{res:B07}. Its
parametrization by an ordered product was proved in the
same result using the positive coordinate grading.
Compare \cite[Section 21.d]{milne} and \cite[Theorem 5.1.13]{conrad}.
\end{proof}

For a split pair, fix \(B\supset T\) and write \(\Delta\) for the simple roots of the chosen positive system. A \emph{pinning} specifies, for each \(\alpha\in\Delta\), an isomorphism \(x_\alpha:\Ga\to U_\alpha\) equivariant for the character \(\alpha\). Equivalently it specifies a nonzero vector in each simple root line. Changing the parametrization multiplies the vector by a nonzero scalar. Over an arbitrary field a connected reductive group with a \(k\)-Borel is \emph{quasi-split}. Its maximal torus in that Borel need not be split: Result~\ref{res:D10} will give an example.

\begin{result}{C07}{Central isogenies and the two lattices of type \(A_1\)}
Suppose \(f:(G_1,T_1)\to(G_2,T_2)\) is a central isogeny of split connected reductive pairs carrying \(T_1\) onto \(T_2\). Then \(f^*:X^*(T_2)\to X^*(T_1)\) is a finite-index injection. It identifies the roots bijectively, the corresponding root-group restrictions are isomorphisms, and the dual map \(f_*:X_*(T_1)\to X_*(T_2)\) carries each coroot to its corresponding coroot. For a semisimple split group, the adjoint form has character lattice equal to the root lattice, and the simply connected form has cocharacter lattice equal to the coroot lattice.
\end{result}

\begin{proof}
The central kernel is contained in \(T_1\) by the scheme
center formula (Result~\ref{res:C05}). Each torus restriction
is an isogeny, so the character anti-equivalence of
Result~\ref{res:B01} gives the finite-index injection.
The kernel of \(df\) is contained in \(\Lie(T_1)\) by
Result~\ref{res:A12}. Equivariance sends a nonzero root
space of \(G_1\) injectively into the weight space of the
pulled-back \(T_2\)-action on \(\Lie(G_2)\). Thus its root
has the form \(f^*\beta\). Distinct roots give distinct
\(\beta\)'s, and both groups have the same dimension and
torus rank, so the count in Result~\ref{res:C04} makes
this a bijection.

For \(\alpha=f^*\beta\), \(f\) sends the maximal subtorus
of \(\ker\alpha\) onto that of \(\ker\beta\), from the
finite-index identification of the two rational character
spaces. Hence the rank-one centralizer at \(\alpha\) maps to
that at \(\beta\), and taking the limit along the image of a
positive cocharacter sends \(U_\alpha\) to \(U_\beta\).
The restriction has trivial scheme kernel since \(U_\alpha\)
meets \(T_1\) trivially. By Result~\ref{res:A07} its nontrivial
image in the target \(\Ga\) is a closed subgroup isomorphic
to \(U_\alpha\). A closed one-dimensional subgroup of that
smooth connected curve is the whole group. Thus the restriction
is an isomorphism in every characteristic.

The rank-one representative for \(\alpha\) lies in its
centralizer and induces the nontrivial rank-one Weyl action.
Its image is in the corresponding normalizer and induces
the nontrivial Weyl action for \(\beta\): on rational
cocharacters it is conjugate to the original reflection by
the isomorphism \(f_* \otimes \QQ\).
Comparing the reflection formulas of Result~\ref{res:C04}
and using \(f^*\beta=\alpha\) gives
\(f_*(\alpha^\vee)=\beta^\vee\).
The scheme center formula yields the adjoint criterion
\(X=\ZZ\Phi\).

For the simply connected criterion, use only the realization
and central lifting in Result~\ref{res:C08}. The argument for
those two assertions below uses the root and coroot formulas
of C04 and C05, independently of the existence of the
simply connected cover here. For a semisimple group
the common annihilator of the roots is zero by
Result~\ref{res:C03}, so the coroots span \(Y\otimes\QQ\).
The sublattice \(Y_{\mathrm{sc}}=\ZZ\Phi^\vee\)
is therefore of finite index in \(Y\). With its dual lattice
and the same root and coroot vectors this is a root datum:
the reflection formulas restrict to these lattices and
the pairings stay integral. The inclusion has power \(1\)
on each root and coroot. The realization and isogeny lifting
produce a split group with this datum and a central isogeny
onto the given group. If the original group is simply
connected the index must be one.
Conversely every central isogeny has source cocharacter
lattice a sublattice of \(Y\) containing the coroot lattice,
by the coroot calculation above. When \(Y=\ZZ\Phi^\vee\)
such a torus map is an isomorphism, so the central kernel
inside the torus is trivial and the isogeny is an isomorphism
by Result~\ref{res:A07}. This argument is geometric in the
converse, where we may split all source tori after field
extension. Compare \cite[Proposition 23.59]{milne}.
\end{proof}

For \(\SL_2\), let \(x(\operatorname{diag}(t,t^{-1}))=t\), and write \(y\) for the dual generator. The root and coroot are \(2x\) and \(y\). For \(\operatorname{PGL}_2\), use the ratio of diagonal entries as the generator \(x'\), with dual \(y'\). Its root and coroot are \(x'\) and \(2y'\). The central quotient sends \(x'\) to \(2x\) by pullback and \(y\) to \(2y'\) by pushforward. Its kernel \(\mu_2\) has a one-dimensional tangent space in characteristic two. More explicitly, the derivative \(\mathfrak{sl}_2\to\mathfrak{gl}_2/kI\) sends a matrix to its class and kills exactly \(kI\) in characteristic two. The analogous map \(\SL_n\to\operatorname{PGL}_n\) has kernel \(\mu_n\).

An isogeny need not be surjective on rational points. For \(\SL_n\to\operatorname{PGL}_n\), take a representative \(A\in\GL_n(k)\) for a projective class (Hilbert 90 for the scalar torus). Adjusting the determinant to \(1\) by a scalar is possible exactly when \(\det A\in(k^\times)^n\). Over \(\RR\) with \(n=2\), a class with negative determinant gives an explicit failure of lifting. Result~\ref{res:D02} identifies this as a torsor obstruction.

\begin{result}{C08}{Existence and isomorphism for split pinned groups}
The groupoid of split pinned connected reductive groups over \(k\) is equivalent to the groupoid of based reduced root data. Every such datum is realized. A specified based root-datum isomorphism lifts to a unique pinning-preserving group isomorphism. On forgetting the pinning, the isomorphism classes are still classified by reduced root data up to isomorphism.
\end{result}

\begin{proof}
There is a based root-datum functor by Results~\ref{res:C04}
and~\ref{res:C05}. We first prove faithfulness, also after
any scalar algebra extension. It suffices to treat an
automorphism acting identically on the torus and the pinned
simple root coordinates. For a simple root \(\alpha\), it
preserves the rank-one centralizer and its two dynamic groups
\(P^\pm\), since those are the centralizer and limit functors
of subtori and cocharacters of the fixed torus.
On the rank-one quotient \(H/P^+\simeq\mathbb P^1\)
from Result~\ref{res:C04} it fixes \(0,\infty\).
It also commutes with the upper transvections, which are the
action of the pinned positive root group in those coordinates.
A projective automorphism fixing \(0,\infty\) is a scaling
over every test algebra, and a scaling commuting with all
translations is the identity, as seen by setting the translation
parameter to \(1\). Hence it acts trivially on this projective
line. Its equivariance for the \(H\)-action shows that it fixes
the negative root parametrization as well, since the negative
root group acts faithfully by lower transvections
(Result~\ref{res:C04}). It fixes the resulting simple
Weyl representative expressed as a word in the two root
subgroups. All root groups are conjugate to simple ones
up to sign by these representatives, by the chamber argument
of Result~\ref{res:D04}. The automorphism is thus the
identity on \(T\) and the root groups and hence on the open
product of Result~\ref{res:D03}. That open remains
schematically dense over every scalar algebra, giving
identity everywhere. Applying this to the ratio of two lifts
proves uniqueness of a pinned isomorphism.

Here is a characteristic-zero construction of both existence
assertions. We first use a dimension observation. Let \(H_i\)
be finitely many smooth geometrically connected subgroups of an
affine group in characteristic zero. Suppose their Lie algebras
span a subspace \(L\) preserved by every \(H_i\) under conjugation.
In a multiplication map from a product of the \(H_i\), after
translation to the identity, the derivative is a sum of
conjugates of the subalgebras \(\Lie(H_i)\), so it takes values in
\(L\). In characteristic zero the dimension of the closure of
the image is its generic differential rank. Indeed, at the
generic point of the image the extension of function fields is
separable, and the differential of a transcendence basis from
the image stays linearly independent. Among all words choose
one with maximal image dimension and let \(D\) be its reduced
closure over an algebraic closure. Adding \(H_i\) on the left
gives an irreducible closure of the same dimension containing
\(D\), so \(H_iD\subset D\) for every \(i\). It follows that \(D\)
is exactly the closure of the subgroup generated by the \(H_i\).
Inversion and multiplication preserve this closure by density,
including density in the product of two copies. The construction
descends to the original field since the word and its
scheme-theoretic closure are defined there and closure
commutes with field extension on affine schemes. Thus the
generated group is smooth, geometrically connected, and has
Lie algebra \(L\): its dimension is at most \(\dim L\) and its
Lie algebra contains the summands spanning \(L\).

Given \(R=(X,\Phi,Y,\Phi^\vee)\), first construct the model
over \(\QQ\). For \(\Phi=\varnothing\) take \(D(X)\).
Otherwise take the split semisimple Lie algebra with the
Cartan matrix of a base of \(\Phi\), and add the vector
space of central directions
\[
 Z_{\QQ}=\{v\in Y\otimes\QQ:\alpha(v)=0\text{ for every }\alpha\in\Phi\}.
\]
The Cartan subspace of the enlarged algebra is \(Y\otimes\QQ\):
the simple coroots span a direct complement to \(Z_{\QQ}\)
because the Cartan matrix of a finite root system is
invertible. Use the Serre generators and relations over \(\QQ\).
Their scalar extension to \(\mathbb C\) is the construction
in \emph{Lie Algebras and Finite-Dimensional Representations},
Chapter~6, Results F02--F05. Since scalar extension is faithful,
the triangular decomposition and the one-dimensional root
spaces over \(\QQ\) follow from those results. Choose the
positive simple vectors and their linked negatives in this
presentation. For a nonsimple root, a rational root vector
and its opposite can also be linked by rescaling the opposite
one: their bracket is a nonzero rational multiple of the
displayed coroot in characteristic zero.

Choose finitely many dominant characters \(\lambda\in X\)
generating \(X\) as a group. For instance, if there are roots,
add sufficiently large multiples of \(2\rho=\sum_{\alpha>0}\alpha\)
to a lattice basis of \(X\), and include \(2\rho\) itself.
For each \(\lambda\), form the
rational highest-weight module by imposing the Verma
relations and \(f_i^{\lambda(\alpha_i^\vee)+1}v_\lambda=0\),
and let the central directions act through \(\lambda\).
The formation of this module commutes with field extension:
its presentation is a quotient of a rational enveloping
algebra by the left ideal of the stated relations. The
finite dimension, the highest line, and the weight support
in \(\lambda-\ZZ_{\geq0}\Delta\) follow by extension to
\(\mathbb C\), using Chapter~7, Results G03--G06 of the same
notes. In particular all weights lie in \(X\).
On the direct sum \(V\) let \(T=D(X)\) act by these weights.
The action is a closed embedding because the weights contain
the chosen generators of \(X\): the coordinate characters
of the image then generate the Laurent coordinate ring of
\(T\). Each root vector acts on \(V\) nilpotently, since it
shifts weights in a finite set. The Lie representation is
faithful: its Cartan action is faithful by the weight
generators, and a root vector acting by zero would, with
its opposite, make its nonzero coroot act by zero. The
kernel is stable under \(T\), so the Cartan and root-space
tests suffice.

Put \(x_\alpha(t)=\exp(t e_\alpha)\) in \(\GL(V)\), with
\(e_\alpha\) now acting on \(V\). The sum is finite
and the binomial formula gives the additive group law.
Its differential is nonzero, so it is an embedding
of \(\Ga\): an additive group in characteristic zero
has no nontrivial finite subgroup, and the scheme image
identifies the trivial-kernel quotient by Result~\ref{res:A07}. Conjugation by this exponential
on the image Lie algebra is the exponential of
\(\operatorname{ad}(t e_\alpha)\). Both \(T\) and
these root subgroups therefore normalize that Lie algebra.
Apply the dimension observation to generate \(G_\QQ\).
It has exactly the chosen split Lie algebra. An ideal
consisting of nilpotent endomorphisms cannot contain
a simple ideal, which contains a coroot acting nontrivially
and diagonally. It cannot have a nonzero projection onto a
simple ideal either, since bracketing with that ideal would
generate it. The remaining central directions act diagonally,
so such an ideal is zero. By B06 the Lie algebra of the
unipotent radical consists of nilpotent endomorphisms.
Thus \(G_\QQ\) is reductive. Its torus is maximal because
the zero weight space of its Lie algebra is exactly
\(\Lie(T)\).

The root characters are those in the multiplication
description of \(V\). For each root the Lie triple integrates
inside \(\GL(V)\) by the rank-one calculation:
in a weight string of highest weight \(m\), a highest
vector and the successive divided powers of the lowering
operator give the Lie matrices of \(\operatorname{Sym}^m\QQ^2\).
The strings span after extension to \(\mathbb C\), by the
rank-one decomposition. They can be chosen over \(\QQ\):
in each highest weight the kernel of the raising
operator on its weight space is defined over \(\QQ\),
and any rational basis yields the same decomposition
after extension. Consequently the resulting \(\SL_2\)
action has upper and lower root exponentials
\(x_\alpha,x_{-\alpha}\), and its diagonal acts on
\(V_\mu\) by \(t^{\langle\mu,\alpha^\vee\rangle}\).
The \(\SL_2\) action has image in \(G_\QQ\) since
its opposite elementary subgroups generate a dense
open cell. Its diagonal is the coroot of C04, whose
rank-one description characterizes it without changing
the torus lattice. This gives the prescribed pinned datum
over \(\QQ\), hence over every characteristic-zero field.

\paragraph{Integral realization.}
First suppose the datum is semisimple with cocharacter
lattice \(\ZZ\Phi^\vee\). Then the fundamental weights
form a basis of \(X\). In the rational construction
take their highest-weight modules, and choose the
root vectors to form a Chevalley basis. To make this
choice over \(\QQ\), take the rational span of the
integer bracket table in Chapter~6, Result F05 of
\emph{Lie Algebras and Finite-Dimensional Representations}.
Its simple root triples give the rational Serre
presentation: the map from that presentation becomes
an isomorphism over \(\mathbb C\) by F01--F03.
The coroots are unchanged. We first construct lattices
in these rational modules that are preserved by
every root subgroup polynomial. Only rational matrix
identities will be used to do so, but we track their
integer coefficients.

Use \(x_a(x)=\exp(xe_a)\) as in the rational model.
If \(a,b\) are independent roots, give their cone a
positive order and commute \(x_a(x)\) past \(x_b(y)\).
Other factors have roots \(ia+jb\), \(i,j>0\),
and parameters integer multiples of \(x^iy^j\).
Here is the calculation of integrality. Conjugation
by \(x_a(x)\) takes \(\exp(ye_b)\) to the exponential
of
\[
 y\sum_{j\geq0}x^j(\operatorname{ad}e_a)^j e_b/j!.
\]
If the root string starts at \(b-pa\), the coefficients
in its root lines have magnitude \(\binom{p+j}{j}\):
successive brackets have magnitudes \(p+1,\ldots,p+j\).
The possibilities for the subalgebra on these lines
can be checked in the four two-dimensional root systems
by listing their vectors in the simple basis. In
\(A_1\times A_1\) and \(A_2\) the string lines commute.
For \(B_2\), two long roots commute, and for two
short roots whose sum is a root the one extra bracket
of the ordered pair is central of magnitude \(2\).
When the first root is short and the second long,
the only noncommuting orientation is the simple
pair up to the Weyl group. The string lives on
\(b,a+b,2a+b\), which commute with one another.

For \(G_2\), the long roots give the \(A_2\) calculation.
Two short roots at angle \(60\) degrees have a single
extra root in their positive cone, their long sum;
the bracket coefficient is of magnitude \(3\).
At angle \(120\) degrees their string adds \(a+b\)
and \(2a+b\), and the remaining root in the cone
is \(a+2b\). The only bracket between these string
lines is between \(b\) and \(a+b\), of magnitude
\(3\); the coefficient of the \(a+b\) line already
has a factor \(2\). In the short-long case the
angles \(30\) and \(90\) degrees give commuting
groups. At \(150\) degrees use the simple pair
\(a,b\) up to the Weyl group. With
\(v_j=(\operatorname{ad}e_a)^j e_b/j!\) for
\(0\leq j\leq 3\), the only brackets among the
\(v_j\) lie in the central line for \(3a+2b\).
They satisfy
\([v_1,v_2]=-3[v_0,v_3]\), by applying
\(\operatorname{ad}e_a\) to \([v_0,v_2]=0\).
The coefficient of \([v_0,v_3]\) is \(1\) or \(-1\).
For completeness, the short and long lists used
here are \(a,a+b,2a+b\) and \(b,3a+b,3a+2b\),
and their negatives. In \(B_2\) they are \(a,a+b\)
and \(b,2a+b\), and their negatives. These lists
also show there are no further bracket pairs in
the strings. The orientations with the long root
first can be obtained by inverting the
conjugation product for the short root first.
For a central bracket,
\(\exp(A)\exp(B)=\exp(A+B+[A,B]/2)\);
this follows by differentiating the two
sides with \(A,B\) scaled by a common parameter.
Consequently all the listed string exponentials
separate into ordered root factors over \(\ZZ\).
In the two nonabelian \(G_2\) strings the
halves are integral, since the coefficients
are respectively \(2\cdot3\) and \(1-3\)
up to the stated signs. This proves the
required integer commutation assertion.

For opposite roots use the rank-one calculation
\[
 x_a(x)x_{-a}(y)
 =
 x_{-a}\bigl(y/(1+xy)\bigr)\,
 a^\vee(1+xy)\,
 x_a\bigl(x/(1+xy)\bigr).
\]
It holds by multiplication in \(\SL_2\) and the
rank-one action constructed above. At the identity
all denominators here have constant term one.
Together with torus conjugation, these formulas
give the big-cell product in coordinates as formal
power series over \(\ZZ\). Explicitly, use the
parameters for root groups and \(1+s_i\) in a
cocharacter basis for the torus, and sort words
into the order of the negative roots, torus, and
positive roots. Terms with parameters of orders
\(d,e\) commute modulo order \(d+e\), since
interchanging them in the displayed formulas
introduces only such corrections. The torus
corrections, including \(a^\vee(1+xy)\), have
coefficients given by integral binomial series.
To work through degree \(m\), retain all corrections of
orders at most \(m\) in the word, and discard higher
orders modulo degree \(m+1\). First order the original
factors while placing corrections at the end of the
word. Each correction starts at order two. When a
correction of order \(d\) crosses any factor, all the
new corrections start at order \(d+1\). Insert and
combine the corrections in increasing order until
order \(m\) is reached. At each order there are only
finitely many factors and crossings. This gives
consistent sorted coordinates in all truncations.
The compatible reductions give the integral
formal coordinates. They agree with the product
over \(\QQ\), by uniqueness of the big-cell
coordinates around the identity.

For a rational highest-weight module, take the
\(\ZZ\)-span \(M\) of the coefficients obtained by
applying an ordered big-cell product to the
chosen highest vector \(v_\lambda\).
Positive root factors are on the right and fix
this vector; torus coordinates multiply it by
an integral binomial series. Only the finitely
many coefficients of the product of the
negative-root polynomials are needed, so \(M\)
is finite free and spans the rational module.
Spanning follows because the ordered negative
monomials span by PBW, and factorials are
invertible over \(\QQ\). Each generator has
a specified \(X\)-weight, and the highest
weight lattice is exactly \(\ZZ v_\lambda\).
To check stability, multiply the product on
the left by \(x_a(z)\) and collect formal
coefficients. The just-proved integral product
in the big-cell parameters expresses each
coefficient as an integral linear combination
of those defining \(M\). Hence all coefficients
of \(x_a(z)\) preserve \(M\), and so does the
split torus by its weight decomposition. The
same holds for \(x_a(-z)\). On the sum of the
fundamental lattices the torus is still a closed
matrix subgroup, since its weights include the
fundamental weights.

We will need a linear coefficient equal to a unit
for each root operator on this sum. For a simple
root \(a_i\), use the \(\omega_i\) module and
\(e_{a_i} e_{-a_i} v_{\omega_i}=v_{\omega_i}\).
Use the coefficient functional with value 1
on \(v_{\omega_i}\) in its highest-weight
lattice. For another root conjugate this test
by a product of simple Weyl representatives.
These act integrally on all the \(M\)'s, by
their root word expressions. They also act
integrally and invertibly on the Chevalley
Lie lattice: on the root strings this follows
from the divided coefficients above, and on
the rank-one triple it follows from
\([e_a,e_{-a}]=a^\vee\) and
\([e_a,a^\vee]=-2e_a\). On other Cartan
directions the exponential adds only a
multiple of \(e_a\) with integral coefficient.
Thus these representatives take one primitive
root vector to another up to sign, so the
transported test still gives a unit. In
particular, in every root direction there is
an integral vector of some weight \(\mu\)
and an integral functional of weight
\(\mu+a\) whose matrix coefficient on \(e_a\)
is 1.

Choose an ordering of each sign of the root
system and a weight basis ordered by a regular
cocharacter. The products of the positive and
the negative root matrices are respectively
upper and lower unitriangular. Each ordered
product is a closed immersion into the
corresponding triangular group over \(\ZZ\).
Indeed the coefficient test just constructed
for a root \(a\) reads its parameter plus a
polynomial in parameters of roots of strictly
smaller positive height (use negative height
for the negative group). A matrix term
contributing to that coefficient must shift
the weight by exactly \(a\); every root term
in the product has positive height in the
chosen sign. A term involving an equal-height
root can therefore only be the term linear
in \(a\). Induction in height expresses the
inverse coordinate functions as polynomials
in the matrix entries, over \(\ZZ\). This is
the ring surjectivity for closed immersion.
The two images are subgroups, since the
closed-subscheme equations for products and
inverses hold over \(\QQ\) and the source
coordinate algebras of those operations
are torsion-free.

The open subset of the full matrix group
where the leading principal minors are units
admits unique lower--diagonal--upper
decomposition over any ring, by row and
column elimination one pivot at a time. In
these coordinates the product
\(U^-\times T\times U^+\) just constructed
is a smooth closed subscheme of that open
subset. Its generic fibre is the big cell
of \(G_\QQ\), by D03. Since it is dense in
the intersection of \(G_\QQ\) with the
matrix open and closed there, it is the
whole intersection. Take the schematic
closure \(H\) of \(G_\QQ\) in the matrix
group over \(\ZZ\). The closure is flat,
since its coordinate algebra is the image
of the matrix coordinate algebra in
\(\QQ[G_\QQ]\). It is a subgroup, since
the group identities defining the image
persist by torsion-freeness on products
and inverses. Formation of schematic
closure commutes with this open
restriction, as can be checked by
localizing the defining ideal. Hence
the integral big cell \(\Omega\)
constructed above is an open subscheme
of \(H\), containing the identity in
every fibre. Each geometric fibre of
\(H\) is smooth by translating this
open neighborhood in the fibre; flatness
and finite presentation give smoothness of
\(H\) over \(\ZZ\).

One can retain only the geometrically connected fibre components,
staying affine. The image of \(\Omega\times\Omega\) under multiplication
is the open subgroup of fibrewise identity components: multiplication is
open by smoothness, and every element of an identity component is a
product of two elements of the dense open cell in that irreducible
component. The complement is closed with empty generic fibre, hence lies
over finitely many primes. To remove the other components at one such
prime \(p\), write \(R_H\) for the coordinate ring and \(I\) for the
defining ideal of the retained component in the special fibre, viewed as
an ideal containing \(p\). The algebra \(R_H[I/p]\) inside \(R_H[1/p]\)
represents exactly the desired open restriction. In fact \(I=(p)\)
locally off the removed components, since the retained component is open
and closed in the special fibre, so the algebra does not change there.
Over a removed point a function \(i\in I\) is nonzero in the residue
field, whereas the equation \(p(i/p)=i\) would make it zero. There are
no other points above the complement. Repeating at the finitely many
primes gives the affine smooth connected-fibre model.

Its torus and tangent weights are visible on \(\Omega\), as are the root
and coroot relations. The rank-one map from \(\SL_2\) exists over
\(\ZZ\): the pivot cell formulas use \(x_a,x_{-a}\) and the integral
cocharacter, and the other pivot cell is obtained by its integral Weyl
representative. These maps glue since they agree on the generic fibre of
the overlap. The group law agrees over \(\QQ\) and therefore over
\(\ZZ\) on the flat product. Each coroot is primitive in
\(\ZZ\Phi^\vee\), being Weyl conjugate to a simple basis element. Its
tangent vector is consequently nonzero in every fibre. If the Lie
algebra of the unipotent radical were nonzero, it would contain a
nonzero diagonal vector or a root line, since it is stable under the
split torus. The former cannot be nilpotent, and bracketing the latter
with its opposite also gives a nonzero diagonal coroot. This argument
applied to the radical itself proves its triviality and hence
reductivity in each fibre. The zero-weight space proves maximality of
the torus. The rank-one maps identify the coroots as in C04. The
integral model therefore has exactly the desired simply connected
semisimple datum.

For a general datum, put \(Y_{\rm der}=\ZZ\Phi^\vee\) and
\(Y_0=\bigcap_{\alpha\in\Phi}\ker\alpha\) in \(Y\). The map \(Y_{\rm
der}\oplus Y_0\to Y\) is a finite-index injection, by the direct-sum
decomposition of its rationalization. Form the integral product of the
simply connected model for \(Y_{\rm der}\) and the split torus for
\(Y_0\). Its torus maps by this lattice injection to the specified
torus, with finite diagonalizable kernel \(Z\). This kernel is central:
every root character is the pullback of the corresponding character of
the specified torus, so \(Z\) commutes with all root groups and with the
torus. Multiplication from two copies of the open cell is smooth and
surjective on this connected-fibre model. The equality expressing
centrality therefore descends from its product with the kernel under
this faithfully flat cover.

Here the central quotient can be constructed directly. Write
\(A=\bigoplus_{m\in M}A_m\) for the grading for right translation by
\(Z=D(M)\). The map that sends \((g,z)\) to \((g,gz)\) is a closed
immersion: use the isomorphism \((g,h)\mapsto(g,g^{-1}h)\) of the full
product. On its coordinate rings this surjectivity represents \(1\otimes
e_m\) as a sum of tensors, whose components of degrees \((-m,m)\) give
\(1\in A_{-m}A_m\). These identities make multiplication
\(A_m\otimes_{A_0}A_{-m}\to A_0\) an isomorphism: if \(\sum_\nu u_\nu
v_\nu=1\) in opposite degrees, insertion of this sum gives the inverse
and also the dual-basis relations. Thus the \(A_m\) are invertible
\(A_0\)-modules, and \(A\) is finite faithfully flat over \(A_0\); the
same identities give the torsor isomorphism \(A\otimes_{A_0}A\simeq
A\otimes\ZZ[M]\). Centrality makes the Hopf maps descend to \(A_0\). For
finite type, choose a finite list of \(A_0\)-module generators of \(A\),
including 1. Adjoin to \(\ZZ\) the finitely many coefficients expressing
their products and the original algebra generators. Then \(A\) is
already finite over the resulting noetherian subalgebra of \(A_0\). Its
submodule \(A_0\) is also finite over this subalgebra. It follows that
\(\Spec A_0\) is an affine flat group of finite presentation.

The product cell is invariant under \(Z\), and its fppf quotient is the
open cell \(U^-\times(T'/Z)\times U^+\), where \(T'\) denotes the source
torus. Indeed an invariant open descends under the open surjective
torsor, and scalar extension identifies its coordinates with the
asserted product. The quotient therefore has smooth, connected geometric
fibres. It is reductive: a nontrivial smooth unipotent radical of a
geometric fibre would have a preimage whose reduced identity component
is a positive-dimensional connected solvable normal subgroup of the
source. Its tori map trivially to the radical and hence are trivial, as
the kernel is finite. By B12 it is unipotent, contradicting source
reductivity. The torus is the specified one, the root coordinates pass
through the quotient isomorphically, and each rank-one map from the
source has cocharacter the prescribed image of the coroot. This
completes the integral realization without using the covering existence
assertion of C07.

To prove lifting in characteristic zero, start with
two existing split pinned groups and a root bijection
\(d:\Phi_1\to\Phi_2\), compatible with the bases, and
a finite-index character injection \(h:X_2\to X_1\)
satisfying
\[
 h(d(\alpha))=\alpha,\qquad
 h^\vee(\alpha^\vee)=d(\alpha)^\vee .
\]
The same Serre relations hold on their linked simple
root triples. The mixed brackets vanish since a difference
of distinct simple roots is not a root, and the
Serre powers vanish by the root strings. These triples
generate all root spaces: conjugation by the simple
Weyl representatives is expressed by products of root
exponentials on the adjoint module, and it reaches every
root space from a simple one. The Serre algebra and the
central directions describe the Lie algebras by dimension,
so the displayed torus map and corresponding pinned triples
give a Lie isomorphism \(\ell\). For each \(\alpha\), let
the map \(U_\alpha\to U_{d(\alpha)}\) be the unique root
group isomorphism with derivative \(\ell\) there. The
adjoint actions of the paired root groups intertwine
\(\ell\) since an algebraic \(\Ga\)-action in characteristic
zero is its exponential, by differentiating its group law.
The paired torus actions intertwine it by the character
identities. Thus in \(G_1\times G_2\) every graph subgroup
normalizes the graph of \(\ell\). The dimension observation
gives a connected group \(P\) with precisely this graph
Lie algebra.

Both projections from \(P\) are surjective: the images
contain \(T_i\) and the root groups, hence the dense
cell of D03. The projections have finite kernels and these
kernels are central, since a smooth connected group acts
trivially on a finite \'{e}tale group by conjugation. A
unipotent radical of \(P\) would map trivially to \(G_1\)
and then have dimension zero. Thus \(P\) is reductive,
and the graph torus is maximal by its zero weight space.
The first projection is an isomorphism on the graph
torus, and by C05 its central kernel lies in that torus,
so the first projection is an isomorphism by A07.
The second projection therefore gives the required
central isogeny. If \(h\) is an isomorphism then the
second projection is also an isomorphism on the graph
torus, and the same kernel argument applies.

\paragraph{Recognizing integral models.}
Suppose two smooth affine groups over \(\ZZ\) have connected
reductive geometric fibres and split fibrewise maximal tori
whose roots form the same fixed based datum. A pinning here
chooses an integral basis of each simple root line, and the
root-datum identification matches these choices. The two
integral groups are isomorphic under that identification.

To construct the necessary coordinates, let \(A\) be the
coordinate algebra of either group, graded for conjugation
by \(T\), and \(I\) its augmentation ideal. The torus
characters on \(I/I^2\) are the tangent weights in the
convention \(f(tgt^{-1})=\chi(t)f(g)\). This is a finite
free weight-graded module since the group is smooth.
For a root \(\alpha\), restrict first to the centralizer of
the largest subtorus of \(\ker\alpha\) using its weight
equations. Then impose contraction to the identity for
a cocharacter positive on \(\alpha\), just as in C02.
On the coordinate algebra this kills the functions of
nonzero weight for that subtorus, all functions of negative
cocharacter weight, and all augmentation functions of
cocharacter weight zero. These are equations for the
indicated limit functor because a Laurent series along the
cocharacter extends to zero precisely when its negative
coefficients vanish, and its constant term then gives
the limit. They commute with base change by their
homogeneous construction. The resulting augmentation
quotient of the cotangent module keeps just the
\(\alpha\)-line. Its coordinate ring has positive
cocharacter grading off the constants, so a homogeneous
lift of a basis of the cotangent line generates it, by
induction on degree. Over \(\QQ\), the coordinate is a
polynomial coordinate by C04. Thus the surjection
\(\ZZ[x]\) onto that ring has zero kernel. Its coproduct
is primitive by grading and the counit identities, since
a product of two positive-degree variables has higher
degree than \(x\). This gives the integral root group
with its specified tangent basis.

To link opposite generators, fix the simple \(x_+\) and
temporarily choose any integral \(x_-\), with tangent
vectors \(X_+\) and \(X_-\). The coefficient of
\(u^2X_+\) in \(\operatorname{Ad}(x_+(u))X_-\)
is an integer. In each residue field, compare the adjoint
action with the projective line in the proof of C04.
The two root groups map isomorphically to the upper and
lower transvections, so if \(x_+\) maps with scale one and
\(x_-\) with scale \(c\), the corresponding matrix
calculation gives coefficient \(-c\). Here the action
on the positive root line of the Lie algebra is read
faithfully in the projective group, even in characteristic
two. Thus the coefficient is nonzero in every residue
field, so is a unit in \(\ZZ\). Rescale \(x_-\) to make it
\(-1\). The linked Weyl word from C04 now defines a
\(\ZZ\)-point normalizing \(T\): the normalizing equality
holds over \(\QQ\) and extends by flatness of \(T\).
Conjugating simple root coordinates by these words gives
integral coordinates on all root groups. The conjugate
maps are isomorphisms by the same weight and contraction
description, or by applying the inverse Weyl word.

Multiplication of the ordered integral root groups and
the torus gives the integral open cell. The morphism is
\'etale because its fibre maps are the open immersions of
D03 and the source and target are smooth over \(\ZZ\);
equivalently, the relative differentials are invertible
everywhere. Its geometric fibres over points of the
target have at most one point, so it is an \'etale
radicial morphism, hence an open immersion. The pinned
generic isomorphism just proved matches the linked simple
coordinates and their Weyl conjugates. It therefore
matches the ordered integral cells by an integral
isomorphism \(f_\Omega:\Omega_1\to\Omega_2\).

The map \(\Omega_1\times\Omega_1\to G_1\) given by
\((a,b)\mapsto ab^{-1}\) is a smooth surjection:
it is the restriction of a projection up to change of
coordinates on \(G_1\times G_1\), and is onto in
each geometric fibre by intersection of two dense opens.
On this cover use the morphism
\((a,b)\mapsto f_\Omega(a)f_\Omega(b)^{-1}\).
The two pullbacks to the overlap are equal over
\(\QQ\), hence everywhere, since the overlap is smooth
over \(\ZZ\) and the target is separated. Indeed
the overlap is smooth over \(\Omega_1\times\Omega_1\)
by base change, so its generic fibre is schematically
dense. Faithfully flat descent gives the integral morphism.
The homomorphism identity follows by the same density on
\(G_1\times G_1\). Applying the construction in reverse
gives its inverse. The central-isogeny version follows too.
For the indicated torus map, its finite diagonalizable
kernel is central in the source by the root-character
condition: it commutes with all root coordinates and hence
with the group by the smooth cell cover above. Form its
affine quotient by the grading argument in the integral
realization. The cell of the quotient retains its root
coordinates and has the quotient torus. The rank-one
root words show that the coroots are the images of the
source coroots, so the quotient and the specified target
have the same pinned datum. The argument of this paragraph
identifies them. In particular neither version for these
integral models uses recognition of arbitrary groups over
positive-characteristic fields.

\paragraph{Lifting in arbitrary characteristic in type \texorpdfstring{\(A_1\)}{A1}.}
The preceding lattice conditions also suffice for products of type \(A_1\)
when the torus map is inseparable. Use the linked coordinates
of Result~\ref{res:C04} and put
\(w(z)=x_+(z)x_-(-z^{-1})x_+(z)=\alpha^\vee(z)w(1)\).
The rank-one identity
\[
 x_+(u)x_-(v)
 =x_-\bigl(v/c\bigr)\alpha^\vee(c)x_+\bigl(u/c\bigr),
 \qquad c=1+uv\ne0 ,
\]
holds in every such group. Indeed, for \(v\ne0\) rewrite
the left side as \(x_+(c/v)w(-1/v)x_+(1/v)\),
and rewrite \(x_-(v/c)\) as
\(x_+(c/v)w(-c/v)x_+(c/v)\).
Since \(w(-c/v)^{-1}w(-1/v)=\alpha^\vee(c)\),
cancellation on the left and the torus action on \(x_+\)
give the displayed formula. At \(v=0\) the identity
extends by density and separatedness.

Root groups of distinct factors commute, including with
either choice of sign. For roots \(a,b\) from distinct factors,
choose a regular cocharacter making both positive. The
commutator is then in the positive unipotent group, whose
ordered root coordinates from C05 each transform by the
corresponding character. Its coordinate functions in \(u,v\)
vanish when either parameter is zero. Each term \(u^iv^j\)
would therefore have \(i,j>0\) and would require the root
\(ia+jb\) by torus equivariance. There is no such root in this
system, so all the commutator coordinates vanish.

Map the open cell of \(G_1\) to that of \(G_2\) by the
given torus homomorphism and by the matched linked root
coordinates. The equality of roots intertwines the
torus commutations, and the equality of coroots intertwines
the displayed rank-one identity. Roots from different
factors commute by the preceding calculation. Thus the cell
map respects multiplication on a dense open of the product
of the cells: multiply the two factor words, commute the
distinct factors, and apply the displayed identities where the
finitely many denominators \(1+uv\) are units. All denominators
are units at the pair of identity elements. For the rootless
case the map is just the given torus homomorphism.

A rational homomorphism obtained in this way extends
everywhere. Over an algebraic closure, for each \(a\) in
the source cell \(\Omega\), use \(g\mapsto f(ga)f(a)^{-1}\)
on \(\Omega a^{-1}\). The cell identity holds whenever
\(g,a,ga\) are in \(\Omega\), because the two morphisms
agree on a dense open of that domain. The formulas all agree
with \(f\) on their intersections with \(\Omega\), hence with
one another. The translated opens cover since
\(\Omega\cap g^{-1}\Omega\) is nonempty in the irreducible group.
The extended morphism descends: both pullbacks to the tensor
square of the field extension agree on the schematically dense
cell. Its homomorphism identity extends from the dense open
of the product. Its image contains the target torus and root
groups, so it is the whole target by A07 and D03. The source
cell is the inverse image of the target cell by the Bruhat
indices of D04, since the map sends the simple Weyl representatives
and the two Borel subgroups to their counterparts. This identifies
the open subschemes by their geometric points. In the cell
coordinates its kernel is just the finite torus kernel, which
is central by C05. This proves the claimed lift and, when the
lattices are identified, the pinned isomorphism.

\paragraph{Lifting in rank two in types \texorpdfstring{\(A_2\) and \(B_2\)}{A2 and B2}.}
In the following calculations all groups and root coordinates are
those of an existing split group over the given field. For two
independent roots, choose a chamber in which both are positive.
Their commutator belongs to the positive unipotent group. Its
coordinate in a root \(\gamma\) is a polynomial in the two
parameters \(u,v\); by equivariance under \(T\) and vanishing on
the two axes its monomials must have weights \(ia+jb=\gamma\)
with \(i,j>0\). Thus the positive-cone root list prescribes all
possible correction factors and the powers in their parameters.
The same weight test applied to conjugation in a parabolic
radical will be used below. In a Levi of semisimple rank one,
write \(n_a=x_a(1)x_{-a}(-1)x_a(1)\) in its linked coordinates.
Then \(n_a^2=a^\vee(-1)\) because \(w(-1)=w(1)^{-1}\).
Conjugation sends \(x_a(u)\) to \(x_{-a}(-u)\), as can be
read faithfully on the corresponding transvections of C04.

First consider an abelian two-root string in a rank-one
parabolic radical, or in a quotient of that radical, for which
the two roots have weights \(-1,1\) under the simple coroot.
Order the coordinates by these weights and normalize the second
root coordinate as the translate of the first by \(n_a\).
Conjugation by \(x_a(u)\) and \(x_{-a}(v)\) has matrices
\[
 L(u)=\begin{pmatrix}1&0\\ Au&1\end{pmatrix},\qquad
 R(v)=\begin{pmatrix}1&Dv\\0&1\end{pmatrix}.
\]
In the instances below the string roots have the same
nonzero coefficient on the other simple root. The full torus grading then forces
linearity in the string coordinates, and forces the indicated
powers of \(u,v\). The matrix \(L(1)R(-1)L(1)\) interchanges
the two weight lines, so its diagonal entries \(1-AD\) vanish.
Its lower left entry is \(A(2-AD)=A\), which is \(1\) by
the normalization. Thus \(A=D=1\) in any characteristic.

In type \(A_2\) put \(c=a+b\) and choose \(x_c\) by
conjugating \(x_b\) by \(n_a\). The two positive roots \(b,c\)
are the abelian radical for the Levi on \(a\). The two-root
calculation gives
\[
 x_a(u)x_b(v)x_a(-u)=x_b(v)x_c(uv).
\]
The other two nonopposite pairs of positive roots commute.
In the radical on \(a,c\) for the Levi on \(b\), reversing
the displayed commutator gives the coefficient \(-1\) for
the lower triangular action of \(x_b(v)\). The same
two-by-two product, now without changing the already
chosen \(c\) coordinate, has lower left entry \(-1\).
Hence \(n_b\) sends \(x_a(u)\) to \(x_c(-u)\).
The square identity for \(n_a,n_b\) gives the return scales.
All normalizer scales on the remaining opposite coordinates
are now determined: conjugation preserves the linked pair,
and scaling the positive generator scales the negative
generator by the reciprocal. The simple reflected root
uses the rank-one sign above. Thus both the positive
commutations and the simple representatives' action are fixed
by the two pins.

In type \(B_2\) take \(a\) short, \(b\) long, and write
\(c=a+b,\ d=2a+b\). Choose \(x_c=n_b x_a n_b^{-1}\)
and \(x_d=n_a x_b n_a^{-1}\). The radical on \(a,c,d\)
for the Levi on \(b\) has central \(U_d\), and its
quotient by \(U_d\) is the abelian two-root string on
\(a,c\). The centrality follows from the positive-cone
test, and \(U_d\) commutes with both \(U_b\) and \(U_{-b}\)
by the corresponding signed cone tests. The two-root
calculation therefore gives the coefficient \(+1\) for
the commutator of \(x_b(v)\) with \(x_a(u)\) in \(U_c\).
It also gives the scales for \(n_b\) on \(a,c\); \(n_b\)
fixes \(x_d\).

The radical on \(b,c,d\) for the Levi on \(a\) is abelian
and these roots have coroot weights \(-2,0,2\).
In the same order the action of \(x_a(u)\) has matrix
\[
 L(u)=\begin{pmatrix}
 1&0&0\\ Au&1&0\\ Bu^2&Cu&1
 \end{pmatrix}.
\]
All three roots have \(b\)-coefficient one, so equivariance
again forces linearity in the radical coordinates and
gives the displayed \(u\)-powers. The coefficient \(A\)
is \(-1\), as is seen by reversing the preceding commutator
modulo its \(d\)-coordinate. In this basis \(n_a\) exchanges
the first and last axes with scale \(1\) by the normalization
and the square identity. It acts by some \(s\) on the middle
axis, with \(s^2=1\). Since it conjugates \(x_a(1)\) to
\(x_{-a}(-1)\), it gives the negative action matrix
\[
 R(-1)=\begin{pmatrix}1&C/s&B\\0&1&sA\\0&0&1\end{pmatrix}.
\]
In \(L(1)R(-1)L(1)\) the top-right entry is \(B\), hence
\(B=1\). Its top-left entry is \(2+AC/s\), which is zero.
Its middle-right is \(A(1+s)\), also zero. Since \(A=-1\)
is a unit, \(s=-1\) and \(AC=2\), so \(C=-2\).
There is no division by \(2\) in this calculation.
This matrix specifies both correction factors for the pair
\(a,b\) and the only correction factor for the pair \(a,c\).
Every other distinct positive pair commutes by its root
list. The action of the two representatives on the positive
roots is fixed by the preceding calculations and the
rank-one sign; on opposite roots use the linked pairs.
Thus the whole list is forced, including in characteristic
two.

\paragraph{The rank-two group of type \texorpdfstring{\(G_2\)}{G2}.}
For the short simple root \(a\) and long simple root \(b\)
write \(c=a+b,d=2a+b,e=3a+b,f=3a+2b\).
Choose the positive coordinates for \(c,d,e,f\) by
conjugating the coordinates for \(a,c,b,e\) by
\(n_b,n_a,n_a,n_b\), respectively. Choose the linked
negative coordinates. The Levi on \(b\) acts on the
abelian quotient of its radical on \(a,c\) obtained
by discarding the roots of \(a\)-coefficient at least two.
It also acts on the abelian string on \(e,f\) in its radical.
The two-root calculation gives the coefficient \(+1\)
for \([x_b(v),x_a(u)]\) in \(U_c\), and gives
\([x_b(u),x_e(v)]=x_f(uv)\). The root \(d\)
commutes with both \(b\)-root groups by the signed
positive-cone test. This determines the action of \(n_b\)
on positive coordinates other than \(b\), using its square
on \(c\) and \(f\) to determine the return scales.

The radical \(N\) for the Levi on \(a\) has roots \(b,c,d,e,f\).
The subgroup \(U_f\) is central in \(N\) and is fixed by
the \(a\)-root groups by the positive-cone test.
Modulo \(U_f\) the radical is abelian, since the only
positive-cone corrections among \(b,c,d,e\) have root \(f\).
The four roots of the quotient have \(b\)-coefficient one.
Equivariance makes the action of \(x_a(u)\) on their
ordered coordinates linear, of the form
\[
 L(u)=\begin{pmatrix}
  1&0&0&0\\
  Au&1&0&0\\
  Bu^2&Cu&1&0\\
  Du^3&Eu^2&Fu&1
 \end{pmatrix},\qquad
 J=\begin{pmatrix}
  0&0&0&-1\\
  0&0&-1&0\\
  0&1&0&0\\
  1&0&0&0
 \end{pmatrix}
\]
for the \(n_a\)-action. The signs of \(J\) are given by
the chosen coordinates and \(n_a^2=a^\vee(-1)\),
since the four weights under \(a^\vee\) are \(-3,-1,1,3\).
Reversing the \(c\)-coordinate in the commutator of \(b,a\)
gives \(A=-1\). Conjugation by \(J\) sends \(L(1)\)
to the action of \(x_{-a}(-1)\). Thus
\(L(1)J L(1)J^{-1}L(1)=J\).
The top three entries of the last column on the left
are \(-D,-B+D,-1-BC-BD\). Equality forces
\(D=B=1,\ C=-2\). The coefficient identities in
\(L(u)L(v)=L(u+v)\) give \(EA=3D=FB\), so
\(E=-3,\ F=3\). These computations use no division
by the small coefficients.

The only unordered pairs in \(N/U_f\) whose commutators
can be nontrivial in \(N\) are \(b,e\) and \(c,d\).
The commutator is alternating and bilinear in the
quotient coordinates because \(U_f\) is central.
The coefficient on \(b,e\) is \(1\). If the coefficient on
\(c,d\) is \(K\), apply \(L(u)\) to the pair \(b,d\).
In the notation of their coordinate axes their images are
\(b-uc+u^2d+u^3e\) and \(d+3ue\).
Preservation of the zero commutator gives \(3u-Ku=0\),
and \(K=3\), valid in every characteristic.
It follows also that \(n_a\) fixes \(U_f\) pointwise,
since it preserves the commutator of \(b,e\), replacing
those axes by \(e,-b\).

It remains to recover the central part of the short-root
action. Denote the coordinates in order by \((b,c,d,e,z)\),
reusing root letters for scalars and denoting the \(f\)-coordinate
by \(z\). When multiplying two such coordinate vectors the
central correction to \(z+z'\) is \(-eb'-3dc'\),
by the two commutators just computed. In the action of
\(x_a(u)\), torus equivariance permits a central term
\(Pu^3v^2\) on \(x_b(v)\), and \(Quv^2\) on \(x_c(v)\).
It permits none on \(x_d(v),x_e(v)\), since the corresponding
\(u\)-powers would be negative. Let \(\varphi\) be the action
of \(x_a(1)\), and let \(J_N\) be the action of \(n_a\) on \(N\).
The product correction and the already computed linear parts
give
\[
 \begin{split}
 \varphi(b,c,d,e,z)&=(b,c-b,d+b-2c,e+b-3c+3d,
                      z+Pb^2+Qc^2-3bc),\\
 J_N(b,c,d,e,z)&=(-e,-d,c,b,z+eb+3dc).
 \end{split}
\]
The identity \(\varphi J_N\varphi J_N^{-1}\varphi=J_N\)
on the unit \(b\)- and \(c\)-axes gives the central
coefficients \(2P+Q-5=0\) and \(9P+5Q-24=0\).
For the first axis the coordinates right after
\(J_N^{-1}\varphi\) are \((1,1,1,-1,P-2)\).
For the second they are \((-3,-2,-1,0,Q-6)\),
which yields the second equation by the same substitution.
The linear system has determinant \(1\), so \(P=1,Q=3\)
in every characteristic. This forces the positive
commutations involving \(a\), including the corrections
from both \(b\) and \(c\) to \(f\). The only other
nontrivial positive commutations are the two already
computed in \(N\). The normalizer scales are known on
all positive and linked negative coordinates by \(J_N\),
the \(n_b\)-action, and the rank-one opposite signs.

In these rank-two systems the signed nonopposite relations are determined
by these positive ones. Indeed a Weyl word carries any
chosen positive chamber to the base chamber, and its action
on coordinates is the product of the simple representative
actions just computed. The opposite-root relation was
proved above in every characteristic. To compare two based
data, use the prescribed map on the torus and match all
these root coordinates in their common order. The lattice
conditions intertwine the torus actions and the opposite
root identities. All signed nonopposite commutations now
match as well. They determine formal multiplication at the
identity by the sorting argument in the integral realization,
which works over the field in question. For the product of
two formal cells each correction of order \(m\) introduces
only higher corrections when moved across a factor. Uniqueness
of the cell coordinates identifies the resulting series with
the actual product. The cell map therefore preserves
multiplication as a formal identity, and hence as a rational
identity. Indeed the completed local ring at the pair of
identity points contains its local ring faithfully, and
multiplication stays in the cell on an open neighborhood
of that pair. The rational homomorphism extends and has
the prescribed central kernel by the cell gluing and Bruhat
argument above. This proves both pinned identification and
central lifting in all rank-two cases, including the orthogonal case above.

\paragraph{Combining the rank-two identifications.}
Start with two split pinned groups having the root and
coroot lattice compatibility above. The bijection identifies
their real root hyperplane arrangements. For each pair of
simple roots, take the centralizer of the largest subtorus
of the common kernel of the two roots. It is smooth connected
reductive by C02 and has the given full maximal torus.
Its roots are exactly the roots in the rational span of the
pair, by the weight description of the centralizer. Its
base is the pair, because a positive root has nonnegative
simple-root coefficients. The rank-two calculation gives a
map of these centralizers that matches the two pins and the
prescribed torus map. For a single simple root use its
rank-one centralizer. The maps agree on each rank-one overlap
with the same simple root, by uniqueness on its open product.

We first compare normalizers. The normalizer is the union
of translates of its torus by the words in the simple
representatives, since their Weyl elements generate by
D04. Conjugation of the torus by each representative is
given by its reflection, and the square of each representative
is its coroot value at \(-1\). For a pair \(a,b\) of simple
roots let \(m\) be the order of \(s_as_b\). The element
\((n_an_b)^m\) belongs to the torus, and the comparison
just constructed in the pair's centralizer matches that
torus value in the two groups.

These relations suffice to define a homomorphism of the
normalizers. Here is the needed chamber fact. Any loop of
chambers, stepping across successive walls, is generated by
crossing a wall and back and by the cyclic walks around
codimension-two faces. Indeed realize the gallery by a
polygonal path crossing walls in their interiors. Fill the
loop by a piecewise-linear disk in the real cocharacter
space. A generic choice of its interior vertices makes
it avoid strata of codimension at least three and meet
codimension-two strata at finitely many transverse interior
points. It may be triangulated so that these points lie
inside triangles and its wall arcs avoid the vertices.
Sliding paths across the disk changes the sequence only by
backtracking at one wall or by going around one of those
points the other way. Locally at such a point the walls
are those of a rank-two arrangement. Starting in any
of its adjacent chambers the cyclic gallery alternates
the two simple wall steps for that chamber, of length
\(2m\). Thus a word in simple Weyl reflections
representing the identity is a product of conjugates of
these rank-one and rank-two relations. The same applies
after removing the common linear subspace when the
arrangement has a central direction.

Given a word in the normalizer, first move its torus
factors to one side using the conjugation relations.
By the chamber fact any two words representing the same
Weyl element differ by products of the square and pair
relations, conjugated by words. Their exact torus
discrepancies agree under the prescribed torus map, by
the rank-one and rank-two comparisons and the compatibility
of the torus actions. Therefore sending each simple
representative to its counterpart and applying the given
torus map defines the claimed homomorphism. It also
defines a scheme map on each component \(n_wT\),
since these components have fixed rational representatives
and the torus maps are morphisms.

For each root \(\gamma\) we next compare its two root
groups by transporting the positive pin from a chamber
having \(\gamma\) as one of its simple positive roots.
Use a representative of the passage from the base
chamber in both groups, matched by the normalizer map.
The choice of representative differs only by a torus
element, whose action matches by the root-character
condition. We check independence of the chamber.
The chambers in question correspond to the open wall
facets of the arrangement cut out on the hyperplane
\(\gamma=0\), with the adjacent chamber chosen on the
\(\gamma>0\) side. A generic polygonal path within that
hyperplane joins any two such facets, crossing the
codimension-one intersections in succession and avoiding
their smaller strata. At a crossing, the two chosen
chambers are incident to the same codimension-two face
of the full arrangement. In the first chamber exactly
two of its simple walls contain that face; after moving
back to the base chamber the chambers around the face
are the translates by the rank-two group of that
simple pair. This follows by going around the face across
its successive walls, each step using one of the two
simple reflections in that moving chamber. The two
pin transports therefore agree by the rank-two comparison
already made for that pair. This proves independence
along the path and hence in general.

The comparisons on root groups now intertwine the
normalizer actions, by transporting the chosen chamber
and its representative along the normalizer element.
The opposite comparisons respect the linked coordinates:
one may choose the neighboring chamber across the
root's own wall, and its representative differs by the
rank-one Weyl word, whose negative scale is the same
on both groups. In particular every Weyl representative
defined by a linked root pair is mapped to its counterpart.

Take now any independent roots \(a,b\), with either sign.
Their centralizer of the common kernel subtorus has
root system \(\Phi\cap(\mathbb Q a+\mathbb Q b)\).
Choose its two simple roots and prescribe as pins on
the target the images under the root comparisons just
constructed. The rank-two theorem applies to these
pins. Its map agrees on every root in this subsystem,
because each is transported from a local simple root
by the local Weyl words, and our global comparisons
already intertwine those words. The full commutation
relation for the original pair thus matches between the
two groups, by applying this local map to it. Opposite
pairs and torus commutations match as before. The
formal sorting and open-cell argument therefore gives a
homomorphism of the two groups over the given field.
It is a central isogeny with the specified torus map,
and an isomorphism if that torus map is an isomorphism.
Together with pinned uniqueness and integral realization,
this proves the stated equivalence of groupoids and gives
the central lifts used in C07 in every characteristic.

For comparison, the general isogeny theorem also treats
root bijections with powers \(q(\alpha)\) in the two
lattice equalities instead of \(1\), where the root-group
map is the corresponding Frobenius power. See
\cite[Theorem 6.1.16 and Section 6.2]{conrad}.
The ordinary central lifts just constructed have all
powers equal to \(1\), including for torus kernels that
are not smooth.

Finally an unpinned isomorphism can be adjusted to match
split Borel pairs. First align the split maximal tori
by rational conjugacy, then align positive systems by a
rational Weyl representative, using
Results~\ref{res:C02} and~\ref{res:C06}.
Any reduced root-datum isomorphism can likewise be
adjusted by Weyl conjugacy of chambers to match bases.
This proves classification without the pinning.
The full lattices remain part of the invariant, as in
the two type-\(A_1\) examples of Result~\ref{res:C07}.
\end{proof}
\chapter{Homogeneous spaces, Bruhat cells, and flags}
\label{ch:flags}

For a closed subgroup scheme \(H\subset G\), \(G/H\) denotes the fppf sheaf associated to \(R\mapsto G(R)/H(R)\). Its points are locally represented by cosets \(gH\), with \(H\) acting on \(G\) from the right. By a parabolic subgroup of a connected reductive group we mean a smooth closed subgroup whose geometric fibre contains a Borel. Properness of \(G/P\) is an equivalent condition by Result~\ref{res:C06}.

\section{Quotients and rational lifting}

\begin{result}{D01}{Representability of a homogeneous quotient}
Let \(G\) be a finite-type affine group scheme over a field \(k\) and \(H\) a closed subgroup scheme. Then \(G/H\) is a quasi-projective \(k\)-scheme and the projection \(q:G\to G/H\) is an fppf right \(H\)-torsor. In particular,
\[
  G\times H \xrightarrow{\ \sim\ } G\times_{G/H}G,
  \qquad (g,h)\longmapsto(g,gh).
\]
Formation of \(G/H\) commutes with extension of the ground field. If \(H\) is smooth then \(q\) is smooth. If \(H\) is normal, the quotient is the affine group quotient of Result~\ref{res:A07}.
\end{result}

\begin{proof}
By the scheme-stabilizer version of Chevalley's theorem (Result~\ref{res:A15}), there is a finite-dimensional representation \(V\) of \(G\) and a line \(L\subset V\) with scheme stabilizer \(H\). For smooth \(G\) we can construct the quotient directly. Write
\(\alpha(g)=gL\) and let \(Z\subset\mathbb P(V)\) be the
scheme-theoretic closure of its image. In a standard affine chart
the inverse image in \(G\) is affine: choose a vector spanning \(L\)
and invert the corresponding coordinate of its translate.
Thus the image ideal on each chart is the kernel of a ring map.
It commutes with every field extension by exactness of tensoring
over \(k\). In particular \(Z_{\bar k}\) is reduced, since every
chart ring of this closure embeds into the ring of an open
subscheme of the smooth \(G_{\bar k}\).

Over \(\bar k\) choose a nonempty affine open \(U\) of the closure
contained in one chart and in the locus of exactly one irreducible
component. It can be chosen principal in that chart. Its ring is a
domain, and the ring of its inverse image under \(\alpha\) contains
it by the definition of the closure. This inverse image is affine:
the extra principal-open condition is another localization of the
chart inverse image. Apply the generic freeness argument in
Result~\ref{res:A07} to that finite-type algebra. On a smaller
nonempty principal open of \(U\), its inverse image is faithfully
flat over it, and takes geometric points onto those of the open.

Translation by any geometric group element preserves the closure:
it preserves the orbit and hence its reduced closure. It carries
the restriction of the orbit map over this open to an isomorphic
restriction over its translate. The union \(O\) of these translates
is open in \(Z_{\bar k}\) and contains precisely the geometric orbit
on geometric points. Indeed, it contains at least one point of
the orbit, all translates remain in the image, and the group is
transitive on its own orbit. The inverse image of \(O\) is all of
\(G_{\bar k}\), since its closed complement has no closed points.
Thus over \(\bar k\) the orbit map is faithfully flat onto \(O\).

For descent, the quasi-compact open \(O\) can be defined over some
finite extension \(k'/k\) by writing it as a finite union of
principal opens in the finitely many charts of \(Z_{\bar k}\) and
retaining their coefficients. That open is exactly the set image
of \(\alpha_{k'}\): set images of scheme morphisms commute with
base change, as nonemptiness of a fibre survives field extension.
The finite surjection \(Z_{k'}\to Z\) is a closed quotient map of
underlying spaces, so the set image of \(\alpha\), whose inverse
image is the descended open, is itself open in \(Z\). The map onto
it is faithfully flat by checking after field extension. In affine
coordinates flatness descends because injectivity of a map after
tensoring can be checked after a faithfully flat field extension.
The map is of finite presentation since both schemes have finite
type over a field. The fibre-product identity below then
identifies the open with the sheaf quotient: a section of the open
lifts fppf locally, and two lifts differ exactly by \(H\).
Being an open of a closed subscheme of projective space, it is
quasi-projective. Compare the smooth orbit theorem in
\cite[Proposition 7.17 and Theorem 7.18]{milne}.

For an arbitrary \(G\), embed it as a closed subgroup of a
smooth group \(M=\GL_n\) by Result~\ref{res:A06}. Both \(M/H\)
and \(M/G\) are represented by the smooth case just proved.
The inverse image of the point \(eG\) under \(M/H\to M/G\)
represents \(G/H\): a locally chosen lift \(m\) of a point
of this fibre satisfies \(mG=eG\), which is exactly the condition
\(m\in G\) checked fppf locally. Equality of these sheaf maps
follows by descent. Since \(eG\) is a closed \(k\)-point and
\(M/H\) is quasi-projective, this fibre is quasi-projective.

For any represented quotient, equality of the two image
sections \(q(g)\) and \(q(g')\) is equivalent to
\(g^{-1}g'\in H\): membership can be checked fppf locally
since \(H\) is closed. This proves the displayed isomorphism
with inverse \((g,g')\mapsto(g,g^{-1}g')\).
Fppf local sections are part of the quotient-sheaf definition.
They give local products with the finite-type group scheme \(H\),
so faithful flatness and finite presentation descend. If \(H\)
is smooth, smoothness follows by the same local products.
The quotient sheaf commutes with field extensions because fppf
covers of test schemes over the extension are covers in either
description.

For the affineness assertion, suppose \(H\) is normal, and
first take the field algebraically closed. Use the pair
\((V,L)\) from Result~\ref{res:A15}, and let \(\chi\) be
the character through which \(H\) acts on \(L\). In characteristic
\(p>0\) choose \(m=p^r\) so that the image \(I\) of the iterated
relative Frobenius \(G\to G^{(p^r)}\) is smooth. This is possible:
the ring of the image is the subring of \(m\)-th powers, which is
reduced when \(m\) kills the nilradical of the noetherian
coordinate algebra. Use Result~\ref{res:A09}. In characteristic
zero take \(m=1\) and \(I=G\).
The image \(J\) of \(H\) in \(I\) is normal since \(G\to I\)
is faithfully flat by Result~\ref{res:A07}; both image
maps can be tested by lifting the points fppf locally.

In the corresponding Frobenius twist of \(V\) (or \(V\)
when \(m=1\)), the twist of \(L\) is a \(J\)-eigenline
whose \(H\)-character is \(\chi^m\).
Indeed on pulling this representation back to \(G\),
the matrix coefficients are the \(m\)-th powers of those
of the original representation in the twisted basis, and the
line and its character descend from \(H\) to \(J\) by faithful
flatness. Let \(W\) be the sum of all the \(J\)-eigenlines.
Normality of \(J\) makes this sum invariant under \(I(k)\),
so the density of \(I(k)\) in the smooth \(I\) makes it
an \(I\)-subrepresentation. As a \(J\)-representation it is
a sum of character spaces, a direct sum for distinct characters
since distinct group-like elements of a coalgebra are linearly
independent. Thus \(W^\vee\) has an \(H\)-eigenline \(L'\)
with character \(\chi^{-m}\), obtained by a functional on the
corresponding character space and extended by zero on the others.

The line \(L^{\otimes m}\otimes L'\) in
\(V^{\otimes m}\otimes W^\vee\) is fixed pointwise by \(H\).
Any element of \(G\) stabilizing this tensor line must stabilize
\(L\), on every test algebra. To see this choose bases \(e_i,f_j\)
whose first vectors span \(L,L'\). If \(g e_1=\sum a_i e_i\)
and \(g f_1=\sum b_j f_j\), the coefficient \(a_1^m b_1\)
is a unit, so \(a_1,b_1\) are units. The mixed coefficients
\(a_i a_1^{m-1}b_1\) for \(i>1\) vanish, forcing \(a_i=0\).
Consequently the action of \(G\) on the \(H\)-invariants of
this tensor representation has kernel exactly \(H\).
The invariant subspace is \(G\)-stable by normality on
test algebras, and commutes with scalar extension by
Result~\ref{res:A05}. The kernel contains \(H\) by its definition
and is contained in \(H\) since it fixes the displayed tensor
line. All the coefficients of this representation and the
kernel equality descend to some finite extension \(K/k\)
by finite presentation and faithful flatness. Use the induced
map from \(G\) to the Weil restriction of its target general
linear group (Result~\ref{res:B04}). Its kernel is \(H\),
as membership can be tested after \(R\to R\otimes_k K\).
Result~\ref{res:A07} now realizes \(G/H\) as its affine scheme
image, proving the normal case. Compare
\cite[Propositions 5.15--5.18 and Section 7.e]{milne}.
\end{proof}

\begin{result}{D02}{Cosets, points, and the obstruction torsor}
There is an injection \(G(k)/H(k)\to(G/H)(k)\). For \(x\in(G/H)(k)\), let \(\delta(x)\) be the class of the \(H\)-torsor \(q^{-1}(x)\) in \(H^1_{\mathrm{fppf}}(k,H)\). The image of the injection is exactly the neutral fibre of \(\delta\), and the image of \(\delta\) is the neutral fibre of the map
\[
 H^1_{\mathrm{fppf}}(k,H)\longrightarrow H^1_{\mathrm{fppf}}(k,G).
\]
This is a sequence of pointed sets. Two quotient points have the same torsor class exactly when they lie in the same \(G(k)\)-orbit.
\end{result}

\begin{proof}
The fibre product identity in Result~\ref{res:D01} proves the injection. A torsor under \(H\) is trivial if and only if it has a rational point: that point and its translates identify the torsor with \(H\). For the next assertion, put \(T=q^{-1}(x)\). Extending the structure group gives \(T\times^H G\), the quotient of \(T\times G\) under \((t,g)\cdot h=(th,h^{-1}g)\). The multiplication map \([t,g]\mapsto tg\) is a trivialization as a right \(G\)-torsor, as can be seen after any fppf cover on which \(T\) has a section. Conversely, if an \(H\)-torsor \(T\) becomes trivial after this extension of structure group, choose a right \(G\)-equivariant identification \(T\times^H G\simeq G\). The induced map \(T\to G\) is \(H\)-equivariant. Its composite with \(q\) descends to a point \(x\) over \(k\). The induced map \(T\to q^{-1}(x)\) is an equivariant map of \(H\)-torsors, and therefore an isomorphism by local trivialization. Finally an \(H\)-equivariant isomorphism \(f:q^{-1}(x)\to q^{-1}(y)\) makes \(f(t)t^{-1}\) invariant under the right \(H\)-action. It descends to a rational \(g\in G(k)\) with \(gx=y\). Left translation by any such \(g\) supplies the converse. This is the torsor version of the homogeneous-space exact sequence \cite[Section 5.b and Appendix B]{milne}, \cite[Example A.1.2]{conradfeng}.
\end{proof}

For smooth \(H\), Galois cohomology computes these torsor classes by passage to a separable closure. With \(k=\FF_p(t)\), the quotient \(\Gm/\mu_p\simeq\Gm\) is given by \(p\)-th power, since adjoining a \(p\)-th root of an invertible element is finite faithfully flat. The point \(t\) has no lift over \(k\), as its \(t\)-adic valuation is \(1\) and the valuation of every \(p\)-th power is a multiple of \(p\). Its fibre is a nontrivial fppf \(\mu_p\)-torsor with no point over a separable closure.

\section{The split open cell and the Bruhat decomposition}

Fix a split connected reductive pair \((G,T)\) and a Borel \(B\supset T\). Write \(\Phi^+\) for the positive roots determined by \(B\), \(U^+=R_u(B)\), and \(U^-\) for the unipotent radical of the opposite Borel. Write \(W=W(G,T)\) and let \(w_0\) be its longest element relative to the simple reflections of \(B\).

\begin{result}{D03}{The open cell}
Multiplication gives an open immersion with dense image
\[
  U^- \times T\times U^+\longrightarrow G.
\]
If \(n_{w_0}\in N_G(T)(k)\) represents \(w_0\), left translation by \(n_{w_0}\) identifies the image with the open double orbit \(Bn_{w_0}B\). In \(\GL_n\) with \(B\) upper triangular, the first open cell consists of the invertible matrices whose leading principal minors are all invertible.
\end{result}

\begin{proof}
The dynamic open immersion proved in Result~\ref{res:C02},
for a cocharacter regular in the positive chamber, gives this
product, using the positive and negative root products of
Result~\ref{res:C05}. Its image is dense because the smooth
geometrically connected \(G\) is geometrically irreducible.
Conjugation by \(n_{w_0}\) interchanges \(U^-\) and \(U^+\)
and normalizes \(T\), so left translation gives
\(U^+n_{w_0}B=Bn_{w_0}B\), absorbing the torus across \(n_{w_0}\).
The translated subscheme is the fppf double orbit, since every
element in the indicated product lies in that orbit and every
pair of factors of \(B\) reduces to this product.

For \(\GL_n\), this is the factorization \(LDU\) with lower and
upper unitriangular factors and an invertible diagonal factor.
The leading \(j\)-th principal minor is the product of the first
\(j\) diagonal entries. Conversely elimination with the first
pivot splits off the invertible top-left entry and leaves the
Schur complement. Its leading minors are the corresponding
successive minors of the original matrix divided by that pivot.
Induction gives the factors. Each multiplier is forced by
the pivot and its row or column, so the formulas give the inverse
morphism over the locus where those minors are units.
Compare \cite[Theorem 21.84 and Example 21.87]{milne}.
\end{proof}

\begin{result}{D04}{Split Bruhat cells}
Choose \(n_w\in N_G(T)(k)\) above each \(w\in W\) and set \(U_w=U^+\cap n_wU^-n_w^{-1}\). The group \(U_w\) is an ordered product of the root groups for \(\Phi^+\cap w\Phi^-\). The locally closed double orbit \(C(w)=Bn_wB\) is isomorphic to \(U_w\times B\) by \((u,b)\mapsto un_wb\). The orbits give a disjoint locally closed stratification of \(G\). On \(G/B\) the corresponding cells are isomorphic to \(\mathbb A^{\ell(w)}_k\), and \(B(k)n_wB(k)=C(w)(k)\).
\end{result}

\begin{proof}
Let \(\lambda\) be regular in the positive chamber and put
\(\mu=w\lambda\). In the ordered coordinates of \(U^+\) from
Result~\ref{res:C05}, conjugation by \(\mu(t)\) scales the
\(\gamma\)-coordinate by \(t^{\langle\gamma,\mu\rangle}\).
The orbit stays in the closed \(U^+\), so its limit can be
tested in these coordinates on every test algebra. The elements
having limit \(1\) for \(\mu\), respectively \(-\mu\),
are exactly the products supported on
\(\Phi^+\cap w\Phi^+\), respectively \(\Phi^+\cap w\Phi^-\).
Write \(U'_w\) for the first intersection.
Choosing the order with the negative-\(\mu\) roots first gives
\(U^+=U_w U'_w\) as a product of schemes.
The same limit test in \(B=T\ltimes U^+\) gives
\[
 B\cap n_w B n_w^{-1}=T\ltimes U'_w,\qquad
 U_w\cap n_w B n_w^{-1}=1.
\]
Thus \(U_w\) represents the \(B\)-homogeneous space of the
coset \(n_wB\): the factorization of \(B\) by \(U_w\) and
the stabilizer gives inverse maps on test algebras.
This homogeneous space is its locally closed orbit in \(G/B\):
use the projective orbit embedding of \(G/B\) from
Result~\ref{res:C02} and apply the smooth-group orbit argument
of Result~\ref{res:D01} to the \(B\)-action in that projective
space. Pulling the orbit back along \(G\to G/B\) gives
\(C(w)\). Its inverse maps are supplied by the section
\(u\mapsto un_w\) above \(U_w\) and the \(B\)-torsor structure
from Result~\ref{res:D01}. Explicitly the isomorphism is
\((u,b)\mapsto un_wb\). Its restriction to \(k\)-points
gives the asserted rational double orbit.

The torus action on the orbit \(U_w\) scales its root
coordinates by nontrivial characters, so its only geometric
fixed point is \(n_wB\). If two \(B\)-orbits meet they coincide,
hence have the same fixed point. The representatives are then
equal modulo \(N_G(T)\cap B=T\), by the torus quotient
of \(B\) as in Result~\ref{res:C05}. Thus the orbits are disjoint.

For coverage, take a simple root \(\alpha\) and its
reflection \(s=s_\alpha\). Choose an integral cocharacter
\(\eta\) in the relative interior of this wall, not on any
other root hyperplane. Then the only vanishing roots are
\(\pm\alpha\), and all other positive roots pair positively.
The dynamic Levi \(Z_G(\eta)\) is \(H_\alpha=Z_G(S)\) of
Result~\ref{res:C04}: the inclusion follows because
\(\eta\) takes values in \(S\), and the two centralizers are
smooth connected with equal tangent weight spaces. The
positive-coordinate argument of Result~\ref{res:C05} applies
also to \(\eta\), giving \(U_G(\eta)\) as a product of all
the remaining positive root groups, hence inside \(B\).
The quotient for \(P_s=P_G(\eta)\) is therefore
\(P_s/B=H_\alpha/(B\cap H_\alpha)\).
The intersection is its positive Borel, since it is
\(Z_B(S)\) with the torus and positive \(\alpha\)-line as
tangent weights. Result~\ref{res:C04} identifies that flag
variety with \(\mathbb P^1\), with its positive root group
acting by upper transvections. On geometric points the two
orbits are the fixed point and the affine-line orbit, so
\[
 P_s=B\ \cup\ Bn_sB,\qquad Bn_sB=U_\alpha n_sB
\]
on these points. If \(v(\alpha)>0\), moving \(U_\alpha\)
past \(n_v\) into \(B\) gives \(C(v)Bn_sB=C(vs)\).
If \(v(\alpha)<0\), apply the positive case to \(vs\):
\(C(vs)P_s=C(vs)\cup C(v)\), whence \(C(v)P_s\) is
contained in this same union. Both cells occur because
\(1,n_s\in P_s\). In both cases,
\[
 C(v)P_s=C(v)\cup C(vs)
\]
on geometric points. This calculation also gives the
rank-one product rule for use in Result~\ref{res:D06}.

The simple reflections generate the Weyl group: take a
gallery of adjacent chambers starting from the chosen chamber,
and at each step the reflection in the crossed wall is the
conjugate of the corresponding simple reflection by the
element already reached. Every root hyperplane is a wall
of some chamber, by taking a point on it avoiding the other
hyperplanes. Thus every root is conjugate up to sign to a
simple root, and the sign can be changed by its reflection.
Consequently the \(P_s(\bar k)\) and \(B(\bar k)\) generate
all the root subgroups under conjugation by simple Weyl
representatives and include \(T(\bar k)\). They contain
the points in the dense open cell of Result~\ref{res:D03}.
A subgroup containing this open is the whole \(G(\bar k)\),
since it meets every translate of the open. The union of
the cells, stable under right multiplication by all these
generators and containing the identity, therefore covers
\(G(\bar k)\). The locally closed subschemes constructed over
\(k\) give the resulting stratification over \(k\).

Finally the number of root coordinates in \(U_w\) is the
number of hyperplanes separating the base chamber from its
\(w\)-translate: a positive root contributes exactly when it
becomes negative under \(w^{-1}\). Each gallery crosses all
separating hyperplanes, and a generic straight segment between
interior points crosses just those, one at a time. Since a
gallery of \(r\) steps is a word of \(r\) simple reflections
and conversely, the number is \(\ell(w)\).
Each root group is a \(k\)-affine line, proving the cell
description. Compare \cite[Proposition 10.3.3 and Appendix N.2]{conradfeng}.
\end{proof}

\section{Parabolics, Schubert varieties, and partial flags}

\begin{result}{D05}{Standard parabolics and their Levi subgroups}
Let \(\Delta\) be the simple roots and \(I\subseteq\Delta\). Set \(\Phi_I=\Phi\cap\ZZ I\), and let \(T_I\) be the maximal subtorus in \(\bigcap_{\alpha\in I}\ker\alpha\). There is a unique standard parabolic \(P_I\) containing \(B\) with Weyl indices \(W_I=\langle s_\alpha:\alpha\in I\rangle\). Its Levi \(L_I=Z_G(T_I)\) has absolute roots \(\Phi_I\), and
\[
    P_I = R_u(P_I)\rtimes L_I,\qquad
    R_u(P_I)=\prod_{\alpha\in\Phi^+\setminus\Phi_I} U_\alpha ,
\]
where the product denotes an ordered scheme parametrization. Every geometric parabolic is its own scheme-theoretic normalizer.
\end{result}

\begin{proof}
Fix a cocharacter \(\lambda\) pairing to zero on \(I\) and
positively on \(\Delta\setminus I\). It exists over \(\QQ\)
by the independence of the simple roots proved in
Result~\ref{res:C05}, and can be multiplied to be integral.
Put \(P_I=P_G(\lambda)\). Its dynamic construction and
smoothness are proved in Result~\ref{res:C02}.
The torus \(T_I\) contains \(\lambda(\Gm)\).
The split connected centralizers \(Z_G(T_I)\subseteq Z_G(\lambda)\)
have the same tangent weight spaces: the roots trivial on \(T_I\)
are exactly the roots in the rational span of \(I\), and the
positive simple-coordinate description in Result~\ref{res:C05}
shows these are precisely the roots pairing to zero with \(\lambda\).
Thus the centralizers coincide, and they are reductive by
Result~\ref{res:C02}. The dynamic contracting group is
smooth connected normal unipotent with this reductive quotient,
so it is \(R_u(P_I)\): it is contained in the unipotent
radical, and the image of that radical would otherwise be a
nontrivial smooth connected normal unipotent subgroup of the
reductive quotient. The root-product parametrization for the
strictly positive weights, including cocharacters on walls,
was proved in Result~\ref{res:C05}.
Every positive root group thus belongs to \(P_I\), either
in its Levi or in its radical. The group contains \(B\) and
is parabolic in the stated convention.

The Borel \(B\cap L_I\) of \(L_I\) is given by the
torus-centralizer result in Result~\ref{res:C02}.
Its positive roots have simple basis \(I\), by the
simple-coordinate description of Result~\ref{res:C05}.
Bruhat decomposition of \(L_I\) by Result~\ref{res:D04}
therefore gives indices \(W_I\).
Passing from this Levi to \(P_I\) adds the radical contained
in \(B\), since \(B=(B\cap L_I)\ltimes R_u(P_I)\).
Thus \(P_I\) is the union of \(C(w)\) for \(w\in W_I\).

Conversely let \(P\) be any geometric parabolic containing \(B\).
It is a closed subgroup and a union of Bruhat cells.
Let \(I\) be the simple roots with \(n_{s_\alpha}\in P\).
If \(C(w)\subset P\), its closure is also in \(P\). By the
subword proof in Result~\ref{res:D06}, every simple letter
in a reduced expression for \(w\) belongs to \(I\). This use
of Result~\ref{res:D06} only needs the rank-one calculation
of Result~\ref{res:D04}. Thus \(w\in W_I\). Conversely every
product of these simple reflections admits a representative
in \(P\) by the subgroup property and \(T\subset P\).
This gives equality of \(P\) and \(P_I\) on geometric
points and hence as smooth closed subgroups. The same
comparison gives uniqueness over \(k\) by field descent.

To check its normalizer geometrically, use the Bruhat expression
\(g=b_1n_wb_2\) for a point of the normalizer.
Since \(b_1,b_2\in P_I\), the element \(n_w\) normalizes
\(P_I\) and preserves its tangent root set, namely the
roots with nonnegative pairing with \(\lambda\).
It also preserves their negatives, so it preserves the
full sign pattern. Hence the \(w\)-translate of the base
chamber has the same face, the one whose relative interior
contains \(\lambda\), in its closure. A small ball at that
point meets no hyperplanes except those through the face.
A gallery within the ball joins the two chambers by
crossing just those hyperplanes, one at a time.
Starting from the base chamber, each step is a translate
of a simple wall through \(\lambda\) under the previously
obtained Weyl element. Inductively that element belongs to
\(W_I\) and fixes \(\lambda\), so each next simple letter
lies in \(I\). By the free chamber action from
Result~\ref{res:C05}, \(w\in W_I\), and \(g\in P_I\).
The tangent space to the scheme normalizer is contained
in \(\Lie(P_I)\) by the same first-order normalizer argument
as in Result~\ref{res:C05}: the quotient of Lie algebras has
no zero \(T\)-weights. The resulting equality of local
and tangent dimensions implies smoothness. Agreement on
geometric points is therefore equality as schemes, and
the assertion for every geometric parabolic follows by
Borel conjugacy. Compare \cite[Section 21.h]{milne}
and \cite[Corollary 5.2.8]{conrad}.
\end{proof}

\begin{result}{D06}{Schubert closures and the subword order}
Let \(C_w=C(w)/B\subset G/B\) in the split case. The closure of \(C_w\) is the union of \(C_v\) for \(v\le w\), where \(v\le w\) means that \(v\) is represented by a subword of a reduced simple-reflection expression for \(w\). Equivalently the expression for \(v\) can be required to be reduced. The closure formula for \(C(w)\subset G\) uses the same order.
\end{result}

\begin{proof}
Write \(w=s_1\cdots s_m\) reduced. Form the iterated quotient
\[
  Z=P_{s_1}\times^B P_{s_2}\times^B\cdots\times^B P_{s_m}/B
\]
by coupling each right \(B\)-action with the next left \(B\)-action and taking the right \(B\)-quotient at the end. Equivalently, start with the flag \(B\) and choose a sequence of flags whose consecutive entries have the same image in \(G/P_{s_i}\) at step \(i\). This realizes the iterated quotient as a scheme by fibre products. Each step is a \(\mathbb P^1\)-bundle in the fppf topology, since \(G/B\to G/P_{s_i}\) becomes the product with \(P_{s_i}/B\) after the \(P_{s_i}\)-torsor trivializes. Thus the successive projections are proper, smooth and open with geometrically irreducible fibres of dimension one. It follows by induction that \(Z\) is proper and geometrically irreducible of dimension \(m\). The multiplication map \(Z\to G/B\) takes the last flag and is proper since \(Z\) is proper and \(G/B\) separated.

By the rank-one product rule from Result~\ref{res:D04}, the geometric image is precisely the union of the cells indexed by products of subwords of \(s_1,\ldots,s_m\). It contains \(C_w\), which has dimension \(m\). The irreducible image is closed and of dimension at most \(m\), so it equals the closure of that cell. We give the deletion step for the equivalence of the two
subword descriptions. A reduced gallery crosses exactly
the separating hyperplanes, each once, by the length
calculation in Result~\ref{res:D04}. If appending a simple
letter \(s\) to a reduced word \(u=s_1\cdots s_r\) shortens
its length, the hyperplane \(uH_s\) has already been crossed.
Suppose it was crossed at the \(i\)-th step. Its reflection
is both \(usu^{-1}\) and \(u_{i-1}s_i u_{i-1}^{-1}\), where
\(u_{i-1}=s_1\cdots s_{i-1}\). Multiplying this identity by
\(u\) gives \(us=s_1\cdots\widehat{s_i}\cdots s_r\).
In any word which is not reduced, take the first prefix
which fails to be reduced. That prefix has just appended
a letter across a hyperplane already crossed by the preceding
reduced gallery; delete the two letters as above. Repeating
produces a reduced subword with the same product.
This proves the equivalence. The geometric closure is
intrinsic, so the criterion is independent of the chosen
reduced expression. Compare \cite[Sections 8.5.4--8.5.6]{springer}.

The geometric image determines the closed support over the split field, with its reduced variety structure. Finally \(G\to G/B\) is smooth and surjective, hence open. For an open map, inverse image commutes with topological closure, because an open neighbourhood of a point maps to a neighbourhood of its image. Apply this to \(C_w\) to obtain the formula for \(C(w)\).
\end{proof}

\begin{result}{D07}{Partial flag varieties and their cells}
For any parabolic \(P\) of a connected reductive group, \(G/P\) is smooth and projective. In the split case \(P=P_I\), the \(B\)-orbits on \(G/P_I\) are indexed by the unique shortest representatives \(W^I\) of the cosets \(wW_I\) in \(W/W_I\). The orbit of \(n_wP_I\) for \(w\in W^I\) is \(\mathbb A_k^{\ell(w)}\). Its closure is the union of the orbits indexed by \(v\in W^I\) with \(v\le w\).
\end{result}

\begin{proof}
The quotient exists and is quasi-projective by
Result~\ref{res:D01}, and is proper by Result~\ref{res:C06}.
A proper immersion into projective space is closed, so it is
projective. Smoothness descends from smooth \(G\) under its
smooth surjective quotient map.

For \(P=P_I\), the rank-one rule of Result~\ref{res:D04}
and the cell decomposition of \(P_I\) in Result~\ref{res:D05}
show that multiplying \(C(w)\) on the right by \(P_I\)
gives the union of cells indexed by \(wW_I\).
Indeed, successively multiply by the \(P_s\) for words in
the simple letters of \(I\). These subgroups generate
\(P_I(\bar k)\) together with \(B(\bar k)\), because
the products of their double-orbit representatives cover its
Bruhat cells. This gives exactly those indices in the fibre of
a partial-flag orbit.

For the length assertions, if \(w\) minimizes length in its
coset, it sends every simple root of \(I\) to a positive root.
Otherwise right multiplication by that simple reflection
would cross a separating hyperplane in reverse, shortening
the length by one, by the gallery distance calculation
in Result~\ref{res:D04}. By positive simple coordinates in
Result~\ref{res:C05} it follows that \(w(\Phi_I^+)\subset
\Phi^+\). An element \(u\in W_I\) preserves the sign of
every root outside \(\Phi_I\): each generator preserves
its coordinates outside \(I\), and at least one of these
is a nonzero coordinate determining the sign. Thus the
hyperplanes separating the base chamber from its \(u\)-translate
are all from \(\Phi_I\). Their translates by \(w\) separate
\(w\) from \(wu\), and they do not separate the base chamber
from its \(w\)-translate, since \(w(\Phi_I^+)\subset\Phi^+\).
The two sets of separating hyperplanes are disjoint, and
their union is the separating set for \(wu\).
Consequently \(\ell(wu)=\ell(w)+\ell(u)\). This proves
uniqueness of the minimal representative as well.

Fix such a minimal \(w\), and choose the cocharacter
\(\lambda\) of Result~\ref{res:D05} representing \(P_I\).
Testing limits of the ordered coordinates of \(B\) along
\(w\lambda\), as in Result~\ref{res:D04}, gives the scheme
intersection \(B\cap n_wP_In_w^{-1}\) with root factors
\(\Phi^+\cap w(\Phi^+\cup\Phi_I)\) and torus \(T\).
Minimality makes this root set exactly
\(\Phi^+\cap w\Phi^+\): the extra negative Levi roots
remain negative under \(w\). This is the same stabilizer
as for \(n_wB\). Both orbit schemes are their homogeneous
spaces by the smooth orbit argument of Result~\ref{res:D01}
in the projective line embeddings of the flag varieties
from Result~\ref{res:A15}. Thus projection identifies the
full-flag cell of \(w\) with this partial-flag cell,
isomorphic to \(\mathbb A^{\ell(w)}_k\).

The projection of full to partial flags is proper,
as the source is proper and the target separated.
It maps the closure of that full-flag cell onto the
closure of its image, by closedness and continuity.
If \(v\le w\) occurs in the full-flag closure, then its
minimal representative \(\bar v\) satisfies \(\bar v\le v\):
write \(v=\bar v u\) with \(u\in W_I\) and concatenate
reduced words using length additivity. The reduced
subword criterion in Result~\ref{res:D06} is transitive
by taking a reduced subword for \(v\) and then one for
\(\bar v\) in it. It follows that \(\bar v\le w\).
Conversely every minimal \(v\le w\) belongs to the full-flag
closure and supplies its partial-flag cell. This gives the
asserted closure description.
\end{proof}

For \(\GL_n\), the standard \(P_I\) stabilizes a flag of selected dimensions in the coordinate flag \(E_j=k^j\subset k^n\). The quotient represents such flags in families, since a basis adapted to the successive subbundles exists locally. For the Grassmannian of \(d\)-planes, the indices are increasing pivot positions \(1\le i_1<\cdots<i_d\le n\). The cell has
\[
  \dim(L\cap E_j)=\#\{r:i_r\le j\}\qquad(1\le j\le n).
\]
Indeed a reduced echelon basis has \(r\)-th vector equal to \(e_{i_r}\) plus non-pivot terms in positions less than \(i_r\). Prescribe those columns in an upper triangular unipotent matrix. There are \(\sum_r(i_r-r)\) free entries. On cell closure the equalities become lower bounds on \(\dim(L\cap E_j)\), closed determinantal conditions on \(L\to k^n/E_j\). The bounds allow exactly pivot tuples componentwise at most \((i_1,\ldots,i_d)\). This is the Bruhat order on the Grassmannian representatives, as follows from its rank criterion \cite[Section 8.5]{springer}. One can get from the larger tuple to any allowed smaller one by moving pivots left past unoccupied positions, each an adjacent-reflection shortening. These are the rank conditions on the reduced closure.

\section{Relative structure over an arbitrary field}

\begin{result}{D08}{Rational parabolics, Levis, and flags}
Let \(G\) be connected reductive over any field \(k\). Every \(k\)-parabolic \(P\) has the form \(P_G(\lambda)\) for a \(k\)-cocharacter \(\lambda\). Its unipotent radical is \(k\)-split, and it has a \(k\)-Levi. Any two \(k\)-Levis are conjugate by a unique element of \(R_u(P)(k)\). One has \((G/P)(k)=G(k)/P(k)\). Minimal \(k\)-parabolics are conjugate under \(G(k)\), and maximal \(k\)-split tori are conjugate under \(G(k)\). The quotient \(G/P\) parametrizes parabolics of the geometric type of \(P\).
\end{result}

\begin{proof}
Choose a maximal \(k\)-torus \(T\) of \(P\) by
Result~\ref{res:C02}. It is geometrically maximal in \(G\)
since \(P_{\bar k}\) contains a Borel. Over a finite Galois
splitting field for \(T\), Result~\ref{res:D05} gives the cone
of cocharacters of \(T\) determining \(P\): the pairings are zero
on the roots of the Levi and positive on those of the unipotent
radical. Equivalently the roots are tested by membership of
their two root groups in \(P\), so the cone is Galois-stable.
An integral point of this cone averaged by summing its Galois
translates gives a \(k\)-cocharacter by Result~\ref{res:B02}.
Thus \(P=P_G(\lambda)\). The dynamic construction in
Result~\ref{res:C02} gives \(P=U\rtimes L\), where
\(U=U_G(\lambda)\) and \(L=Z_G(\lambda)\) is reductive.
Consequently \(U\) is the geometric unipotent radical,
by the quotient argument of Result~\ref{res:D05}.

In fact \(U\) is split over \(k\). The coordinate grading
from the contracting action is positive away from constants.
Lift a homogeneous basis of the cotangent space at the
identity. The chosen functions generate the algebra by
degree induction, since the error term is a sum of products
of positive functions of smaller degree. They are algebraically
independent by smoothness, dimension, and geometric
irreducibility of \(U\). The coproduct of a generator of
degree \(m\) is its two linear terms plus products with
both factor degrees strictly below \(m\), by the counit
identities and the grading. Thus the subalgebras generated by
all variables up to any occurring degree are Hopf subalgebras:
the antipode preserves degrees and sends the augmentation ideal
to itself. Successive kernels of their quotient maps give a
filtration, whose successive quotients are vector groups on the
variables of one degree. Indeed the coproduct on these
variables becomes primitive after the lower variables are set
to zero. The quotients and kernels here are the scheme quotients
of Result~\ref{res:A07}; the kernel at each stage has polynomial
coordinates on the remaining variables. Refining the vector
groups by coordinate flags gives the \(\Ga\)-series.

For conjugacy of Levis work first geometrically. Each Levi
contains a maximal torus of \(P\), since a torus of the
standard Levi has the full geometric rank of \(P\) and all
Levis map isomorphically to \(P/U\).
Conjugate a torus of the other Levi into \(L\) by an
element of \(P\), using Result~\ref{res:C02}.
Writing that element in \(LU\) shows that this conjugacy
into \(L\) can be achieved by \(U\) alone.
A Levi containing a given maximal torus \(T\) of \(L\)
is unique: the quotient isomorphism identifies its \(T\)-weights
with those of \(L\), and the corresponding root groups in
\(G\) are unique by Result~\ref{res:C04}.
Together with \(T\) they generate the Levi by the open
product in Result~\ref{res:C05}.
A point of \(U\) normalizing a Levi centralizes it on every
test algebra, since it acts trivially on \(P/U\).
As \(L\) contains \(\lambda(\Gm)\), its centralizer in \(U\)
is trivial by the contraction equations. The transporter in
\(U\), represented as in Result~\ref{res:A08}, is therefore
a single reduced point geometrically. Faithfully flat descent
makes it a \(k\)-point, proving uniqueness and existence
over \(k\).

We next justify rational density for reductive groups over
infinite fields. In a smooth connected subgroup \(H\) of
\(G_{\bar k}\) containing a geometric maximal torus \(T\),
suppose its Lie algebra has a symmetric subset of the
root lines and \(H\) contains the corresponding root groups.
Then \(H\) is reductive. For a regular cocharacter of \(T\),
the centralizer in \(H\) is \(T\), and the positive and negative
limit subgroups have equally many root lines in their tangent
spaces. Any Borel of \(H\) containing \(T\) has at most one
line from each opposite pair: applying the smooth dynamic
construction inside that Borel and the relevant rank-one
centralizer gives the root group whenever it gives its tangent
line. Both opposite root groups and \(T\) fill an open
subset of the corresponding rank-one group of \(G\),
which is nonsolvable by Result~\ref{res:C04}.
It follows by dimension, as in Result~\ref{res:C05}, that
the two opposite limit parabolics of \(H\) are Borels of \(H\).
Its unipotent radical lies in both and maps trivially to \(T\)
in both, so lies in the intersection of the two opposite
contracting groups, which is trivial.

Apply this to the smooth connected centralizer of a matrix
semisimple Lie element \(Y\) of \(G\), using the centralizer
smoothness proof of Result~\ref{res:C02}. To see it contains a
geometric maximal torus of \(G\), first choose a maximal torus
\(T\) of its own. The element \(Y\) belongs to the Lie
algebra of \(Z_H(T)\), since \(H\) centralizes it.
As \(Z_H(T)=T\times U_0\) over \(\bar k\), the matrix
Jordan decomposition puts \(Y\) in \(\Lie(T)\).
Any larger geometric torus of \(G\) containing \(T\)
would centralize \(Y\) and belong to \(H\), contradicting
maximality there. Now the tangent roots of \(H\) are exactly
those with \(d\alpha(Y)=0\). They come in opposite pairs,
and their root groups centralize \(Y\): differentiate the
identity for conjugation of \(x_\alpha(a)\) by \(T\) in the
direction \(Y\). They belong to the identity component since
they are connected. The preceding symmetric-root argument
therefore proves that \(H\) is reductive.

For generation by \(k\)-tori, induct on the dimension of a
connected reductive \(G\), with the torus case immediate.
If there are geometric noncentral semisimple Lie elements,
take the nonempty open in \(\Lie(G)\) from the proof of
Result~\ref{res:C02}. For each \(X\) in its rational points,
the rational semisimple part in characteristic zero, or its
large \(p\)-power in characteristic \(p\), gives \(Y\)
commuting with \(X\). So \(X\) belongs to the Lie algebra
of the proper reductive centralizer component of \(Y\).
These Lie algebras span \(\Lie(G)\) because the rational
points of that affine open are dense. Each centralizer
component is generated by its \(k\)-tori by induction.
Consequently the \(k\)-tori generate \(G\): take finitely
many of the centralizers spanning the Lie algebra and the
closure of their repeated products. The nested irreducible
product closures stabilize by dimension and are stable
under products and inversion, using density and swapping
the order or including the inverse products. They define
a smooth subgroup whose tangent space contains the chosen
sum. Equality of dimensions makes the subgroup all of \(G\).

In the remaining characteristic-\(p\) case take the central
quotients of Result~\ref{res:C02}. Those quotients are reductive:
the reduced identity component of the geometric inverse image
of a smooth connected normal unipotent subgroup is smooth
connected normal and has no nontrivial torus (its torus image
would be trivial and the kernel is finite). The geometric
rank-zero argument in Result~\ref{res:C02} makes it unipotent,
so it is trivial. Finiteness of the quotient forces the original
unipotent subgroup to be trivial. The quotients eventually
reduce to the preceding generation argument or to a torus.
The \(k\)-tori generating a quotient lift to \(k\)-tori
by Result~\ref{res:C02}. The subgroup they generate has
full dimension, since the quotient is finite, so is \(G\).
This completes the induction.

Finally the \(k\)-points of any \(k\)-torus \(T'\) are
geometrically dense. Take a finite Galois splitting field
\(K\). The norm from \(\operatorname{Res}_{K/k}(T'_K)\)
to \(T'\) is the multiplication of conjugate factors
after splitting; it is compatible with their Galois
permutation and hence descends by Result~\ref{res:F02}.
It is surjective after splitting. The Weil restriction
of this split \(K\)-torus is the open subset of affine
space where the coordinate multiplication determinants
are nonzero (Result~\ref{res:B04}). Its \(k\)-points
are geometrically dense over infinite \(k\). Indeed
a polynomial over \(\bar k\) vanishing on them is zero
by expanding its finitely many coefficients in a
\(k\)-linearly independent family and using polynomial
density over \(k\). The same applies to the resulting
torus image. Thus the geometric closure of \(G(k)\)
contains each \(k\)-torus and hence equals \(G_{\bar k}\),
as just proved. Compare \cite[Section 17.j]{milne}.

The open product \(U_G(-\lambda)\times P\subset G\)
gives an open of \(G/P\) with a rational section.
Over infinite \(k\), its \(G(k)\)-translates cover
geometrically: for any geometric coset the elements
moving it into the open form a nonempty open in \(G\),
and such an open contains a \(k\)-point by density.
These sections lift every rational point of \(G/P\).
Over finite \(k\) every quotient fibre of a rational point
is a torsor under the smooth connected \(P\), so has a
rational point by Result~\ref{res:F08}.
Injectivity is Result~\ref{res:D02}.

We record the needed intersection argument. Geometrically
any two parabolics \(P,Q\) share a maximal torus.
Choose a Borel in each. If the second is \(gBg^{-1}\)
for the first, write \(g=b_1n_wb_2\) by
Result~\ref{res:D04}. Then \(b_1Tb_1^{-1}\) lies
in both Borels. A cocharacter \(\lambda\) of this torus
defining \(P\) gives \(P\cap Q=P_Q(\lambda)\), since a
limit of elements of the closed \(Q\) stays in \(Q\).
The intersection is smooth connected by the dynamic
calculation in Result~\ref{res:C02}.
Its projection to \(L=P/U\) is parabolic: describe
\(Q=P_G(\mu)\) for a cocharacter of the common torus.
Projection from \(P\) is \(\mu\)-equivariant and sends
\(P\cap Q\) onto \(P_L(\mu)\) because this latter group
already has a section in \(L=Z_G(\lambda)\).
The inverse image of \(P_L(\mu)\) in \(P\) is parabolic
in \(G\). Indeed it contains \(U\) and a Borel of \(L\),
whose product is the Borel of \(G\) for a small regular
perturbation of \(\lambda\) toward its chamber in \(L\).
The descriptions descend, using the faithfully flat
image factorization of Result~\ref{res:A07} for the
projection of \(P\cap Q\).

For minimal \(P\), let its split solvable radical \(U\)
act on the proper \(G/Q\) for any \(k\)-parabolic \(Q\).
Result~\ref{res:B14} gives a rational fixed point,
and the rational lifting just proved gives its
representative in \(G(k)\). After conjugating \(Q\)
it therefore contains \(U\).
The parabolic \((P\cap Q)U=P\cap Q\) contained
in \(P\) equals \(P\) by minimality. This proves
rational minimal-parabolic conjugacy.

Any split torus is conjugate into \(P\) by the same
fixed-point argument, acting on \(G/P\).
A split torus \(S'\) in \(P\) lies in a maximal
\(k\)-torus of \(P\): choose a maximal torus of
\(Z_P(S')\) using Result~\ref{res:C02}. It contains
\(S'\) by centrality and has the full geometric rank
because over \(\bar k\) one can extend \(S'\) to a maximal
torus of \(P\) in that centralizer.
The Galois averaging construction of \(P=P_G(\lambda)\)
with this torus produces a Levi containing \(S'\).
Hence \(S'\) can be conjugated into the chosen \(L\)
by \(U(k)\).
If \(P\) is minimal, every split torus in \(L\) is central:
a noncentral cocharacter would define a proper \(k\)-parabolic
of \(L\) and its inverse image would be proper in \(P\).
Every noncentral split torus has such a cocharacter since
its lattice spans it and contains a lattice point outside
the central subspace. The central split tori have a unique
largest product, so the maximal split tori in \(P\) are
conjugate by \(U(k)\). Together with rational conjugacy
into \(P\) this proves the assertion in \(G\).
In particular if \(S\) is maximal split in a minimal \(P\),
one can take \(L\) containing \(S\). The defining cocharacter
then lies in \(S\), the largest split torus of \(L\).
As \(S\) is central in \(L\), it follows that
\(Z_G(S)=Z_G(\lambda)=L\).

The parameter interpretation is given by
\(gP\mapsto gPg^{-1}\), since the scheme normalizer
is \(P\) by Result~\ref{res:D05}. We check local
conjugacy also for families. For a smooth closed parabolic
family \(Q\) in \(G_R\) of the same geometric type,
the transporter is represented by the equations
\(gP_Rg^{-1}\subset Q\). Indeed these are coefficient
equations on the fixed \(k\)-algebra \(k[P]\).
The inclusion forces equality after every geometric
fibre, and thus equality: flatness of the closed
subgroup being included makes the fibre of the
quotient ideal inject into the ambient fibre ring,
where it vanishes. Local Nakayama for that finitely
generated ideal gives zero. Finite presentation of
the transporter follows by taking coefficients of a
finite generating set of equations for \(Q\).
Every geometric fibre of the transporter is nonempty.

For its smoothness, consider a square-zero extension with
ideal \(I\) and an existing conjugacy modulo \(I\).
Use smoothness of \(G\) to lift the conjugating point
locally, so the two subgroup schemes then agree modulo \(I\).
The first-order difference of their ideals is a section
of the normal bundle of \(P\) in \(G\) with coefficients
in \(I\). Explicitly, an equation of \(P\) can be
replaced by an equation of \(Q\) congruent to it modulo
\(I\); restrict their difference in \(I\mathcal O_G\)
to \(P\). The ambiguity vanishes since flatness gives
\(I\mathcal O_G\cap\mathcal I_Q=I\mathcal I_Q\).
This construction is linear on the conormal module
and zero exactly when the two ideals agree.
Right translation trivializes the normal bundle with
fibre \(M=\Lie(G)/\Lie(P)\). The displacement is a
regular cocycle \(c\) on \(P\) with values in \(M\otimes I\):
multiplying perturbed elements gives the infinitesimal
displacement \(c(pq)=c(p)+\operatorname{Ad}(p)c(q)\).
This identity follows in matrix coordinates by multiplying
\((1+v)p(1+w)q\) with square-zero coefficients.
A coboundary is the displacement by a single
infinitesimal conjugation of \(G\).

Restrict the cocycle to the defining \(\Gm\) via
\(\lambda\). The constant-coefficient averaging on
the Laurent polynomial algebra in Result~\ref{res:B13}
works over the test ring with coefficients \(M\otimes I\)
and makes the restriction a coboundary. Subtract it.
The adjusted cocycle vanishes on \(\lambda\), hence is
equivariant for conjugation by \(\lambda\).
On \(L\) its image is fixed and zero, since \(M\)
has only strictly negative weights. On \(U\) it is
zero since its weight coordinates would be regular
functions of negative degree in the contracting positive
grading. Thus it vanishes on \(P=UL\).
This lifts the conjugacy over the extension, proving
formal smoothness, and the finitely presented transporter
is smooth. It is therefore an fppf cover of the base
of the family. Descent of the local cosets gives the
inverse parameter map, with equality checked using
\(N_G(P)=P\). Compare \cite[Corollaries 5.2.8--5.2.9]{conrad}
and \cite[Section 25.a]{milne}.
\end{proof}

\begin{result}{D09}{The relative Bruhat theorem}
Let \(S\) be a maximal \(k\)-split torus and \(P\) a minimal \(k\)-parabolic containing it. Write \(Z=Z_G(S)\), \(N=N_G(S)\), and \(W_k=N(k)/Z(k)\). The group scheme \(N/Z\) is finite constant with group \(W_k\). For \(n_w\in N(k)\) representing \(w\), one has
\[
     G(k)=\coprod_{w\in W_k} P(k)n_wP(k).
\]
The orbit schemes \(Pn_wP\) are locally closed and pairwise disjoint. Their geometric points need not cover \(G(\bar k)\).
\end{result}

\begin{proof}
By Result~\ref{res:D08}, \(P=U\rtimes Z\) with
\(Z=Z_G(S)\). We first check that any two minimal
\(k\)-parabolics have a common maximal split torus.
Their intersection \(P\cap Q\) is smooth connected
of full geometric torus rank, by the intersection
argument in Result~\ref{res:D08}. It therefore has
a maximal \(k\)-torus \(T\) of that geometric rank
by Result~\ref{res:C02}. Projection to \(Z\) gives a
maximal torus of \(Z\), and restricts to an isomorphism
onto it, since its kernel in \(T\) is both of
multiplicative type and unipotent
(Result~\ref{res:B07}). This image contains \(S\):
a central torus belongs to every geometric maximal
torus, since multiplying the two gives a torus.
The inverse image of \(S\) in \(T\) is consequently
a split torus of the full split rank, in both \(P\)
and \(Q\). Apply this to \(P\) and \(gPg^{-1}\)
for \(g\in G(k)\). By rational conjugacy of maximal
split tori within a minimal parabolic
(Result~\ref{res:D08}), move this common torus to
\(S\) inside \(P\), and move \(gSg^{-1}\) to the
common torus inside \(gPg^{-1}\). This gives a
representative of the double coset of \(g\) in \(N(k)\).

Conjugation of the split \(S\) gives a map from \(N\)
to the constant lattice automorphisms of \(S\).
Exactly as in Result~\ref{res:C02}, the weights of a
faithful representation generate the lattice and are
permuted, so the possible images lie in a finite
constant group. Its nonempty geometric fibres are
translates of \(Z\), the kernel, and it maps fppf onto
the union of the corresponding components of that
constant group. Thus \(N/Z\) is finite constant.
Moreover \(N\cap P=Z\) on every test algebra:
if a point of \(P\) normalizes \(S\), then projection
to \(Z\) identifies its conjugation action on \(S\)
with that of its \(Z\)-component, which is trivial.
The constant group \(N/Z\) acts on the minimal
\(k\)-parabolics containing \(S\) by using local
representatives and descending conjugation; \(Z\)
lies in every such parabolic by Result~\ref{res:D08}.
This action is free by the scheme normalizer statement
of Result~\ref{res:D05} and \(N\cap P=Z\).
The subgroup \(N(k)/Z(k)\) acts transitively: conjugate
the minimal parabolics by \(G(k)\), then align their
maximal split tori within the target using
Result~\ref{res:D08}. Freeness of the full group and
transitivity of the subgroup give equality. Every
element of \(N/Z\) therefore has a \(k\)-rational
representative.

To separate the double orbits it suffices to compare
their geometric points. Fix a geometric maximal torus
\(T\) in \(Z\). Each representative in \(N\) normalizes
\(Z\), so after multiplying by an element of \(Z(\bar k)\)
it normalizes \(T\), by conjugacy in \(Z_{\bar k}\).
Take a geometric Borel \(B\subset P\) containing \(T\).
In absolute Weyl indices, \(P_{\bar k}\) is standard,
with Levi Weyl group \(W_I=W(Z_{\bar k},T)\).
The absolute double cells in its double orbit with a
representative \(n\in N_G(T)(\bar k)\) are exactly the
indices \(W_I w W_I\). For right multiplication this
is the rank-one product rule of Result~\ref{res:D04}
and the parabolic indices of Result~\ref{res:D05};
left multiplication follows by inverting. Each
\(W_I\)-element acts trivially on \(S\), since it has
a representative in \(Z\). Thus two absolute indices
of representatives of \(N\) lying in that same double
orbit give the same action on \(S\), hence the same
element of \(N/Z\). This proves disjointness.

For the scheme assertion, \(G/P\) embeds equivariantly
in projective space using the line with stabilizer \(P\)
from Result~\ref{res:A15}, the orbit construction of
Result~\ref{res:D01}, and properness. The smooth group
\(P\) has a locally closed orbit through each \(n_wP\)
by that same orbit construction. Its inverse image in
\(G\) is exactly the indicated locally closed double
orbit, using fppf lifts of quotient points.
Disjointness follows from the geometric-point calculation.
Coverage was asserted only on \(k\)-points; the example
below gives the failure of geometric coverage.
Compare \cite[Theorem 10.3.1 and Appendix P]{conradfeng}.
\end{proof}

The relative roots are the nonzero weights of \(S\) on
\(\Lie(G)\). A maximal \(k\)-torus \(T\) containing \(S\)
can be obtained in \(Z_G(S)\) by Result~\ref{res:C02};
the weights follow by restricting the absolute roots over
\(k_s\) and discarding the zero restrictions.
Every \(k\)-parabolic containing \(S\) is \(P_G(\lambda)\)
for a \(k\)-cocharacter of \(S\): in the construction in
Result~\ref{res:D08} take a maximal \(k\)-torus in
\(Z_P(S)\) containing \(S\), and observe that any
\(k\)-cocharacter of that torus has image in its largest
split subtorus, which equals \(S\).
Two such dynamic groups are equal exactly when the sets
of relative roots with nonnegative pairing agree.
Indeed the absolute root sets after extension consist of
all roots restricting to zero and all roots restricting
to that relative subset, hence determine the parabolic
by Result~\ref{res:D05}.

If \(\lambda\) avoids the relative root hyperplanes,
its Levi is \(Z_G(S)\) by equality of the zero weight
spaces and smooth connectedness. Any split torus of
\(Z_G(S)\) is contained in \(S\), since it commutes with
\(S\) and the product is again split. So this Levi
has no proper \(k\)-parabolic: the cocharacter of a
proper parabolic would be noncentral by
Result~\ref{res:D08}. Any parabolic contained in
\(P_G(\lambda)\) contains its radical, as one checks
over the algebraic closure by using a Borel in the
smaller parabolic and the root description of
Result~\ref{res:D05}. Thus \(P_G(\lambda)\) is minimal.
If a nonzero relative root pairs to zero, \(S\) is
noncentral in the Levi, which then has a proper
\(k\)-parabolic defined by a cocharacter of \(S\).
Its inverse image gives a proper parabolic inside
\(P_G(\lambda)\) by the Levi projection argument in
Result~\ref{res:D08}. Thus minimal parabolics through
\(S\) are exactly the relative chambers.

For a relative root \(a\), use the full scheme kernel
\(D=\ker(a:S\to\Gm)\) and put \(H=Z_G(D)^\circ\).
The centralizer smoothness proof in Result~\ref{res:C02}
applies to this split diagonalizable group: its
character grading, even when the character group has
torsion, permits homogeneous cotangent lifts, and the
nonzero-character functions give the fixed-locus ideal.
In particular no reducedness of \(D\) is needed.
The geometric maximal torus \(T\) of \(G\) is in \(H\).
The root spaces in \(\Lie(H)\) are exactly those whose
characters are trivial on \(D\), and each corresponding
root subgroup is in \(H\), by its scalar weight action.
This is a symmetric set, so \(H\) is reductive by the
symmetric-root argument in Result~\ref{res:D08}.
The relative roots of \(H\) lie on the line of \(a\):
a character of \(S\) trivial on \(D\) is an integral
multiple of \(a\), by the character description in
Result~\ref{res:B01}. Both \(a\) and \(-a\) occur.
The two relative chambers for \(H\) therefore give
exactly two minimal parabolics through \(S\).
The action just proved in Result~\ref{res:D09},
applied to \(H\), provides an element interchanging
them. Its action \(r_a\) on \(S\) negates \(a\) and
fixes the full kernel \(D\), since it has a representative
in \(H(k)\). Hence \(t\mapsto t\,r_a(t)^{-1}\)
factors through \(a:S\to\Gm\), giving an integral
cocharacter \(a^\vee\) and the equations
\[
 r_a(x)=x-\langle x,a^\vee\rangle a,\qquad
 \langle a,a^\vee\rangle=2 .
\]
In particular this argument retains the lattice assertion
when both \(a\) and \(2a\) are roots. Conjugation by the
normalizer takes the construction for \(a\) to that for
its image, so the characters and cocharacters are
permuted along with their reflections. The reflection for a
wall is unique in the finite normalizer image: two
involutions fixing the same hyperplane pointwise and
reversing the transverse line have unipotent product, and
a finite-order unipotent rational linear map is the identity.
The coroots are determined by the formula and distinguish
the roots: distinct proportional roots give distinct
inverse scalings of the same nonzero coroot, and
nonproportional roots have different normal lines for an
invariant inner product averaged over the finite Weyl group.
This gives the possibly non-reduced relative root datum.
The reflections interchange adjacent relative chambers, since
they fix their shared wall and preserve the arrangement.
Following a gallery proves that they generate \(N/Z\),
whose action on chambers is free by the result above.
Compare \cite[Sections 11.3--11.4]{conradfeng}.

\begin{result}{D10}{A real group with two relative and six absolute cells}
For the real group \(\operatorname{SU}(2,1)\) the relative root system is \(BC_1\). A minimal real parabolic gives two real rational double cosets. Its complexification is a Borel of \(\SL_3\) with six absolute cells. The two relative orbit schemes complexify to the cells of the identity and the longest permutation, and hence do not cover the complex group.
\end{result}

\begin{proof}
Use the hermitian form on \(\CC^3\) with matrix
\[
 J=\begin{pmatrix}0&0&1\\0&1&0\\1&0&0\end{pmatrix}
\]
in the basis \(e_1,e_2,e_3\). Its signature is \((2,1)\), and the equations are \(\bar g^{\,t}Jg=J\) and \(\det(g)=1\). The real split torus \(S=\{\operatorname{diag}(t,1,t^{-1})\}\) is maximal. Indeed for any split torus the form pairs a character space only with that of the inverse character: on real points in the identity component the character values are real and their product must be \(1\) for a nonzero pairing. For a generic cocharacter the sum of the positive-weight spaces is totally isotropic, of complex dimension at most \(1\) by the signature. Consequently at most one inverse pair of distinct nontrivial characters occurs. Faithfulness implies that the split torus has rank at most one.

Let \(P\) stabilize the isotropic line \(\CC e_1\). It also stabilizes the orthogonal plane \(\CC e_1+\CC e_2\). On complexifying the equations, the two factors of \(\CC\otimes_{\RR}\CC\) are interchanged by conjugation and the equation for the form determines the second matrix from the first. Thus \(G_{\CC}\simeq\SL_3\). Stabilizing the line in both factors becomes stabilizing the line and the plane in the first factor. Hence \(P_{\CC}\) is the upper triangular Borel and \(P\) is a minimal real parabolic.

Both simple roots of the complex diagonal torus restrict to \(a:t\mapsto t\), and their sum restricts to \(2a\). With negatives these give \(BC_1\), whose Weyl group has order two. The relative reflection is represented by the real matrix interchanging \(e_1,e_3\) and sending \(e_2\) to \(-e_2\), which preserves \(J\) and has determinant \(1\). Its complex Weyl index is \((13)=w_0\). Result~\ref{res:D09} gives two real rational double cosets, whereas Result~\ref{res:D04} gives all six permutations over \(\CC\). For example, a representative of \((12)\) belongs to neither of the two complexified relative orbits. This shows the precise range of \cite[Remark 10.3.2 and Appendix Q.4]{conradfeng}. The geometric covering wording of \cite[25.27]{milne} does not apply to this nonsplit example.
\end{proof}

\chapter{Rational representations and highest weights}
\label{ch:representations}

A rational representation uses the action of the group scheme. In particular, when the ground field is finite, the action of the abstract group of rational points can lose information. We begin with operations valid for any affine group scheme. After that let \(G\) be a connected split reductive group over a field \(k\), with split Borel pair \(T\subset B\). The roots in the unipotent radical \(U\) of \(B\) are positive. Write \(B^-=TU^-\) for the opposite Borel, \(\Phi^+\) for the positive roots, \(\Pi\) for the simple roots, and \(W\) for the Weyl group. Let \(X=X^*(T)\), and write \(\mu\leq\lambda\) if \(\lambda-\mu\in\sum_{\alpha\in\Pi}\mathbb Z_{\geq 0}\alpha\). A weight \(\lambda\in X\) is dominant if \(\langle\lambda,\alpha^\vee\rangle\geq 0\) for all \(\alpha\in\Pi\). All representations in the chapter are finite dimensional.

\section{Operations, weights, and primitive vectors}

\begin{result}{E01}{Rational operations and change of field}
Let \(H\) be an affine group scheme over \(k\) and \(A=\mathcal O(H)\). The category of finite-dimensional right \(A\)-comodules has tensor products, duals, restrictions along homomorphisms, and invariants. Equivariant Hom spaces and invariants commute with every field extension \(K/k\). If a submodule \(U_0\subset V\) admits an \(H_K\)-equivariant retraction after extension to \(K\), it already admits an \(H\)-equivariant retraction over \(k\). Consequently semisimplicity descends along field extensions.
\end{result}
\begin{proof}
Use the action-comodule identification of Result~\ref{res:A05}. For \(\delta(v)=\sum v_0\otimes v_1\) and \(\delta(w)=\sum w_0\otimes w_1\), put
\[
 \delta(v\otimes w)=\sum(v_0\otimes w_0)\otimes v_1w_1.
\]
The bialgebra axioms give the comodule axioms. In a basis \((v_j)\), write \(\delta(v_j)=\sum_i v_i\otimes a_{ij}\). Then \(\Delta(a_{ij})=\sum_\ell a_{i\ell}\otimes a_{\ell j}\) and \(\varepsilon(a_{ij})=\delta_{ij}\). The two antipode identities say that the entrywise antipode \(S(a_{ij})\) forms the inverse matrix. On the dual basis \((v^j)\), set \(\delta(v^j)=\sum_i v^i\otimes S(a_{ji})\). This is the inverse-transpose action, so evaluation and coevaluation are equivariant. Restriction is composition of the coaction with the map of coordinate algebras. A vector is invariant precisely when \(\delta(v)=v\otimes1\).

Equivariance of \(f\colon V\to V'\) is the kernel condition
\[
 \delta_{V'}f-(f\otimes1)\delta_V=0
\]
on the vector space of linear maps. Tensoring this kernel, or the invariants kernel, with \(K\) preserves it by flatness. For the inclusion \(i\colon U_0\hookrightarrow V\), a retraction exists precisely when \(\operatorname{id}_{U_0}\) belongs to the image of
\[
 \operatorname{Hom}_H(V,U_0)\longrightarrow\operatorname{End}_H(U_0),
 \qquad f\longmapsto f i .
\]
Its class in the cokernel vanishes if it vanishes after extension because a vector space injects into its scalar extension. Apply this to every submodule to descend semisimplicity. The equations here are identities in \(A\), so smoothness, reducedness, and density of \(H(k)\) play no role. See \cite[Chapter~4, especially 4.3--4.5, 4.19--4.20 and 4.31--4.34]{milne}.
\end{proof}

\begin{result}{E02}{Root subgroups and the movement of weights}
If \(v\in V_\mu\) and \(x_\alpha\colon\mathbb G_a\to U_\alpha\) is a root parametrization, then
\[
 x_\alpha(c)v=v+\sum_{j\geq1}c^j v_j,\qquad v_j\in V_{\mu+j\alpha},
\]
with finitely many nonzero terms. A representative \(n_w\in N_G(T)(k)\) carries \(V_\mu\) isomorphically to \(V_{w\mu}\).
\end{result}
\begin{proof}
The orbit map of \(x_\alpha\) on \(v\) has polynomial coordinates, so write it as \(\sum_j c^jv_j\). Substituting \(c=0\) gives \(v_0=v\). For the generic \(t\) of \(T\), the identity \(t x_\alpha(c)t^{-1}=x_\alpha(\alpha(t)c)\) in the split root data of Result~\ref{res:C04} implies
\[
 \sum_j c^j(t v_j)=\mu(t)\sum_j c^j\alpha(t)^jv_j
\]
over \(\mathcal O(T)[c]\). Equating coefficients proves the assertion in the weight decomposition of Result~\ref{res:B05}. For the normalizer, \(t(n_w v)=n_w(n_w^{-1}tn_w)v=(w\mu)(t)n_wv\). Apply the same calculation to \(n_w^{-1}\) for the inverse. These arguments compare regular functions, including over a finite field. See \cite[22.14--22.15]{milne}.
\end{proof}

A nonzero \(T\)-weight vector fixed by \(U\) is called a \emph{primitive vector}. It spans a \(B\)-stable line, with \(B\) acting through \(B/U=T\). The \(U\)-invariant subspace of any module is \(T\)-stable and decomposes into such weight spaces.

\begin{result}{E03}{A cyclic module generated by a primitive vector}
If \(V\) is generated by a primitive vector \(v_\lambda\), then \(V_\lambda=kv_\lambda\), all other weights are \(<\lambda\), and \(\lambda\) is dominant. The sum of all proper \(G\)-submodules is the unique maximal proper submodule. Each simple module has a unique primitive line. Two simple modules with the same highest weight are isomorphic.
\end{result}
\begin{proof}
Write the coaction of \(v_\lambda\) as \(\sum_i v_i\otimes a_i\) with the \(a_i\) linearly independent. Coassociativity, followed by linear functionals on the last tensor factor selecting each \(a_i\), implies that the span of the \(v_i\) is a subcomodule. It contains \(v_\lambda\) by the counit and is contained in every subcomodule containing \(v_\lambda\). Hence a linear functional annihilates the cyclic module exactly when its pairing with the orbit morphism is zero. It is enough to test this on the open dense cell \(U^-B\) from Result~\ref{res:D03}, since \(G\) is integral. Multiplication on the right by \(B\) only multiplies \(v_\lambda\) by a scalar, so the coefficients of the \(U^-\)-orbit suffice. The negative root groups multiply onto \(U^-\) as in Result~\ref{res:C06}. Applying Result~\ref{res:E02} successively to each factor shows that the constant term is \(v_\lambda\) and every remaining term has strictly smaller weight. The weight decomposition therefore has \(\lambda\)-part \(kv_\lambda\). For each simple root \(\alpha\), the weight \(s_\alpha\lambda\) occurs by Result~\ref{res:E02}. Since it is at most \(\lambda\), the relation \(\lambda-s_\alpha\lambda=\langle\lambda,\alpha^\vee\rangle\alpha\) gives dominance. A submodule meeting \(kv_\lambda\) nontrivially contains the generator, so every proper submodule misses the \(\lambda\)-space. Their sum misses it as well by the weight decomposition. The quotient by this sum is simple.

A simple module contains a primitive vector: any nonzero vector in a maximal weight space is fixed by all positive root groups, by the weight calculation. Each such vector generates the simple module. Two primitive vectors thus have weights at most each other's weights, which coincide, and the corresponding space is one dimensional. For simple \(V,V'\) of highest weight \(\lambda\), take the cyclic submodule \(M\subset V\oplus V'\) generated by \(v_\lambda+v'_\lambda\). Its projection to either summand is surjective, and its intersection with the other summand misses that summand's primitive vector by the one-dimensional \(\lambda\)-space in \(M\). The intersection is a proper submodule of a simple module, hence zero. Both projections identify \(M\) with the corresponding simple module. See \cite[22.16--22.19]{milne}.
\end{proof}

\section{Induced modules and classification}

\begin{result}{E04}{The flag line bundle and induction}
Write \(k_\lambda\) for the one-dimensional \(B^-\)-module coming from \(\lambda\colon T\to\mathbb G_m\). On \(Y=G/B^-\), use the line bundle
\[
 \mathcal L_\lambda=G\times^{B^-}k_\lambda,\qquad
 (g,z)b=(gb,\lambda(b)^{-1}z).
\]
Its global sections \(\nabla(\lambda)=H^0(Y,\mathcal L_\lambda)\) are the regular functions satisfying \(f(gb)=\lambda(b)^{-1} f(g)\). Left translation \((h f)(g)=f(h^{-1}g)\) gives a rational action. Evaluation at \(1\) gives an isomorphism
\[
 \operatorname{Hom}_G(V,\nabla(\lambda))
 \simeq \operatorname{Hom}_{B^-}(V,k_\lambda).
\]
\end{result}
\begin{proof}
Results~\ref{res:D01} and \ref{res:D07} give the \(B^-\)-torsor \(G\to Y\) and a projective \(Y\). A section of the pullback is a regular function on \(G\), and its two pullbacks to \(G\times B^-\) agree after the prescribed identification precisely by the displayed covariance condition. Faithfully flat descent of the trivialized line bundle gives the claimed functions. Equivariance of the condition under left translation makes their space a subcomodule of the left regular comodule: the orbit formula, viewed on \(G\times G\), preserves the condition in the second variable. This subspace is finite dimensional by projectivity, and the formula for the action gives \((b f)(1)=f(b^{-1})=\lambda(b)f(1)\). For a \(B^-\)-equivariant linear map \(\ell\colon V\to k_\lambda\), send \(v\) to the function \(g\mapsto\ell(g^{-1}v)\). The regular action and inversion give regularity, and equivariance of \(\ell\) gives the required covariance under \(b\). This construction is \(G\)-equivariant by left translation. Conversely if \(\phi\colon V\to\nabla(\lambda)\) is equivariant, \(\phi(v)(g)=\phi(g^{-1}v)(1)\), proving the two constructions inverse. See \cite[22.21, 22.29--22.30, 22.48]{milne} and \cite[Part~I, Chapter~3]{jantzen}.
\end{proof}

We will use the following elementary comparison for a central isogeny \(\pi\colon G'\to G\) of connected split reductive groups. Suppose \(B'^-\) is the opposite Borel above \(B^-\), so \(\pi^{-1}(B^-)=B'^-\) as schemes. Indeed \(B'^-\) maps faithfully flat onto \(B^-\) (on the torus it is the torus isogeny and on the unipotent radical it is an isomorphism) and contains \(\ker\pi\). The induced map of fppf quotients \(G'/B'^-\to G/B^-\) is an isomorphism, and under this identification \(\mathcal L_{\pi^*\lambda}\) is the pullback of \(\mathcal L_\lambda\). These assertions follow by taking local representatives in the torsors, lifting \(G\)-points fppf locally, and using \(\pi^{-1}(B^-)=B'^-\) for equality of lifted cosets. The same calculation gives the comparison of the associated line bundles. In particular it works if \(\ker\pi\) is non-smooth.

\begin{result}{E05}{Dominance, nonvanishing, and the simple socle}
The space \(\nabla(\lambda)\) is nonzero exactly for dominant \(\lambda\). When nonzero, its simple socle is the unique simple submodule, its primitive line has weight \(\lambda\), and each weight \(\mu\) satisfies \(w_0\lambda\leq\mu\leq\lambda\), where \(w_0\) is the longest Weyl element.
\end{result}
\begin{proof}
On the open dense cell \(UB^-\), right covariance determines a section by its restriction to \(U\). If it is \(U\)-invariant, that restriction is constant by the regular left action of \(U\) on itself. Thus \(f\mapsto f(1)\) is injective on \(\nabla(\lambda)^U\). A nonzero \(\nabla(\lambda)\) has a nonzero \(U\)-invariant weight vector by maximality of a weight and Result~\ref{res:E02}. The invariant line has weight \(\lambda\), since \((t f)(1)=f(t^{-1})=\lambda(t)f(1)\). Result~\ref{res:E03} gives dominance. Every simple submodule has a primitive vector and therefore contains this same line, which generates the simple submodule. Existence of a simple submodule follows from finite dimensionality. Every maximal weight is the weight of a primitive vector, hence is \(\lambda\). This proves the upper bound, and applying \(w_0\) gives the lower bound.

We prove nonvanishing for dominant weights using the construction of \cite[22.23--22.27]{milne}. First suppose \(G\) semisimple. For each simple root \(\alpha_i\), let \(P_i\supset B\) be the maximal proper standard parabolic omitting \(\alpha_i\). By Result~\ref{res:A15}, there is a line \(kv_i\) with scheme stabilizer \(P_i\). Its generator is primitive and has dominant weight \(\eta_i\). In the cyclic module generated by \(v_i\), the \(\eta_i\)-space is precisely \(kv_i\). Simple reflections other than \(s_i=s_{\alpha_i}\) have representatives in \(P_i(k)\), so each fixes the weight \(\eta_i\). The representative of \(s_i\) is outside \(P_i(k)\), hence does not preserve \(kv_i\). If \(s_i\eta_i=\eta_i\), Result~\ref{res:E02} would force it to preserve this line. Thus \(\eta_i=d_i\omega_i\) for an integer \(d_i>0\), where \(\omega_i\) is the \(i\)-th fundamental weight in the rational character space. The multiple \(\eta_i\) is in the actual character lattice even when \(\omega_i\) is not. Tensoring the \(v_i\)'s, with repetitions, produces a primitive vector of weight \(d\eta\) for any dominant \(\eta\) and any common multiple \(d\) of the \(d_i\), using \(\eta=\sum_i \langle\eta,\alpha_i^\vee\rangle\omega_i\).

If \(v\) is primitive of weight \(\eta\), then \(n_{w_0}v\) is a \(B^-\)-eigenvector of weight \(w_0\eta\). Choose a linear functional \(a\) with \(a(n_{w_0}v)\ne0\). The coefficient \(g\mapsto a(g n_{w_0}v)\) lies in \(\nabla(-w_0\eta)\), since right translation by \(b\in B^-\) multiplies it by \((w_0\eta)(b)\). Apply this with \(\eta=d(-w_0\lambda)\), which is dominant and is a weight of the form just constructed. The normalized \(U\)-invariant section of \(\nabla(d\lambda)\) restricts on \(UB^-\) to the function \(ub\mapsto(d\lambda)(b)^{-1}\). It equals the \(d\)-th power of the regular function \(ub\mapsto\lambda(b)^{-1}\) there. The latter is a rational function on \(G\) integral over \(\mathcal O(G)\), since its \(d\)-th power is global. The affine variety \(G\) is smooth and integral, so \(\mathcal O(G)\) is integrally closed. The resulting global function lies in \(\nabla(\lambda)\): on \(UB^-\times B^-\) it satisfies the covariance formula, which then holds on \(G\times B^-\) by density. Its value at \(1\) is \(1\).

In general, take the central isogeny \(Z_0\times G_{\rm sc}\to G\) given by the maximal central torus and the simply connected cover of \(G^{\rm der}\), using Results~\ref{res:C03} and \ref{res:C07}. The torus factor contributes a one-dimensional representation, and the restriction of the pulled-back weight to \(G_{\rm sc}\) is dominant. The semisimple case gives a nonzero section on the cover, and the comparison of flag line bundles preceding the result gives a nonzero section on \(G/B^-\). This applies to the actual character \(\lambda\in X\) chosen at the start and does not require the degree of the isogeny to be invertible in \(k\).
\end{proof}

\begin{result}{E06}{Classification and absolute simplicity}
Dominant weights in \(X\) classify the simple rational \(G\)-modules over \(k\). Denote by \(L(\lambda)\) the simple socle of \(\nabla(\lambda)\). Then \(\operatorname{End}_G L(\lambda)=k\), and \(L(\lambda)\) remains simple with highest weight \(\lambda\) after every field extension. For semisimple \(G\), the entire weight lattice is the character lattice exactly in the simply connected case.
\end{result}
\begin{proof}
Results~\ref{res:E03} and \ref{res:E05} give existence and uniqueness, and an equivariant endomorphism is determined by the scalar on the cyclic highest weight line. Under any extension, the \(U\)-invariant line of \(L(\lambda)\) remains the invariant line by Result~\ref{res:E01}. A nonzero submodule after extension has a primitive vector by its maximal weight. That vector lies in the same invariant line. The line still generates because the linear span of the coefficients of the coaction on a generator commutes with extension. Thus the submodule equals \(L(\lambda)_K\). The assertion about the lattice is the simply connected root datum criterion of Result~\ref{res:C07}. See \cite[22.8--22.12 and 22.27]{milne}.
\end{proof}

For \(\mathrm{SL}_2\), with diagonal parameter \(t\mapsto\operatorname{diag}(t,t^{-1})\), weights are integers, the positive root is \(2\), and the dominant weights are \(m\geq0\). The central isogeny \(\mathrm{SL}_2\to\mathrm{PGL}_2\) pulls the target lattice back to \(2\mathbb Z\). The scheme centre is \(\mu_2\) and acts on \(L(m)\) by \(z\mapsto z^m\), since every weight differs from \(m\) by an even integer. Hence \(L(m)\) descends if and only if \(m\) is even. In characteristic \(2\), the character of \(\mu_2\) of odd exponent is still nontrivial in the coordinate Hopf algebra, even though \(\mu_2(\overline k)\) is trivial. This also illustrates why rational-point checks alone cannot determine the descent. See \cite[22.11--22.12, 22.33 and 22.35]{milne}.

\section{Semisimplicity and formal characters}

\begin{result}{E07}{Complete reducibility in characteristic zero}
Suppose \(\operatorname{char} k=0\). For a connected affine algebraic group \(H\), the following conditions are equivalent: \(H\) is reductive, all its finite-dimensional rational representations are semisimple, and it has a faithful semisimple finite-dimensional rational representation.
\end{result}
\begin{proof}
We first explain the passage from Lie-invariance to group-invariance in characteristic zero. If \(W\subset V\) is stable under \(\operatorname{Lie}(H)\), its scheme stabilizer in \(H\) has Lie algebra \(\operatorname{Lie}(H)\) by differentiating the condition that the composite \(W\to V\to V/W\) vanish. Both the stabilizer and \(H\) are smooth by Cartier's theorem (Result~\ref{res:A09}). Over an algebraic closure the stabilizer therefore has the dimension of \(H\), and a closed subgroup of that dimension in the connected smooth group must be all of \(H\). Equality of these closed subgroup schemes descends to \(k\).

Suppose \(H\) is reductive. Over \(\bar k\), Result~\ref{res:C03}
gives the faithfully flat central multiplication map \(Z_{\rm tor}
\times D\to H_{\bar k}\), with \(D=(H_{\bar k})^{\rm der}\)
connected semisimple. Here is why the Lie algebra of \(D\) is
semisimple in characteristic zero. Choose a maximal torus \(T_D\)
by Result~\ref{res:C02}; it is split over \(\bar k\).
Result~\ref{res:C04} expresses \(\operatorname{Lie}(D)\) as the
direct sum of \(\operatorname{Lie}(T_D)\) and the one-dimensional
root spaces. The roots span the character space over \(\mathbb Q\):
by Results~\ref{res:C03} and \ref{res:C05} the scheme centre is
both finite and diagonalizable with character group the quotient
by the roots. Their differentials therefore span
\(\operatorname{Lie}(T_D)^*\) over \(\bar k\).

Any Lie ideal is stable under \(\operatorname{Lie}(T_D)\), and
decomposes into its zero and root weight spaces. One can project
to these spaces by polynomials in the adjoint action of a toral
element that separates the finitely many weights; it exists
because the torus differential of distinct characters is distinct
in characteristic zero. If a solvable ideal contains a nonzero
root vector, bracket it with a vector in the opposite root space.
The rank-one matrix models of Result~\ref{res:C04} give a
nonzero coroot vector. Bracketing with the opposite root space
then puts that whole root space in the ideal as well, since the
coroot acts there by \(-2\). These three vectors span a copy of
\(\mathfrak{sl}_2\), whose derived algebra is itself, contradicting
solvability. Thus a solvable ideal is contained in the toral space.
Bracketing its elements with each root vector now gives zero,
and the spanning property of the root differentials forces that
ideal to be zero.

The central torus gives weight spaces by Result~\ref{res:B05}.
The action of \(D\) preserves them and a submodule decomposes
by the same weight spaces. On each weight space,
\(\operatorname{Lie}(D)\) acts completely reducibly by the proof
of Theorem C05 in Chapter 3 of the companion volume
\emph{Lie Algebras and Finite-Dimensional Representations}.
Its hypotheses are a semisimple Lie algebra and a finite-dimensional
module over a characteristic-zero field, as above. Choose a
Lie-equivariant complement on each weight space. The stabilizer
argument makes it \(D\)-stable, and the sum is also stable under
the central torus. Pullback on coordinate rings along the
faithfully flat multiplication map is injective, so this is an
\(H_{\bar k}\)-equivariant complement. Result~\ref{res:E01}
then descends the retraction to \(k\).

If all modules are semisimple, the faithful linear realization in Result~\ref{res:A06} is one of them. Conversely, let \(V\) be faithful and semisimple and let \(N=R_u(H)\). The field is perfect and \(N_{\overline k}=R_u(H_{\overline k})\) is smooth, connected and unipotent. In every nonzero simple summand \(S\) of \(V\), the space \(S^N\) is nonzero: after extension to \(\overline k\), unipotent triangularization from Result~\ref{res:B06} gives a fixed vector, and invariants commute with extension by Result~\ref{res:E01}. Normality of \(N\) makes \(S^N\) an \(H\)-submodule, so it equals \(S\). Thus \(N\) acts trivially on \(V\) as a group scheme. Faithfulness gives \(N=1\), proving reductivity. See also \cite[22.39--22.42]{milne}.
\end{proof}

\begin{result}{E08}{Weyl modules, induced modules, and duality}
For dominant \(\lambda\), put \(\lambda^*=-w_0\lambda\) and define the Weyl module \(\Delta(\lambda)=\nabla(\lambda^*)^\vee\). The socle of \(\nabla(\lambda)\) and the head of \(\Delta(\lambda)\) are both \(L(\lambda)\). The highest weight spaces of all three representations are one dimensional, and \(L(\lambda)^\vee\simeq L(\lambda^*)\). In characteristic zero,
\[
 \Delta(\lambda)\simeq L(\lambda)\simeq \nabla(\lambda).
\]
\end{result}
\begin{proof}
A dual replaces each weight \(\mu\) by \(-\mu\) and exchanges simple submodules and simple quotients. The unique lowest weight of \(L(\lambda)\) is \(w_0\lambda\), with multiplicity one: the upper and lower bounds follow from Result~\ref{res:E03} and Weyl symmetry, and \(n_{w_0}\) identifies its weight space with the highest line. The dual is simple, so its highest weight \(\lambda^*\) identifies it by Result~\ref{res:E06}. The bounds and multiplicities for \(\nabla(\lambda)\) follow from Result~\ref{res:E05} and the same Weyl symmetry. Dualizing the inclusion of the simple socle in \(\nabla(\lambda^*)\) gives the quotient of \(\Delta(\lambda)\) by its unique maximal submodule. Indeed, any maximal submodule would dualize to a simple submodule of \(\nabla(\lambda^*)\), and there is only one. In characteristic zero, Result~\ref{res:E07} gives semisimplicity, so a module with a simple socle is simple. See \cite[22.27--22.31]{milne} and \cite[Part~II, Sections~2.6 and~2.13]{jantzen}.
\end{proof}

\begin{result}{E09}{The character ring}
For a rational \(G\)-module \(V\), set
\[
 \operatorname{ch}_T V=\sum_{\mu\in X}(\dim_k V_\mu)e^\mu\in\mathbb Z[X].
\]
It is additive on short exact sequences and multiplicative under tensor products. Restriction to \(T\) induces an isomorphism of rings
\[
 K_0(\operatorname{Rep}_k G)\xrightarrow{\ \sim\ }\mathbb Z[X]^W .
\]
\end{result}
\begin{proof}
The weight decomposition is exact and the tensor product adds weights, proving the first two assertions. The normalizer preserves weight multiplicities by Result~\ref{res:E02}. The finite-length category has Grothendieck group freely generated by the classes of simple modules by the Jordan--H\"older theorem. The invariant group ring has a basis of Weyl orbit sums \(m_\lambda=\sum_{\mu\in W\lambda}e^\mu\), indexed by dominant weights. By Result~\ref{res:E06}, the character of \(L(\lambda)\) is
\[
 m_\lambda+\sum_{\nu<\lambda,\ \nu\ {\rm dominant}} c_{\lambda\nu}m_\nu .
\]
Each sum is finite. For spanning, put
\(2\rho^\vee=\sum_{\alpha\in\Phi^+}\alpha^\vee\).
With a Weyl-invariant positive definite form on the root
span, the coroot corresponding to a root is its positive
scalar multiple under the induced identification of root
and coroot spans. A nonnegative expansion of a positive root
in the simple roots therefore gives a nonnegative expansion
of its coroot in the simple coroots. Hence \(2\rho^\vee\)
pairs nonnegatively with each dominant weight. A simple
reflection \(s_i\) permutes the positive coroots other than
\(\alpha_i^\vee\): for any positive root other than
\(\alpha_i\), its expansion has a positive coefficient at some
other simple root and that coefficient is unchanged by
\(s_i\). Thus
\(s_i(2\rho^\vee)=2\rho^\vee-2\alpha_i^\vee\), and the
reflection formula gives
\(\langle\alpha_i,2\rho^\vee\rangle=2>0\). Thus its pairing with \(\nu\) is a strictly smaller nonnegative integer for \(\nu<\lambda\) dominant. In the rootless case the order has no strict comparisons. For linear independence, a maximal weight among the finitely many terms of any proposed relation contributes its orbit sum with coefficient equal to its own coefficient in the relation. This proves the additive isomorphism, and compatibility with tensor products proves the ring assertion. Formal characters take their coefficients in \(\mathbb Z\), even if \(k\) has positive characteristic, and \(K_0\) identifies extensions with the sums of their composition factors. See \cite[22.36--22.38]{milne}.
\end{proof}

\section{Flag cohomology and the Weyl formula}

Let \(\rho=\frac12\sum_{\alpha\in\Phi^+}\alpha\). For root expressions use \(\mathbb Z[\frac12 X]\) when \(\rho\) is not in \(X\). Put
\[
 A(e^\nu)=\sum_{w\in W}(-1)^{\ell(w)}e^{w\nu}.
\]
The pullback under a central isogeny injects character lattices and does not change the Weyl group or the root correspondence. Thus a character identity can be checked on the cover, including a cover for which \(\rho\) is integral.

\paragraph{Restriction to Frobenius kernels.}
Work over an algebraically closed field of characteristic \(p>0\),
write \(q=p^r\), and put \(G_r=\ker(F_G^r)\).
For a finite group scheme \(K\), its modules are equivalently
modules over the algebra \(k[K]^\vee\), with multiplication
dual to the Hopf coproduct. Indeed the coaction pairs against
linear functionals, and using a basis with its dual recovers it
from the action. We use the notation \(\operatorname{Dist}(K)\)
for this finite-dimensional algebra.

Let \(M=L(\eta)\) for a \(q\)-restricted dominant \(\eta\),
and let \(v\) be its highest vector. For each simple root use
opposite root coordinates \(x_+,x_-\), normalized to map
to the two standard transvections in the rank-one
projective-line action of Result~\ref{res:C04}.
Any nonzero scaling of the positive coordinate can be
matched by the inverse scaling of the negative one.
For invertible \(u\) the definitions of \(w(t)\) and
\(\alpha^\vee(t)=w(t)w(1)^{-1}\) give
\[
 x_-(u)=x_+(u^{-1})\,\alpha^\vee(-u^{-1})\,w(1)\,
                      x_+(u^{-1}).
\]
Indeed expand \(w(-u^{-1})\) as the three root factors.
If \(m=\langle\eta,\alpha^\vee\rangle\), then
\(w(1)v\) has weight \(s_\alpha\eta\), and \(x_+\) fixes \(v\).
Thus
\[
 x_-(u)v=(-u)^m x_+(u^{-1})w(1)v .
\]
The left side is polynomial in \(u\), and on the right no
power exceeds \(m\). Its coefficient of \(u^m\) is
\((-1)^m w(1)v\). Since \(m<q\), all coefficients of
\(x_-(u)v\) are obtained from the action of
\(\operatorname{Dist}({\Ga}_r)\), by taking coefficients in
\(k[u]/(u^q)\).

Put \(M'=\operatorname{Dist}(G_r)v\). Normality of \(G_r\)
makes conjugation by a group point preserve its distribution
algebra. The calculation shows \(x_-(u)M'\subseteq M'\)
for each simple negative root. The torus and the positive
simple root subgroups preserve \(M'\) because they preserve
the highest line. These subgroups generate \(G(k)\):
they give the simple Weyl representatives, their conjugates
give all the root groups by Result~\ref{res:D04}, and the
open root product gives a dense open whose points generate
the connected \(G\). Thus \(M'\) is stable under \(G\)
by geometric density and equals \(M\).

We check that \(M\) remains simple on \(G_r\).
The socle as a \(G_r\)-module is nonzero and is preserved
by \(G(k)\) by normality, hence by \(G\). Thus \(M\) is
semisimple on \(G_r\).
The image algebra of \(\operatorname{Dist}(G_r)\) on
\(M\) is a product of full matrix algebras on the distinct
simple constituents. To see this, on each constituent
the simple-module density argument from
Result~\ref{res:B06} gives its full endomorphism algebra.
Annihilators of nonisomorphic constituents are distinct
maximal two-sided ideals: a matrix algebra has a unique
simple module by the matrix units. Pairwise these ideals
are comaximal; successive products of elements representing
\(1\) modulo one ideal and \(0\) modulo the others give
the simultaneous projections. This proves the product
description. Conjugation by the connected \(G\) permutes
the finite set of its central primitive idempotents and
therefore fixes them. Their images in \(M\) are \(G\)-stable,
so only one type occurs.

Write \(M=S\otimes V\) with the image algebra
\(\operatorname{End}(S)\) on the first factor. The action
of \(T\) on this algebra gives a homomorphism to
\(\operatorname{PGL}(S)\). The matrix-algebra automorphism
description holds as a sheaf: locally use the images of
matrix units to decompose a free module into rank-one
summands, and choose a generator in one and transport it
to the others by the matrix units. This gives conjugation
by a change of basis, unique up to a scalar.
The pullback of \(\operatorname{GL}(S)\to
\operatorname{PGL}(S)\) along \(T\) is a central extension
of \(T\) by \(\Gm\). It is commutative: its commutator
descends to a homomorphism in each variable on
\(T\times T\) with values in \(\Gm\). Such a map is
trivial because a family of torus characters is locally
constant (the group-like coefficient argument of
Result~\ref{res:C02}) and the character at the identity
is trivial. This smooth connected commutative group
is a torus, by Result~\ref{res:B12} and the vanishing
of its unipotent radical under projection and in the
kernel. The exact sequence of its free character
lattices splits, hence \(T\) lifts to
\(\operatorname{GL}(S)\). Remove this lifted action
on the first factor from the action on \(M\).
The remaining action commutes with the matrix algebra
and therefore acts on \(V\), where it is a torus
representation. Its weight lines decompose \(M\)
into \(G_r\)-simple, \(T\)-stable summands.
The highest weight space has dimension one by
Result~\ref{res:E06}, so \(v\) lies in one of them.
Its \(G_r\)-cyclicity now gives simplicity on \(G_r\).

For the special weight \(\eta=(q-1)\rho\) of a simply
connected semisimple group, compute the dimension as follows.
The open root product at the identity contains \(G_r\),
whose underlying geometric space is just the identity.
Since the product commutes with the relative Frobenius,
it identifies \(G_r\) as a scheme with
\(U^-_r\times T_r\times U^+_r\). Thus the ordered products
of the three distribution spaces span
\(\operatorname{Dist}(G_r)\); \(T_r\) acts on the highest
line by its character and \(U^+_r\) acts trivially there.
Consequently \(A=\operatorname{Dist}(U^-_r)\) generates
\(M\) from \(v\). Its coordinate algebra, in any order
of negative roots, is \(k[u_1,\ldots,u_N]/(u_1^q,\ldots,u_N^q)\)
by the root product and the relative Frobenius, where
\(N=|\Phi^+|\).

Give these coordinates their positive degrees under a regular
contracting cocharacter for \(U^-\).
The coproduct of each generator has its two linear terms
and mixed terms using only variables of smaller degree
in both factors, since the coproduct preserves the total
degree. Thus translations on the coordinate algebra send
each variable to itself plus a polynomial in smaller-degree
variables with coefficients in the base algebra.
The top-coefficient functional \(\omega\), extracting the
coefficient of \(\prod_i u_i^{q-1}\), is invariant under
translations. Indeed take the exponents in decreasing
degree order: using a correction term decreases the
first corresponding exponent without increasing any
earlier one. No monomial with all individual exponents
less than \(q\) other than the top one can contribute.
The coefficient on that top one stays \(1\).
The pairing \((f,g)\mapsto\omega(fg)\) is nondegenerate, by
complementary monomials, and invariant under translations.
This identifies the coordinate regular representation with
its dual. Its invariant functions are constants, by evaluating
the translation identity at the identity point, so the
invariants in the left regular \(A\)-module form the
one-dimensional line \(k\omega\). The augmentation ideal
of \(A\) is nilpotent by the positive grading of this
finite-dimensional algebra. Therefore every nonzero left
ideal contains a nonzero invariant, by taking its last
nonzero image under successive augmentation powers.

Choose a reduced expression \(w_0=s_1\cdots s_N\).
Writing \(w_j=s_1\cdots s_j\), the roots
\(\beta_j=w_{j-1}\alpha_j\) for the simple roots
\(\alpha_j\) of \(s_j\) are all positive roots, each once:
they are the distinct crossed hyperplanes of the reduced
gallery, with positive sign from the base chamber.
Use the negative root coordinates obtained by conjugating
the simple coordinate for \(-\alpha_j\) by the chosen
representative \(\dot w_{j-1}\), in reverse order in
the product. The vector \(\dot w_{j-1}v\) is fixed
by the corresponding positive \(\beta_j\)-root group
and has pairing \(q-1\) with \(\beta_j^\vee\), since
\(\langle \rho,\alpha_j^\vee\rangle=1\).
The rank-one polynomial calculation shows the coefficient
of degree \(q-1\) in the next root action gives, up to sign,
\(\dot w_jv\). Extracting successively all these coefficients
in the ordered product gives \(\omega v=\pm\dot w_0v\ne0\).
Here \(\omega\) is the functional for those ordered coordinates,
to which the same grading argument applies.
Hence the left ideal kernel of \(A\to M,\ a\mapsto av\)
is zero. The map is surjective as above, and
\(\dim M=q^N\).
Compare \cite[Part~II, Sections 3.15--3.18]{jantzen}.

\begin{result}{E10}{Euler characteristic and Kempf vanishing}
For every \(\lambda\in X\), the equivariant Euler characteristic on \(Y=G/B^-\) satisfies
\[
 \sum_{i\geq0}(-1)^i\operatorname{ch}_T H^i(Y,\mathcal L_\lambda)
 =\frac{A(e^{\lambda+\rho})}{A(e^\rho)} .
\]
If \(\lambda\) is dominant, then \(H^i(Y,\mathcal L_\lambda)=0\) for \(i>0\), in every characteristic.
\end{result}
\begin{proof}
We can extend to an algebraic closure for cohomology: a finite affine open cover of the separated variety \(Y\) computes the cohomology of a quasi-coherent sheaf by its \v{C}ech complex, and extension of the field commutes with that complex and its cohomology. We may also use the central isogeny \(Z_0\times G_{\rm sc}\to G\) and the flag comparison in this chapter. On the flag variety the torus factor then only changes the linearization by a character, so for vanishing it suffices to treat a split simply connected semisimple group. For the Euler characteristic its weight is retained on the cover. A torus has flag variety a point, for which both assertions are immediate.

For the Euler characteristic let \(E=G_0^T(Y)\) be the
group of \(T\)-equivariant coherent sheaves with the relations
from short exact sequences. It is a \(\ZZ[X]\)-module by
tensoring with a torus character. We use the following
localization consequence of this definition. For an invariant
closed subset \(Z\), restriction to its open complement is
surjective and its kernel is generated by the classes on
\(Z\). In the situation here the assertion can be proved
directly. An equivariant coherent sheaf on an invariant
open extends by taking a coherent submodule of its
quasi-coherent direct image that generates the restriction.
One obtains such a module by choosing finitely many local
generators and clearing denominators on an affine cover.
Then take the equivariant span: on invariant affine charts
the \(T\)-action is a grading, and the weight components
of finitely many generators suffice. These graded spans
glue because weight decompositions commute with restriction
on invariant principal opens. The charts can be taken
from an equivariant projective embedding by inverting an
eigen-coordinate; for locally closed pieces one extends to
their projective closures and restricts.
An exact sequence on the open extends by first extending
its middle term and then taking the inverse image of the
subsheaf in the direct image. This inverse image is coherent
because the middle term is noetherian. Two extensions of
the same sheaf differ in the group by classes supported on
the complement: extend the graph of their common
identification and use its two projections, whose kernels
and cokernels are supported on the complement.
Finally a coherent sheaf supported on \(Z\) has a finite
filtration by powers of the ideal of \(Z\), with factors
on \(Z\). These arguments establish the claimed restriction
sequence on the class groups.

By the split Bruhat cells of Results~\ref{res:D04}
and~\ref{res:D06}, the fixed points of \(T\) are
\(x_w=n_w B^-\), one in each cell, and cell closures
are unions of cells. In the ordered root coordinates
of a cell, the complement of its origin has a finite
stratification by the first nonzero coordinate.
Each such coordinate has a nonzero character and is
invertible on its stratum. Multiplication by it identifies
a sheaf twisted by that character with the original one.
Thus \(1-e^\alpha\) for that weight (or its negative,
depending on the coordinate convention) kills the
equivariant class group of the stratum. Filter by the
Bruhat closures and these coordinate strata, and apply
the restriction sequence repeatedly. After inverting
all \(1-e^\alpha\) for nonzero roots the complement of
the fixed points has zero class group, so \(E\) is
generated by the images of the fixed-point classes.

On this localized group, derived restriction to \(x_w\)
is computed by the alternating \(\operatorname{Tor}\)
groups with \(k(x_w)\). It is additive by the long exact
sequence and finite because the local ring at \(x_w\)
is regular. The contribution of the point class at any
other point is zero; at \(x_w\) the Koszul resolution
of its residue field gives
\(\prod_\nu(1-e^\nu)\), where \(\nu\) ranges over the
cotangent weights. To see the equivariance, choose
homogeneous lifts of a basis of the cotangent space
on an invariant affine chart; they are a regular system
of parameters in the smooth local ring. The weight
calculation on the Koszul factors gives precisely the
alternating exterior powers of the cotangent space.
Thus the coefficient of the point class in
\([\mathcal L_\lambda]\) is its fibre character divided
by this product. Equivariant Euler characteristic
is additive on coherent sheaves by properness and the
cohomology exact sequence, and sends each point class
to \(1\). Applying it proves the fixed-point formula
in this case. Compare \cite[Theorem 22.50]{milne}.

The cotangent weights at the identity coset are
\(-\alpha\) for \(\alpha\in\Phi^+\), by the tangent
quotient \(\Lie(G)/\Lie(B^-)\). At \(x_w\) they are
\(-w\alpha\), and the fibre character is \(w\lambda\).
Put \(D=\prod_{\alpha>0}(1-e^{-\alpha})\).
A simple reflection changes only the factor in its own
root to its negative exponent, and permutes the other
factors. Thus \(wD=(-1)^{\ell(w)}e^{\rho-w\rho}D\).
The fixed-point sum becomes \(A(e^{\lambda+\rho})/(e^\rho D)\).

For completeness, \(e^\rho D=A(e^\rho)\) because \(e^\rho D\) is alternating, has coefficient \(1\) at \(\rho\), and is supported on weights \(\leq\rho\) in \(\rho+\mathbb Z\Phi\). A dominant representative \(\nu\) of an orbit carrying nonzero coefficient must be regular, since a simple reflection fixing \(\nu\) forces its coefficient to be zero in the integral group ring. The simple coroot pairings of \(\nu\) are positive integers, hence \(\nu-\rho\) is dominant. Pairing with \(2\rho^\vee\) rules out any strict inequality \(\nu<\rho\). This proves the denominator identity, and establishes the stated formula.

For vanishing we use the ampleness of \(\mathcal L_\eta\) for strictly dominant \(\eta\) on the simply connected flag variety. One can see this directly from the maximal standard parabolics \(P_i^-\supset B^-\) omitting \(-\alpha_i\). Apply the primitive-vector argument of Result~\ref{res:E05} with \(B^-\) as the positive Borel. A line with scheme stabilizer \(P_i^-\) then has weight \(-d_i\omega_i\) for \(d_i>0\). The orbit map \(G/P_i^-\to\mathbb P(V_i)\) of that line is a monomorphism by the stabilizer condition and is proper, hence a closed immersion. The pullback of \(\mathcal O(1)\) to \(Y\) is \(\mathcal L_{d_i\omega_i}\), the dual of the tautological line with weight \(-d_i\omega_i\). The intersection of the \(P_i^-\)'s is \(B^-\) as subgroup schemes (the standard parabolic description, Result~\ref{res:D05}). It follows that the product orbit map from the projective \(Y\) is again a proper monomorphism, hence a closed immersion. Choose a common multiple of the \(d_i\)'s. A positive tensor power of \(\mathcal L_\eta\) is now the restriction of an external product of positive powers of the hyperplane bundles, which is ample by the Segre and Veronese embeddings. Thus \(\mathcal L_\eta\) is ample.

In positive characteristic put \(q=p^r\) and let \(F_r\) be
the relative Frobenius on \(Y\). Identify the target twist
with the split flag variety. Its pullback of \(\mathcal L_\nu\)
is \(\mathcal L_{q\nu}\) by multiplying the transition characters
by \(q\). Put \(\sigma_r=(q-1)\rho\) and \(N=\dim Y\).
The preceding proof shows \(S_r=L(\sigma_r)\) is simple
on the group Frobenius kernel and has dimension \(q^N\).

Via Result~\ref{res:E05}, regard \(S_r\) as a submodule of \(H^0(Y,\mathcal L_{\sigma_r})\) on the source of \(F_r\). Here is the resulting sheaf argument. The evaluation \(S_r\otimes\mathcal O_Y\to\mathcal L_{\sigma_r}\) is nonzero at every geometric point: the submodule is nonzero and the source group acts transitively on the flag variety. Adjunction gives a map \(S_r\otimes\mathcal O_Y\to F_{r*}\mathcal L_{\sigma_r}\) on the target. The kernel of the group Frobenius acts trivially on the target and the map is equivariant for this kernel. At each geometric point of the target, the induced fiber map is nonzero, since its composite with evaluation at the unique point of the Frobenius fiber is the nonzero evaluation on the source. It is injective by kernel-simplicity. Relative Frobenius is finite flat of degree \(q^N\), as can be checked on translates of the affine-space open cell, where its coordinate functions are \(q\)-th powers. Both fibers therefore have dimension \(q^N\), so the fiber maps are isomorphisms. It follows that the map of vector bundles is an isomorphism.

The projection formula now identifies
\[
 F_{r*}\mathcal L_{\sigma_r+q\lambda}
 \simeq S_r\otimes\mathcal L_\lambda.
\]
Since \(F_r\) is finite, cohomology on the left is that of
\(\mathcal L_{\sigma_r+q\lambda}\) on the source, namely
\(\mathcal L_{-\rho}\otimes\mathcal L_{\lambda+\rho}^{\,q}\).
The twist \(\mathcal L_{\lambda+\rho}\) is ample, since
\(\langle\rho,\alpha^\vee\rangle=1\) on the simple roots:
a simple reflection permutes the positive roots other than its
own, so \(s_\alpha\rho=\rho-\alpha\).
A coherent sheaf tensored by high powers of an ample line
bundle has no positive cohomology here. One may check the
needed assertion by taking a very ample power and using
its projective embedding. On projective space a coherent
sheaf has a surjection from finitely many \(\mathcal O(-a)\)'s:
take finite generators on the standard affine charts and
multiply by high powers of the chart coordinates to clear
denominators. Repeat with the kernels, which are coherent.
The standard affine cover computes cohomology and gives no
terms above the dimension of projective space. The Čech
complex of \(\mathcal O(d)\) splits into monomial complexes
indexed by exponent vectors of total degree \(d\). For
positive cohomology only the top degree can survive, and
then all exponents must be negative: when a coordinate has
nonnegative exponent, adjoining that chart index contracts
the monomial complex. Thus the positive cohomology
vanishes for \(d\ge0\).
Applying the exact sequences of the finite initial part
of the resolution, up to the dimension bound, gives the
required vanishing for every sufficiently high twist.
If a very ample power was used, apply this to the finitely
many residual powers as twists of the original sheaf.
This proves
\(S_r\otimes H^i(Y,\mathcal L_\lambda)=0\) for \(i>0\)
when \(r\) is sufficiently large, hence the positive
characteristic assertion. Compare \cite[Theorems 2.1
and 3.1]{haboush}.

For characteristic zero, reduce the same split flag
calculation modulo large primes. Use the integral model and
the characteristic-zero identification in Result~\ref{res:C08}.
Over \(\QQ\), embed the
flag in a projective representation using the line
stabilizer construction of Results~\ref{res:A15}
and~\ref{res:D01}. Over \(A=\ZZ[1/m]\), after increasing
\(m\), the closure of this projective model is a flat
smooth projective scheme \(\mathcal Y\) with the group
action and the rational base point extended.
Here one spreads the finitely many matrix coefficients
and equations and inverts their denominators. The flat
model is the schematic closure, whose local rings have
no \(A\)-torsion. The nonsmooth locus is closed, avoids
the generic fibre, and its image is a finite set of
primes by properness; remove it. The stabilizer of the
base point agrees with the split negative Borel after
inverting more denominators, since both defining ideals
are finitely generated and agree over \(\QQ\).
The line bundle and its equivariance extend after
inverting finitely many primes as well: on a finite
affine trivializing cover extend the transition functions
and their inverses, and clear denominators in the
compatibility and action equations. The weight on the
base point remains \(\lambda\) by the same comparison.

To see that the special fibres are precisely the flag
varieties, first note that coherent cohomology of this
projective scheme over \(A\) is finite over \(A\).
The projective-space argument just given proves finiteness
as well: for fixed degree, the monomial calculation for
\(\mathcal O(d)\) has only finitely many surviving monomials
(all exponents nonnegative in degree zero or all strictly
negative in top degree), and exact sequences along a
finite initial resolution reduce the general case to these.
A finite affine Čech cover for any locally free sheaf
on \(\mathcal Y\) has \(A\)-flat terms, and its reduction
at a prime \(\ell\nmid m\) computes the fibre cohomology.
The short exact sequence given by multiplication by \(\ell\)
on that complex gives
\[
 0\longrightarrow H^i(\mathcal Y,\mathcal F)/\ell
   \longrightarrow H^i(\mathcal Y_{\mathbf F_\ell},
                          \mathcal F_{\mathbf F_\ell})
   \longrightarrow H^{i+1}(\mathcal Y,\mathcal F)[\ell]
   \longrightarrow0 .
\]
With \(\mathcal F=\mathcal O_{\mathcal Y}\), remove the
finitely many torsion primes in the first two groups.
The generic \(H^0\) is one-dimensional, so every remaining
fibre has \(H^0=\mathbf F_\ell\), also after extension to
the algebraic closure. Smoothness thus makes those
geometric fibres connected and irreducible.
Their dimension is the generic dimension: smooth relative
dimension is locally constant and any part of a different
dimension has image away from the generic point, which
can again be removed by properness.
The orbit of the base point of each geometric fibre is
a closed flag subvariety by its Borel stabilizer,
Result~\ref{res:D01}, and properness.
It has the full dimension by the common root datum,
so this closed immersion is the whole smooth geometric
fibre. The equivariant line bundle on it is the
associated line bundle for its stabilizer character
\(\lambda\), as in Result~\ref{res:E04}.
Thus the positive-characteristic vanishing applies
at every remaining prime. Taking the above cohomology
sequence with the extended line bundle in positive degree
shows that each finite \(A\)-module \(H^i\) has zero
free rank. Its generic fibre vanishes, and field extension
now proves the assertion in characteristic zero.
\end{proof}

\begin{result}{E11}{Weyl character and dimension formulas}
For a dominant weight \(\lambda\), the induced and Weyl modules have the same formal character in every characteristic:
\[
 \operatorname{ch}_T\nabla(\lambda)
 =\operatorname{ch}_T\Delta(\lambda)
 =\frac{A(e^{\lambda+\rho})}{A(e^\rho)}.
\]
Their common dimension is
\[
 \prod_{\alpha\in\Phi^+}
 \frac{\langle\lambda+\rho,\alpha^\vee\rangle}
      {\langle\rho,\alpha^\vee\rangle}.
\]
These are formulas for the simple \(L(\lambda)\) when the induced module is simple, in particular in characteristic zero.
\end{result}
\begin{proof}
Result~\ref{res:E10} identifies the Euler characteristic with the character of \(\nabla(\lambda)\). The opposition involution \(-w_0\) fixes \(\rho\), permutes the positive roots, and conjugates the Weyl group preserving the sign character. It carries the displayed quotient for \(\lambda^*\) to that for \(\lambda\). On the invariant character ring its effect agrees with the duality operation \(e^\mu\mapsto e^{-\mu}\), because the additional action of \(w_0\) is trivial. Result~\ref{res:E08} gives the character of \(\Delta(\lambda)\).

For the dimension, fix a regular rational cocharacter \(h\) and apply \(e^\mu\mapsto\exp(t\langle\mu,h\rangle)\) as a formal power series over \(\mathbb Q\). Put \(N=|\Phi^+|\). At degree \(j\) in \(A(e^\nu)\), the coefficient is
\[
 \frac1{j!}\sum_{w\in W}(-1)^{\ell(w)}\langle w\nu,h\rangle^j ,
\]
an alternating homogeneous polynomial of degree \(j\) in \(\nu\). It is divisible by the defining linear form \(\langle\nu,\alpha^\vee\rangle\) of each root reflection hyperplane. For \(j<N\) the coefficient is therefore zero, and for \(j=N\) it is a scalar (depending on \(h\)) times the product of these linear forms. The scalar is nonzero since \(A(e^\rho)=e^\rho\prod_{\alpha>0}(1-e^{-\alpha})\) has order \(N\) for regular \(h\). Taking the ratio of the degree-\(N\) coefficients yields the dimension product. This computation evaluates the regular character at the identity without substituting directly into its displayed quotient. The empty-root product gives \(1\), as for a torus. See \cite[22.50--22.52]{milne} and \cite[Part~II, Chapter~5]{jantzen}.
\end{proof}

In type \(A_1\) the quotient for \(\mathrm{SL}_2\) and \(m\geq0\) is
\(e^m+e^{m-2}+\cdots+e^{-m}\), of dimension \(m+1\). The computation is valid for \(\nabla(m)\) in every characteristic. Result~\ref{res:E14} will exhibit a simple socle of dimension \(2\) in the case \(m=p>0\), where the induced dimension is \(p+1\).

\section{Matrix examples, differentiation, and Frobenius}

\begin{result}{E12}{Rational and polynomial representations of \(\mathrm{GL}_n\)}
For \(n\geq1\), with upper triangular Borel subgroup, the dominant weights of \(\mathrm{GL}_n\) are integer tuples \(\lambda_1\geq\cdots\geq\lambda_n\). For any finite-dimensional rational representation, some twist by \(\det^a\) with \(a\geq0\) is polynomial. A simple module \(L(\lambda)\) is polynomial precisely when \(\lambda_n\geq0\). Scalar matrices act on \(L(\lambda)\) with degree \(\sum_i\lambda_i\). More generally the central degree-\(d\) summand of a polynomial representation has polynomial coefficients homogeneous of degree \(d\).
\end{result}
\begin{proof}
The characters of the diagonal torus are freely generated by its projections \(\varepsilon_i\). The simple coroot between adjacent coordinates gives the difference \(\lambda_i-\lambda_{i+1}\). The coordinate ring is \(k[x_{ij},\det^{-1}]\), so take a common \(\det^a\) denominator for the finitely many matrix coefficients. The twist adds \(a\) to every coordinate of every weight. A representation with polynomial coefficients extends to a morphism from the full matrix monoid to the endomorphism monoid: the multiplicativity identity holds on the scheme-theoretically dense open of pairs of invertible matrices and hence everywhere. Thus restricting to diagonal matrices exhibits nonnegative exponents on all weights whenever the representation is polynomial.

For the converse, let \(V=k^n\) be the natural module. In \(\bigwedge^j V\), the vector \(e_1\wedge\cdots\wedge e_j\) is primitive of weight \(\varepsilon_1+\cdots+\varepsilon_j\). For \(j<n\) take \(\lambda_j-\lambda_{j+1}\) tensor copies of this module, and take \(\lambda_n\) copies of the determinant line. The tensor product is polynomial and the tensor of the indicated vectors has weight \(\lambda\). Result~\ref{res:E03} gives \(L(\lambda)\) as the simple quotient of its cyclic submodule. Stable subspaces and quotients of a polynomial module remain polynomial: in an adapted basis the off-block coefficients that vanish on \(\mathrm{GL}_n\) vanish as polynomials on all matrices, and the induced diagonal blocks are polynomial. Finally the scalar torus is central and its weight spaces give a decomposition into group-stable summands. On \(L(\lambda)\) the central degree is \(\sum_i\lambda_i\) because every root vanishes on scalars. In degree \(d\), the matrix identity \(r(zX)=z^d r(X)\) over \(k[z,z^{-1}]\) shows by coefficient comparison that each polynomial matrix entry is homogeneous of degree \(d\). See the weight convention in \cite[22.32]{milne}.
\end{proof}

For example, the dual of \(k^n\) has highest weight \((0,\ldots,0,-1)\) when \(n>1\) (weight \(-1\) for \(n=1\)). Twisting by \(\det\) replaces its inverse-matrix coefficients by cofactors. The comparison with Chapter~4 of the companion \emph{Advanced Finite Group Representation Theory} is as follows. Over \(\mathbb C\), in homogeneous degree \(d\), the weights allowed by the polynomial criterion are exactly the partitions of \(d\) of length at most \(n\), the same indexing as for the multiplicity modules in its Schur--Weyl decomposition (Results D04--D05 and D11). Here the classification of the polynomial simple modules and the degree decomposition hold over any field, while the cited complete reducibility and double-centralizer arguments in that volume are stated over \(\mathbb C\). In positive characteristic a homogeneous polynomial representation need not be semisimple, as shown by the same symmetric-power construction used for \(\mathrm{SL}_2\) in Result~\ref{res:E14}.

\begin{result}{E13}{Differentiation and the integral lattice}
For a connected semisimple group \(H\) in characteristic zero, differentiation from rational \(H\)-modules to modules of \(\mathfrak h=\operatorname{Lie}(H)\) is fully faithful. It is essentially surjective if \(H\) is simply connected. For \(H\) split, a simple Lie algebra module integrates precisely when its highest weight belongs to the actual group character lattice. In characteristic \(p>0\), pullback along relative Frobenius can give a nontrivial rational representation with zero differential.
\end{result}
\begin{proof}
If a linear map \(f\colon V\to W\) intertwines Lie actions, its graph in \(V\oplus W\) is Lie-stable. The smooth stabilizer argument in Result~\ref{res:E07} makes it group-stable, so \(f\) intertwines group actions. The converse follows by differentiation.

Suppose first that \(H\) is split and choose \(T\subset B\).
The proof in Result~\ref{res:E07} gives a semisimple Lie algebra
and identifies its weight spaces for \(\operatorname{Lie}(T)\).
The differentials \(h_i\) of the simple coroots form a basis of
\(\operatorname{Lie}(T)\) in characteristic zero, since the
simple coroots span its rational cocharacter space. For each \(i\)
the associated rank-one matrix subgroup gives the
\(\mathfrak{sl}_2\) relations for \(e_i,f_i,h_i\).
Apply the rank-one decomposition proved in Theorem C06
of \emph{Lie Algebras and Finite-Dimensional Representations},
Chapter 3. After algebraic closure \(h_i\) is diagonalizable with
integer eigenvalues. The same minimal polynomial splits with
distinct roots in \(k\), so this diagonalization holds already
over any split field. The \(h_i\) commute, and therefore act by a
simultaneous weight decomposition. A maximal weight in a nonzero
finite-dimensional module is killed by the positive root spaces,
since a root operator adds its root to the weight. Its value on
each simple coroot is a nonnegative integer, by the highest
weight line in the rank-one strings.

For uniqueness in a simple module, use the triangular
Lie-algebra decomposition into the negative root spaces, the
toral space and the positive root spaces. The sum of the positive
spaces with the toral space is a Lie subalgebra, since the
bracket adds weights. Induce the one-dimensional module on
which the positive part is zero and the toral part acts by
the chosen highest weight \(\lambda\). Ordered monomials
in the negative root spaces form its basis by the PBW theorem
proved in Theorem A05 of the same companion volume, Chapter 1.
The \(\lambda\)-weight space is one dimensional, since all
nonconstant monomials lower the weight by a nonzero sum of
positive roots. A proper submodule misses that line because
the line generates the induced module. Each submodule is a sum
of weight spaces: on any vector only finitely many weights
occur, and interpolation in a suitable toral element recovers
each component. The sum of all proper submodules still misses the
highest line and is thus the unique maximal proper submodule.
It follows that a simple finite-dimensional highest weight
module is unique for its highest weight over the given split
field. No integration of Lie modules is used in this argument.

The simple rational \(H\)-module of highest weight \(\lambda\)
has simple differential, by the smooth stabilizer argument
of Result~\ref{res:E07}. Its highest Lie weight is \(d\lambda\),
since the simple coroots of \(T\) and the positive root
parametrizations differentiate to the ones just used. In
characteristic zero the differential embeds the character
lattice in the toral dual and preserves the integer coroot
pairings. Result~\ref{res:E06} supplies the rational simple
module for every dominant weight in this lattice. The uniqueness
argument proves the claimed integrability test. For simply
connected \(H_{\bar k}\), Result~\ref{res:C07} identifies
the character lattice with the full weight lattice. Thus
over \(\bar k\) every simple Lie module integrates, and then
Weyl complete reducibility from Result~\ref{res:E07} gives
integration of any finite-dimensional Lie module.

Over a general characteristic-zero field, an integration on \(V_{\overline k}\) has matrix coefficients defined over a finite Galois extension \(K/k\), since there are finitely many coefficients. Each Galois conjugate of the integration on \(V_K\) has the same differential. By full faithfulness, the identity on \(V_K\) intertwines the two representations, so they coincide. The matrix coefficients are Galois invariant, hence defined over \(k\), and the action identities descend. This proves essential surjectivity in the simply connected case
over \(k\). When \(H\) is split, the preceding simultaneous weight
and uniqueness arguments apply directly over \(k\), so the
criterion uses the actual \(k\)-character lattice. See also
\cite[22.12 and Theorem~22.53]{milne}.

In characteristic \(p\), the relative Frobenius \(F\colon H\to H^{(p)}\) has zero differential, since on functions it raises coordinates to \(p\)-th powers after scalar twist. For a split model the target twist is identified with the same split group. The pullback of its standard \(\mathrm{SL}_2\)-action along \(F\) is nontrivial on the diagonal torus, where its weights are \(p\) and \(-p\) as characters. Its differential is zero. The example is a rational group action even over a finite field.
\end{proof}

\begin{result}{E14}{A nonsplit symmetric power in characteristic \(p\)}
Let \(k\) have characteristic \(p>0\). In \(\operatorname{Sym}^p(k^2)\) for \(\mathrm{SL}_2\), the subspace \(S=kx^p+ky^p\) is the Frobenius pullback of the natural module and has no equivariant complement. This symmetric power is \(\nabla(p)\).
\end{result}
\begin{proof}
Use the standard basis \(x,y\) of \(V=k^2\) with upper unipotent subgroup fixing \(x\) and sending \(y\) to \(y+c x\). The binomial identity \((a x+b y)^p=a^p x^p+b^p y^p\) holds over every \(k\)-algebra, so \(S\) is stable and has the Frobenius matrices. The natural module is simple: any proper nonzero submodule would be a torus weight line, and neither line is fixed by both opposite unipotent subgroups. Relative Frobenius is faithfully flat, so its pullback \(S\) is simple as well.

The monomials \(x^{p-i}y^i\), \(0\leq i\leq p\), have distinct torus characters \(t\mapsto t^{p-2i}\). A torus-equivariant complement to \(S\) must therefore equal the sum of the intermediate weight lines. In particular it contains \(x^{p-1}y\). Its image under the upper unipotent is \(x^{p-1}y+c x^p\), which is outside the complement for \(c=1\). This proves nonsplitting for all \(p\), including \(p=2\). The weight characters are homomorphisms of tori as group schemes, so they remain distinct regardless of the size of \(k\).

The flag variety \(\mathrm{SL}_2/B^-\) is \(\mathbb P(V)\) with base point the line \(ky\), whose weight is \(-1\). Thus \(\mathcal O(1)\) with its natural linearization is \(\mathcal L_1\), and its \(p\)-th power has global sections \(\operatorname{Sym}^p(V^\vee)\). The invariant determinant pairing identifies \(V^\vee\simeq V\), which identifies these sections with the stated module. It follows that \(\nabla(p)\) of dimension \(p+1\) has the simple socle \(S=L(p)\) of dimension \(2\). Compare \cite[12.56 and 22.46]{milne}.
\end{proof}

\begin{result}{E15}{Steinberg's tensor product formula}
Suppose \(k\) is algebraically closed of characteristic \(p>0\) and \(G\) is connected, semisimple and simply connected. A dominant weight has a unique expansion
\[
 \lambda=\sum_{i\geq0}p^i\lambda_i,\qquad
 0\leq\langle\lambda_i,\alpha^\vee\rangle<p
 \quad(\alpha\in\Pi)
\]
with finitely many nonzero dominant digits. Write \(M^{[i]}\) for the pullback of the split-model representation \(M\) along the \(i\)-fold Frobenius. Then
\[
 L(\lambda)\simeq\bigotimes_{i\geq0} L(\lambda_i)^{[i]} .
\]
\end{result}
\begin{proof}
The digits are obtained by taking the base-\(p\) expansion of the coefficient of each fundamental weight. Since \(G\) is simply connected and semisimple, these weights form a basis of \(X\), giving uniqueness including the trivial digit.

The Frobenius-kernel calculation before
Result~\ref{res:E10} proves that \(S=L(\eta)\) for a
\(p\)-restricted dominant \(\eta\) remains simple on \(G_1\).
We use this with \(q=p\).

Let \(M=L(\mu)^{[1]}\). The relative Frobenius \(F\) is finite faithfully flat: over the algebraic closure, translates of \(U^-\times T\times U\) give local coordinates in which \(F\) raises the affine and invertible coordinates to \(p\)-th powers. Thus \(M\) is simple as a \(G\)-module because a subspace stable under its pullback action is stable under the original action. On \(G_1\), \(M\) is trivial. Thus \(S\otimes M\) restricts to a direct sum of copies of the simple \(G_1\)-module \(S\). A \(G\)-submodule \(N\) of \(S\otimes M\) also restricts to such a direct sum. Schur's lemma over the algebraically closed field gives \(\operatorname{End}_{G_1}S=k\), so
\[
 S\otimes \operatorname{Hom}_{G_1}(S,N)\longrightarrow N
\]
is an isomorphism by evaluation. Normality of \(G_1\) makes \(\operatorname{Hom}_{G_1}\) here stable under the natural action of \(G\) on Hom. The identification
\(\operatorname{Hom}_{G_1}(S,S\otimes M)=M\), given by \(m\mapsto(s\mapsto s\otimes m)\), is \(G\)-equivariant. Thus if \(N\ne0\), the Hom space from \(S\) into \(N\) equals \(M\) by simplicity, and \(N=S\otimes M\). The tensor product is therefore simple. A tensor of highest vectors has weight \(\eta+p\mu\), because Frobenius multiplies torus characters by \(p\). It is primitive because both factors are. By Result~\ref{res:E06} this simple tensor product is \(L(\eta+p\mu)\). Repeating the argument recursively proves the displayed formula, as in \cite[Part~II, Sections~3.16--3.17]{jantzen}. In the rank-one example of Result~\ref{res:E14}, the only nonzero digit of the weight \(p\) is the weight \(1\) in place \(1\), recovering the Frobenius submodule there.
\end{proof}
\chapter{Galois descent, forms, and arithmetic}
\label{ch:descent}

Let $k_s$ be a separable closure of $k$ and let $\Gamma_k=\Gal(k_s/k)$, with its Krull topology. If $K/k$ is finite Galois, write $\Gamma=\Gal(K/k)$. A semilinear action on a $K$-space $V$ satisfies $s_\sigma(av)=\sigma(a)s_\sigma(v)$ and $s_{\sigma\tau}=s_\sigma s_\tau$. The first two results use only linear algebra, field theory and the Hopf description in Result~\ref{res:A02}. They can be read before the classification of tori.

\section{Effective descent}

\begin{result}{F01}{Semilinear descent for vector spaces}
Let $K/k$ be finite Galois. For a $K$-vector space $V$ with semilinear $\Gamma$-action, the canonical map
\[
 K\otimes_k V^\Gamma\longrightarrow V,\qquad a\otimes v\longmapsto av
\]
is an isomorphism. Equivariant $K$-linear maps correspond to $k$-linear maps of fixed spaces. This equivalence preserves tensor products and the tensor unit. In finite dimension, $\dim_k V^\Gamma=\dim_K V$. The descent construction is valid without a dimension bound.
\end{result}
\begin{proof}
For $a\in K$ and $v\in V$, the vector $\sum_{\sigma\in\Gamma}\sigma(a)s_\sigma(v)$ is fixed, by reindexing the finite sum. If the $K$-span of the fixed vectors were proper, choose a nonzero $K$-linear functional $\ell$ vanishing on that span. For each $v$ and every $a$ we would have
$\sum_\sigma \ell(s_\sigma(v))\sigma(a)=0$. Independence of the distinct field automorphisms over $K$ gives $\ell(v)=0$ from the coefficient of the identity automorphism. Recall the independence argument: take a relation of smallest nonzero length, evaluate at $ba$ and subtract one of the values on $b$ times the relation at $a$. Two distinct automorphisms have different values at some $b$, giving a shorter nonzero relation. Thus the map is surjective, and no division by $|\Gamma|$ is needed.

To prove injectivity, let $v_1,\dots,v_m$ be $k$-independent fixed vectors. In a putative shortest $K$-linear dependence normalize its first coefficient to $1$. Subtract the relation obtained by applying $s_\sigma$. Minimality forces every coefficient to be fixed by $\Gamma$, contradicting $k$-independence. Conversely, in the extension of a $k$-space, write a vector using a $k$-basis. Invariance of that vector is precisely invariance of each coefficient. An equivariant map is determined by its restriction to the fixed vectors, and every map of fixed spaces extends. Fixed bases for $V$ and $W$ give a fixed basis of pure tensors for $V\otimes_K W$. This identifies $(V\otimes_K W)^\Gamma$ with $V^\Gamma\otimes_k W^\Gamma$ and sends the unit $K$ to $k$. See \cite[Lemma 2.3.8]{gs}, \cite[Proposition 3.36]{milne}, and \cite[Tag 0CDR]{stacks}.
\end{proof}

\begin{result}{F02}{Effective descent for affine groups}
Suppose an affine $K$-group has semilinear Galois descent data. Its coordinate Hopf algebra $A$ descends to the Hopf algebra $A^\Gamma$ over $k$, and equivariant morphisms descend uniquely. Finite type and smoothness descend, and stable Hopf ideals correspond to closed subgroup schemes of the descended group.
\end{result}
\begin{proof}
Write the datum as a semilinear action on $A$ compatible with its multiplication, unit, comultiplication, counit and antipode, using the inverse pullbacks if the datum is given on the group. Result~\ref{res:F01} descends the five maps. In particular, the image of a fixed element under comultiplication belongs to $(A\otimes_K A)^\Gamma=A^\Gamma\otimes_k A^\Gamma$. The Hopf identities descend because extension of scalars is faithful. The same argument descends a morphism of Hopf algebras and reverses its direction upon applying $\Spec$.

For finite type, span the Galois translates of a finite set of algebra generators. The resulting finite-dimensional stable space has a fixed basis, which generates $A$ after extension. The $k$-subalgebra of $A^\Gamma$ generated by that basis equals $A^\Gamma$: the quotient of the underlying vector spaces vanishes after tensoring with $K$. For smoothness apply the geometric regularity criterion over an algebraic closure containing $K$ to the descended algebra of finite presentation. A stable Hopf ideal descends by Result~\ref{res:F01}. Its ideal and coideal identities can be checked after extension, including exactness for the quotient. Conversely, extending a Hopf ideal of $A^\Gamma$ gives a stable ideal. See \cite[Sec.~3.j]{milne}, \cite[Proposition 2.15]{prr}, and \cite[Tag 0CDQ]{stacks}.
\end{proof}

For a discrete $\Gamma_k$-module of groups $M$, a continuous cocycle is a map $a_\sigma\in M$ with $a_{\sigma\tau}=a_\sigma\,\sigma(a_\tau)$. We use the equivalence
\[
 a'_\sigma=b^{-1}a_\sigma\,\sigma(b).
\]
The quotient $H^1(\Gamma_k,M)$ is a pointed set, with base point the class of the constant identity cocycle. For a smooth $k$-group $H$ we write $H^1(k,H)=H^1(\Gamma_k,H(k_s))$.

\begin{result}{F03}{Forms and twisting}
Let $G_0$ be an affine $k$-group of finite type. For finite Galois $K/k$, the $k$-isomorphism classes of groups becoming isomorphic to $(G_0)_K$ are classified by
$H^1(\Gamma,\Aut_K((G_0)_K))$. Replacing $K$ by $k_s$ gives the analogous classification for forms isomorphic to $G_0$ over $k_s$, with continuous cocycles and discrete coefficients.
\end{result}
\begin{proof}
Identify the two $K$-groups by an isomorphism $i\colon (G_0)_K\to G_K$. The semilinear action transported to $(G_0)_K$ has form $a_\sigma s_\sigma^0$ for its standard action $s_\sigma^0$. Composing two such semilinear maps gives the cocycle law for conjugation by $s_\sigma^0$. Replacing $i$ by $ib$ gives the displayed equivalence. Conversely $a_\sigma s_\sigma^0$ gives a datum and Result~\ref{res:F02} descends it. An isomorphism of descents becomes an intertwiner of the two data after extension, precisely the same equivalence on cocycles.

In the $k_s$ case each isomorphism and its inverse are specified by finitely many coefficients in a finite Galois extension. Conversely a continuous cocycle has finite image by compactness. Its identity preimage is an open subgroup, and a common field of definition for its finitely many values can be chosen finite Galois. Shrink to an open normal subgroup of the identity preimage fixing that field. The cocycle is constant on its right cosets and defined over the corresponding finite extension. An equivalence between two cocycles descends over a common finite Galois extension containing the coefficients of the isomorphism and its inverse. Thus both existence and equality of classes reduce to finite-level descent. Compare \cite[Secs.~2.3.2--2.3.3 and Lemma 2.3.9]{gs} and \cite[Secs.~3.43, 3.49]{milne}.
\end{proof}

\section{Torsors and examples of forms}

A \emph{right $H$-torsor} becomes isomorphic, with right $H$-action, to $H$ with its translation action after a faithfully flat base change of finite presentation. Over a field, for an affine group one can check this by nonemptiness, finite type and the isomorphism $X\times H\to X\times X$, $(x,h)\mapsto(x,xh)$.

\begin{result}{F04}{Matrix Hilbert 90 and smooth torsors}
For $n\geq 1$, $H^1(k,\GL_n)=H^1(k,\SL_n)=1$. The analogous finite Galois cohomology sets for $\GL_n(K)$ and $\SL_n(K)$ are singletons. For a smooth affine $k$-group $H$ of finite type, $H^1(k,H)$ classifies its fppf torsors up to isomorphism. A torsor is trivial exactly when it has a $k$-point.
\end{result}
\begin{proof}
For a matrix cocycle $A_\sigma$, act on column space by $v\mapsto A_\sigma\sigma(v)$. A fixed basis by Result~\ref{res:F01} gives an invertible matrix $M$ with $A_\sigma\sigma(M)=M$. Thus $A_\sigma=b^{-1}\sigma(b)$ for $b=M^{-1}$. If $\det A_\sigma=1$, then $\det M\in k^\times$. Multiply $M$ on the right by a $k$-diagonal matrix of determinant $(\det M)^{-1}$ to make $b$ special linear. The same proof applies at a finite level in the separable closure case.

A smooth-group torsor is smooth by faithfully flat descent of smoothness. Over the separably closed field $k_s$, a nonempty smooth scheme has a point: an étale chart has nonempty open image in affine space, which contains a rational point, and its nonempty étale fiber has a rational point. Choose $p$ over $k_s$ and write $\sigma(p)=p h_\sigma$. This determines a continuous cocycle, with the change of $p$ giving the equivalence from Result~\ref{res:F03}. Conversely, act semilinearly on $H_{k_s}$ by $x\mapsto h_\sigma\,\sigma(x)$ and descend the resulting affine scheme by the algebra argument of Result~\ref{res:F02}. Right translation descends, as does the defining torsor isomorphism (test it after extension). A $k$-point gives the equivariant isomorphism $h\mapsto ph$. For a non-smooth $H$ the fppf classification is needed in general. In characteristic $p$, $x^p=a$ for $a\notin k^p$ gives a $\mu_p$-torsor without a $k_s$-point. See \cite[Example 2.3.4]{gs} and \cite[Secs.~3.38--3.40, 3.46--3.50]{milne}.
\end{proof}

\begin{result}{F05}{Central simple algebras and inner forms of type A}
Let $K/k$ be finite Galois and $n\geq 1$. The central simple $k$-algebras of degree $n$ split by $K$ are classified up to algebra isomorphism by $H^1(\Gamma,\operatorname{PGL}_n(K))$. Their twists of $\SL_n$ and $\operatorname{PGL}_n$ by conjugation are the reduced-norm-one group $\operatorname{SL}_1(A)$ and the adjoint group $\operatorname{PGL}_1(A)$. The cocycles remember the inner twisting data, which can be lost when passing to the isomorphism class of the group alone.
\end{result}
\begin{proof}
Every $K$-algebra automorphism of $M_n(K)$ is conjugation by an invertible matrix. Indeed, twisting the natural $n$-dimensional module by such an automorphism gives a simple module. The matrix units show uniqueness of the simple module: for a nonzero vector in the image of the first diagonal idempotent, the first-column matrix units give the natural basis of a simple submodule. An isomorphism of simple modules is an intertwining invertible matrix. Its ambiguity is exactly a scalar, the commutant of all the matrix units. This identifies the automorphism group on $K$-points with $\operatorname{PGL}_n(K)$, compatibly with Galois action.

Descending a matrix algebra uses the multiplication part of the argument in Result~\ref{res:F02}. The center after extension is the extension of the center (take the joint kernel of commutators with a finite basis). A nontrivial proper two-sided ideal would extend to such an ideal in a matrix algebra. This proves central simplicity, and dimension gives the degree. Conversely, a splitting of a central simple algebra gives the trivialization for Result~\ref{res:F03}'s cocycle argument. Determinant is invariant under conjugation, so its coordinate polynomial descends to the reduced norm. Its norm-one locus identifies equivariantly with $\SL_n$ over $K$. Quotienting the units by the scalars as fppf sheaves commutes with extension and descends the adjoint group. For $n\geq3$, transpose-inverse gives an outer symmetry. In particular, inversion into the opposite algebra induces a group isomorphism $\operatorname{SL}_1(A)\to\operatorname{SL}_1(A^{\mathrm{op}})$, illustrating how the group can forget the inner data. See \cite[Secs.~2.4.1--2.4.3]{gs} and \cite[Secs.~24.22--24.25 and 24.f]{milne}.
\end{proof}

\begin{result}{F06}{Outer data and the quasi-split inner form}
Let $G_0$ be split connected reductive over $k$ with pinned based root datum $\Psi^+$. Its fppf automorphism sequence has kernel $(G_0)_{\mathrm{ad}}$ and quotient the constant étale group $\Aut(\Psi^+)$, with a pinning giving a section. A form of $G_0$ has the twisted outer sheaf and the corresponding adjoint kernel. Two classes in $H^1(k,\Aut(G_0))$ are inner forms of each other exactly when their images in $H^1(k,\Aut(\Psi^+))$ agree. Each such fiber has exactly one quasi-split isomorphism class.
\end{result}
\begin{proof}
The adjoint quotient is an affine group, and conjugation
gives it as the sheaf of inner automorphisms: a lift from
\(G_0/Z(G_0)\) acts by conjugation and two lifts differ
by a central point. Its map into the automorphism sheaf
is injective by the definition of the scheme centre.

For an automorphism over a test algebra \(R\), its image
of the pinned Borel pair can be brought back to \((B,T)\)
fppf locally. For the Borel this is the family statement in
Result~\ref{res:D08}. For tori inside \(B\), use the analogous
transporter argument there. The transporter is cut out by
including the image of the fixed split \(T\) into the other
maximal torus; inclusion is equality by flatness and equality
on geometric fibres. The equations on the coefficients of
\(k[T]\) give a finitely presented scheme. Its geometric
fibres are nonempty by conjugacy of maximal tori in
\(B_{\bar k}\) from Result~\ref{res:B13}.
For a square-zero lift, the displacement is a \(T\)-cocycle
with values in \(\Lie(B)/\Lie(T)\) tensored with the
square-zero ideal, by the conormal calculation of
Result~\ref{res:D08}. The regular-cochain averaging in
Result~\ref{res:B13} for the split torus works over the
test ring and kills this cocycle by infinitesimal conjugation.
Thus the transporter is smooth and gives an fppf cover.
The ambiguity in aligning the pair is \(T\):
the Borel is its own normalizer by Result~\ref{res:D05},
and its normalizer of \(T\) equals \(T\) by projection to
\(B/R_u(B)\) and Result~\ref{res:C02}.

An aligned automorphism induces an automorphism of the
based datum by Results~\ref{res:C04} and~\ref{res:C05}.
It is locally constant, since the automorphisms of a
split torus are locally lattice automorphisms by the
group-like coefficient argument of Result~\ref{res:C02}.
The ambiguity \(T\) acts trivially on the datum, hence
the map to the constant group descends.
By the pinned lifting of Result~\ref{res:C08}, each
based automorphism has a unique pinned automorphism
over the field and therefore defines one by scalar extension
on each component. Uniqueness makes this a section
of sheaves of groups. After composing with its inverse,
the aligned automorphism fixes \(T\) and acts on each
simple root coordinate by a unit scale. Indeed a
\(T\)-equivariant map \(\Ga\to\Ga\) in the same root
character is linear: the character \(T\to\Gm\)
is faithfully flat, and equivariance for scalar
multiplication forces the polynomial to have degree one
on every test algebra. The simple root characters form
a basis of the adjoint torus character lattice by
Result~\ref{res:C05}, so a unique point of that torus
corrects all these scales by conjugation. The resulting
automorphism is the identity by the pinned faithfulness
proof in Result~\ref{res:C08}. This proves exactness
of the fppf sequence. Twisting by descent twists
both kernel and quotient, and gives the asserted
outer sheaf of a form. The constant group can be
infinite because of the central torus.
Compare \cite[Theorem 7.1.9]{conrad}.

To compare cohomology classes, after adjusting outer
trivializations suppose their cocycles \(a,b\) have the
same outer image. The section supplies the adjustment.
Write \(d_\sigma=b_\sigma a_\sigma^{-1}\) in the adjoint
kernel. The equality
\[
 d_{\sigma\tau}=d_\sigma a_\sigma \sigma(d_\tau)a_\sigma^{-1}
\]
makes \(d\) a cocycle for the twisted adjoint action.
Conversely each such cocycle multiplied by \(a\)
gives a class with the same outer image.
This is exactly the inner-form relation, and does not
assert uniqueness of a preimage in twisted adjoint
cohomology. The lift of the outer cocycle via the
pinned section fixes \(B,T\) as subgroups, so their
descent gives a \(k\)-Borel and hence a quasi-split
form \(G_q\).

For uniqueness, let an inner twist of \(G_q\) also
have a \(k\)-Borel. Over the separable closure
align its Borel with that of \(G_q\) using
Result~\ref{res:D08}. Its inner cocycle now has values
in the Borel of the adjoint group, by the scheme
normalizer statement of Result~\ref{res:D05}.
In that Borel the quotient by the unipotent radical
is the torus descending from the pinned adjoint
torus. Its character basis of simple roots is permuted
by Galois, so it is the product of the Weil restrictions
indexed by the orbits of simple roots, by
Results~\ref{res:B02} and~\ref{res:B04}. Its first
cohomology is trivial by Result~\ref{res:F07}.
The unipotent radical has a \(k\)-split series by
Result~\ref{res:D08}. The first cohomology of each
additive term is zero: on a finite Galois layer pick
\(x\) with trace one and average the values \(v_\sigma\)
of an additive cocycle as \(t=\sum_\sigma\sigma(x)v_\sigma\);
the reindexing identity gives \(v_\tau=t-\tau(t)\).
Trace is nonzero by independence of embeddings.
Induct along the series, trivializing each image
cocycle and taking adjusting lifts over the separable
closure by smoothness.
Thus the adjoint Borel cocycle is trivial, first in
the torus quotient and then in the unipotent radical.
This proves uniqueness of the quasi-split class.
Compare \cite[Section 7.2]{conrad}.
\end{proof}

\begin{result}{F07}{A norm-one torus}
For a separable quadratic extension $K/k$, let $T=\ker(N\colon\Res_{K/k}\Gm\to\Gm)$. It is a $k$-torus and $X^*(T_K)=\ZZ$ with the nonidentity Galois element acting by negation. There is a canonical isomorphism
\[
 H^1(k,T)\simeq k^\times/N_{K/k}K^\times.
\]
For $\CC/\RR$ the resulting $H^1$ has order two.
\end{result}
\begin{proof}
On the two split factors the norm from Result~\ref{res:B04} is $(x,y)\mapsto xy$. Its kernel has elements $(x,x^{-1})$, and the permutation of the factors inverts this coordinate. This proves the character statement and surjectivity of the norm on $k_s$-points.

For a finite separable field extension $E/k$ there is a useful vanishing $H^1(k,\Res_{E/k}\Gm)=1$. Over a finite Galois field $K'$ splitting $E$ and the cocycle, a cocycle descends a product of one-dimensional spaces indexed by the embeddings of $E$. More precisely, the twisted semilinear action on $E\otimes_k K'$ is multiplication by the cocycle followed by the standard semilinear action, so it commutes with the action of $E$. The fixed module has $k$-dimension $[E:k]$ by Result~\ref{res:F01} and thus is a one-dimensional $E$-module. An $E$-generator extends to a generator in every factor of $E\otimes_k K'$. Its fixedness makes the cocycle a coboundary. This argument also covers the factors used in Result~\ref{res:F06}.

Lift $a\in k^\times$ to $u\in(\Res_{K/k}\Gm)(k_s)$ with $N(u)=a$. Then $u^{-1}\sigma(u)$ lies in $T(k_s)$, and its class is independent of the lift. It is trivial exactly when $u$ can be adjusted by a $T(k_s)$-element to be fixed, which says $a$ is a norm from $K^\times$. Conversely a $T$-cocycle becomes a coboundary in the Weil restriction by the vanishing above. The norm of a witnessing element is fixed and gives the inverse. Multiplicativity follows by taking products of lifts. For $\CC/\RR$ the nonzero norms are exactly the positive reals. See \cite[Lemma 2.22 and Sec.~2.2.4]{prr} and \cite[Sec.~12.i]{milne}.
\end{proof}

\section{Finite fields and rational points}

\begin{result}{F08}{Lang's theorem}
Let $G$ be a connected group scheme of finite type over $\FF_q$. A nonempty finite-type $\FF_q$-scheme $X$ with a $G$-action transitive on geometric points has an $\FF_q$-point. If $G$ is smooth, the map on geometric points
\[
 L(g)=F(g)g^{-1}
\]
is surjective, where $F$ is the $q$-power Frobenius endomorphism of $G$ over $\FF_q$. Its fiber at $L(g)$ is $gG(\FF_q)$. Consequently a torsor under such a connected group is trivial.
\end{result}
\begin{proof}
First suppose $G$ is smooth. Over an algebraic closure it is connected: the geometric identity component descends and is open and closed. Consider the action $g*x=F(g)xg^{-1}$ on $G$. For every $x$ the orbit map has invertible differential at the identity, because $dF=0$ and the remaining term is inversion followed by translation. The orbit map is \'{e}tale there by the smooth local criterion, hence its image contains an open neighborhood of $x$. Thus every orbit of geometric points is open, and connectedness leaves one orbit. The orbit of the identity is the image of $L$. If $L(g')=L(g)$, substituting $g'=gh$ gives $F(h)=h$, so $h\in G(\FF_q)$, and these $h$ all give the same value.

For $X$ choose a geometric point $x$ and an element $h$ with $F(x)=hx$. Since $g\mapsto g^{-1}F(g)=L(g^{-1})^{-1}$ is surjective, choose $g$ with $g^{-1}F(g)=h^{-1}$. Then $F(gx)=gx$, hence $gx$ is rational over $\FF_q$. For a general $G$, replace it by its reduction. Over the perfect field $\FF_q$ the reduction is a smooth closed subgroup: products of geometrically reduced finite-type schemes are reduced, so the structure maps factor through the reduction, which is smooth by the smoothness criterion for group schemes. The underlying geometric points and connectedness are preserved. A torsor has transitive geometric points, and any rational point trivializes it. See \cite[Theorem A.1.1 and Proposition A.1.3]{conradfeng}.
\end{proof}

\begin{result}{F09}{Rational liftings and Borels over finite fields}
For a smooth affine $k$-group $G$ and a smooth closed subgroup $H$, there is a sequence
\[
 G(k)\longrightarrow (G/H)(k)\longrightarrow H^1(k,H)\longrightarrow H^1(k,G)
\]
exact at the indicated base points. Two points of $(G/H)(k)$ give the same $H^1(k,H)$ class exactly when they are in the same $G(k)$-orbit. If $k$ is finite and $H$ connected, every rational point lifts. Over finite $k$, any smooth connected affine group admits a $k$-Borel and a geometrically maximal $k$-torus. In particular a connected reductive group is quasi-split.
\end{result}
\begin{proof}
The fiber of $G\to G/H$ over a rational point is a right $H$-torsor by Result~\ref{res:D01}. For a geometric lift $g$ over $k_s$ the class is given by $g^{-1}\sigma(g)\in H(k_s)$. A trivial class gives a rational lift by Result~\ref{res:F04}. Two lifts with identical cohomology class can be adjusted inside their fibers so that their cocycles agree. Then the ratio of the adjusted lifts on the left is fixed by Galois, giving the claimed orbit description. The image in $H^1(k,G)$ is trivial since $g$ is a witness. Conversely, if an $H$-cocycle is $g^{-1}\sigma(g)$ in $G(k_s)$ then $gH$ is fixed and descends to a $k$-point. This proves exactness and the orbit claim. For $k$ finite and $H$ connected use Result~\ref{res:F08}. Disconnectedness matters, as $\Gm\to\Gm$, $t\mapsto t^2$, fails to lift a nonsquare over $\FF_q$ with $q$ odd.

Here is geometric conjugacy in the generality needed for this
finite-field application. Work first over an algebraically closed
field. Embed the smooth connected affine group faithfully into
\(\GL(V)\) by Result~\ref{res:A06}, and let it act on the full flag
variety of \(V\). This variety is projective: it is the incidence
locus defined by inclusion of consecutive subspaces in the product
of the Pl\"ucker Grassmannians.

Each orbit of geometric points is open in its closure. Indeed the
group is irreducible: its regular local rings are domains, so its
irreducible components are disjoint and its connectedness gives just
one. The closure of an orbit is therefore irreducible. On a nonempty
affine open of that closure take a nonempty affine open of its
inverse image under the orbit map. This is a dominant map of affine
varieties. Its coordinate inclusion has a nonzero free localization
by the monomial-division argument of Result~\ref{res:A07}. The
argument applies to any domain inclusion with a finitely generated
algebra. After localization, each geometric fiber is nonempty,
and a nonempty finite-type fiber over the algebraically closed
field has a rational point. Thus the orbit contains a nonempty open
subset of its closure. Its translates are open in that same closure
and cover the orbit. The boundary has strictly smaller dimension,
so an orbit of minimum dimension is closed.

For a geometric point of this closed orbit its scheme stabilizer
lies inside a triangular subgroup of \(\GL(V)\). Its reduced
identity component is consequently smooth, connected, and solvable.
A maximal smooth connected solvable subgroup of \(G\) exists by
maximizing dimension. Each such subgroup fixes a point of the
closed projective orbit by Result~\ref{res:B14}. Since it is
smooth and connected, it lies in the reduced identity component of
that point's stabilizer. Maximality gives equality, and all these
stabilizers' reduced identity components are conjugate because the
flags lie in one orbit. This proves geometric Borel conjugacy.
Likewise a geometric torus is contained in a Borel by maximizing
the dimension of a smooth connected solvable group containing it.
A maximal torus of \(G\) is maximal in its Borel. Conjugacy of
Borels and conjugacy of maximal tori inside a geometric connected
solvable group (Result~\ref{res:B13}) give geometric conjugacy
of maximal tori of \(G\).

Now return to the finite field. The base change of \(G\) is smooth
connected. Start with a geometric Borel \(B\), and let \(F\) be
the \(q\)-power Galois Frobenius on geometric points. Since
\(F(B)=hBh^{-1}\) for some \(h\), choose \(g\) with
\(g^{-1}F(g)=h^{-1}\) by Result~\ref{res:F08}. Then
\(F(gBg^{-1})=gBg^{-1}\), so \(gBg^{-1}\) descends by
Result~\ref{res:F02}. Stability under \(F\) suffices because the
subgroup equations are defined over a finite layer and Frobenius
generates its Galois group. Apply the same argument to a maximal
torus, using the geometric conjugacy just proved. A connected
reductive group having a rational Borel is quasi-split by definition.
Compare \cite[Example A.1.2 and Corollary A.1.5]{conradfeng}
and \cite[Sec.~25.a]{milne}.
\end{proof}

\section{Relative rank and arithmetic}

\begin{result}{F10}{Relative rank and anisotropy}
For a connected reductive $k$-group $G$, the $k$-rank is the dimension of a maximal $k$-split torus. For a split group it equals the geometric rank. For a semisimple group, positive $k$-rank is equivalent to the existence of a proper $k$-parabolic. For a reductive group the criterion uses the $k$-rank of $G_{\mathrm{ad}}$. A norm-one quadratic torus has geometric rank one and $k$-rank zero.
\end{result}
\begin{proof}
Conjugacy in Result~\ref{res:D08} makes the split rank
well-defined. A split torus is contained in a geometric maximal
torus, so the rank does not exceed the geometric rank, and equality
holds for a split group by definition. A cocharacter of \(G\)
gives a proper parabolic exactly if it is noncentral: nontriviality
of a root character on it gives a positive and a negative root
weight geometrically, and hence a proper dynamic subgroup.
If all root weights vanish it is central by
Result~\ref{res:C05}. Every proper \(k\)-parabolic arises from
a \(k\)-cocharacter by Result~\ref{res:D08}.

We verify the lift from the adjoint quotient, including its
possibly nonsmooth central kernel. Put \(Q=G/Z(G)\), an
affine group by Result~\ref{res:D01}; it is smooth and
connected geometrically since its coordinate algebra embeds
in that of \(G\) and the quotient is surjective.
The reduced identity component of the geometric inverse
image of a smooth connected normal unipotent group of \(Q\)
is smooth connected normal in \(G\). It projects onto the
unipotent group: the finitely many components of the reduced
inverse image are translates of the identity component,
and the irreducible unipotent group cannot be covered by
finitely many translates of a proper subgroup image.
This inverse-image component is solvable, as its quotient
is unipotent and its kernel central. In a reductive group
such a normal subgroup is a central torus by
Results~\ref{res:C01} and~\ref{res:C03}.
Its unipotent image is trivial, proving reductivity of \(Q\).

For a geometric maximal torus \(T'\) of \(Q\) its
inverse image is a torus in \(G\).
The inverse image is commutative: its commutator
descends to a morphism \(T'\times T'\to Z(G)\) by centrality,
and is a homomorphism in each argument by the commutator
identities for central commutators. A family of homomorphisms
from a split torus into a diagonalizable group is locally
constant on character lattices, by the group-like coefficient
calculation of Result~\ref{res:C02}. The parameter torus is
connected, and the family equals the trivial map at the identity.
The centre is diagonalizable here by Result~\ref{res:C05}.
Thus the commutator vanishes. The reduced identity component
of the inverse image is smooth connected commutative and projects
onto \(T'\) by the same component-image argument.
Its unipotent radical maps trivially to \(T'\) and meets the
central multiplicative-type kernel trivially by
Result~\ref{res:B07}, so this component is a torus.
Extend it to a maximal torus \(T\) of \(G\); its image is
then \(T'\) by maximality. The full centre is in \(T\)
scheme-theoretically by Result~\ref{res:C05}. Consequently the
full inverse image of \(T'\) is \(T\), since one can lift
along \(T\to T'\) fppf locally and then divide by a kernel point.
This calculation descends: the inverse image of a geometric
maximal \(k\)-torus of \(Q\) is a maximal \(k\)-torus of \(G\).

If \(Q\) has a nonzero split cocharacter, extend its image
to a geometric maximal \(k\)-torus in \(Q\) by
Result~\ref{res:C02} applied in the smooth connected
centralizer of its image. The map from the inverse-image
torus onto that torus is surjective on rational cocharacter
spaces. Averaging over a finite splitting Galois group
preserves surjectivity of this map on invariants, so
a positive multiple of the split cocharacter has a
\(k\)-cocharacter lift by Result~\ref{res:B02}.
The lift is noncentral and gives a proper parabolic.
Conversely the cocharacter of any proper parabolic has
nontrivial image in \(Q\), giving positive split rank there.
In the semisimple case the centre is finite by
Result~\ref{res:C03}. The corresponding maps on maximal-torus
rational cocharacter spaces and their invariants are then
isomorphisms, so the split ranks of \(G\) and \(Q\) agree.
Finally the norm-one torus of Result~\ref{res:F07} has
geometric rank one and its nontrivial Galois action is
negation on the cocharacter group, so its split rank is
zero. Compare \cite[Sections 3.3 and 11.1]{conradfeng}.
\end{proof}

\begin{result}{F11}{Arithmetic subgroups and integral models}
For an affine algebraic $\QQ$-group $G$, fix a closed faithful matrix realization $i\colon G\hookrightarrow\GL_N$. The commensurability class in $G(\QQ)$ of $i^{-1}(\GL_N(\ZZ))$ is independent of $i$. The analogous statement over a number field holds with a fixed ring of $S$-integers $\OO_{K,S}$ for finite $S$ containing the infinite places.
\end{result}
\begin{proof}
For another closed realization $j$, each matrix entry of $j$ and its inverse determinant pulls back to a polynomial in the matrix entries of $i$ and the inverse determinant of $i$, with rational coefficients. Choose an integer $m\geq 1$ clearing the denominators of all these polynomials. Inside $i^{-1}(\GL_N(\ZZ))$ take the subgroup congruent to the identity modulo $m$. It has finite index by reduction to the finite group over $\ZZ/m\ZZ$. The inverse determinant is congruent to $1$ as well. For each of the polynomials $f$ the integral polynomial $mf$ takes values congruent to $mf(1)$ modulo $m$, so the actual value of $f$ differs from its integral value at the identity by an integer. This puts the congruence subgroup inside $j^{-1}(\GL_{N'}(\ZZ))$. Reversing $i$ and $j$ proves finite index on both sides. Over $\OO_{K,S}$ choose a nonzero common denominator in the ring and reduce modulo the ideal it generates, whose quotient is finite. The same argument applies.

An arithmetic subgroup over $\QQ$ is a subgroup of $G(\QQ)$ in this commensurability class. For example $\SL_n(\ZZ)$ and its principal congruence subgroups are arithmetic. The latter are kernels of reduction on the former. The groups also embed in $G(\RR)$ under extension to the real completion. The term \emph{lattice} in $G(\RR)$ also imposes finite covolume, a separate assertion. For $\SL_n$ the $\QQ$-rank calculated in Result~\ref{res:F10} is $n-1$. Rank alone neither defines arithmeticity nor proves finite covolume. See \cite[Secs.~5.1 and 9.1]{morris} and \cite[Ch.~4]{prr}.
\end{proof}

\begin{result}{F12}{Adelic topology and the strong approximation criterion}
Fix a number field $K$ and a finite set of places $S$. In a closed matrix realization of $G$, use the subgroups $G(\OO_v)=G(K_v)\cap\GL_N(\OO_v)$ at finite places. The group of $S$-adeles of $G$ is the restricted product of $G(K_v)$ for $v\notin S$ with respect to these subgroups, with any archimedean factors outside $S$ given their ordinary topology. The topology is independent of the chosen closed realization and there is a diagonal homomorphism from $G(K)$. For connected absolutely almost simple $G$ over a global field, strong approximation outside a finite nonempty $S$ holds exactly when $G$ is algebraically simply connected and $\prod_{v\in S}G(K_v)$ is noncompact. The general criterion is the strong approximation theorem \cite[Theorem 2.3]{rapinchuk}. This section constructs the topology and illustrates both hypotheses.
\end{result}
\begin{proof}[Construction and approximation arguments]
The subgroup $G(\OO_v)$ is compact open at any finite place because $\GL_N(\OO_v)$ is compact open in $\GL_N(K_v)$ and $G(K_v)$ is closed. The basic open sets prescribe open subsets at finitely many places and $G(\OO_v)$ at the remaining finite places. The local operations and the subgroup conditions define a group topology. For $g\in G(K)$ the finitely many coordinates of $g$ and its inverse determinant are integral at all but finitely many places. Thus the diagonal map lands in the restricted product. For another realization the change-of-coordinates polynomials in the proof of Result~\ref{res:F11} are integral at all finite places outside one finite set. Applying the argument in both directions identifies the two prescribed local subgroups away from that set, proving independence. Weak approximation concerns projection to a finite product, while a basic open neighborhood for strong approximation also imposes integrality at all other finite places.

Here is the case $\SL_n/\QQ$ for $n\geq2$ and $S=\{\infty\}$. Fix the finite set of primes $T$ on which a neighborhood prescribes a local target. If $N=\prod_{p\in T}p$, the diagonal ring $\ZZ[1/N]$ is dense in $\prod_{p\in T}\QQ_p$: clear denominators by a high power of $N$, then solve the resulting finite system of congruences by the Chinese remainder theorem. Over any field, row reduction generates $\SL_n$ by the elementary matrices $I+aE_{ij}$, since elementary matrices also realize the required row interchanges and determinant-one diagonal steps. For the diagonal step in two coordinates, the product of transvections
\[
 w(t)=
 \begin{pmatrix}1&t\\0&1\end{pmatrix}
 \begin{pmatrix}1&0\\-t^{-1}&1\end{pmatrix}
 \begin{pmatrix}1&t\\0&1\end{pmatrix}
 \quad\text{satisfies}\quad
 w(t)w(-1)=\operatorname{diag}(t,t^{-1}).
\]
In the product of the local groups, concatenate the elementary words for each target using identity factors in the other coordinates. Approximate each tuple of parameters by a common rational in $\ZZ[1/N]$. The products approximate the local targets by continuity and are integral at all primes outside $T$. This proves strong approximation in the example and $\SL_n(\RR)$ is noncompact.

For comparison, $\operatorname{PGL}_2/\QQ$ has a noncompact real group. Consider the open set at finite adeles with local class represented by an invertible integral matrix at every prime, and determinant of nonsquare unit class at $5$. It is nonempty by taking $\operatorname{diag}(2,1)$ at $5$ and identities elsewhere. Square class of the determinant is independent of a local scalar, and on the image of integral invertible matrices the $5$-adic condition is open. This gives an open set for the restricted product, using the integral closed model of $\operatorname{PGL}_2$ acting by conjugation on $M_2$ (it is the automorphism group of the matrix algebra). In this model a $\ZZ_p$-point is represented by an invertible integral matrix since a scalar torsor over the local ring $\ZZ_p$ is trivial. A rational point in this open set would have a determinant $d\in\QQ^\times$ with even valuation at every prime, hence square class $\pm1$. Both are squares in $\QQ_5$, since $-1$ is a square modulo $5$ and a square root lifts by Hensel's lemma. This contradicts the condition at $5$. Here the central isogeny $\SL_2\to\operatorname{PGL}_2$ exhibits the missing simply connectedness.

For the role of noncompactness, if $\prod_{v\in S}G(K_v)$ is
compact, the projection of the diagonal to the outside places
is closed and discrete. For number fields, the additive
diagonal is discrete by requiring integrality at every finite
place and small coordinates at the infinite places, where the
ring of integers is a lattice. For function fields, an element
integral at all places is a constant. For example, choose a
separating parameter over the finite constants. The coefficients
of the characteristic polynomial of such an element are
integral at every place of the rational function field by
integrality at every extension of each place, so they are
constants. Now require the value at one place to have
strictly positive valuation, excluding all nonzero constants.
A discrete subgroup of a Hausdorff topological group is closed:
an identity neighborhood whose pairwise differences contain no
nonzero group element has at most one group element in any
translate. Apply additive discreteness to matrix entries and
the inverse determinant. Any product of the compact inside
group with a compact neighborhood outside contains only
finitely many diagonal points, which proves the projection
assertion. A positive-dimensional smooth group has a
nondiscrete outside local group, since an \'etale chart
gives local coordinates near its identity. Density fails.

We prove the other necessity by central isogenies, including
when their kernels are not smooth. Suppose $G$ is not
algebraically simply connected. The split central cover
of C07--C08 descends to a finite faithfully flat
$K$-homomorphism $\pi:\widetilde G\to G$ with nontrivial
central kernel $Z$. Here are the descent details. Choose
a finite Galois splitting extension $E/K$. The Galois
action on $G_E$ is a product of adjoint conjugations and
pinned automorphisms by F06. A pinned automorphism acts
on the simply connected datum by permuting the simple
coroot basis, so it lifts by C08, compatibly with $\pi$
by pinned uniqueness there. Adjoint conjugation lifts
as well: the isogeny induces an isomorphism on adjoint
groups. To check the latter claim on schemes, the
preimage of the target torus is the source torus,
since the torus isogeny is faithfully flat and the
kernel is inside that torus. The root-character
condition from C07 and the centre formula C05 then
identify the preimage of the target centre with the
source centre. Faithfully flat quotient gives the
adjoint isomorphism. Two lifts of the same group
automorphism differ by a morphism from the smooth
connected cover to $Z_E$. This morphism is constant,
as is seen after extension to an algebraic closure
(a reduced connected scheme has constant image in a
finite scheme). Its value at the identity is the
identity. The Galois cocycle therefore lifts, and
F02 gives the asserted cover and isogeny over $K$.

Over $E$, choose the decomposition $Z_E=\prod_i\mu_{n_i}$
from the Smith bases of the split torus isogeny, using B01,
with at least one $n_i>1$. Enlarge the finite set
of places under consideration so that the groups,
the matrix models, the tori, and the isogeny with its
kernel are defined over the ring of integers of $E$
away from that set. For the extension of the isogeny
one uses only finite presentation: extend a presentation
of the generic coordinate algebra as a finite projective
module over the target coordinate algebra, including
its multiplication and a finite dual basis. Clear
denominators in the dual basis identities and in the
torsor isomorphism and its inverse. After also clearing
the group and action identities these give the same
finite flat torsor over the localized ring. The rank is
positive, and hence it is faithfully flat: inverting
finitely many primes eliminates the rank-zero locus,
which has empty generic fibre. The torsor lies over
the equation model of $G$ used for the integrality
conditions, because the two finite lists of model
equations agree over $E$ and hence after clearing
denominators. Enlarge the set in $K$ by all places below
the removed places in $E$ and by the original $S$.

For $g\in G(K)$ integral outside this enlarged set,
the pullback of $\pi$ gives a $Z_E$-torsor. Its
$i$-th component over $E$ corresponds to
$a_i\in E^\times/E^{\times n_i}$. This description follows
from the $\ZZ/n_i$-grading of the torsor coordinate algebra:
the degree-one part is a line with trivialized $n_i$-th
tensor power, as one sees after the faithfully flat
base change trivializing the torsor. Multiplication
identifies each other degree with the corresponding
tensor power. The same argument works over a local
valuation ring, where the line is free and its
trivialization is specified by a unit. Thus all
valuations of $a_i$ outside the enlarged set in $E$
are multiples of $n_i$. This uses the finite flat
description and remains valid in characteristic
dividing $n_i$.

There are finitely many such classes for each $i$.
For clarity, the needed arithmetic finiteness works
uniformly over either kind of global field. Let $R$
be the integral closure of $\ZZ$ in the number-field
case or of $\FF_q[t]$ for a separating $t$ in the
function-field case. It is finite free over the
base principal ideal domain: take an integral field
basis and bound its trace-dual lattice using the nonzero
discriminant. Its nonzero ideals are invertible since
it is the noetherian integrally closed domain of
dimension one obtained by this normalization.
For an ideal $I$, two short elements of $R$
with the same image in the finite quotient $R/I$
give a nonzero $a\in I$ with
$|R/(a)|\leq C|R/I|$ for a constant $C$ depending
only on the chosen basis. In degree $d$ over $\QQ$,
choose the coefficients of the short elements
in $\{0,\ldots,m\}$ with $(m+1)^d>|R/I|$ and
$m\leq |R/I|^{1/d}$. The degree-$d$ norm polynomial
bounds $|N(a)|$ by $Cm^d$. Over $\FF_q(t)$ choose
coefficients of degree below $m$ with
$q^{md}>|R/I|$ and $q^{md}\leq q^d|R/I|$.
The determinant of multiplication by $a$ has
degree at most $md+C'$ since each entry has degree
bounded by $m$ plus a constant. This gives the
required index bound, as can also be read from
diagonalizing this multiplication matrix over the
base principal ideal domain. The ideal $(a)I^{-1}$
is integral and represents the inverse class of $I$.
There are finitely many ideals of bounded index:
over $\ZZ$ they contain their integer index times
$R$, and over $\FF_q[t]$ they contain the multiple
of $R$ given by the characteristic polynomial of
multiplication by $t$ on $R/I$, of degree
$\log_q|R/I|$. Hence the ideal class group is finite.
Localizing at more places only quotients this class
group.

The unit group is finitely generated. In the number
case, the logarithms of the archimedean absolute
values of units have discrete image with finite kernel:
a bounded range of logarithms bounds all the
embeddings, and the integral lattice has only finitely
many points there. To see that a discrete subgroup
of a real vector space is finitely generated, choose
a real basis of its span from its elements, then
subtract integer combinations of the basis to put
each element in a bounded parallelepiped. In the
function case, the valuations at the finitely many
places above infinity give a homomorphism to a finite
rank integer lattice with finite constant kernel.
In both cases allowing poles at finitely many more
places extends the group by a subgroup of a finite
rank integer lattice via the corresponding valuations.
The enlarged set may include all the places above
infinity of the chosen separating parameter. Now set
$A$ equal to the integers of $E$ outside that set.
For the class of $a_i$ choose a fractional ideal
$I_i$ of $A$ with $(a_i)=I_i^{n_i}$.
Its class belongs to the finite $n_i$-torsion of the
class group. If two such ideals give the same class,
their representatives $a_i$ differ by an $n_i$-th
power and a unit. The latter classes are finite
by finite generation. This proves the finiteness
used here without a finite generation assertion for
$G(\OO_{K,S})$.

We will use the elementary existence of infinitely many
places split in any finite Galois extension $L/K$.
Choose an integral primitive element with monic minimal
polynomial $f(x)$ over the initial integer ring of $K$;
take $f(0)=b\ne 0$ by a translation if needed. For any
finite excluded set of primes take a nonzero $c$ in
their product. The values $f(bcx)/b$ are integral
and congruent to $1$ at all excluded primes.
Their norms grow without bound along positive integers
$x$ for number fields, or have unbounded degree along
$x=t^r$ for function fields. In the second case the
norm is a polynomial in $x$ over $\FF_q[t]$ of
positive degree, and for large $r$ its leading term
has strictly greatest $t$-degree. Consequently at
least one value is not a unit and gives a prime outside
the excluded set where $f$ has a root. Exclude also
the primes dividing the discriminant. The root is
simple and Hensel lifting embeds $L$ into $K_v$;
since $L/K$ is Galois, it splits at $v$.

First suppose one chosen $n_i>1$ is prime to the
characteristic of $K$ (in characteristic zero any
will do). Choose a finite Galois extension of $K$
containing $E$, the $n_i$-th roots of unity, and
$n_i$-th roots of representatives of the finitely
many relevant classes of $a_i$. The preceding
argument gives a finite place $v$ split in this
extension and outside the enlarged set, with residue
characteristic prime to $n_i$. Fix an embedding
$E\hookrightarrow K_v$. All the corresponding
global $i$-th obstruction classes vanish locally.
But the residue field has a unit of nontrivial class
modulo $n_i$-th powers: its multiplicative group is
cyclic and contains all $n_i$-th roots of unity.
The split torus isogeny realizes this unit obstruction
on a point of $T(\OO_v)$. Indeed its character map
diagonalizes in integer bases, giving the coordinate
power maps and their Kummer classes. The quotient by
$n_i$-th powers on local units is discrete, by
Hensel lifting on principal units.

Otherwise in positive characteristic $p$ use the
quotient modulo $p$-th powers of an obstruction
coordinate with $p\mid n_i$. Choose $v$ split in $E$
and outside the enlarged set. The $i$-th obstruction
of global integral points still has finite image.
The local unit quotient modulo $p$-th powers is
infinite and Hausdorff. In the completed valuation
ring the residue field embeds by Hensel lifts of
the roots of $x^{q_v}-x$, giving the series
description $\OO_v=k_v[[\varpi]]$. Here $p$-th
powers in $K_v$ have precisely the exponents
divisible by $p$, a closed condition. The units
$1+\varpi^j$ for positive $j$ prime to $p$ give
distinct classes, since the first differing term
of their quotient has degree not divisible by $p$.
The split torus isogeny realizes all these classes
after passing to the quotient modulo $p$-th powers.
Choose one outside the finite global image.

In either case the obstruction is continuous on
the open product cell $U^-TU^+$ around the chosen
torus point: both root factors lift under $\pi$,
and on the torus it is the quotient of a coordinate
of the split torus by the specified powers. An open
neighborhood of this point in $G(\OO_v)$ therefore
avoids the images of all global points integral
outside the enlarged set. Use this neighborhood as
the local condition at $v$, impose integrality at
all other finite outside places beyond the enlarged
set, and put no restrictions at the finitely many
remaining outside places. This is a nonempty open
subset of the original restricted product, containing
no diagonal point. This proves necessity of algebraic
simple connectedness in both characteristics.

For sufficiency in the split simply connected case,
we use additive approximation away from a single place.
We give the additive argument because it must allow $S$
to omit some or all of the infinite places of a number
field. Fix a place $w\in S$ to be omitted. For number
fields, all finite congruences away from $w$ can be met,
with integrality at unspecified finite places, by the
Chinese remainder theorem applied to fractional ideals.
In detail, allow the required denominators at the finitely
many specified places in a fractional ideal $I$. Their
precision conditions select a coset $x_0+J$ of a smaller
fractional ideal $J\subset I$: the integral ideal
$J I^{-1}$ is a product of powers of distinct prime
ideals, and the quotients of the invertible module $I$
at these powers reduce independently. When $w$ is
finite we may also add any element of
$\bigcup_{m\geq0}J\mathfrak p_w^{-m}$, preserving
all the required finite-place conditions.

The remaining archimedean conditions can be met in each
case. Write $V=K\otimes_{\QQ}\RR$, a product over the
infinite places, and use its real trace pairing.
A closed proper subgroup of a finite-dimensional real
vector space admits a nonzero linear functional taking
integer values on it. To verify this, let $M$ be the
largest vector subspace contained in it. In an
orthogonal complement to $M$ the subgroup is closed
and discrete: a sequence of nonzero elements tending
to zero would have normalized directions tending to
some unit direction, and suitable integer multiples
would put its whole line in the subgroup. The discrete
part is a lattice in its span by the bounded
parallelepiped argument. A nonzero integer-valued
functional on that lattice, or a functional on an
unused complement when its span is proper, proves
the assertion.

In the archimedean case for $w$, if the projection of
$J$ to the other components were not dense, such a
functional would pull back to a nonzero form on $V$
vanishing on the $w$-component and integer valued on
$J$. An integer-valued form on $J$ comes from an
element of $K$ under the trace pairing: it first
extends to a rational linear form on $J\otimes\QQ=K$,
and the field trace is nondegenerate. Its restriction
to any whole component cannot vanish unless that
element is zero. In the finite case for $w$, a form
integer valued on the union would give an element
of the intersection of all the trace-dual ideals
$(J\mathfrak p_w^{-m})^\#=J^\#\mathfrak p_w^m$.
This intersection is zero by valuations. Here the
formula for the dual follows by viewing $J^\#$ as
the quotient of the trace dual of $\OO_K$ by the
invertible fractional ideal $J$ under ideal
multiplication. Thus both projected subgroups are
dense in the relevant real vector spaces.

For function fields, let $u$ be an initial separating
parameter, $R$ the integral closure of $\FF_q[u]$, and
$V_\infty=\prod_{v\mid\infty_u}K_v$. With an
$\FF_q[u]$-basis of $R$, completion gives
$V_\infty\simeq\FF_q((1/u))^{[K:\FF_q(u)]}$.
This follows by factoring the minimal polynomial of
a primitive element over the completion, hence
decomposing the tensor algebra into its completed
field factors. By subtracting polynomial parts
in the chosen coordinates one sees that
$V_\infty/(R+\prod_{v\mid\infty_u}\OO_v)$
is finite-dimensional over $\FF_q$; the polynomial
coordinate lattice plus a fixed bounded product
lattice already covers $V_\infty$, and two such
open product lattices are commensurable.
If $w$ is finite over $u$, the quotient
$\mathfrak p_w^{-n}/R$ has dimension
$n[k_w:\FF_q]$. The condition of being integral
at the original infinite places removes at most
a bounded number of dimensions from this quotient,
by the displayed finite-dimensional quotient.
If $w$ is infinite over $u$, apply the same
argument to the quotient obtained by increasing
only the pole bound at $w$ in
$\prod_{v\mid\infty_u}\OO_v$. In both cases the
space of rational functions with poles bounded
by $n$ at $w$ and no other poles has dimension at
least $n[k_w:\FF_q]-C$ for a fixed $C$.

There is consequently a separating parameter
$t$ with sole pole at $w$. Indeed if all actual
positive pole orders there were divisible by
the characteristic $p$, the dimension up to order
$n$ would be at most $\dim_{\FF_q}K_{\rm const}
+\lfloor n/p\rfloor[k_w:\FF_q]$: the increase at
each pole order is measured by its leading
coefficient in $k_w$. Choose a function of pole
order not divisible by $p$. It is not a $p$-th
power. Since $[K:K^p]=p$ by the initial separating
parameter, this means that $dt$ is a basis of
the differentials of $K$ over $\FF_q$, so the finite
extension $K/\FF_q(t)$ is separable. The integer
ring for this $t$ is exactly the ring of functions
integral away from $w$. Finite congruences away
from $w$, with integrality at all other places,
are now solved by its fractional-ideal Chinese
remainder theorem as before. This completes the
additive approximation argument in both cases.

In a split simply connected semisimple group,
the rational points over any local field $F$
are generated by the root groups. Indeed the
Bruhat cells of D04 express a point by root
factors and normalizer representatives. The
simple Weyl representatives are root words
by C04; the simple coroot basis of the cocharacter
lattice makes every torus point a product of
coroot values, and each such value is a product
of root elements by the rank-one matrix identity
displayed above. Fix now a finite set of places
outside $S$ with prescribed target neighborhoods.
Include among them all the other outside places
where the chosen root parametrizations have
nonintegral coefficients, prescribing the identity
there if no other target was specified. Concatenate
local root words realizing the desired targets,
using zero parameters at the other coordinates
for each word. Approximate every resulting tuple
of root parameters by a global parameter by the
additive argument. At the prescribed places the
products approximate the targets by continuity.
Elsewhere outside $S$ the parameters and the
coefficients are integral, including the inverse
root matrices, so the products are integral.
This proves the split sufficiency in the full
adelic topology.

The same additive argument proves sufficiency for all
quasi-split simply connected groups. It remains to show
that their rational points over each local field \(F\) are
generated by the two unipotent radicals \(U^\pm(F)\) of
the opposite rational Borels. Work in the model descended
from the split pinning by the diagram action, as in F06.
The torus \(T\) of the Borel pair has its simple coroots
as a cocharacter basis permuted by Galois. Thus \(T\) is
a product of Weil restrictions of \(\Gm\), one for each
Galois orbit \(I\) of simple roots, by B02 and B04.

For such an orbit, use the derived group \(H_I\) of the
standard Levi on \(I\). The Levi descends because \(I\)
is stable. Its derived torus has cocharacter lattice
given by the intersection of \(X_*(T)\) with the rational
span of the coroots in \(I\), by the central product in
C03 and the coroot comparison in C07. This intersection
is the span of the simple coroots in \(I\) over \(\ZZ\),
since those coroots are part of the full basis. The group
\(H_I\) is therefore simply connected, and its torus
factor is precisely the indicated factor of \(T\).

Here is the calculation in these orbit groups. The possible
orbits in an irreducible Dynkin diagram consist of
mutually orthogonal roots, except for the middle adjacent
pair exchanged in type \(A_{2r}\). This follows by
inspecting the diagram symmetries: the reflection of
\(A_n\), the permutations of the end vertices in
\(D_n\) and \(D_4\), and the reflection of \(E_6\);
the other diagrams have no nonidentity length-preserving
symmetry. An orthogonal orbit gives the product of
\(\SL_2\)'s over the splitting field, descended as
\(\Res_{L/F}\SL_2\) over the orbit field \(L\), by the
pinned identification in C08. Its two unipotent radicals
generate its torus points and its longest Weyl
representative, by the rank-one elementary matrix words.

The adjacent pair gives the quasi-split special unitary
group of the hermitian form with anti-diagonal entries
\(1,1,1\) over a separable quadratic extension \(L/F\).
Indeed this group has split form \(\SL_3\) after base
change to \(L\), its \(F\)-Borel is upper triangular,
and conjugation exchanges the two simple coroots, as
seen by the rule for its diagonal entries. It is the
pinned outer form for that diagram action by F06.
Write \(\bar z\) for the involution of \(L/F\). Its torus
is
\[
 t(b)=\operatorname{diag}(b,\bar b/b,\bar b^{-1})
 \quad(b\in L^\times).
\]
Use the upper unipotent matrices
\[
 u(a,b)=\begin{pmatrix}1&a&b\\0&1&-\bar a\\0&0&1\end{pmatrix},
 \qquad b+\bar b=-a\bar a,
\]
and the lower unipotents \(l(a,b)=w u(a,b)w^{-1}\), where
\[
 w=\begin{pmatrix}0&0&1\\0&-1&0\\1&0&0\end{pmatrix}.
\]
For admissible \(a,b\) with \(b\ne0\), multiplication gives
\[
 l(a/\bar b,1/\bar b)\,u(a,b)\,l(a/b,1/\bar b)=t(b)w.
\]
Both lower matrices are admissible by the equation on
\(a,b\). On the right of \(u(a,b)\) the chosen lower
matrix makes the top row \((0,0,b)\), and on the left the
other makes the right column \((b,0,0)^{\mathsf T}\).
The other entries are forced by the hermitian form and
determinant, proving the identity also in characteristic
two. The ratios of admissible nonzero \(b\)'s exhaust
\(L^\times\): for \(s\in F^\times\), take nonzero \(b\) of
trace zero and both \(b,sb\) are admissible with \(a=0\).
For \(s\in L^\times\setminus F^\times\), the two linear
forms \(\operatorname{Tr}(b)\), \(\operatorname{Tr}(sb)\)
are independent by the nondegenerate trace pairing.
Solve for both to be \(-1\); then both values are
admissible with \(a=1\). Taking ratios of the displayed
products gives every \(t(s)\), and then also \(w\).
All these matrices are in \(H_I\) and its unipotent
factors lie in \(U^\pm\).

For the rational Weyl group, put \(S_T\) equal to the
largest split subtorus of \(T\). Its centralizer is \(T\):
the sum of all positive coroots is a Galois-fixed
cocharacter strictly positive on the simple roots, so
no nonzero root vanishes on \(S_T\). The fixed subgroup
of the absolute Weyl group identifies with the relative
Weyl group. In fact a fixed Weyl element preserves the
invariant rational cocharacter subspace of \(T\), hence
normalizes \(S_T\), and every element of this normalizer
over the separable closure gives a rational relative
Weyl element by the constant quotient of D09. Conversely
the rational relative representative normalizes \(T\)
as the centralizer of \(S_T\), so gives a fixed absolute
element. The fixed Weyl group is generated by the longest
elements of the orbit subsystems: for a nonidentity fixed
element choose a simple right descent in the absolute
Weyl group. The entire orbit \(I\) is a right descent,
so all positives in the subsystem on \(I\) are sent
negative. Its longest element reverses those and permutes
the remaining positive roots. Right multiplication by it
therefore decreases length by the number of subsystem
positives. Induction proves the generation statement.

The orbit group calculations now place torus points and
the generators of the rational Weyl group inside the
subgroup generated by \(U^\pm(F)\). The relative Bruhat
decomposition D09 proves the asserted generation of \(G(F)\).
Apply it at each desired outside place and concatenate
the words there. By the positive coordinate grading
argument of D08, both global \(U^\pm\) are affine spaces
with the identity at the origin. Their coordinate maps
to the matrix model, including the inverse determinant,
are integral polynomials outside a fixed finite set of
places. Approximate each tuple of affine parameters as
in the split case, prescribing identity at the additional
places in this finite set. This yields the required
local approximations with integrality everywhere else.

For a globally isotropic group the approximation step has
a uniform reduction to local generation, even when the
group is not quasi-split. Fix a noncentral \(K\)-cocharacter
\(\lambda\), and use the opposite parabolics \(P_G(\lambda)\)
and \(P_G(-\lambda)\) with radicals \(V^+,V^-\).
At a completion \(F\) let \(S_F\) be a maximal split torus
in their common Levi containing \(\lambda\). Such a torus
is maximal split in \(G_F\). To see this, choose a minimal
\(F\)-parabolic inside \(P_G(\lambda)\). Its maximal split
torus lies in a Levi by the construction in D08. Conjugate
that Levi into the given one by its unipotent radical.
This shows that the given Levi has the full split rank.
Every maximal split torus of it contains its central split
torus \(\lambda(\Gm)\), since multiplying by that torus
would otherwise enlarge it. Choose a relative chamber
by a small regular perturbation of \(\lambda\).
Its minimal parabolic is contained in \(P_G(\lambda)\),
by the sign criterion following D09.

Use the indivisible roots of the relative root system for
\(S_F\). For one such root \(a\), let \(U_{(a)}\) denote the
positive contracting group in \(Z_G(\ker a)^\circ\) for a
cocharacter positive on \(a\), as in D09. Its roots along
the relative ray are \(a\) and possibly \(2a\).
This includes all positive relative multiples: the two
integral reflection pairings for proportional roots force
their ratio to be \(1,2\) or \(1/2\). If \(a\) pairs
positively with \(\lambda\), then \(U_{(a)}\subset V^+\)
and \(U_{(-a)}\subset V^-\); the same holds with signs
reversed when the pairing is negative. The normalizer
acts on these groups through its action on the relative
root rays, by conjugating the kernel and limit definitions.

In the rank-one centralizer for \(a\), apply D09 to any
nonidentity point of \(U_{(-a)}(F)\). This group has
nonidentity points by the affine coordinates of D08.
Such a point is outside the positive parabolic and is in
the nontrivial double coset. Writing that double coset as
\(U_{(a)}(F)Z(F)n_aU_{(a)}(F)\) gives a representative
of the relative reflection in the subgroup generated by
\(U_{(\pm a)}(F)\); the factor in \(Z(F)\) is part of
the representative. Thus whenever both \(a\)-rays are
present the subgroup can conjugate by that reflection.

The relative root system here is irreducible. Indeed,
if it had two components, take the subgroup over
\(\overline F\) generated by the absolute root groups
whose restrictions belong to the first component.
The closure is smooth connected by the product-closure
argument of C03. An absolute root with zero restriction
normalizes this subgroup by the root commutation test
of C08: any resulting root \(i\alpha+j\beta\), where
\(\beta\) has restriction zero, still restricts into the
same component as \(\alpha\). An absolute root with
restriction in another component commutes with this
subgroup, since a positive combination of the two
restrictions cannot be a relative root. The torus
normalizes it as well. The open product therefore makes
it normal in \(G_{\overline F}\). Absolute almost
simplicity forces it to be all of \(G_{\overline F}\),
contradicting the noncentral root groups from the
other component that commute with it.

Let \(I\) consist of the relative simple roots pairing
to zero with \(\lambda\). Since \(P_G(\lambda)\) is
proper, there is a simple root outside \(I\).
The group generated by \(V^+(F),V^-(F)\) contains the
two rays of every positive root outside the span of
\(I\). To get the remaining simple rays, follow the
connected relative simple diagram from a simple root
\(b\notin I\). If \(a\) is an adjacent simple root, the
root \(s_b(a)\) has positive \(b\)-coefficient, so it and
its negative are outside the span of \(I\). Reflecting
with the representative for \(b\) gives the two rays of
\(a\). This propagation also works when \(b\) has just
been obtained inside \(I\): by induction maintain the
claim for all relative roots outside the span of the
still unreached simple roots. In fact for a root in the
span of the unreached simple roots, reflection across
a reached simple root has nonzero coefficient there
whenever the root pairs nontrivially with its coroot.
Its two rays then follow by conjugation. On adjoining
an adjacent simple root \(a\) to the reached set, any
positive root newly outside the unreached span either
pairs nontrivially with a previously reached simple
coroot, or has support disjoint from the neighbors of
every previously reached root. The latter is impossible,
since its support contains \(a\), which is adjacent to
one of them, and its other coefficients are
nonnegative. Thus the induction holds. It obtains
all the simple rays and hence all rays by Weyl conjugacy.

The rational groups generated by the relative rays
also contain the full radical of the chosen minimal
parabolic. Indeed the ordered product of the positive
ray groups is its scheme product. This can be checked
after extending to a splitting field, where it simply
groups the absolute root coordinates from D05 by their
relative rays. The centralizer of \(S_F\) and all its
normalizer representatives normalize the group generated
by the rays. The Bruhat decomposition D09 then shows
that group is normal in \(G(F)\). It contains every
rational parabolic radical by conjugacy of the minimal
parabolics and containment of radicals under parabolic
inclusion. Consequently the subgroup generated by the
two fixed radicals is exactly \(G(F)^+\), the subgroup
generated by the \(F\)-points of all parabolic radicals.

Both fixed radicals are affine spaces over \(K\) by
D08, with the zero coordinates at the identity.
The additive argument used for the quasi-split group
therefore approximates any finite list of elements in
the local \(G(K_v)^+\)'s with the required integrality at
all other outside places. In particular, for globally
isotropic groups, local density of \(G(K_v)^+\) in
\(G(K_v)\) at each outside place suffices for strong
approximation. This reduces the remaining part of this
case to the local generation theorem.

Another case of sufficiency, including nonsplit groups, is
\(G=\SL_m(D)\) for a central division algebra \(D/K\)
and \(m\geq2\). Here \(\SL_m(D)\) means the reduced-norm-one
group of \(M_m(D)\). We prove local density of the elementary
matrices in this case. The division-algebra and reduced-norm
facts used below are Results C04.02, C04.05, C04.07--C04.10
of \emph{Rings, Modules, and Associative Algebras}.

First let \(B\) be a central division algebra of degree \(n\)
over a nonarchimedean local field \(F\) with valuation \(v\)
and finite residue field \(k_F\). Write
\[
 w(x)=v(\operatorname{Nrd}_B(x))/n
 \qquad(x\ne0).
\]
This gives a valuation on \(B\). It is multiplicative by the
reduced norm, and on a commutative subfield \(F(x)\) it is
the unique prolongation of \(v\). To see the latter claim,
the reduced characteristic polynomial of \(x\) has only
roots of its irreducible minimal polynomial and is a power
of that polynomial, by the split matrix representation.
Consequently the reduced norm is the field norm raised
to the power \(n/[F(x):F]\). For a finite separable extension
of \(F\) its integral closure is finite over the complete
valuation ring, by the trace-lattice bound. It is a
Dedekind domain and has only one maximal ideal: its
idempotents modulo the uniformizer lift by completeness,
whereas a domain has no nontrivial idempotents. This
gives the unique extended valuation and the norm formula
by conjugating in a finite normal extension. In a finite
inseparable extension apply this to the separable part
and recover the value of any element by taking its
\(p^r\)-th power there and dividing by \(p^r\).
Now \(w(z)\geq0\) implies \(w(1+z)\geq0\), checked in
\(F(z)\). Applying this to \(z=x^{-1}y\) when
\(w(x)\le w(y)\) gives the ultrametric inequality.

The ring \(\mathcal O_B=\{w\geq0\}\) is a compact open
\(\mathcal O_F\)-lattice. Indeed the reduced norm has no
zeros on the compact unit sphere of any coordinate
maximum norm for \(B\), so its absolute value is bounded
above and bounded away from zero there; scaling by
powers of a uniformizer compares the two norms.
The ideal \(\mathfrak p_B=\{w>0\}\) has a finite
division-ring quotient. This quotient is a field.
For the finite division-ring argument, a smallest
noncommutative finite division ring would have field
centralizers for noncentral elements. If its center
has \(q\) elements and its dimension is \(d>1\), each
noncentral conjugacy class has size
\((q^d-1)/(q^{d_i}-1)\) for a proper divisor \(d_i\) of \(d\).
The cyclotomic value \(\Phi_d(q)\) divides these indices
and \(q^d-1\), so the class equation would make it divide
\(q-1\). But \(\Phi_d(q)>q-1\) by its product over
primitive complex roots of unity. This proves the assertion.

Let \(\ell=\mathcal O_B/\mathfrak p_B\) have degree \(f\)
over \(k_F\), and write \(w(B^\times)=(1/e)\ZZ\).
Since this group is contained in \((1/n)\ZZ\), one has
\(e\le n\). Multiplication by an element of value \(1/e\)
identifies each successive valuation layer with \(\ell\).
The \(e\) layers modulo a base uniformizer have total
\(k_F\)-dimension \(n^2\) by the lattice rank, so \(ef=n^2\).
Choose an element of \(\mathcal O_B\) whose residue generates
\(\ell/k_F\). Its commutative subfield is complete and has
that residue field. Lift the residue generator there
by Hensel's lemma applied to a monic lift of its separable
minimal polynomial. The resulting subfield \(L/F\) has
degree \(f\), residue degree \(f\), and is unramified.
Indeed the polynomial is irreducible by its reduction,
and each of its distinct residue roots lifts in \(L\).
It is Galois with the cyclic residue-field group.
The degree of this separable subfield is at most \(n\)
by the centralizer dimension formula in C04.08 of the
algebra notes. Thus \(e=f=n\), and \(L\) is maximal.

By Skolem--Noether, each \(\sigma\in\operatorname{Gal}(L/F)\)
has a conjugating representative in \(B^\times\), unique
up to \(L^\times\), since the centralizer of the maximal
subfield is \(L\). Its valuation gives a homomorphism
\(\operatorname{Gal}(L/F)\to(1/n)\ZZ/\ZZ\). It is injective.
If the value is integral, scale the representative by a
central scalar to value zero. Its conjugation is then
trivial on the commutative residue field, so \(\sigma\)
is trivial. Both groups have order \(n\).
Choose a representative \(\Pi\) of value \(1/n\)
after a central scaling, and write \(\tau\) for its
action on \(L\). This action generates the Galois group.
The element \(t=\Pi^n\) is in \(L\) by centralization
and is \(\tau\)-fixed, so is a uniformizer of \(F\).
Distinct valuation classes give
\(\mathcal O_B=\bigoplus_{j=0}^{n-1}\Pi^j\mathcal O_L\).
On the right \(L\)-vector space with basis \(\Pi^j\),
left multiplication is the split representation after
extension to \(L\). The coefficient \(a\in \mathcal O_L\)
acts diagonally by \(\tau^{-j}(a)\), and terms with
positive powers of \(\Pi\) act strictly triangularly
modulo \(t\), since a wrap in the basis introduces \(t\).
Hence for \(y\in \mathcal O_B\) the reduced norm and trace
reduce to \(N_{\ell/k_F}(\bar y)\) and
\(\operatorname{Tr}_{\ell/k_F}(\bar y)\).

The norm-one group of \(B\) is in the closure of its
multiplicative commutator group. A norm-one element has
value zero, and its residue has norm one. In a finite
cyclic field extension the elements of norm one are
\(\bar u/\tau(\bar u)\): the two groups have the same
order, as follows from the cyclic multiplicative group.
The commutator \([u,\Pi]\) of a lift of \(\bar u\)
has that residue. Divide by it to start in
\(1+\mathfrak p_B\). Suppose a norm-one element has
been reduced to \(1+\mathfrak p_B^r\). If \(n\) does
not divide \(r\), conjugating \(1+\Pi^r x\) by an
appropriate \(u\in\mathcal O_L^\times\) gives the
leading commutator coefficient
\((\overline{\tau^{-r}(u)/u}-1)\bar x\).
Choose \(u\) for which this scalar is nonzero, and then
\(x\) to match any desired coefficient.
If \(r=nm\), the element is \(1+t^m y\), where \(m\ge1\).
The determinant expansion modulo \(t^{m+1}\) shows that
\(\operatorname{Tr}_{\ell/k_F}(\bar y)=0\).
Here \(2m\ge m+1\), so only the trace term contributes.
The map \(\tau-1\) on \(\ell\) has image the trace-zero
space, since its kernel is \(k_F\) and trace is nonzero
by separability. The commutator \([\Pi,1+t^m x]\)
realizes that leading coefficient by choosing \(\bar x\).
In either case divide by the chosen commutator to proceed
one layer further. Any fixed precision therefore uses a
finite product of commutators. This is exactly the
density assertion.

At a complex place the local division factor is the field.
At a real place it is the field or the quaternions, by
C04.12 of the algebra notes. A quaternion of norm one
belongs to a complex subfield, choosing any such subfield
for a scalar quaternion. There it has the form
\(u/\overline u\): one may take \(u=1+x\) for \(x\ne-1\),
and \(u=\mathrm i\) otherwise. Complex conjugation on
this subfield is inner by Skolem--Noether, so this
quotient is a commutator. A scalar field has trivial
norm-one kernel. This completes all the needed local
division-algebra cases.

For each completion \(F=K_v\), write \(D_F=M_s(B)\) with
\(B\) a division algebra. Elementary row and column
operations over \(B\) reduce a matrix in \(\SL_{ms}(B)\)
to a diagonal matrix. Such operations have reduced norm
one, by passing to a split matrix representation, and
the usual two-coordinate word \(w(a)w(-1)\) gives
\(\operatorname{diag}(a,a^{-1})\) for any \(a\in B^\times\).
Use these diagonal pairs to move every nonidentity
diagonal entry into the final coordinate, leaving a
single entry of reduced norm one. Also, writing
\(D(a)=\operatorname{diag}(a,a^{-1})\) in the last pair,
\(D(a)D(b)D(ab)^{-1}=\operatorname{diag}(1,[a^{-1},b^{-1}])\).
The density of commutators proved above therefore makes
the elementary group dense in \(\SL_{ms}(B)\).
Its elementary generators with indices in distinct
\(D_F\)-blocks lie in the off-diagonal block groups.
A generator within one block is obtained by a commutator
through an index in another block, since
\([1+xE_{ij},1+E_{jk}]=1+xE_{ik}\) for distinct indices.
There is another block because \(m\ge2\).
Thus the original block groups \(1+aE_{ij}\), \(a\in D_F\),
already generate a dense local subgroup.

Choose a \(K\)-basis of \(D\) to regard each off-diagonal
block parameter as an additive tuple. The block maps and
their inverses are polynomial in these coordinates and
integral outside finitely many places. Approximate the
local words at the prescribed places, concatenate them
using identity parameters at the other places, and apply
the additive approximation proved above to every
coordinate. This proves strong approximation for
\(\SL_m(D)\) outside any nonempty \(S\). These groups
are noncompact at every omitted local place because
they contain the additive off-diagonal blocks.

We can also prove local generation for the isotropic forms of
type \(A\) with nontrivial diagram action. We first explain
the algebra description, to make its use independent of a
local classification of forms. In split type \(A_{d-1}\),
\(d\geq3\), the group is \(\SL_d\) by C08. Its automorphisms
are inner or the composite of an inner automorphism with
\(g\mapsto(g^{\mathsf t})^{-1}\), by F06: the latter has the
nontrivial symmetry of the chain as its based-datum action.
Choose a finite Galois splitting field. In a split matrix
model, write each semilinear group action as
\(\beta_\sigma(g^{(-1)^{\epsilon_\sigma}})\), where
\(\epsilon_\sigma\) is its diagram character and
\(\beta_\sigma\) is respectively a semilinear algebra
automorphism or antiautomorphism. In the outer case
\(\beta_\sigma\) is coefficient conjugation and transposition,
followed by inner conjugation. Inversion commutes with
all these maps on the group, so the cocycle identity gives
\(\beta_\sigma\beta_\tau=\beta_{\sigma\tau}\) on \(\SL_d\).
It gives equality on the entire matrix algebra: the two
sides have the same semilinearity and parity, and their
values on \(1+E_{ij}\) determine the off-diagonal matrix
units, which generate the algebra.

For trivial diagram action, Galois descent gives a central
simple algebra and the reduced-norm-one group. For nontrivial
action let \(E\) be the quadratic subfield fixed by its kernel.
Descend the matrix algebra to a central simple algebra \(A/E\)
using the even subgroup. The other coset gives an involutory
antiautomorphism \(\ast\) of \(A\) inducing the nonidentity
automorphism of \(E/F\). The original group is
\[
 \{g\in A^\times:g^\ast g=1,\ \operatorname{Nrd}_A(g)=1\}.
\]
Indeed in the split coordinates over the splitting field
the two copies of the algebra are exchanged by the involution;
the norm-one unitary equations express one copy of \(g\)
as the inverse adjoint of the other, and the remaining equation
is determinant one. The descent action in these equations
is precisely the original action. This uses just the descent
of algebras and Skolem--Noether as proved in C04.13 and
C04.07 of \emph{Rings, Modules, and Associative Algebras}.

Over a nonarchimedean local field \(F\), every such algebra
\(A/E\) with \(E/F\) a quadratic field is split. Here is the
argument. Its division factor \(B/E\) inherits an involution
of the second kind. To verify this directly, choose a primitive
matrix idempotent \(e\) in \(A=M_s(B)\). For \(x\in Ae\), all
elements have matrix rank at most one over \(B\). If every
\(x^\ast x\) vanished, polarize with \(x+y\) and with \(zx+y\)
for a central \(z\in E\) not fixed by \(\ast\). Subtraction
would give \(x^\ast y=0\) for all \(x,y\in Ae\),
contradicting \(e^\ast A e\ne0\). Thus there is a nonzero
self-adjoint element \(b=x^\ast x\) of rank one.
Likewise \(\operatorname{Trd}(b c c^\ast)\) is nonzero for
some \(c\in A\). Otherwise the same polarization and central
scaling give \(\operatorname{Trd}(b c d^\ast)=0\) for all
\(c,d\), contrary to the nondegenerate reduced-trace
pairing, which is seen on matrix units after splitting.
Now \(b'=c^\ast b c\) is self-adjoint of rank one and is
not nilpotent, since its reduced trace is nonzero. Its
kernel and image are complementary, with \(b'\) invertible
on its image. When the kernel is nonzero the minimal
polynomial is \(t p(t)\) with \(p(0)\ne0\), and has fixed
coefficients by self-adjointness and uniqueness. A polynomial
over \(F\) that is \(0\) modulo \(t\) and \(1\) modulo
\(p(t)\) therefore gives a self-adjoint primitive
idempotent \(e'\) on evaluation at \(b'\). If the kernel
is zero, use \(e'=1\). Restriction to the corner \(e'Ae'\)
gives the desired involution of the division factor.

Apply the division-algebra valuation description proved above
to \(B/E\), of degree \(n\). The involution preserves valuation
because it conjugates reduced norms, so it acts on the finite
residue field. Conjugation by a uniformizer \(\Pi\) induces a
generator \(\tau\) of the residue extension of degree \(n\).
Every uniformizer induces the same residue automorphism:
two differ by a unit, which acts trivially by conjugation on
the commutative residue field. Applying \(\ast\) to
\(\Pi x\Pi^{-1}\) thus replaces \(\tau\) by \(\tau^{-1}\)
after conjugating by the residue involution. Finite-field
automorphisms commute, so \(n\leq2\).

If \(n=2\), the canonical quaternion antiautomorphism
\(j(x)=\operatorname{Trd}(x)-x\) is involutory and commutes
with \(\ast\); these properties follow from the adjugate
formula on split two-by-two matrices. The composite
\(j\ast\) is a semilinear algebra automorphism of order two,
so \(B\) descends to a degree-two central simple algebra
over \(F\). Any such algebra is split by \(E\).
In the division case the local description above writes it
using an unramified quadratic subfield \(L/F\) and a
uniformizer generator of square \(t\), a uniformizer of
\(F\). If \(E/F\) is unramified, it is isomorphic to \(L\)
and splits that algebra. If it is ramified, \(LE/E\) is
unramified quadratic and \(v_E(t)=2\).
The norm of \(LE/E\) takes all units onto units. This follows
on residue fields from cyclicity, and on principal units
from successive trace maps
\[
  N(1+\varpi^r a)\equiv
     1+\varpi^r\operatorname{Tr}(\bar a)\pmod{\varpi^{r+1}},
  \qquad r\geq1,
\]
using completeness and surjectivity of the finite-field
trace. Thus \(t\) is a norm, by first matching its valuation.
Rescaling the conjugating generator by an element of \(LE\)
now makes its square one. Multiplication by \(LE\) together
with the generating field automorphism gives a split matrix
representation on \(LE\) over \(E\). This contradicts \(B\)
being a division algebra of degree two. Hence \(B=E\).
At an archimedean place with a quadratic field extension
only \(\mathbb C/\mathbb R\) is possible and the algebra
over \(\mathbb C\) is split as well.

In the split algebra \(M_d(E)\) the involution is adjoint to
a Hermitian form \(h\). Indeed compose \(\ast\) with conjugate
transposition and use Skolem--Noether to write
\(X^\ast=H^{-1}\bar X^{\mathsf t}H\). Its square being one
makes \(\bar H^{\mathsf t}=uH\) for a scalar \(u\) of norm
one. Rescale \(H\) to make \(u=1\) by Hilbert 90, as in F07.
An \(F\)-split cocharacter of the unitary group pairs its
opposite weights under \(h\). Thus if the group is isotropic
the Hermitian space has a nonzero isotropic vector in a
nontrivial weight, and a hyperbolic pair \(e,f\) can be
chosen with \(h(e,f)=1\). Indeed choose \(f\) of that pairing,
then change it by a multiple of \(e\) to make its norm zero
using the surjectivity of \(\operatorname{Tr}_{E/F}\).
Let \(W\) be the orthogonal complement of this plane.

The elementary unitary calculation on any nonisotropic line
\(Ew\subset W\) works over every infinite \(F\), including
characteristic two. Put \(q=h(w,w)\in F^\times\).
In the coordinates \(e,w,f\) use
\[
 u_q(a,b)=\begin{pmatrix}
 1&a&b\\0&1&-\bar a/q\\0&0&1
 \end{pmatrix},\qquad b+\bar b=-a\bar a/q.
\]
These matrices fix the orthogonal complement of the three
coordinates, and lie in the unipotent radical at \(Ee\).
Take \(w_0\) to interchange \(e,f\) and send \(w\) to
\(-w\), as in the three-dimensional calculation of
Result~\ref{res:D10}. These operations preserve the
hyperbolic pairing and the middle coefficient \(q\).
With this \(w_0\),
put \(l_q(a,b)=w_0u_q(a,b)w_0^{-1}\) for the opposite
radical. The identical row and column operations, now
with the indicated \(q\)-factor, give
\[
 l_q(a/\bar b,1/\bar b)u_q(a,b)l_q(a/b,1/\bar b)
       =\operatorname{diag}(b,\bar b/b,\bar b^{-1})w_0 .
\]
In detail both lower factors satisfy the displayed trace
condition. The right factor turns the top row into
\((0,0,b)\), the left one turns the right column into
\((b,0,0)^{\mathsf T}\), and unitarity and determinant
force the other entries. For \(s\in F^\times\), choose
\(b\) of trace zero to obtain both \(b,sb\) on the
right. For \(s\in E^\times\setminus F^\times\), use the
independence of the two trace forms to arrange
\(\operatorname{Tr}(b)=\operatorname{Tr}(sb)=-1/q\);
both are allowed with \(a=1\). Ratios yield every
\(\operatorname{diag}(s,\bar s/s,\bar s^{-1})\) in
the subgroup of the two parabolic radicals. On \(W\)
these realize the Hermitian reflections scaling a
nonisotropic \(w\) by any scalar of norm one and
fixing its orthogonal complement.

For completeness these reflections generate the unitary
group of \(W\). Take \(v\) of nonzero norm \(q_0\), and
\(u\) of the same norm which is to be taken to \(v\).
If \(z=u-\eta v\) has nonzero norm for some norm-one
\(\eta\), then \(h(z,u)\ne0\), since
\(h(z,z)=\operatorname{Tr}_{E/F}h(z,u)\).
The reflection about \(z\) of scale
\(1-h(z,z)/h(z,u)\) has norm-one scale and takes \(u\)
to \(\eta v\). A reflection about \(v\) then removes
\(\eta\). If \(h(v,u)\ne0\), the bad \(\eta\)'s are
finite, since the equation for zero norm becomes
quadratic on multiplying by \(\eta\) and using
\(\bar\eta=\eta^{-1}\). The norm-one group is infinite:
the ratios \((1+as)/(1+a\bar s)\), \(a\in F\),
are pairwise distinct for \(s\notin F\).
If \(h(v,u)=0\), the norm of \(u-v\) is nonzero when
the characteristic is not two. In characteristic two
use the intermediate vector \(a u+(1+a)v\) with
\(a\in F\setminus\{0,1\}\). It has the same norm by
\(a^2+(1+a)^2=1\) and pairs nontrivially with both
vectors. The preceding procedure works in two steps.
After aligning \(v\), reduce to its orthogonal
complement. The induction applies since a nonzero
nondegenerate Hermitian space has a nonisotropic
vector: otherwise trace polarization with a scalar
not in \(F\) forces its pairing to vanish.

The Levi preserving both isotropic lines acts by
\((t,A,\bar t^{-1})\) with \(A\in U(W)\).
The just constructed words remove \(A\) using
unitary reflections; the determinant condition
then forces the remaining \(t\) into \(F^\times\).
This scale is furnished by the same three-dimensional
diagonal words, with trivial action on \(W\).
If \(W=0\), the group is split of rank one
and the preceding split calculation applies.
By the relative root argument above the remaining
Bruhat factors lie in the generated group. Hence
\(\operatorname{SU}(h)(F)^+=\operatorname{SU}(h)(F)\).

For a global outer form at a place where the
quadratic algebra splits, the involution exchanges
the two simple algebra factors and eliminates
one as an inverse adjoint. The local group is
therefore the reduced-norm-one group in a matrix
algebra over a division factor. If isotropic it
has at least two matrix coordinates: with only
one coordinate the norm-one subgroup of the
division algebra is compact by the valuation
description above (and by the quaternion norm
at a real place). The local elementary density
already proved for that case applies; the
off-diagonal transvections lie in parabolic
radicals by contraction under a diagonal
cocharacter. Together these arguments establish
local density for all isotropic groups of type \(A\)
at their completions. The global additive argument
then gives sufficiency for all globally isotropic
groups of this type.

The local generation requirement is also elementary for
an isotropic group in the Clifford spin model of a regular
quadratic form. In characteristic two the polar form can
be nondegenerate or have a one-dimensional radical on which
the quadratic form is nonzero. We work over the local field
\(F\), using \(V=Fe\oplus W\oplus Ff\) with a hyperbolic
pair \(e,f\), both isotropic and of polar pairing one. In the Clifford algebra use
\(v^2=q(v)\), hence \(ef+fe=1\).
Write \(\widetilde g\) for the reversion of \(g\).
The spin group consists of even Clifford units normalizing
\(V\) by conjugation and with \(\widetilde g g=1\).

The elements \(1+ew\) and \(1+fw\), \(w\in W\), belong to
this group. For example \((ew)^2=0\), its reversion is
\(-ew\), and conjugation by \(1+ew\) fixes \(e\), takes
\(f\) to \(f-w-q(w)e\), and takes \(v\in W\) to
\(v+\langle w,v\rangle e\), by multiplication in the
Clifford algebra. These elements contract along the
appropriate signs of the cocharacter acting by \(s^2\)
on \(e\), \(s^{-2}\) on \(f\), and trivially on \(W\).
This is a spin cocharacter, given by \(s\,ef+s^{-1}fe\).
Hence the two families lie in the respective opposite
parabolic radicals. In characteristic zero they give the
whole radicals. In the orthogonal action the first radical
fixes \(e\) and is the identity on the associated graded space
of the flag \(Fe\subset Fe\oplus W\subset V\). The image
of \(f\) has the form \(f-w+ce\); isotropy forces
\(c=-q(w)\), and preservation of its pairing with \(W\)
forces the displayed action on \(W\). This determines the
orthogonal element, and a contracting radical meets the
finite central kernel of the spin cover trivially. The other
radical is obtained by interchanging \(e,f\).
For \(q(w)\ne0\) their product
\[
 n_w=(1+ew)(1+fw/q(w))(1+ew)=(e+f/q(w))w
\]
belongs to the generated group. For the equality, the two
nilpotent terms \(E=ew\) and \(F_0=fw/q(w)\) satisfy
\(EF_0=-ef\), \(F_0E=-fe\), \(EF_0E=-E\).
Conjugation by the displayed even product induces the
two orthogonal reflections in \(e+f/q(w)\) and in \(w\).
It acts on \(W\) by the reflection in \(w\) and
interchanges the lines \(Fe,Ff\), sending
\(e\mapsto-f/q(w)\) and \(f\mapsto-q(w)e\).

Let \(L\) be the Levi centralizing this cocharacter.
In characteristic different from two its points act on
\(W\) by elements of \(\operatorname{SO}(W)\), since
the even Clifford elements commute with the volume
element, on which a normalizing orthogonal action has
its determinant. Every orthogonal isometry is a product
of reflections. To check this when the characteristic is
not two, choose a nonisotropic vector \(v\). If its image
differs, reflection in the difference aligns it when this
difference is nonisotropic. Otherwise the sum is
nonisotropic since the two norms add to \(4q(v)\);
reflection in the sum takes the image to \(-v\), and a
reflection in \(v\) finishes. Induct on \(v^\perp\).
The determinant makes the number of reflections even
for the action of \(L\) on \(W\).

In characteristic two one aligns a nondegenerate
two-dimensional plane at each induction step. For
vectors \(u,v\) of equal norm outside the polar radical,
if their polar pairing is nonzero, symmetry in \(u+v\)
takes \(u\) to \(v\). If their pairing is zero, choose
\(z\) with \(q(z)\ne0\) and with both pairings nonzero.
This is possible because the three nonvanishing
conditions define nonempty opens in a vector space
over the infinite field \(F\). The vector obtained
from \(u\) by symmetry in \(z\) pairs nontrivially with
both \(u\) and \(v\), so two symmetries align them.
Choose the first vector \(v\) of nonzero norm outside
the radical. After aligning it, choose a partner \(w\)
of polar pairing one with \(v\). Its image has the same
pairing. Run the same alignment for \(w\) and its image,
now taking the intervening \(z\) in \(v^\perp\); both
pairing functionals on \(v^\perp\) are nonzero because
\(v\) belongs to that space and pairs nontrivially
with both. The quadratic form there is nonzero since
\(q(v)\ne0\). All these symmetries fix \(v\), and after
aligning \(w\) one inducts in the plane's perpendicular
space. The possible one-dimensional radical with
nonzero quadratic form has trivial isometry group
over the field. This proves generation by symmetries.

We need the even-number assertion in characteristic
two for the action of \(L\), or a way to arrange it.
For a quadratic space of nondegenerate polar form the
center of its Clifford algebra is scalar. This follows
after splitting into hyperbolic planes over an algebraic
closure: exterior multiplication and contraction by
paired basis vectors generate all matrix units on the
exterior algebra, by first projecting to the vacuum and
then changing monomials. The defining Clifford relations
give the matching dimension bound by ordered monomials.
The claim then descends. If the polar radical is the
line \(Fr\) with \(q(r)\ne0\), split it off from a
complement of nondegenerate polar form. In this
characteristic its Clifford algebra is the tensor
product with \(F[r]/(r^2-q(r))\), since the radical
vector commutes with the complement. Its center is
therefore that commutative factor, whose even part is
scalar.

In the nondegenerate polar case, a product of vectors
of \(W\) realizing the symmetry action of an element
of \(L\) acts on \(Fe+Ff\) as the identity in characteristic
two. After multiplying by \(t\,ef+fe\) it has the same
action on \(V\) as the given element of \(L\).
The ratio centralizes \(V\) and is scalar by the center
calculation, so the number of vector factors is even.
When the polar radical is one-dimensional, the symmetry
of its nonisotropic generator \(r\) acts trivially on
\(W\). Adjoin this symmetry if necessary to make the
number even. In all cases the even product of the
corresponding \(n_w\)'s removes the action on \(W\)
and preserves \(Fe,Ff\).

The element left in \(L\) induces a scale \(t,t^{-1}\)
on the two lines and the identity on \(W\).
The Clifford element \(x_t=t\,ef+fe\) has exactly
that conjugation action, and \(\widetilde x_t x_t=t\)
since \(ef,fe\) are orthogonal idempotents interchanged
by reversion. An even Clifford element centralizing
every vector is scalar. In characteristic two this
was verified by the center calculation. In the other
characteristics use an orthogonal basis: ordered
monomials form a basis by the relations and their
leading exterior terms acting on \(1\).
Conjugation by a nonisotropic basis vector changes
the sign of exactly the even monomials containing it,
so only the scalar monomial commutes with all vectors.
Thus any spin lift of the scale is a
scalar multiple of \(x_t\), and its norm forces \(t\)
to be a square. The product \(n_{sw}n_w^{-1}\) realizes
scale \(s^2\), using any nonisotropic \(w\in W\).
The remaining scalar kernel is \(\{1,-1\}\), and
\(n_{-w}n_w^{-1}=-1\). This gives all Levi points in the
subgroup generated by the parabolic radicals. The
dimension-two torus case has no absolutely almost
simple group and is not used here.

By the relative root argument above, the subgroup
of the two radicals already contains the unipotent
radical of a minimal parabolic and the rational Weyl
representatives. The centralizer of its maximal split
torus lies in \(L\), and D09 then gives the full group.
Applying additive approximation proves sufficiency
for every globally isotropic group in these spin models
under the stated absolute simplicity hypothesis.

This includes all globally isotropic groups of type \(B\)
in characteristic different from two. In the split odd-dimensional
quadratic model, the special orthogonal group is the adjoint group
of the simply connected spin group. For the lattice check, take
the ordinary coordinates \(t_i\) on the split orthogonal torus.
In each hyperbolic plane the element \(t_i e_i f_i+f_i e_i\)
implements the indicated scale and has Clifford norm \(t_i\).
Their product gives the spin torus by adjoining a scalar with
square \((\prod_i t_i)^{-1}\), since any two implementers differ
by a centralizing scalar as above. The lift is a connected
torus and includes the full scalar kernel. Its cocharacter
lattice inside the orthogonal lattice is the even-sum lattice,
which is the coroot lattice of the roots
\(\pm\varepsilon_i,\ \pm\varepsilon_i\pm\varepsilon_j\).
The split root groups of the orthogonal model can be read
from the hyperbolic shears and the elementary transformations
between two planes. The lift is a central double cover:
over an algebraic closure every even product of reflections
lifts after normalizing its Clifford norm, and the kernel is
the étale scalar group of order two. The cover is smooth by
the tangent dimension test, and connected because its identity
component surjects with the scalar kernel already in its torus.
Thus C07 gives the simply connected form. Since the diagram
has no symmetry, the twist of the spin group is by its adjoint
orthogonal group, by F06. That group acts on the split quadratic
space itself, so the same cocycle descends the vector space and
quadratic form. The Clifford construction is equivariant for
this action and commutes with base change, giving the spin
group of the descended form. A split cocharacter acting
nontrivially has a nonzero weight on the quadratic space,
and the corresponding vectors have norm zero. Hence an
isotropic group gives the isotropic quadratic model just used.

In characteristic two the same type-\(B\) descent can be
checked in the regular odd-dimensional model. Start with the
split \(q_0\) on \(V=V_0+Fr\), where \(V_0\) is a sum of
\(n\) hyperbolic planes, \(Fr\) is the polar radical,
and \(q(r)=1\). The orthogonal group fixing \(r\) is
the group of pairs \((A,\ell)\) with \(A\) symplectic on
\(V_0\) and \(\ell\) linear, subject to
\[
 q_0(Ax)+\ell(x)^2=q_0(x) \qquad(x\in V_0).
\]
The difference of the quadratic forms has zero polar form,
so in a basis it is a sum of squared-coordinate terms.
Projection to the symplectic group is finite faithfully
flat and radicial: its algebra adjoins the \(2n\) square
roots of those coefficients. The group of pairs is smooth.
At the identity its tangent equations on \(X\in\mathfrak{sp}\)
are the vanishing of the diagonal of \(JX\) for the
polar matrix \(J\); these are \(2n\) independent equations
on the symmetric matrices \(JX\). The tangent \(\ell\)
is unrestricted and has dimension \(2n\), so tangent
and scheme dimensions agree. Translation gives smoothness
everywhere. The group is connected by radiciality and
is reductive, since a smooth unipotent radical would have
normal unipotent image in the symplectic group and cannot
lie in a finite kernel. Its diagonal torus has the
weights \(\pm\varepsilon_i\pm\varepsilon_j\) from the
off-diagonal part of \(\mathfrak{sp}\) and
\(\pm\varepsilon_i\) from the \(\ell\)-coordinates.
It is maximal by comparison with the symplectic group.
These are the \(B_n\) roots, and they span the character
lattice, so this is the split adjoint group by C05.

Embed \(V\) into the split even quadratic space with one
additional hyperbolic plane \(Fe_0+Ff_0\), taking
\(r=e_0+f_0\) and \(V=r^\perp\). The even Clifford spin
group is smooth by the alternating-shear open charts
in Appendix C of \emph{Nonassociative Algebras}.
Its orbit map at \(r\) is smooth. Indeed the tangent
space to the smooth norm-one quadric at \(r\) is \(V\).
The shears with Clifford terms \(e_i e_0,f_i e_0\)
give the \(V_0\) tangent directions, by their commutators
with \(r\). The torus cocharacter scaling \(e_0,e_1\)
by \(s\) and \(f_0,f_1\) by \(s^{-1}\) lifts to spin by
the Clifford scalar \(s^{-1}\), and differentiates to
the remaining direction \(r\). The tangent test at the
identity, followed by translation, makes the stabilizer
\(H\) of \(r\) smooth.

In the full Clifford algebra the centralizer of \(r\)
is \(\operatorname{Cl}(V)\). Indeed in the ordered
monomial basis write every element as \(a+bf_0\)
with \(a,b\in\operatorname{Cl}(V)\). Commutation with
\(r\) detects \(b\), since \(rf_0+f_0r=1\) and
\(r\) commutes with \(V\) in characteristic two.
Thus the restriction \(H\) to \(V\) has scalar
norm-one kernel \(\mu_2\), by the even center computation
above. It surjects geometrically onto the adjoint group
of pairs. An isometry of \(V\) is a product of symmetries
in vectors of \(V\) by the plane argument above.
Pad the product to even length by symmetry in \(r\),
which is trivial on \(V\). In the even Clifford space
this even product implements the isometry and preserves
\(r\). Normalize its norm by a scalar over the algebraic
closure. The resulting element belongs to \(H\).
The morphism \(H\to O(V)_r\) is therefore faithfully
flat by A07, with kernel \(\mu_2\). Its identity component
surjects by dimension, and the kernel has one geometric
point, so \(H\) is connected. It is reductive by the
same unipotent radical argument as for the group of pairs.
Its torus maps by the same norm equation
\(u^2\prod_i t_i=1\) as above. The even-sum lattice
is the coroot lattice, hence \(H\) is the simply
connected split group by C07.

This stabilizer is the intrinsic Clifford spin model.
Any even norm-one unit in \(\operatorname{Cl}(V)\)
normalizing \(V\) induces an isometry fixing \(r\),
which is central in that algebra. Lift this isometry
to \(H\) after an fppf extension by the preceding
surjection. The remaining implementer is a scalar of
norm one and belongs to \(H\). The scalar claim also
holds on test algebras, since commuting with the vector
generators is a linear system over the field, and its
kernel commutes with scalar extension. This proves
equality of the two functors. By F06, type \(B\) has
only inner twists; its adjoint group of pairs acts
linearly on the odd quadratic space. Descending that
space and its Clifford model therefore gives the
odd spin description over the original field.
A noncentral split cocharacter supplies a nonzero
isotropic weight vector as above. We have thus
proved sufficiency for globally isotropic type \(B\)
also in characteristic two.

There is an analogous proof for all globally isotropic groups
of type \(C\). Over a splitting field the group is the
symplectic group, with the
diagonal hyperbolic torus and roots of type \(C\); its
cocharacter lattice has the basis of simple coroots.
By F06 the form is an inner twist. The adjoint symplectic
group acts on the split matrix algebra by conjugation
preserving its symplectic adjoint involution, since scalars
act trivially. Descent gives a central simple algebra with
a symplectic involution of the first kind, and the descended
group is the unitary group of that involution. These statements
can be checked over the splitting field, where they are the
ordinary symplectic group equations.

At a nonarchimedean completion write the algebra as \(M_r(B)\).
The involution takes a primitive idempotent to a primitive
idempotent. Conjugate it back by an invertible matrix, which
exists by choosing bases of images and kernels. This gives
a center-preserving antiautomorphism on its division corner;
an involution of the corner is not needed for this argument.
As above the division valuation is invariant and the residue
automorphisms of the finite field commute, while the
antiautomorphism interchanges the action of a uniformizer with
its inverse. Thus the degree of \(B\) is at most two.
At real and complex completions the same alternatives
are the field or the real quaternions, as above.
If the division factor is the field, the group is the split
symplectic group. In the other case use canonical quaternion
conjugation, denoted by a bar. Any involution on \(M_r(B)\)
has the form
\[
 X\longmapsto H^{-1}\bar X^{\mathsf t}H
\]
by Skolem--Noether applied to its composite with the
standard adjoint. Its square forces
\(\bar H^{\mathsf t}=\epsilon H\) with \(\epsilon=\pm1\).
For a symplectic involution \(\epsilon=1\).
Indeed on split quaternion matrices conjugation is the
two-dimensional alternating adjoint, so the standard
matrix adjoint has a block alternating matrix \(J\).
The new form has matrix \(JH\) and
\((JH)^{\mathsf t}=-\epsilon JH\).
In characteristic different from two its being alternating
forces \(\epsilon=1\), and in characteristic two there is
only that choice of sign. In both cases each diagonal block
\(J_2 H_{ii}\) is alternating, as a restriction of \(JH\).
A two-by-two alternating matrix is a scalar multiple of
\(J_2\); hence each \(H_{ii}\) is a central scalar, also
before splitting. Thus the group acts on a nondegenerate
quaternionic Hermitian space with \(h(v,v)\in F\),
since each off-diagonal pair contributes a reduced trace.
A nontrivial split cocharacter gives a nonzero weight vector
of norm zero as before. Choose a hyperbolic pair \(e,f\)
with pairing one, adjusting the second vector by a multiple
of \(e\) using surjectivity of reduced trace. Let \(W\) be
its orthogonal complement.

We check the Levi generation by the same small matrices,
now over a noncommutative quaternion division ring \(B\).
The norm \(q=h(w,w)\) of any nonisotropic \(w\in W\)
is central by the symplectic form calculation above. On
\(e,w,f\) use the matrices \(u_q(a,b),l_q(a,b)\) above
with \(b+\bar b=-a\bar a/q\). They fix the orthogonal
complement and lie in the two parabolic radicals.
For \(b\ne0\), carry out the multiplication in the
indicated orders:
\[
 l_q(\bar b^{-1}a,\bar b^{-1})u_q(a,b)
       l_q(b^{-1}a,\bar b^{-1})
 =
 \begin{pmatrix}
  0&0&b\\
  0&1+\bar a b^{-1}a/q&0\\
  \bar b^{-1}&0&0
 \end{pmatrix}.
\]
The two lower factors are admissible since norms of
their first parameters are \(N(a)/N(b)\), and
\(\operatorname{Trd}(\bar b^{-1})=\operatorname{Trd}(b)/N(b)\).
The right and left operations remove the other entries
from the top row and right column as before. The middle
entry is obtained by adding \((\bar a/q)b^{-1}a\) to
\(1\), and preservation of the Hermitian pairing forces
the remaining entries. This computation uses no
commutation of \(a,b\).

If a quaternion \(r\) of norm one differs from \(1\),
take \(a=1\) and \(b=-q^{-1}(1-r)^{-1}\).
Then \(\bar b=-rb\) and \(b+\bar b=-q^{-1}\), so
the middle entry is \(r\). For \(a=0\) any nonzero
trace-zero \(b\) is admissible, and the middle entry
is \(1\). Taking the ratio with such an interchange
gives a Levi element acting on \(W\) as scaling by
\(r\) on \(w\) and fixing its perpendicular space.

These quaternionic reflections generate the full
isometry group of \(W\). For two vectors \(u,v\) of
the same nonzero norm, try \(z=u-vr_0\) with
\(N(r_0)=1\). If \(z\) is nonisotropic, then
\(a=h(z,u)\) satisfies \(a+\bar a=h(z,z)\).
The reflection sending \(z b\) to \(z r b\), with
\(r=1-h(z,z)a^{-1}=-\bar a a^{-1}\), has norm-one
scale and takes \(u\) to \(vr_0\). A reflection on
\(v\) then removes \(r_0\). In characteristic different
from two at least one of \(u-v\) and \(u+v\) is
nonisotropic, since their norms add to \(4h(u,u)\).
In characteristic two, if \(c=h(u,v)\ne0\), there
is an \(r_0\) of norm one with
\(\operatorname{Trd}(c r_0)\ne0\), which makes
\(z\) nonisotropic. To see this, choose \(b\in B\)
with \(\operatorname{Trd}(b)\ne0\) and
\(\operatorname{Trd}(cb)\ne0\) by nondegeneracy of
trace and infinitude of \(F\). Such a \(b\) is
noncentral and generates a separable quadratic
subfield. The norm-one elements of this subfield
span it over \(F\), by the field ratios used above.
Thus the nonzero functional \(\operatorname{Trd}(c-)\)
on that subfield is nonzero on a norm-one element.
If \(h(u,v)=0\), the same-norm intermediate
\(a u+(1+a)v\) for \(a\in F\setminus\{0,1\}\)
reduces to the nonorthogonal case twice.
After aligning \(v\), induct on its orthogonal
complement. A nonisotropic vector exists in a
nonzero nondegenerate space, by trace polarization
and the nondegenerate trace pairing.

We can thus remove the entire \(W\)-component from any
Levi element, leaving \(\operatorname{diag}(t,1,\bar t^{-1})\)
on \(e,W,f\). Any \(t\in B^\times\) is a ratio \(b c^{-1}\)
of nonzero trace-zero quaternions. Choose \(c\ne0\) in
the simultaneous kernels of \(\operatorname{Trd}(c)\)
and \(\operatorname{Trd}(tc)\), possible by
dimension four; take \(b=tc\). The two interchange
words with zero \(a\) and parameters \(b,c\) furnish
exactly this scale as their ratio. They also work for
\(W=0\), since then they are products of the upper and
lower two-dimensional shears with trace-zero
parameters. Explicitly for a pure \(b\) the two-by-two
product of upper shear \(b\), lower shear \(-b^{-1}\),
and upper shear \(b\) is the desired interchange.
This proves Levi generation in all these local models.
The other Bruhat factors are supplied by the relative
root argument above. Additive approximation proves
sufficiency over the global field.

Finally, in every characteristic an isotropic group
of type \(G_2\) is split. The diagram
has no symmetry and the split center is trivial, so
by F06 the form is an inner twist by the split group
itself. This is the automorphism group of the split
octonions by \emph{Nonassociative Algebras},
Result NA13-C04.08. The same cocycle therefore
descends that algebra and its norm, and identifies the
given group with its automorphism group. A nontrivial
split cocharacter acts on the octonions with a nonzero
weight because this representation is faithful. Its
weighted vectors have isotropic norm, by norm
invariance. The trace-one idempotent and Zorn
coordinate construction in the proof of NA13-C04.08
gives the split octonion algebra from a nonzero
isotropic vector in every characteristic. The split
approximation argument then applies.

\medskip\noindent
\textit{Globally anisotropic inner type \(A\) over a number field.}
Write \(k\) for the number field, \(D/k\) for the central division
algebra of degree \(n\), \(G=\operatorname{SL}_1(D)\), and
\(\Gamma=G(k)\). Use the reduced norm \(\operatorname{Nrd}\) of
Result~\ref{res:F06}. First, \(G(\mathbb A_k)/\Gamma\) is compact.
Choose an order \(O\) in \(D\) that is free over \(\mathcal O_k\).
To obtain one, include \(1\) in a \(k\)-basis and scale the other
basis vectors by a common integer clearing the multiplication
denominators. At a finite place \(v\), left multiplication by
\(x\in O_v\) has determinant \(\operatorname{Nrd}(x)^n\) on \(D_v\),
as follows after passing to a matrix splitting. This determinant is
integral, so \(\operatorname{Nrd}(x)\in\mathcal O_v\).
View \(D\) over \(\mathbb Q\) in a fixed integral basis of \(O\).
For \(g\in G(\mathbb A_k)\) take the full rational lattice
\[
 L_g=D\cap\prod_{v\text{ finite}}g_v^{-1}O_v.
\]
The local lattices lie between \(M O_v\) and \(M^{-1}O_v\) for a
single positive integer \(M\), and equal \(O_v\) for almost all
\(v\). The Chinese remainder theorem on \(M^{-1}O/MO\) shows that
the completions of \(L_g\) are the prescribed lattices. Its
covolume in the real coordinates of \(O\) is one: the determinant
of a rational basis matrix has absolute \(p\)-adic value one at
each rational prime \(p\), since \(g_v\) has determinant one for
every \(v\mid p\). The real lattice \(g_\infty L_g\) also has
covolume one. Any nonzero \(x\in L_g\) satisfies
\[
  |N_{k/\mathbb Q}\operatorname{Nrd}(g_\infty x)|
  =|N_{k/\mathbb Q}\operatorname{Nrd}(x)|\geq 1,
\]
where the left side means the product at the archimedean places,
with the complex factors squared. Indeed \(\operatorname{Nrd}(x)\)
is a nonzero algebraic integer. This archimedean product is
bounded by a fixed constant times the Euclidean norm of
\(g_\infty x\) to the power \(n[k:\mathbb Q]\), giving a uniform
positive lower bound on lengths of nonzero lattice vectors.

Here is the bounded-basis argument for these real lattices.
A ball of radius depending only on an upper covolume bound has
volume larger than the covolume. Two of its points agree modulo
the lattice, giving a nonzero vector of bounded length. Choose a
shortest vector \(b\). It is primitive, and orthogonal projection to
\(b^\perp\) gives a lattice of covolume divided by \(\|b\|\). Any
nonzero vector of the projected lattice has length at least
\(\sqrt{3}\|b\|/2\), since one can lift it with the component along
\(b\) of absolute length at most \(\|b\|/2\). Induction therefore
produces a bounded basis, lifting and adjusting each projected
basis vector by a multiple of \(b\). For \(g_\infty L_g\) the inverse
basis matrices are bounded too, since their determinant is
\(\pm1\). If \(h_g\) is the basis matrix before multiplying by \(g\),
then \(g h_g\) lies in a fixed compact subset of
\(\operatorname{GL}_{\dim_{\mathbb Q}D}(\mathbb A_{\mathbb Q})\);
its finite components lie in \(\prod_p\operatorname{GL}_{\dim_{\mathbb Q}D}
(\mathbb Z_p)\).

For \(x\) in a rational basis of \(D\), let \(R_x\) denote right
multiplication. The rational matrices
\(h_g^{-1}R_xh_g=(g h_g)^{-1}R_x(g h_g)\) lie in a fixed compact
set, so there are only finitely many tuples of such matrices.
Here the rational diagonal is closed and discrete, as in the
topology argument above. Choose one \(h_i\) per tuple. Then
\(h_g h_i^{-1}\) commutes with every \(R_x\), hence equals left
multiplication by some \(d\in D^\times\), by evaluation at \(1\).
Thus \(gd\) lies in a fixed compact set in \(D^\times(\mathbb A_k)\):
the left multiplication subspace is closed in the ambient matrix
space, and an invertible left multiplication matrix has its
inverse in the same subspace. The finite integral components of
both matrices preserve \(O\), so the restricted product topologies
agree there. The rational numbers \(\operatorname{Nrd}(d)\in
k^\times\) have only finitely many possibilities in the compact
norm image. Choose one \(d_a\in D^\times\) of each such norm
\(a\). Now \(gdd_a^{-1}\) is in a fixed compact subset of
\(G(\mathbb A_k)\), and \(dd_a^{-1}\in\Gamma\). This proves
compactness of the quotient.

Put \(Z=\overline{\Gamma}\) in \(G(\mathbb A_k^S)\). Take a finite
\(T\supseteq S\) containing all the noncompact archimedean places,
and a compact open subgroup \(U^T\) of the adelic product outside
\(T\), using the full groups at the compact archimedean places.
The orbit of the identity under \(G_T\times U^T\) in the compact
quotient just proved is closed, since every orbit of this open
subgroup is open. Hence it is compact. Quotient by \(U^T\) shows
that \(G_T/\Lambda\) is compact, where \(\Lambda\) is the projection
of \(\{\gamma\in\Gamma:\gamma^T\in U^T\}\) to \(G_T\).
This projection is injective, and its image is discrete because
a compact subset of \(G_T\) times \(U^T\) meets the discrete closed
diagonal in finitely many points. There is a compact subset
\(C\subset G_T\) with \(C\Lambda=G_T\), by covering the compact
quotient with images of relatively compact open sets. The
projection of \(\Lambda\) to any nonempty subproduct \(J\) of the
places \(T\setminus S\) is nondiscrete. Otherwise, in decompositions
\(g=c\lambda\) of elements \(g\) supported on \(S\), the projection
of \(\lambda\) to \(J\) could take only finitely many values, being
discrete and lying in a fixed compact set. Rational points embed
into each completion. Thus only finitely many \(\lambda\) could
occur, contradicting the noncompactness of \(G_S\).

We record an invariance argument for such a lattice. Let \(J\)
be a single finite split factor \(\operatorname{SL}_n(k_v)\), or
a subproduct of the noncompact archimedean factors. The latter
are \(\operatorname{SL}_n(\mathbb R)\),
\(\operatorname{SL}_n(\mathbb C)\), or
\(\operatorname{SL}_m(\mathbb H)\) with \(m\geq2\), by the real
division algebra calculation of the algebra notes
(Result C04.12). Work in the adjoint representation over
\(\mathbb Q_p\) for the finite split factor and over \(\mathbb R\)
for the archimedean product. Let \(W\) be a subspace invariant
under the projected \(\Lambda\), of dimension \(r\). Send
\(G_T/\Lambda\) into the projective space of the \(r\)-th exterior
power by sending \(g\Lambda\) to the line \(\bigwedge^r gW\).
The groups here are unimodular: the determinant of the adjoint
action is a character of the semisimple group, hence trivial
by Results~\ref{res:C03} and \ref{res:C04}; the Jacobian of
conjugation for an invariant volume is its local absolute value.
Local Haar measure consequently descends along the discrete
quotient, using neighborhoods disjoint from their nonidentity
right \(\Lambda\)-translates. It gives a finite invariant measure
of full support. Its pushforward has support the whole compact
image of the displayed map.

In any factor choose a diagonal element \(a\) whose entries are
distinct powers of \(p\) at a finite place or distinct powers of
\(2\) at an archimedean place, the exponents summing to zero.
In the exterior representation its eigenvalues are powers of
the same scalar, with distinct absolute values whenever the
eigenvalues differ. An orbit of a projective vector having two
or more nonzero weight components has forward limit in a
different projective line, obtained by retaining the largest
absolute-value component. Such a vector is not recurrent. In a
finite invariant measure almost every point is recurrent: for
each open set the points in it with no future return have
pairwise disjoint forward images, hence measure zero, and a
countable base gives the claim. Thus the full support is in the
closed fixed locus. This also applies to every conjugate of
\(a\), proving \(W\) invariant under the normal subgroup they
generate. For a transvection \(u_{ij}(t)\) the commutator with
\(a\) is \(u_{ij}((a_i/a_j-1)t)\). The coefficient is invertible,
so this normal subgroup contains all transvections in the
factor. They generate each factor. Over \(\mathbb R\) and
\(\mathbb C\) this is elementary row reduction. Over \(\mathbb H\)
row reduction leaves only a unit quaternion in one entry, and
the commutator calculation for unit quaternions above completes
the reduction. We conclude that \(W\) is \(J\)-invariant.

At a finite split factor let \(H\) be the closure of the
projection of \(\Lambda\). On a sufficiently small matrix
neighborhood of \(1\) the logarithm and exponential series
converge and are inverse, and \(\det\exp X=\exp(\operatorname{tr}X)\).
For \(x\in H\) in that neighborhood,
\(\exp(s\log x)\in H\) for \(s\in\mathbb Z_p\), by passage from
integer powers. The \(\mathbb Q_p\)-span of these logarithms is
nonzero by nondiscreteness and is invariant under the projected
\(\Lambda\): taking high \(p\)-powers makes a conjugated
logarithm small without changing its span. By the invariance
argument it is a Lie ideal, upon differentiating the adjoint
action. This Lie algebra is simple over \(\mathbb Q_p\).
After scalar extension to an algebraic closure it is a sum of
simple \(\mathfrak{sl}_n\) factors permuted transitively by the
Galois group. Simplicity of \(\mathfrak{sl}_n\) in characteristic
zero follows by taking brackets with diagonal and elementary
matrices. A basis of logarithms now gives a neighborhood in \(H\):
the product of the exponentials in basis directions has
invertible derivative at zero. In scaled coordinates on a small
ball the nonlinear terms after removing the linear part have
Lipschitz constant less than one, so the inverse exists on a
ball by contraction. Thus \(H\) is open. It is also cocompact
as a quotient image of \(G_T/\Lambda\), and therefore has finite
index. The action on its finite coset space kills every
transvection since the additive group of \(k_v\) is divisible.
It follows that \(H=J\).

For the product \(J\) of the archimedean factors in \(T\setminus
S\), the same argument uses real exponential lines. For the
closed subgroup \(H\) generated by the projected closure, set
\(\mathfrak h=\{X:\exp(\mathbb R X)\subset H\}\).
This is a real vector space by the product limit
\((\exp(tX/N)\exp(tY/N))^N\to\exp(t(X+Y))\). If \(M\) is a
linear complement, the map
\((X,Y)\mapsto\exp X\exp Y\) from \(\mathfrak h\times M\) is a
coordinate chart near zero by its derivative. There are no
nonzero arbitrarily small \(Y\in M\) with \(\exp Y\in H\):
otherwise normalize their directions, take a convergent
subsequence, and take integer powers with integer times the
length tending to any prescribed real time. The limiting
direction would lie in \(\mathfrak h\cap M\). Thus \(H\) near
the identity consists of exponentials of \(\mathfrak h\).
It is invariant under the projected \(\Lambda\), hence by the
invariance argument \(\mathfrak h\) is a Lie ideal of the product.
Each real factor is simple: its complexification is simple
or is two simple factors interchanged by conjugation.
Consequently \(\mathfrak h\) is the sum of some full factor
algebras. The corresponding factors are connected, being
generated by their connected real transvection groups, so they
lie in \(H\). If any factor remains, the projection to the
remaining product is discrete by the exponential chart
argument. This contradicts the nondiscreteness proved for each
such nonempty subproduct. Thus the archimedean projection is
dense as well.

We can now isolate these dense projections from the compact
factors. For \(J\) as above and \(T=S\cup\{\text{places of }J\}\),
with other archimedean factors compact or already omitted, let
\(P=Z\cap(J\times U^T)\). Its projection to \(J\) is closed by
compactness of \(U^T\) and is onto by the lattice argument.
Let \(U'\) be the closure of its projection to \(U^T\), and
let \(U_0=\{u:(1,u)\in P\}\). Then \(U_0\) is normal in \(U'\)
by conjugation in \(P\) and continuity. There is a continuous
homomorphism \(J\to U'/U_0\), sending \(j\) to the second
component of any lift modulo \(U_0\). Continuity follows
because the closed projection \(P\to J\) is a quotient map.
Every transvection is killed: it is contracted to \(1\) by
positive powers of some diagonal \(a\), and a subsequence of
the powers of the image of \(a\) converges in the compact
target. Hence the image of the transvection is conjugate to
\(1\). Since the transvections generate \(J\), we have
\(J\times\{1\}\subset Z\). Apply this first to all noncompact
archimedean factors outside \(S\), if any, and then add them to
the omitted set. The quotient image of \(Z\) is closed and is
the closure of the rational diagonal, so density there lifts.
Apply the isolation argument at each finite split place
outside the enlarged omitted set.

Only finitely many places remain after the split finite factors
have been isolated. Indeed for all but finitely many \(v\) the
map \(O_v\otimes O_v^{\mathrm{op}}\to\operatorname{End}(O_v)\)
is an isomorphism by the central simple enveloping algebra
calculation in the algebra notes (Result C04.03) and clearing
the denominators of its inverse. This property persists on
reduction and forces the residue algebra to be central simple:
two-sided ideals are invariant under all endomorphisms, and
central left multiplications commute with all endomorphisms.
It splits over the finite residue field by the finite division
ring calculation above. Lift a primitive matrix idempotent of
the residue algebra using Newton iteration for \(X^2-X\).
The derivative \(2X-1\) is invertible modulo the maximal ideal
near an idempotent, and the iteration takes place among
commuting functions of the initial lift. The resulting left
ideal is a free module of rank \(n\), since it is a direct
summand and has that rank on reduction. Left multiplication on
it gives \(D_v\simeq M_n(k_v)\).

The isolated factors give their entire restricted product in
\(Z\), since finite supports are dense. At the finite set of
remaining completions the image of the rational diagonal is
dense. For this, \(D^\times(k)\) has simultaneous approximation
in \(D_v^\times\) by additive approximation on the matrix
coordinates and openness of invertibility. The closure of the
commutator subgroup of \(D_v^\times\) contains \(G(k_v)\):
for division factors this is the valuation-layer and quaternion
calculation above. For matrix factors combine those calculations
with row reduction and write each transvection as a commutator
with a diagonal scalar in one position. Approximate the
finitely many entries of local products of commutators at once
by rational units. Their global commutators have reduced norm
one. The projection of \(Z\) to the remaining finite product
is closed since \(Z\) contains the kernel, and is thus the
whole product. This proves strong approximation for the
specified anisotropic inner forms over number fields.

\medskip\noindent
\textit{Globally anisotropic type \(G_2\) over a number field.}
By the octonion identification above, this group is
\(G=\operatorname{Aut}(A)\) for an octonion \(k\)-algebra \(A\)
with anisotropic norm \(q\). The compactness argument just given
also applies to this representation. Choose a free
\(\mathcal O_k\)-lattice containing \(1\) that is closed under
multiplication and on which \(q\) has integral coefficients.
One obtains it by rescaling the other members of a basis
containing \(1\) to clear denominators in the multiplication
and norm tables. The action preserves \(q\) and has determinant
one by connectedness. The lattices \(g_\infty L_g\) from this
representation again have covolume one. Now the nonzero
algebraic integer \(q(x)\), rather than a reduced norm,
gives the uniform lower bound on lengths, using the
archimedean product of the quadratic norm. Use the same
bounded-basis argument to form \(h_g\). Instead of the
right-multiplication matrices, use the bilinear multiplication
tensor and the matrices of scalar multiplication by a
\(\mathbb Q\)-basis of \(k\). Their transforms by \(h_g\)
are rational and range in compact sets, since their transforms
by \(h_g\) equal their transforms by \(gh_g\).
For two choices with the same finite tuple of transformed
tensors, \(h_gh_i^{-1}\) is \(k\)-linear and preserves
multiplication. Thus it already belongs to \(G(k)\).
The compact intersection with the closed multiplication and
scalar equations yields \(G(\mathbb A_k)/G(k)\) compact.

Every finite completion of \(A\) is split. Indeed
Cayley--Dickson doubling writes its norm as
\(N_B(x)-cN_B(y)\) for a quaternion algebra \(B\) and
\(c\ne0\). If \(B\) splits, its norm is the matrix
determinant and is surjective on units. If \(B\) is
division, use the unramified quadratic field \(L/F\) and
the element \(\Pi\) of value \(1/2\) constructed above
for a local division algebra. On \(L^\times\) the reduced
norm is the field norm. That norm is surjective on the
units: the norm on the finite residue fields is surjective,
and on the layer \(1+t^j\mathcal O_L\), \(j\geq1\), the
coefficient of \(t^j\) in the norm is the residue trace.
Lift a prescribed unit residue and then correct one
coefficient at a time, which converges by completeness.
The reduced norm of \(\Pi\) has value one, so together these
norms give all of \(F^\times\). Taking \(N_B(x)=c\)
and \(y=1\) gives an isotropic vector, proving the assertion.
The complex factors are split. At a real place the norm
either is isotropic and the group split, or is positive
definite and the group compact.

The lattice, projective recurrence and isolation arguments
above apply now to each finite split factor and to the
product of the noncompact archimedean factors. In the
recurrence step choose a cocharacter of the split torus
having nonzero pairing with every root, evaluated at \(p\)
at a finite place or \(2\) at an archimedean place.
The eigenvalue and fixed-line argument is unchanged,
and its conjugates generate a subgroup containing all root
groups by commutation. In characteristic zero the split
Lie algebra is simple by its root decomposition:
commutation with a generic Cartan element separates the
one-dimensional root spaces of an ideal. A nonzero Cartan
element brackets nontrivially with some root vector, and
an ideal containing a root space contains its coroot by
bracketing with the opposite root. This gives a simple
root space with nonzero coroot pairing and then all the
simple root spaces by connectedness of the diagram.
Opposite simple root spaces and the Cartan follow by
the same brackets, and the simple root vectors generate the
other root spaces. After restriction of scalars the simple
factors are permuted transitively, as above. Consequently
the small-exponential arguments give an open cocompact
projection at a finite place, hence all of the group since
root groups are divisible there. At the noncompact
archimedean places the same real Lie ideal argument applies,
using connectedness because the groups are generated by
real or complex root groups. The contraction argument for
the projection against a compact factor uses those root
groups in exactly the same way. This isolates the full
restricted product of split finite factors and the
noncompact archimedean factors outside \(S\).

For the remaining compact real factors we give the needed
finite-set density argument. Their groups are connected.
In a positive-definite real octonion algebra choose a
unit imaginary vector \(a\), a unit imaginary vector \(b\)
perpendicular to \(a\), and a unit vector \(c\)
perpendicular to the quaternion span of \(1,a,b,ab\).
The Cayley--Dickson multiplication table on the resulting
basis is independent of the frame. The automorphisms are
therefore identified with this space of frames, starting
from one fixed frame. Successively the spheres of choices
have dimensions \(6,5,3\). These sphere bundles are locally
trivial by choosing local orthonormal bases of the
orthogonal complements, so the space is path connected.

The diagonal of \(G(k)\) is dense near the identity in
the product of these compact factors. Here is the rational
parameter argument. Over an algebraic closure the map
\((g,X)\mapsto\operatorname{Ad}(g)X\) from \(G\) times a
Cartan Lie algebra has surjective derivative at the
identity and any Cartan element on which all roots are
nonzero: the root directions come from the bracket and
the remaining directions from the Cartan. Restricted to
these Cartan elements it therefore has dense constructible
image, containing a nonempty open subset. The image is
independent of the chosen Cartan and is Galois invariant.
Its interior descends to a nonempty \(k\)-open subset.
Since \(k\) is infinite, choose a \(k\)-basis of the Lie
algebra from this open subset. For each basis vector the
identity component of its centralizer is a \(k\)-torus,
as is seen in the Cartan root decomposition, and its Lie
algebra contains the vector.

For each of these tori choose cocharacters forming a
basis over a finite separable splitting extension \(L/k\).
For each cocharacter, multiply its Galois-conjugate
evaluations on the factors of
\(\operatorname{Res}_{L/k}\mathbb G_m\).
This descends to a \(k\)-morphism to the torus. Over an
algebraic closure its differential includes the chosen
cocharacter direction by using just its own factor.
Taking the products of these maps first in each torus and
then in \(G\) gives a map from an open set in affine
\(k\)-space with differential onto \(\operatorname{Lie}G\)
at the parameters for \(1\). Over each real completion
its image contains an identity neighborhood by the
inverse function theorem on a complementary coordinate
slice. Simultaneous additive approximation to the
parameters gives the claimed density near the identity.
The resulting closure is an open subgroup of a connected
product and hence is the whole product. Since the
previous isolation contains the kernel of the projection
to these real factors, this completes strong approximation
for the anisotropic \(G_2\) over number fields.

\medskip\noindent
\textit{Globally anisotropic quadratic spin groups over a number field.}
Let \(q\) be an anisotropic nondegenerate quadratic form over
\(k\), with dimension at least five. First
\(\operatorname{SO}(q)(\mathbb A_k)/\operatorname{SO}(q)(k)\)
is compact by the same lattice argument. On a free integral
lattice choose coordinates with integral norm coefficients.
The archimedean product of \(q(x)\) is bounded away from zero
for nonzero \(x\) in each rational lattice \(L_g\), exactly
as for the octonion norm. The tensors to be normalized by the
bounded bases are the quadratic form, an oriented \(k\)-volume
form and multiplication by a basis of scalars. Equality
of their transforms makes the rational change of basis a
\(k\)-linear determinant-one isometry. There are finitely many
tuples of transforms, which gives compact representatives.

Write \(\pi:\operatorname{Spin}(q)\to\operatorname{SO}(q)\)
for the central covering. Over any field of characteristic
zero, the reflection argument above defines the spinor norm
on \(\operatorname{SO}(q)\) with values in the group of
square classes: take the product of the vector norms in
a reflection word of even length. Two Clifford implementers
of the same isometry differ by a scalar, so this class is
well defined, is multiplicative, and is trivial precisely
when the isometry lifts to the spin group. At a finite place
of odd residue characteristic where \(q\) is unimodular on
the lattice, every integral isometry there has unit spinor
norm class. In fact one can diagonalize the unimodular
form on its lattice by choosing a vector of unit norm
modulo the maximal ideal and taking its perpendicular
summand successively. When aligning one such vector \(v\)
and its image, one of their difference and sum has unit
norm because their two norms add to \(4q(v)\). Reflection in
that vector is integral and aligns the pair, with at most a
further reflection in \(v\). Repeating on the perpendicular
summand gives a word in reflections with unit norms.

For \(g\in\operatorname{Spin}(q)(\mathbb A_k)\) choose
\(d\in\operatorname{SO}(q)(k)\) for which \(\pi(g)d\) lies in
the fixed compact representative set. The local spinor
norm of \(d\) equals that of \(\pi(g)d\). Hence outside
one fixed finite set its valuation is even, since a
compact adelic set is integral outside finitely many places.
The arithmetic finiteness argument for square classes in
the necessity proof above shows that there are only finitely
many global possibilities for the spinor norm of \(d\).
Choose a rational isometry \(d_a\) of each occurring class.
Then \(dd_a^{-1}=\pi(\gamma)\) for a rational spin element
\(\gamma\), and \(\pi(g\gamma)\) lies in a fixed compact set.
The finite covering \(\pi\) has compact inverse images of
compact adelic sets. To check this, its affine coordinate
generators satisfy monic polynomials over the target algebra.
Those equations bound the preimages of compact local sets
and imply integrality of lifts at all places where the
coefficients and the target coordinates are integral.
Together with the closed equations this gives local
compactness of the inverse images and a compact restricted
product bound, after excluding finitely many places.
We obtain compactness of the spin adelic quotient.

We verify that \(q\) is isotropic at every finite completion
\(F=k_v\). A nondegenerate ternary form is a nonzero scalar
multiple of the pure quaternion norm: for
\(\langle a,b,c\rangle\) use the scalar \(\lambda=ab/c\)
and quaternion generators of squares \(-a/\lambda,-b/\lambda\).
If the ternary form is anisotropic the quaternion algebra is
division, since a split quaternion algebra has a nonzero
trace-zero matrix of determinant zero. Every quadratic
extension \(E/F\) embeds into the division quaternion
algebra \(B\). To see this, use the unramified quadratic
subfield \(L/F\) and the generator \(\Pi\) from the valuation
argument above. The unramified quadratic extension is unique,
since it is obtained by lifting the finite residue extension.
If \(E/F\) is ramified, \(LE/E\) is unramified quadratic,
and \(\Pi^2\) has even valuation in \(E\). The unit norm
surjectivity proved above supplies \(z\in(LE)^\times\) with
norm \(\Pi^2\). Represent \(B_E\) on \(LE\) by multiplication
and \(\Pi(y)=z\tau(y)\), with \(\tau\) the nontrivial
automorphism of \(LE/E\). This splits \(B_E\). The unramified
case also splits, since \(L\) is a maximal separable subfield.
In either case the split representation on \(E^2\), viewed
over \(F\), is a one-dimensional left \(B\)-module.
The scalar action of \(E\) embeds into its \(B\)-linear
endomorphism algebra \(B^{\mathrm{op}}\), and hence into \(B\).
Thus a nonzero scalar \(u\) is the norm of a pure quaternion
exactly when \(-u\) is nonsquare: the pure element must square
to \(-u\), and conversely use the embedding of
\(F(\sqrt{-u})\). A nonzero pure element in the division
algebra cannot have a nonzero square as its square, by
factoring that polynomial. The anisotropic ternary form
therefore represents all but one of the nonzero square classes.

An anisotropic binary form is a scalar multiple of a
quadratic extension norm. A ramified quadratic extension
has norms of odd valuation, by taking a uniformizer. For
an unramified quadratic extension the norm on units is
onto, and there is a nonsquare unit: a Kummer defining
class for that unramified extension has even valuation.
Thus the anisotropic binary form represents at least two
square classes. Diagonalize the form \(q\) and take a
ternary and a binary orthogonal subform. If either is
isotropic, so is \(q\). Otherwise the binary values meet
the negatives of the ternary values in some square class,
and scaling within the subforms gives an isotropic vector.

The lattice projection and isolation arguments now apply
to all finite factors and to the noncompact archimedean
factors, using their two generating opposite Clifford
unipotent groups \(1+ew,\ 1+fw\) proved above. The displayed
spin cocharacter from the Clifford construction contracts
each in one direction and acts in the adjoint representation
with eigenvalues powers of its parameter. In the corresponding
radical the commutator of the cocharacter at \(s\) with the
element of parameter \(w\) has parameter
\((s^{\pm 2}-1)w\). Thus projective
recurrence and conjugation by the cocharacter give the same
Lie ideal argument. This uses absolute simplicity, by
the root decomposition after scalar extension, exactly as
in the preceding case. The small logarithms used there
lie in the spin Lie algebra since all their \(\mathbb Z_p\)
exponentials are limits of integer powers in the closed
group. The analytic coordinate argument applies to its
smooth identity neighborhood. Each finite factor is
generated by additive groups over a characteristic-zero
field and has no proper subgroup of finite index. Each
noncompact real factor and complex factor is connected
by its unipotent generating groups. So the full products
of these factors outside \(S\) are isolated as before.

The compact real spin factors are connected as well.
Plane rotations generate the definite orthogonal group
of determinant one, by successively aligning basis vectors.
Their spin lifts lie on the circles
\(\cos(\theta/2)+\sin(\theta/2)e_ie_j\) in an orthonormal
basis, where \((e_ie_j)^2=-1\). These connected circles
also contain the kernel element \(-1\), so they generate
the real spin group. The rational torus parameter argument
used for \(G_2\) applies verbatim, since its characteristic
zero Cartan and centralizer argument takes place after
scalar extension. It gives density in the remaining
compact real product. Together with the isolated kernel
this finishes the quadratic spin case. The absolutely
simple dimension-three spin case is the inner type \(A_1\)
already proved, and dimension four is excluded by absolute
simplicity.

\medskip\noindent
\textit{Globally anisotropic outer type \(A\) over a number field.}
Write \(E/k\) for the quadratic diagram extension and
\((A,\ast)\) for the central simple \(E\)-algebra with
unitary involution obtained above. Let \(n\geq3\) be its
degree. The construction of a self-adjoint primitive
idempotent above applies over the global field \(E\) as
well. The polarization uses only an \(E\)-scalar moved
by the involution, the reduced-trace pairing is
nondegenerate after matrix splitting, and the polynomial
giving the projection to the rank-one image has coefficients
fixed by the involution. Thus the same global division
factor \(D=e'Ae'\) has a unitary involution \(\sigma\).
Write \(A=\operatorname{End}_D(V)\). In a \(D\)-basis
of \(V\), let \(a^\dagger\) be transpose followed by
entrywise \(\sigma\). The composite of two involutions
inducing the same action on \(E\) is an \(E\)-algebra
automorphism. By Skolem--Noether there is an invertible
matrix \(H\) with \(a^\ast=H^{-1}a^\dagger H\).
Applying the involution twice gives \(H^\dagger=\lambda H\)
for \(\lambda\in E^\times\), with \(\bar\lambda\lambda=1\).
Replacing \(H\) by \(\mu H\) makes it Hermitian, taking
\(\mu/\bar\mu=\lambda\). For \(\lambda\ne-1\) take
\(\mu=1+\lambda\), and for \(\lambda=-1\) take any
nonzero anti-invariant scalar. Thus \(\ast\) is the
adjoint involution of the nondegenerate \(D\)-Hermitian
form \(h(x,y)=x^\dagger H y\).

If \(h\) has an isotropic vector \(e\), pick \(f\) with
\(h(e,f)=1\). Choose \(z\in E\) with \(z+\bar z=1\)
and replace \(f\) by \(f-e z h(f,f)\) to make it
isotropic, since \(h(f,f)\) is self-adjoint. The central
scales \(t,t^{-1}\) on these two \(D\)-lines, with the
identity on the perpendicular space, give a split torus
in \(G=\operatorname{SU}(A,\ast)\). Indeed the two
reduced norms are reciprocal after splitting \(D\).
Global anisotropy therefore implies \(h(x,x)\ne0\)
for all nonzero \(x\in V\).

We get compactness of \(G(\mathbb A_k)/G(k)\) by the
bounded-basis argument above. Choose a free integral
lattice in the underlying \(k\)-space of \(V\) on which
the coefficients of \(h\), in a fixed \(k\)-basis of
\(D\), are integral. If \(d=\deg D\), the determinant
over \(E\) of the action on \(V\) is
\(\operatorname{Nrd}_A(g)^d\): after splitting \(D\)
the action is left multiplication on \(d\) columns.
Thus \(G\) preserves covolume. For a nonzero rational
vector in the lattice \(L_g\), some coefficient of
\(h(x,x)\) is a nonzero algebraic integer. Its product
at the archimedean places has absolute value at least
one. The corresponding coefficients of
\(h(g_\infty x,g_\infty x)\) are bounded by a constant
times the square of the real vector length at each
place. Using a single bound for the finite list of
coefficients gives the same uniform lower bound on
nonzero vectors of \(g_\infty L_g\).

Normalize the matrices of right multiplication by
a rational basis of \(D\), the \(D\)-valued form \(h\),
and an \(E\)-valued top exterior form on \(V\), viewed
as rational tensors. Their changes by \(h_g\) and by
\(g h_g\) coincide. For each of the finitely many
possible transformed tuples, compare \(h_g\) with a
fixed representative \(h_i\). Then \(h_gh_i^{-1}\)
comes from a rational element \(u\) of the full unitary
group \(U\), and it preserves the \(E\)-volume. Its
reduced norm is therefore a \(d\)-th root of unity
in \(E\). Choose one rational unitary element \(u_a\)
for each of the finitely many reduced norms that
occur. The products \(u u_a^{-1}\) lie in \(G(k)\)
and \(g u u_a^{-1}\) ranges in a fixed compact set.
Here one takes the intersection with the closed unitary
and reduced-norm equations in the bounded matrix set,
identifying the restricted product topologies as above.
This proves compactness.

Consider a local place \(v\) with \(F=k_v\). If the
quadratic extension remains a field and \(v\) is finite,
the local division argument above shows that \(A\)
splits over \(E_v\). Its Hermitian form then has dimension
\(n\) over \(E_v\), with \(n\geq3\). It admits a diagonal
basis, since in a nonzero nondegenerate Hermitian space
polarization with a scalar moved by the involution gives
a vector of nonzero norm. For three diagonal entries
\(a,b,c\in F^\times\), use the quaternion algebra of
the separable quadratic field \(E_v/F\) with parameter
\(-b/a\). Its reduced norm on two copies of \(E_v\)
is \(N(x)+(b/a)N(y)\). It represents \(-c/a\)
by the local quaternion norm surjectivity above.
This gives an isotropic vector on adding the third
coordinate \(1\). At a real place with \(E_v=\mathbb C\)
the group is compact for definite signature and isotropic
otherwise. At a place where \(E_v\) splits into two
factors, \(U(F)\) identifies with the unit group of
one simple algebra by eliminating the other factor
using \(\ast\). Then \(G(F)\) is the reduced-norm-one
group of \(M_s(B)\) for a local division algebra \(B\).
For \(s\geq2\) it is isotropic. For \(s=1\) it is compact,
and this last case can occur only at finitely many finite
places: the splitting-at-almost-all-places argument above
applies to \(D/E\), while real and complex division factors
have degree at most two, less than \(n\).

At each isotropic completion the generation argument for
type \(A\) above gives a dense subgroup generated by
unipotent groups contracted by split cocharacters. This
is enough for the lattice projection argument. Each
cocharacter evaluated at \(p\) or \(2\) again gives
the projective recurrence fixed-line conclusion. A
closed normal group containing such a diagonal element
\(a\) also contains the group it contracts: if
\(a^j u a^{-j}\to1\), the commutators
\(a^j u a^{-j}u^{-1}\) converge to \(u^{-1}\).
The recurrence invariance therefore extends to the full
local group. In characteristic zero the Lie algebra
is simple over the underlying \(\mathbb Q_p\) or
\(\mathbb R\), by absolute simplicity and the scalar
extension argument above. The real groups in question
are connected since the real unipotent groups are
connected and generate densely. At the finite places
each unipotent has \(m\)-th roots in its unipotent
group for every positive integer \(m\): the polynomial
\(\sum_j\binom{s}{j}(u-1)^j\) agrees with powers for
integers \(s\), hence lies in the group for all \(s\)
and gives an \(m\)-th root at \(s=1/m\).
Consequently a continuous finite quotient kills these
generators. The small-exponential arguments give the
same full projections and the contraction argument
isolates their factors from the compact remainder.
The order is as above: first isolate the noncompact
archimedean product outside \(S\), then the isotropic
finite factors outside the enlarged omitted set.

For the finite set of compact factors use weak
approximation in the full unitary group \(U\).
The translated Cayley chart is valid in all
characteristics. On the open locus where \(1-u\)
is invertible, it is the pair of inverse maps
\[
 y=(1-u)^{-1},\qquad u=1-y^{-1},
 \qquad y+y^\ast=1 .
\]
Indeed \((1-u^{-1})^{-1}=1-(1-u)^{-1}\), and
conversely \(y^\ast=1-y\) is invertible when \(y\)
is, giving \(u^\ast=u^{-1}\). Thus \(U_0\) is
isomorphic to the invertible subset of the affine
space of solutions of \(y+y^\ast=1\), which contains
the scalar \(z\) chosen above. Over an algebraic
closure \(U\) is \(\operatorname{GL}_n\), so it is
irreducible, and \(U_0\) is nonempty by choosing
diagonal entries different from \(1\). Over a completion \(F\) the
rational points of the Cayley open are Zariski dense,
as the source is open in an affine space over the
infinite field \(F\). For any \(u\in U(F)\), the
intersection \(U_0\cap uU_0\) therefore contains an
\(F\)-point, expressing \(u\) as a quotient of two
Cayley values. Simultaneous additive approximation
to these parameters gives weak approximation for
\(U(k)\) at any finite list of completions.

At a finite compact factor the local description
\(U(F)=B^\times\) gives density of commutators in
\(G(F)\) by the valuation-layer calculation above.
At a compact real factor \(U(F)=U(n)\), every
element of \(\operatorname{SU}(n)\) is a commutator
in \(U(n)\). Conjugate it to a diagonal matrix
with entries \(\lambda_i\) of modulus one and
product one. Choose unit scalars \(\eta_i\) with
\(\eta_i/\eta_{i-1}=\lambda_i\), with indices
interpreted cyclically. The diagonal matrix is the
commutator of the diagonal \((\eta_i)\) with the
cyclic permutation. Approximate finite products of
the local commutators using \(U(k)\); the rational
commutators have reduced norm one by multiplicativity.
Thus \(G(k)\) has dense projection to the remaining
compact factors. The adelic closure already contains
the kernel of that projection by isolation of the other
factors, so the closure is the full adelic group
outside \(S\). This completes outer type \(A\)
over number fields.

\medskip\noindent
\textit{Globally anisotropic type \(C\) over a number field.}
Use the symplectic involution model \(G=U(A,\ast)\)
constructed above. We first obtain a Hermitian
description over the global division factor, valid
also in characteristic two. There is \(z\in A\) with
\(z+z^\ast=1\): for a split alternating adjoint take
the projection onto one of two paired Lagrangian
spaces, and solubility of this linear equation
descends to \(k\). Choose a primitive idempotent
\(e\) of \(A\). If \(e^\ast e=0\), set \(a=eze^\ast\)
and \(p=e-a\). The equations
\[
 a+a^\ast=ee^\ast,\qquad a^2=0,\qquad ea=ae^\ast=a,
 \qquad ae=e^\ast a=0
\]
give \(p^2=p\), \(p^\ast p=pp^\ast=0\), and
\(p\ne0\). Then
\(t p+t^{-1}p^\ast+(1-p-p^\ast)\) is a nontrivial
split torus in \(G\): multiplication follows from
the three orthogonal idempotents, and the involution
interchanges the two inverse scales. Thus global
anisotropy makes \(b=e^\ast e\) a nonzero self-adjoint
element of rank one over the division factor. It cannot
be nilpotent. Otherwise \(b^2=0\) because its image is a
line over a division ring. A projector \(e_0\) onto
that line can be written \(e_0=ba_0\) by lifting its
values along \(b\). Then \(e_0^\ast e_0=a_0^\ast b^2a_0=0\),
contradicting the same split-torus argument. The polynomial
projection onto the image of \(b\), as above, yields a
self-adjoint primitive
idempotent \(e'\). Its \(\ast\)-stable corner
\(D=e'Ae'\) is the global division factor, with
involution \(\sigma=\ast|_D\) of the first kind.
As in the outer type \(A\) calculation, write
\(A=\operatorname{End}_D(V)\) and let \(\dagger\)
denote transpose with \(\sigma\) on entries.
Skolem--Noether gives \(a^\ast=H^{-1}a^\dagger H\)
with \(H^\dagger=\epsilon H\), where \(\epsilon=\pm1\)
since the involution fixes the centre. The form
\(h(x,y)=x^\dagger H y\) is nondegenerate and is
\(\epsilon\)-Hermitian.

If \(h(x,x)=0\) for a nonzero vector \(x\), a projector
\(e\) onto its \(D\)-line satisfies \(e^\ast e=0\):
the form pairs any two vectors in that line to zero.
This again gives a split torus. Thus \(h(x,x)\ne0\)
for every nonzero rational vector. The action of \(G\)
on \(V\) has determinant one after splitting \(D\):
it is the standard symplectic action repeated in
\(\deg D\) columns. Choose a free integral lattice
making the components of \(h\) integral. As in the
unitary case just proved, some coefficient of
\(h(x,x)\) gives the archimedean product bound for
each nonzero \(x\in L_g\). In bounded bases normalize
the right multiplications by \(D\) and the \(D\)-valued
form \(h\). Comparing bases with the same transformed
tensors gives a rational \(D\)-linear isometry, which
already belongs to \(G=U(A,\ast)\). This proves the
compactness of the adelic quotient by the same
compact representative argument.

Take absolute type \(C_r\) with \(r\geq2\), since
\(C_1\) is inner type \(A_1\). At every finite place
the local calculation of the involution above gives
either the split symplectic group or a quaternionic
Hermitian form of dimension \(r\).
A nondegenerate quaternionic Hermitian form in
characteristic zero is diagonalizable by successively
taking the perpendicular of a vector of nonzero norm.
For its first two diagonal scalars \(b_1,b_2\),
the equation \(b_1 N(x)+b_2=0\) has a solution
by the surjectivity of the local quaternion norm
proved above. It gives a nonzero isotropic vector.
At the real places split symplectic groups are
isotropic, and quaternionic Hermitian groups are
compact in the definite case and isotropic in the
indefinite case, by the same two-coordinate equation
using positive quaternion norms. The complex groups
are split.

At every isotropic completion the local generation
proved above for \(C\) gives the contracted unipotent
groups. The lattice projection and isolation arguments
apply as in outer type \(A\), using the characteristic
zero Lie simplicity after scalar extension and the
same polynomial roots for unipotents. Thus the adelic
closure contains the product of all outside noncompact
archimedean factors and the restricted product of all
outside finite factors. At the remaining compact real
factors weak approximation follows directly from the
Cayley parameter argument above, since \(G\) is the
full isometry group of the symplectic involution. The
affine space of solutions of \(y+y^\ast=1\) is nonempty
by the construction of \(z\), and over the algebraic
closure the split symplectic torus has elements with
no eigenvalue \(1\). Thus the same open chart in the
geometrically connected group is nonempty.
This completes type \(C\) over number fields.

\medskip\noindent
\textit{Local quotients and anisotropic cases over function fields.}
The small logarithms used above are specific to characteristic
zero. Here is a replacement for certain local factors in any
characteristic, using the rescaling and mixing argument of
\href{https://doi.org/10.1007/BF01081890}{Margulis, Theorem~1}.
Let \(F\) be a nonarchimedean local field. We use a split
simply connected absolutely almost simple group \(J/F\),
the special unitary group of an isotropic Hermitian
form over a separable quadratic field extension of \(F\),
the spin group of an isotropic regular quadratic
space of dimension at least five as above, or the
quaternionic spin group treated below in characteristic
different from two.
Suppose \(H\subseteq J(F)\) is closed, nondiscrete, and
\(H\backslash J(F)\) has a finite nonzero invariant measure.

First a regular diagonal element \(a\) of a maximal split torus,
obtained from a cocharacter nontrivial on every root by
evaluation at a uniformizer, acts mixingly on the quotient.
In any continuous unitary representation, a vector fixed
by the upper unipotent subgroup of a root \(\operatorname{SL}_2\)
is fixed by that entire rank-one group. In fact, writing \(u\)
and \(v\) for upper and lower transvections, respectively,
\[
 u((t-1)/s)\,v(s)\,u(-(t-1)/(st))
   =\begin{pmatrix}t&0\\s&t^{-1}\end{pmatrix}
     \ \longrightarrow\ \operatorname{diag}(t,t^{-1})
 \quad(s\to0).
\]
The matrix coefficient on the given vector tends to its squared
norm because the outer factors fix it. Thus the diagonal fixes
it. Choosing \(t\) with absolute value different from one and
using diagonal conjugation shows that the opposite transvections
fix it too. The argument applies after the root homomorphism
of Results~\ref{res:C07} and \ref{res:C08}.

For the right translation representation on
\(L^2(H\backslash J(F))\), let \(B\) be any weak limit of a
subsequence of the operators \(a^m\), \(m\to\infty\).
Such limits along subsequences exist by boundedness and
separability. If \(a^{-m}u a^m\to1\), then \(uB=B\), by
strong continuity applied before \(a^m\). In the split case, for each pair
of opposite roots one of their two root groups has this
contraction property. The rank-one argument makes every
vector in the image of \(B\) fixed by all the root groups,
hence by \(J(F)\), using the split Bruhat decomposition
and the simple-coroot basis as above.
For the unitary group choose a hyperbolic pair \(e,f\)
paired to \(1\), and include its inverse-scaling
torus in the maximal split torus, extending the pair
to a Witt basis. Replacing \(f\) by \(f\eta\) for
\(\bar\eta=-\eta\ne0\) makes the Gram matrix on the
pair a scalar alternating matrix on \(F\)-coordinates.
In characteristic two take \(\eta=1\).
Thus the \(F\)-matrices of the pair give a root
\(\operatorname{SL}_2\) inside the unitary group,
with determinant one over the quadratic field.
One of its upper and lower groups is contracted by
\(a^{-m}\) as \(m\to\infty\). The rank-one argument
shows that \(B\) has image fixed by the pair's diagonal
torus. Its two signs contract the two opposite line
parabolic radicals. They generate \(J(F)\) by the
relative-root and unitary computations above, so again
the image is \(J(F)\)-invariant.
For the spin case choose the hyperbolic pair \(e,f\)
in a Witt basis and \(w\) in its perpendicular space
with \(q(w)\ne0\). The even Clifford elements
\[
 ef,\quad ew,\quad -fw/q(w),\quad fe
\]
form the four matrix units of a copy of
\(\operatorname{Mat}_2(F)\), in the order
\(E_{11},E_{12},E_{21},E_{22}\), by the Clifford
relations. Its determinant-one subgroup lies in the
spin group because its elementary generators are the
Clifford shears displayed above. Its diagonal is the
spin cocharacter for the pair, and one of its two
transvection groups is contracted by \(a^{-m}\).
As before the Mautner calculation then makes the
image of \(B\) fixed by both opposite line radicals.
They generate the spin group by the preceding
Clifford generation proof. In the quaternionic
case the scalar matrices on a skew-Hermitian
hyperbolic pair give an \(\operatorname{SL}_2\)
in the orthogonal involution group. Its
identity-component preimage under the central
double cover is isomorphic to it, since
\(\operatorname{SL}_2\) is already simply
connected; thus it has a canonical spin lift.
Its upper and lower transvections lie in the
opposite line radicals, and its diagonal
contracts the two radicals with opposite signs.
The rank-one fixed-vector argument applies
to this subgroup and the two radicals
generate densely by the proof below.
In a transitive
quotient these invariant functions are constant almost
everywhere. For example, apply Fubini to the equality
\(f(xg)=f(x)\) and Haar measure in \(g\). The orbit map
for each \(x\) carries the Haar null sets of inverse
images to the same null-set criterion on the quotient,
by transitivity and left translation of Haar measure.
That criterion agrees with quotient measure zero by
invariance and Fubini against a positive integrable
Haar density. This proves the constancy. On the
orthogonal complement of constants every weak limit
is therefore zero, which gives mixing.

Choose noncentral \(h_i\in H\) tending to \(1\).
Embed \(J\) as a closed subgroup of \(\operatorname{SL}_N\)
and use weight coordinates for conjugation by the
split torus on matrices. Choose a relatively compact
identity neighborhood with compact closure \(C\) in
\(J(F)\). Every nonempty open subset of \(J(F)\) is
Zariski dense. Smooth local \'{e}tale coordinates give
an analytic open subset in affine space by the inverse
function theorem. The image of a proper algebraic
closed subset in those coordinates has lower dimension,
so cannot contain such an open subset. For each coordinate consider the
function of \(y\in C\) given by the corresponding
coefficient of \(y^{-1}(h_i-1)y\). These functions all
belong to a fixed finite-dimensional space of regular
functions restricted to \(C\). Its supremum norm is
equivalent to a coordinate norm: finitely many evaluation
maps separate the functions, by choosing each new
evaluation to cut down their common kernel. After
passing to a subsequence, each coordinate function is
either identically zero or always nonzero in that space.
Divide the nonzero ones by a scalar of absolute value
equal to their supremum norm. The resulting functions
converge, after another subsequence, to nonzero functions
in that space. This uses compactness of its unit sphere
over \(F\); the norms are attained on \(C\) and have
values in the absolute-value group. Choose \(x\) in the
interior of \(C\) avoiding the zero sets of the finitely
many limits. As \(y_i\to x\), the change from \(x\) to
\(y_i\) in each coordinate function is then little-oh
of its value at \(x\), by uniform convergence and the
nonvanishing of the limit at \(x\).

Some nonzero-weight coordinate is nonzero along this
subsequence. Otherwise each \(h_i\) centralizes every
conjugate of the split torus, first for conjugations
in \(C\), then by density for all conjugations. These
conjugates topologically generate \(J(F)\): commutators
of a split torus element with a root group lie in the
group they generate, and the closure contains every
contracted radical by taking commutator limits with
powers of the contracting element. The local
generation results above apply. This would
make \(h_i\) central. Choose the regular cocharacter
defining \(a\) so that at least one of the nonzero
coordinates expands under conjugation by \(a^{-1}\).
Such integral cocharacters exist by avoiding finitely
many root hyperplanes in the desired open half-space.
In the matrix coordinate maximum norm take a ball of
radius \(r\) about \(1\) excluding all other central
points. Let \(m_i\) be the last integer before the
first exit of \(a^{-m}x^{-1}h_i x a^m\) from this ball.
Then \(m_i\to\infty\). At time \(m_i\) its distance
from \(1\) is between \(r/M\) and \(r\), where \(M\)
is the largest absolute-value multiplier of conjugation
by \(a^{-1}\). A subsequence converges to a noncentral
\(u\). For any \(y_i\to x\),
\(a^{-m_i}y_i^{-1}h_i y_i a^{m_i}\) has the same limit:
the coordinate changes are little-oh of the corresponding
coordinates, whose rescaled values are bounded.

Fix \(g\in J(F)\). Mixing applied to the images of a
neighborhood of \(x\) and its translate by \(g\) in
the quotient shows that, for every such neighborhood,
\(y a^{m_i}=b z g\) for large \(i\), with \(y,z\)
in the neighborhood and \(b\in H\). The images are
open and of positive measure, since the invariant
measure has full support. Use a sequence of shrinking
neighborhoods and pass to a subsequence to have
\(y_i,z_i\to x\). Then
\[
 b_i^{-1}h_i b_i
   = z_i g a^{-m_i}y_i^{-1}h_i y_i a^{m_i}g^{-1}z_i^{-1}
   \longrightarrow xgug^{-1}x^{-1}.
\]
Thus \(x^{-1}Hx\) contains all conjugates of the noncentral \(u\).

For completeness their closed normal closure is all of \(J(F)\).
Modulo its center, \(J(F)\) acts faithfully and primitively
on the cosets of any maximal proper standard parabolic.
To verify maximality as an abstract subgroup, any subgroup
containing the minimal rational parabolic contains the
relative Weyl representative of each Bruhat cell it meets.
If it contains the representative of \(w\ne1\), choose a
simple relative root \(\alpha\) with \(w^{-1}\alpha<0\).
Conjugating the positive ray group of \(-w^{-1}\alpha\)
supplies the negative ray group of \(-\alpha\).
The two rays give a representative of \(s_\alpha\) by
the rank-one Bruhat calculation above, and \(s_\alpha w\)
has smaller length.
Induction gives exactly the standard parabolic associated
with the simple representatives it contains. The kernel
of the action is central: density makes its elements
belong to every geometric conjugate of the parabolic,
whose intersection is a proper normal algebraic
subgroup and so has only finite, central geometric points.

The center of the parabolic's unipotent radical is
an abelian normal subgroup of the stabilizer. In the
split case it contains a root group of maximal weight
for a cocharacter giving that parabolic, since the
commutator formulas give higher weights against all
other positive weights. The normal closure of one
full root group contains its opposite by rank-one
conjugation, and thus its coroot torus by the rank-one
word. Commuting with that torus gives every root group
with nonzero coroot pairing. Some simple root has such
pairing, and connectedness of the diagram then gives
all the simple and opposite simple groups.
In the unitary case take the maximal proper parabolic
stabilizing the line of \(e\). The center of its
radical contains the upper transvections of the
pair's \(\operatorname{SL}_2\): the matrix with only
its far off-diagonal entry commutes with the other
upper block shears. Its normal closure therefore
contains the opposite transvections and the
inverse-scaling torus. The closed normal closure
contains the radicals contracted by that torus
and by its inverse, giving the full unitary group
as above. For the spin group also use the maximal
proper parabolic of the isotropic line \(Fe\).
Its radical is abelian, by multiplication of
the matrices \(1+ew\) in the Clifford algebra.
It contains the upper transvections of the pair's
\(\operatorname{SL}_2\), so its closed normal closure
again contains the pair cocharacter and both
radicals. In the quaternionic case the
maximal parabolic of a quaternionic isotropic
line likewise has the upper scalar transvections
in the center of its radical, as one sees
from the block shears below. It is maximal
because after splitting it is the parabolic
of an isotropic two-plane in dimension at
least eight. So the same generation by
conjugates applies. In every case the stabilizer's abelian
subgroup topologically generates \(J(F)\) under
conjugation. The group is topologically perfect:
every contracted radical lies in the closure of
the commutator subgroup by the same contraction
formula. A nontrivial closed normal subgroup of the faithful
primitive quotient acts transitively, its orbits being
an invariant partition. Modulo that subgroup every
conjugate of the abelian subgroup in the stabilizer has
the same image, since the group is the product of the
normal subgroup and the stabilizer. The further quotient
is then abelian and must be trivial by topological
perfection.
The closed normal closure of \(u\) therefore surjects to the
central quotient. The group is its product with the
center, so another application of topological perfection makes
the normal closure the full group. This proves \(H=J(F)\).

We apply this to the anisotropic groups of types \(A,C,G_2\) over
a global function field \(k\). For their compactness
argument, replace \(\mathbb Z\) in the bounded-basis proof
by \(R_0=\mathbb F_q[t]\), choosing a separating
parameter \(t\) and writing \(k_0=\mathbb F_q(t)\).
Take the free order or norm lattice over the integral
closure of \(R_0\) in \(k\). It is free over \(R_0\)
by the trace-dual lattice argument above. The finite
local lattices between \(M O_v\) and \(M^{-1}O_v\)
are again recovered from their rational intersection
by the Chinese remainder theorem. A rational basis
matrix has determinant of absolute value one at
the infinity of \(k_0\), since its determinant has
valuation zero at every finite place of \(k_0\).
The determinant is unaffected by \(g\), as before.
The archimedean norm bound is replaced by the absolute
value at this single infinity of a nonzero norm polynomial.
For the division algebra and the octonion algebra the
polynomial is \(N_{k/k_0}(\operatorname{Nrd}(x))\) or
\(N_{k/k_0}(q(x))\), respectively. For a Hermitian form
use \(N_{k/k_0}(c(x))\) where \(c(x)\) is a nonzero
coordinate of \(h(x,x)\), from the finite list of norm
coordinates over \(k\).
That absolute value is at least one and the invariant
polynomial is bounded by a constant times a fixed
power of the maximum coordinate norm. Here all the
places of \(k\) over infinity are included together
in the representation over \(k_0\).

The bounded-basis proof over \(\mathbb F_q((1/t))\)
is the following. The polynomial ring has compact
fundamental domain \(t^{-1}\mathbb F_q[[1/t]]\)
for addition, by taking the polynomial part. A ball
of volume greater than the lattice covolume gives
a bound for the shortest vector, just as over
\(\mathbb R\). A shortest vector is primitive since
dividing by a nonconstant polynomial would shorten
it, and integral coordinate operations over
\(\mathbb F_q[[1/t]]\) put it in the first coordinate
with the same maximum norm. The projected lattice
in the other coordinates has determinant norm
divided by the length of the first vector.
Every nonzero projected vector has at least that
length: lift it and subtract a polynomial multiple
of the first vector so that the absolute value in
the first coordinate is strictly smaller. This
would contradict minimality if the projected length
were smaller. Induction now gives bounded bases and
their bounded inverses, using determinant norm one.
The finite local basis matrices after multiplication
by \(g\) are integral with integral inverse at all
finite places of \(k_0\). Thus the same tensor
comparison takes place in a compact adelic matrix set
over \(k_0\).

For \(\operatorname{SL}_1(D)\) use the right-multiplication
matrices exactly as in the number-field proof. Finitely
many rational reduced norms in the compact norm image
allow the same correction. For outer type \(A\), the
self-adjoint-corner and Hermitian tensor argument above
works over the function field: the trace polarization
uses a central scalar moved by the involution, the
hyperbolic adjustment uses a trace-one scalar, and the
reduced-norm correction has finitely many roots of unity.
For type \(C\) use its Hermitian tensor argument, whose
construction of the global division factor also works
in characteristic two. For \(\operatorname{Aut}(A)\)
with anisotropic octonion \(A\), use its multiplication
and scalar matrices, as above. All these comparisons
give compact adelic quotients, using discreteness of a
global field in its adeles and the displayed bounds.

There are no archimedean places of the function field.
For every split finite factor, and for each special
unitary factor of an isotropic Hermitian form over a
quadratic field, take \(v\) outside \(S\) and use
\(T=S\cup\{v\}\) and the compact-open product elsewhere.
The resulting lattice quotient \(G_T/\Lambda\) is
compact, by the earlier open-orbit argument, and
its projection to \(G(k_v)/\overline{\Lambda_v}\)
carries the finite invariant measure. The projection
is nondiscrete by the same compactness contradiction
using the noncompactness of \(G_S\) and injection of
rational points into the \(v\)-factor. The local
argument just proved therefore makes the
projection dense. Isolation from the compact
product uses the same contracting root groups as
before, in any characteristic.

For inner type \(A\) the local algebra is split at
all but finitely many places by the order and lifted
idempotent argument. The remaining finite set is
handled by weak approximation of units and products
of their local commutators, which was proved above
also in positive characteristic.
For outer type \(A\), all places where the quadratic
extension is a field give split central algebras and
isotropic Hermitian forms by the local calculations
above. They have just been isolated, as have the
split factors at the places where the extension splits.
Only finitely many factors remain, each with full
unitary group the units of a local central simple
algebra. Its reduced-norm-one group is in the closure
of its commutator group by the earlier local proof,
including for matrix algebras over a division factor.
The translated Cayley chart gives simultaneous weak
approximation in the full unitary group, so those
commutators again give density in the remaining factors.
For type \(C\), the defining central simple algebra
splits outside finitely many places, by the same order
and idempotent argument. There the symplectic
involution group is split. At the finitely many other
factors its rational points have simultaneous weak
approximation by the translated Cayley chart, since
it is the full isometry group. This does not require
division by two or that the remaining factors be compact.
For \(G_2\) every completion is split. Here the local
octonion calculation is valid also in characteristic
two: pick a trace-one nonscalar element, so its
quadratic equation gives a separable quadratic
subalgebra in characteristic two. In characteristic
different from two one can instead choose a
nonzero trace-zero vector of nonzero norm.
In either case an anisotropic octonion would give
a quadratic field with nondegenerate norm. Pick a
nonzero-norm vector in its perpendicular space.
The internal doubling identities of
\emph{Nonassociative Algebras}, Result NA13-C03.07,
apply here in every characteristic: they use the
polarized inverse and quadratic equations, orthogonality,
and the Moufang identities, with no division by two.
The resulting four-dimensional subalgebra is the
quaternion algebra of that quadratic field and the
square of the perpendicular vector. Doubling once
more gives the norm \(N_B(x)-cN_B(y)\). The
local quaternion norm is onto by the cyclic and
principal-unit computation above, which uses a
separable residue extension and works in all
characteristics. Thus the octonion norm is
isotropic, a contradiction. The split-octonion
argument in NA13-C04.08 then applies. This leaves
no remaining factors for \(G_2\) and completes the
function-field cases under consideration.

\medskip\noindent
\textit{Quadratic spin groups over function fields.}
Let \(q\) be a globally anisotropic regular quadratic
form of dimension at least five in the spin models
above. In characteristic two we allow either
nondegenerate polar form or its one-dimensional
radical \(kr\) with \(q(r)\ne0\). For the quadratic
isometry group the polynomial-lattice compactness
argument applies with the form and scalar tensors.
In the polar-radical case include the fixed vector
\(r\) as an additional tensor and use its stabilizer.
Thus, writing \(O\) for the corresponding isometry
group, \(O(\mathbb A_k)/O(k)\) is compact. In this
argument determinants have absolute value one:
preservation of a nondegenerate polar form gives
determinant \(\pm1\), and in the radical case the
determinant on the quotient is one and \(r\) is fixed.
The same holds at the auxiliary infinity places.

The map from \(\operatorname{Spin}(q)\) to the
identity component \(O^\circ\) is the finite
central covering with kernel \(\mu_2\).
For the regular odd characteristic-two model this
was checked above using the stabilizer of \(r\).
For nondegenerate polar form it may be checked over
an algebraic closure in the split Clifford model.
The even spin group there is smooth by the shear
charts in Appendix C of \emph{Nonassociative Algebras}.
Every even product of orthogonal symmetries has a
spin lift after adjoining a scalar square root, and
the kernel of the vector action is the scalar
norm-one group by the Clifford center calculation.
In characteristic two the orthogonal group with
nondegenerate polar form is smooth: its orbit on
nondegenerate quadratic forms over an algebraic
closure is the entire open set, and the derivative
of change of basis is onto the space of quadratic
forms. Indeed with polar matrix \(B\) the derivative
is \(x\mapsto x^{\mathsf t} B X x\), and \(BX\)
ranges over all matrices. The connected spin image
has the dimension of \(O^\circ\) by the split charts,
so it is \(O^\circ\). These finite flat central
coverings descend to \(k\).

Apply the compact representative argument to
\(\pi(g)\in O^\circ(\mathbb A_k)\), choosing a rational
isometry \(d\in O(k)\) with \(\pi(g)d\) in a fixed
compact set. There are only finitely many rational
components modulo \(O^\circ(k)\), since they inject
into the finite set of geometric components.
Adjusting by a fixed rational representative per
component therefore gives \(d\in O^\circ(k)\) with a
new fixed compact set, intersected with the closed
subgroup \(O^\circ(\mathbb A_k)\).
The lift obstruction for \(d\)
is a class in \(k^\times/k^{\times2}\) by the
finite flat scalar torsor calculation above.
The obstruction is multiplicative: products of
local lifts multiply their central scalar descent
data. Its local value equals that for \(\pi(g)d\).
All but finitely many places have integral
representatives for the latter points and for the
finite flat covering model, exactly as in the
necessity argument. Thus each obstruction has even
valuation outside one fixed finite set. The square
class finiteness proved there works over function
fields and gives finitely many obstructions. Correct
by a fixed rational isometry of each class and lift.
Compactness of the inverse image under the finite
adelic covering follows from the monic-coordinate
argument above, including for an inseparable
covering. This makes the spin adelic quotient compact.

Every completion is isotropic in this dimension range.
In odd characteristic the ternary and binary
quaternion argument for number fields uses only a
nonarchimedean local field of characteristic different
from two, so applies unchanged. We give the
characteristic-two argument separately. For a
separable quadratic extension \(E/F\) of such local
fields its norm subgroup has index two in \(F^\times\).
If the extension is unramified this is the even
valuation condition, since residue and principal
unit norms are surjective. Suppose it is ramified
and choose a uniformizer \(\pi\) of \(E\), with
nontrivial conjugation \(\sigma\). Write \(d\) for
the normalized \(E\)-valuation of \(\sigma(\pi)-\pi\).
In characteristic two this difference is the trace
of \(\pi\), so \(d\) is even. The basis \(1,\pi\)
has coefficients integral exactly for elements of
\(\mathcal O_E\), since the two contribution valuations
have different parity. The trace of \(a+b\pi\) is
\(b\operatorname{Tr}(\pi)\). Consequently
\[
 \operatorname{Tr}(\mathfrak p_E^M)
   =\mathfrak p_F^{\lfloor(M+d)/2\rfloor}.
\]
For \(N>d\) and \(M=2N-d\), the norms of
\(1+\mathfrak p_E^M\) form \(1+\mathfrak p_F^N\).
In fact \(N(1+x)=1+\operatorname{Tr}(x)+N(x)\),
and \(v_F(N(x))=v_E(x)\). For each \(j\geq N\)
choose \(x\in\mathfrak p_E^{2j-d}\) to hit the
specified trace coefficient at \(j\); the norm term
has larger valuation. Successive corrections
converge and prove surjectivity.

It remains to count the unit norm index modulo
these deep unit groups. A norm-one element is of
the form \(\sigma(x)/x\) by the quadratic identity
using \(x=1+u^{-1}\) for a norm-one element \(u\ne1\).
If \(x\) has odd valuation, scale it by \(F^\times\)
to write \(x=\pi v\) with \(v\) a unit. Then
\(\sigma(v)-v\) has valuation at least \(d\),
by the basis expression, and
\(\sigma(x)-x=(\sigma(\pi)-\pi)\sigma(v)+
\pi(\sigma(v)-v)\) has valuation exactly \(d\).
Hence its ratio difference has valuation \(d-1<M\).
If \(x\) has even valuation, scale it to a unit and
write it as \(a+b\pi\) with \(a\) a unit.
The condition \(\sigma(x)/x\in1+\mathfrak p_E^M\)
is then \(b\in\mathfrak p_F^{(M-d)/2}\).
For \(c=(M-d)/2\), the number of unit classes modulo
these elements is
\[
 \frac{|(\mathcal O_E/\mathfrak p_F^c\mathcal O_E)^\times|}
 {|(\mathcal O_F/\mathfrak p_F^c)^\times|}
 =q_F^c ,
\]
where \(q_F\) is the residue field size.
Thus the norm-one group has
\(2q_F^{(M-d)/2}\) distinct images in
\(\mathcal O_E^\times/(1+\mathfrak p_E^M)\).
These are exactly the kernel of the norm to
\(\mathcal O_F^\times/(1+\mathfrak p_F^N)\),
by the deep norm surjectivity. The index of its
image is therefore
\(2q_F^{N+(M-d)/2-M}=2\).
Norm valuations are all integers in a ramified
quadratic extension. This proves the norm index.

Now split off two nondegenerate binary planes in
the regular quadratic form. If they have no isotropic
vector, each is a scalar multiple of the norm of a
separable quadratic field: the mixed coefficient
in a binary plane is nonzero, so the defining
quadratic equation is separable. Their nonzero
value sets are cosets of index-two subgroups of
\(F^\times\). If the cosets intersect, the sum of
the planes is isotropic in characteristic two.
Otherwise the two underlying subgroups are the
same, since distinct index-two subgroups have
product \(F^\times\). The disjoint cosets then
together represent every nonzero value. Taking a
vector of nonzero norm in the remaining perpendicular
space, which can include the regular radical line,
gives an isotropic vector of the full form. This
proves all the required characteristic-two local
isotropy assertions.

Thus every completion outside \(S\) is one of the
isotropic spin groups of the local quotient argument.
Its projected arithmetic lattice has compact quotient
and nondiscrete projection, using the compact spin
adelic quotient and noncompactness at \(S\).
The mixing and closed-quotient argument gives its
full local projection. Isolation by the contracting
Clifford radicals now supplies each outside factor,
and their finite supports are dense in the restricted
product. This completes the quadratic spin models
over function fields.

\medskip\noindent
\textit{The anisotropic compactness reduction over number fields.}
There is a general characteristic-zero compactness argument
that does not require a particular representation of the
original group. Let \(G\) be a connected semisimple
anisotropic group over a characteristic-zero field \(k\),
and write \(\mathfrak g=\operatorname{Lie}G\).
We first show that no nonzero \(e\in\mathfrak g(k)\) has
nilpotent adjoint action.

The Killing form \(B\) is nondegenerate, as follows
after scalar extension from the absolute root
decomposition in characteristic zero. If
\(\operatorname{ad}e\) is nilpotent and \([e,y]=0\),
then \(\operatorname{ad}e\) and \(\operatorname{ad}y\)
commute, so their product is nilpotent and
\(B(e,y)=0\). The annihilator of the kernel of
\(\operatorname{ad}e\) under \(B\) is its image,
by invariance and dimension. Thus there is a
rational \(h\) with \([h,e]=2e\).

Let \(D_s\) be the semisimple part of
\(\operatorname{ad}h\) as a linear operator.
It is defined over \(k\), since the characteristic
is zero. Write \(\mathfrak g_i\) for its
eigenspaces over a finite Galois splitting field,
with distinct eigenvalues \(\lambda_i\).
These are the generalized eigenspaces of
\(\operatorname{ad}h\). The bracket of a
\(\lambda_i\)-space and a \(\lambda_j\)-space lies
in the \(\lambda_i+\lambda_j\)-space, or is zero
if that value is absent: expand iterates of
\(\operatorname{ad}h-\lambda_i-\lambda_j\)
using the derivation rule. Define the torus of
scalars \(z_i\) on these spaces by the equations
\[
 \prod_i z_i^{m_i}=1
 \qquad\text{whenever}\qquad
 \sum_i m_i\lambda_i=0,\quad m_i\in\mathbb Z.
\]
The relation lattice is saturated, since the
additive group of the splitting field has no
torsion. It is invariant under the Galois
permutation of the \(\lambda_i\), so this is a
torus over \(k\). It preserves brackets by the
displayed weight rule. Its tangent space contains
the scalar tuple \((\lambda_i)\).

The connected automorphism group of the Lie
algebra is \(G_{\mathrm{ad}}\). To see this,
every derivation \(D\) differs from an inner one
by a derivation \(D_0\) orthogonal to all
\(\operatorname{ad}y\) under the trace pairing,
by nondegeneracy of \(B\). Then
\[
 \begin{aligned}
 B(D_0x,y)
 &=\operatorname{Tr}([D_0,\operatorname{ad}x]\operatorname{ad}y)\\
 &=\operatorname{Tr}(D_0\operatorname{ad}[x,y])=0 .
 \end{aligned}
\]
Hence \(D_0=0\). The tangent dimension of the
automorphism group is therefore \(\dim G\),
which is already the dimension of the adjoint
image. The latter lies in the identity component
and equals it. This is also the inner-derivation
calculation in \emph{Lie Algebras and Finite-Dimensional
Representations}, Result B07. The scalar torus just constructed
is consequently a \(k\)-torus in \(G_{\mathrm{ad}}\).
The eigenvalue \(2\) is present, since
\([h,e]=2e\) and \(e\ne0\). Its coordinate
character is Galois invariant and nontrivial,
since its differential takes value \(2\) on
\((\lambda_i)\). Averaging a cocharacter with
nonzero pairing against it over the finite
Galois group gives a nonzero invariant
cocharacter. Thus \(G_{\mathrm{ad}}\) has a
nontrivial split torus. Its inverse image under
the central isogeny has an isogenous torus
as identity component, with the same rational
cocharacter space, so \(G\) is isotropic.
This contradiction proves the claim.

Now take \(k\) to be a number field and apply
the real bounded-basis argument to the adjoint
representation of \(G_{\mathrm{ad}}\) on
\(\mathfrak g\). Take a free integral lattice
with integral bracket constants. For a nonzero
rational \(x\) in \(L_g\), some nonleading
coefficient of the characteristic polynomial of
\(\operatorname{ad}x\) is nonzero by the preceding
claim. Every such coefficient is integral at
the finite places: compute it on \(g_vx\)
in the integral bracket lattice. These coefficients
are invariant under \(g\), are homogeneous of
positive degrees, and have the archimedean
polynomial bounds used above. The product
formula bound for the nonzero integral coefficient
therefore gives a common positive lower bound
on lengths of nonzero vectors of \(g_\infty L_g\).
The determinant of the adjoint representation
is one, so the covolume is one. Use the bracket
tensor and scalar multiplications by a basis of
\(k\) as the bounded-basis comparison tensors.
The rational comparisons are \(k\)-linear Lie
automorphisms. There are only finitely many
components of the algebraic automorphism group,
and rational points in the geometric identity
component lie in \(G_{\mathrm{ad}}(k)\).
Adjusting by one rational representative per
occurring component gives compact representatives
for \(G_{\mathrm{ad}}(\mathbb A_k)/G_{\mathrm{ad}}(k)\).

The finite central covering \(\pi:G\to G_{\mathrm{ad}}\)
allows the same compactness conclusion for \(G\).
For \(g\in G(\mathbb A_k)\) choose
\(d\in G_{\mathrm{ad}}(k)\) with \(\pi(g)d\)
in the fixed compact set. The obstruction
\(\delta(d)\) to lifting is a torsor for
\(Z=\ker\pi\), and the obstruction is
multiplicative because its descent data are
central. Choose one finite Galois extension
\(E/k\) where the kernel is a product of
\(\mu_{n_i}\) with all their geometric points
rational. Outside a fixed finite set the
integral torsor argument in the necessity
proof shows that the obstruction coordinates over
\(E\) have valuations divisible by \(n_i\):
locally they equal those of the compact
representatives \(\pi(g)d\). The arithmetic
finiteness argument there gives finitely many
possible restricted torsor classes. The kernel of
restriction on the global classes is finite.
Indeed, a cocycle trivial on restriction can be
adjusted by a coboundary to be trivial on the
absolute Galois group of \(E\), and then factors
through the finite Galois group of \(E/k\) with
values in the finite \(Z(E)\). Thus only finitely
many \(\delta(d)\) occur. Correct by a fixed
rational adjoint point of each class and lift
to \(G(k)\). The finite adelic covering has
compact inverse images of compact sets by the
monic-coordinate argument above. This proves
compactness of \(G(\mathbb A_k)/G(k)\) without
a restriction on the absolute type.

The same invariant polynomials also give the
local compactness test in characteristic zero.
At an anisotropic completion their maximum
is bounded away from zero on the compact
coordinate unit sphere of the Lie algebra.
If an adjoint orbit were unbounded, scale an
unbounded sequence of orbit vectors to that
sphere. The positive-degree invariant values
then tend to zero, a contradiction. Hence
the images of basis vectors are bounded.
The adjoint determinant is one, so the
inverse matrices are bounded too. The
finite covering then gives compactness
of \(G(k_v)\). Conversely, a nontrivial
split cocharacter gives an unbounded
absolute weight and a noncompact group.

For example, if \(S\) contains the archimedean
places, the compactness reduction leaves only
local generation and finitely many places of
weak approximation to check. To see this
precisely, take an outside noncompact place
\(v\) and \(T=S\cup\{v\}\) in the earlier
lattice construction. The projected closure
\(H\subset G(k_v)\) is nondiscrete and
cocompact. The proof by projective recurrence
applies using a split cocharacter \(\lambda\):
its conjugates fix the space of small
logarithms of \(H\), and their closed
normal closure contains the contracting
radicals by contraction of commutators.
One such nonzero radical and its conjugates
span the whole Lie algebra. Indeed their
Lie span is an invariant ideal by Zariski
density of local points, and absolute
simplicity in characteristic zero. After
restriction from \(k_v\) to \(\mathbb Q_p\)
the Lie algebra is simple by the transitive
factor argument above. Thus the logarithms
of \(H\) span it and \(H\) is open by the
same inverse-function argument. It has
finite index by cocompactness. Every
unipotent element of a parabolic radical
has roots of every positive order by its
polynomial powers, so the action on the
finite coset set kills it. Hence \(H\)
contains the subgroup generated by the
radicals. In the isolation argument the
projection of \(P\) contains \(H\).
For each contracting element some positive
power belongs to \(H\), and that power
still contracts. The compact quotient
argument therefore isolates every
outside local subgroup generated by
the parabolic radicals.

At almost all finite places the group is
quasi-split. The split model and its variety
of Borel subgroups, together with descent
over a finite Galois splitting field,
extend over the integers outside finitely
many primes, by clearing denominators in
the finite descent identities and removing
the ramified primes. The models are smooth
there by the split charts and \'{e}tale
descent. Over a finite residue field a
Borel exists by Lang's theorem,
Result~\ref{res:F08}: conjugate a geometric
Borel by a Lang preimage of its Frobenius
conjugating element. A residue Borel
lifts by smoothness over the complete
valuation ring. At every such place
the local generation for the quasi-split
case above isolates the full factor.
Its restricted product is contained in
the adelic closure by density of finite
supports. Thus simultaneous weak
approximation at the remaining finite
set of places would finish this
number-field case with all archimedean
places omitted.

\medskip\noindent
\textit{Strong approximation in type \(F_4\).}
By F06 and the split automorphism calculation of
\emph{Nonassociative Algebras}, Result NA13-C04.09,
every form of the split \(F_4\) is the automorphism
group of an Albert algebra \(J\): the center and
diagram automorphisms are trivial, so the cocycle
descends the split algebra and its cubic data.
Use its norm \(N\), adjoint \(x^\#\), trace \(t\),
and trace pairing \(T\) from NA13-C04.02.

If \(J\) has a nonzero vector of norm zero,
it has a nonzero \(y\) with \(y^\#=0\).
Indeed use either that vector itself or its
nonzero adjoint, since
\((x^\#)^\#=N(x)x\).
The squares in \(J\) span it over any infinite
base field. Over a splitting field this follows
from the squares of the diagonal idempotents
and, for each off-diagonal vector \(v\) incident
with \(e_i\), from
\((e_i+v)^2-e_i^2-v^2=v\).
An annihilating \(k\)-linear functional would make
a polynomial vanish at every rational argument,
and hence after splitting too. Nondegeneracy of
\(T\) shows that \(T(y,a^2)\ne0\) on a nonempty
open subset of possible \(a\). Take \(a\) there
with \(N(a)\ne0\). The identities
\((U_a y)^\#=U_{a^\#}(y^\#)\) and
\(t(U_a y)=T(y,a^2)\) follow from NA13-C04.02
and self-adjointness of \(U_a\). Thus \(U_a y\)
has zero adjoint and nonzero trace. Dividing
by its trace gives a point \(e\) with \(t(e)=1\),
\(e^\#=0\), and \(e^2=e\) by the cubic square
formula. In particular an isotropic \(F_4\)
supplies such an \(e\): its nontrivial split
cocharacter has a nonzero weight on the
faithful Albert representation, and a vector
of nonzero weight has zero invariant cubic norm.

We next show that a second rank-one idempotent
can be found over the fields needed here.
The scheme of \(e\)'s with \(t(e)=1,\ e^\#=0\)
is geometrically smooth and irreducible by
the explicit charts in Appendix C of
\emph{Nonassociative Algebras}. The geometric orbit
of the standard \(e_1\) is open: the incident
short root directions in that calculation give
a surjective differential. Let \(J_0\) be the
image of \(U_{1-e}\), and \(1_0=1-e\).
This operator is an idempotent by the
quadratic Jordan identities. Its rank is the
constant rank ten computed at \(e_1\), since its
image is a vector bundle on the irreducible
scheme. Put \(Q(z)=t(z^\#)\) for \(z\in J_0\).
The identities
\[
 t(1_0)=2,\quad Q(1_0)=1,\quad
 z^\#=Q(z)e,\qquad z^2=t(z)z-Q(z)1_0
\]
can be checked on the standard orbit, where
\(J_0\) is the last two diagonal coordinates and
their common octonion block. Because they are
polynomial identities on the vector bundle,
they extend over the whole scheme. The polar
pairing of \(Q\) is \(t(z)t(w)-T(z,w)\),
by the cubic trace identities in NA13-C04.02.
The self-adjointness of \(U_{1-e}\) makes
\(T|_{J_0}\) nondegenerate and
\(T(1_0,z)=t(z)\), since \(U_{1-e}e=0\).
Thus a polar-radical vector of \(Q\) must be
\(z=t(z)1_0\). Taking traces gives
\(t(z)=2t(z)\), hence \(z=0\) in every
characteristic. Also
\(B_Q(1_0,z)=t(z)\), where \(B_Q\) is the
polar pairing.

If \(Q\) is isotropic, its isotropic vectors
span \(J_0\). Indeed for a hyperbolic pair
\(u,v\) of pairing one any perpendicular vector
\(w\) is in the span of \(u,v\) and the
isotropic vector \(w+u-Q(w)v\).
Since \(t|_{J_0}\) is nonzero, choose
an isotropic vector of trace one by scaling.
Call it \(e_2\); the displayed identities
give \(e_2^2=e_2\), \(e_2^\#=0\). Put
\(e_3=1_0-e_2\), which also has trace one,
square itself, and zero adjoint since
\(Q(1_0-e_2)=1-t(e_2)=0\).

Here is the coordinate comparison for such a
frame. Geometrically the incidence scheme
of pairs \(e,e_2\) just described is smooth
and irreducible. Its fibers over \(e\) are
given by \(t(z)=1,\ Q(z)=0\) in \(J_0\).
They are smooth because proportionality of
the two differentials would make \(z\) a
multiple of \(1_0\), inconsistent with both
equations. They are irreducible because the
quadratic part restricted to \(\ker t\)
has polar rank at least eight, whereas a
product of two linear polynomials has polar
rank at most two. They are nonempty over
an algebraic closure. Smoothness makes the
projection open, so two nonempty open
subsets meet in some irreducible fiber,
proving irreducibility of the total scheme.
The standard frame orbit is open by the
surjective differential onto the \(24\)
frame directions in the same Appendix C.

Use the off-diagonal projectors
\(U_{e_j,e_\ell}\) with remaining index \(i\),
as in NA13-C04.05, and write \(V_i\) for
their images. The rank-eight projections
and the direct-sum decomposition with
the three diagonal lines hold on the
dense standard frame orbit, so the
polynomial projection identities hold
on the full incidence scheme. Their
ranks are constant there. Set
\(q_i(x)=-t(x^\#)\) in block \(V_i\).
On the open orbit \(q_i\) is the
octonion norm and its polar form is
the trace pairing on the block.
The equality holds everywhere, and
these polar forms remain nondegenerate
since the projections are self-adjoint
and \(T\) is nondegenerate.
On \(V_1\times V_2\), let \(m(x,y)=x\times y\).
The split coordinates give its values
in \(V_3\) and
\(q_3(m(x,y))=q_1(x)q_2(y)\).
These, and the expression of \(N\) in
the three diagonal and three block
coordinates, again extend as polynomial
identities from the dense orbit. Its
three diagonal-block terms are
\(-\alpha_i q_i(x_i)\), and its
triple term is \(T(m(x_1,x_2),x_3)\).

Choose \(a\in V_1,b\in V_2\) with
\(A=q_1(a)\ne0\), \(B=q_2(b)\ne0\).
The maps \(R_b=m(-,b)\) and \(L_a=m(a,-)\)
are similarities and therefore invertible.
On \(V_3\) put
\[
 z\cdot w=m(R_b^{-1}z,L_a^{-1}w),
 \qquad n(z)=q_3(z)/(AB).
\]
This is unital with unit \(m(a,b)\)
and has multiplicative nondegenerate
norm. For clarity, the usual
composition identities apply in
characteristic two too. Polarize
\(L_x^{\mathsf t}L_x=n(x)I\) with the unit,
where the transpose is for the polar
form. One gets \(L_x^{\mathsf t}+L_x=t_C(x)I\);
the right multiplication identity is the
same. Hence \(\bar x=t_C(x)1-x\) acts
as either transpose and
\(\bar x x=x\bar x=n(x)1\).
The quadratic identity gives
\(L_{x^2}=L_x^2\) and \(R_{x^2}=R_x^2\),
so the algebra is alternative. The
polarized quadratic identity and
\(t_C(xy)=b_C(y,\bar x)\) give
\(\overline{xy}=\bar y\bar x\).
Thus
this eight-dimensional unital
composition algebra is an octonion
algebra \(C\). Giving the three
blocks coordinates \(R_b x_1,L_a x_2,\bar x_3\)
identifies \(J\) with \(H_3(C,\Gamma)\)
for \(\Gamma=(B,A,1)\). Indeed the
three block norms have the factors
\(A,B,AB\), and their triple term
becomes \(AB\,t_C((c_1c_2)c_3)\)
because \(T|_{V_3}\) is the polar form
of \(q_3\). This is exactly the
weighted cubic norm and preserves
the unit. The trace polarizations
at the unit and the adjoint derivative
identify the Jordan structures,
as in NA13-C04.04.

At a nonarchimedean completion the original
Albert cubic has a nonzero zero even without an
isotropy assumption on the group. We prove the
local cubic bound by the valuation argument of
Springer, as in
\href{https://www.numdam.org/item/SDPP_1965-1966__7_2_A6_0/}
{Pourchet's exposition, Theorem~3}.
Suppose a cubic form \(f\) over a complete local
field with finite residue field has no nonzero
zero. For independent vectors \(x,y\), the cubic
\(f(Tx+y)\) is irreducible: a reducible cubic has
a root over the field. The three roots, with
multiplicities, have the same absolute value
because the absolute value has a unique extension.
Their common absolute value is
\((|f(y)|/|f(x)|)^{1/3}\). Vieta bounds the two
mixed coefficients by
\(|f(x)|^{2/3}|f(y)|^{1/3}\) and
\(|f(x)|^{1/3}|f(y)|^{2/3}\), respectively.
These estimates for dependent vectors follow
by expanding a cube and using \(|3|\leq1\).
It follows that \(|f(x)|^{1/3}\) is an
ultrametric norm, and if \(|f(y)|<|f(x)|\)
then \(|f(x+y)-f(x)|<|f(x)|\).

For a uniformizer \(\varpi\), let
\(M_j=\{x:|f(x)|\leq|\varpi|^j\}\).
The norm is equivalent to a coordinate norm,
by positivity and compactness on the coordinate
unit sphere. Hence the \(M_j\) are full
lattices over the valuation ring, with
\(M_{j+3}=\varpi M_j\).
On each of the three successive residue layers
between \(M_0\) and \(\varpi M_0\), the
reduction of \(\varpi^{-j} f\) is an
anisotropic cubic form. To verify that it is a
form, choose lifts of a basis of the layer.
The single-coordinate and two-coordinate terms
in \(f\) on those lifts are integral after
division by \(\varpi^j\), by the preceding
coefficient bounds. The three-coordinate
polar coefficient is a signed sum of values
of \(f\) on sums of the lifts, so is integral
by the same ultrametric bound. Every nonzero
residue vector lifts outside \(M_{j+1}\),
so its value is nonzero. Each layer has
dimension at most three, by the finite-field
coefficient sum argument below. Since their
dimensions sum to the number of variables,
an anisotropic local cubic has at most nine
variables. The Albert cubic has \(27\).

The resulting \(e\) gives the Peirce zero
space, whose ten-dimensional \(Q\) is isotropic
by the local quadratic arguments
above, in arbitrary characteristic.
The octonion \(C\) is split by the
local doubling and quaternion norm
argument. A frame then makes \(J\)
split by the norm identity and the
parameter removal in NA13-C04.04.
Thus every \(F_4\) over a
nonarchimedean local field is split.

Over a global function field the
same argument starts without assuming
isotropy of the group. In fact over
a finite extension \(k/\mathbb F_q(t)\)
of degree \(s\), any degree-\(d\)
homogeneous form with \(d\ge 1\) in \(n>d^2\)
variables has a nonzero zero. In a
fixed extension basis take each
coordinate to have polynomial
coefficients of degree at most \(M\)
in \(t\). After fixed denominator
clearing, every coefficient in the
expanded form must vanish, giving
at most \(s(dM+O(1))\) homogeneous
equations of degree \(d\) over
\(\mathbb F_q\) in \(ns(M+1)\)
coefficient variables. For large
\(M\) the sum of the degrees is
less than the number of variables.
To count common zeros modulo the
characteristic, sum the product
of \(1-f_i^{q-1}\) over all coefficient
tuples. Each monomial has some
variable exponent less than \(q-1\).
Summing that variable contributes
zero in \(\mathbb F_q\), also for
exponent zero. Thus the number of
common zeros is divisible by the
characteristic. The zero tuple
already counts, so there is a
nonzero tuple. The basis expansion
keeps it nonzero over \(k\).
Apply this first to the
twenty-seven-dimensional Albert
cubic to obtain \(e\), next to the
ten-dimensional \(Q\) to complete
the frame, and then to the
eight-dimensional octonion norm.
The latter norm is isotropic and
the octonion algebra is split by
NA13-C04.08. Thus every \(F_4\)
over a global function field is
split, and its strong approximation
follows from the split proof above.

We next check the real topology needed for an
arbitrary omitted set over a number field.
Every real \(F_4\) has connected real points.
Let \(\omega\) be the compact conjugation on the
complex simple Lie algebra and use the positive
Hermitian form \(-B(x,\omega y)\).
Results J03--J04 of \emph{Lie Algebras and
Finite-Dimensional Representations} construct
this form and conjugate any real structure into
\(\sigma=\theta\omega\), with \(\theta\)
a unitary Lie involution commuting with
\(\omega\). We can identify the complex algebraic
group itself with the complex Lie automorphism
group: its center is trivial, its connected image
is the full connected Lie automorphism group by
the derivation argument above, and the diagram
has no automorphisms. In this identification
write \(K\) for the group of unitary complex
automorphisms. These are precisely the
\(\omega\)-fixed group elements, since for any
complex bracket automorphism
\(g^\dagger=\omega g^{-1}\omega\).
The polar map
\(g\mapsto g(g^\dagger g)^{-1/2}\)
retracts the complex group to \(K\).
The powers used here preserve brackets: the
positive eigenspaces of \(g^\dagger g\) have
bracket eigenvalues given by their products.
They lie in the identity component by their
path of positive powers. Thus \(K\) is compact
and connected. Since \(\theta\) is inner
and commutes with \(\omega\), its implementing
element \(t\) has \(\omega(t)=t\) by triviality
of the center. In particular \(t\in K\).

The compact centralizer \(C_K(t)\) is connected.
Here is a Weyl calculation for this step.
Every element of \(K\) lies in a maximal
compact torus: equip \(K\) with its positive
bi-invariant metric and take a minimizing
geodesic from \(1\), by the Hopf--Rinow
theorem of \emph{Algebraic Topology and
Riemannian Geometry}, P2-02, Theorem~6.10.
The Koszul formula for left-invariant fields
gives \(\nabla_X X=0\) because the metric is
adjoint-invariant, so these geodesics are
one-parameter subgroups. Their compact
abelian connected closures are tori.
Maximal tori in a compact connected group
are conjugate. For the step used here,
choose \(X\) regular in the Lie algebra of
one torus and \(Y\) with dense exponential
in the other. An extremum of
\(g\mapsto \langle\operatorname{Ad}(g)Y,X\rangle\)
forces the two vectors to commute by
differentiation. The centralizer of a
regular toral vector is the torus Lie
algebra, by the root decomposition, so this
conjugates the second torus into the first.
This also applies in the identity component
of the centralizer, which is compact.

Take \(t\in T\), a maximal torus of \(K\).
Each component of \(C_K(t)\) has a
representative normalizing \(T\):
it sends \(T\) to a maximal torus in
the identity component, and the latter
conjugates these tori. The component group
is therefore the stabilizer of \(t\) in
the Weyl group \(W\), modulo the Weyl
group of the identity component. The
compact normalizer has the complex Weyl
group. Indeed the complexification of
\(T\) is maximal, since the Lie
centralizer of its compact torus is its
own Lie algebra and centralizing
commutes with complexification. The
rank-one Weyl representatives
for simple roots lie in the compact
root \(\operatorname{SU}_2\)'s, and
the complex normalizer quotient is \(W\).
Write \(t=\exp(2\pi i x)\) in the
cocharacter coordinates of \(T\).
The exponential lattice is the coroot
lattice \(Q^\vee\), since \(F_4\) is
simply connected, by C07. Thus the
stabilizer in question is the image of
the stabilizer of \(x\) in
\(Q^\vee\rtimes W\).

For the latter stabilizer, consider the
wall arrangement \(\alpha(z)=m\) for
roots \(\alpha\) and integers \(m\).
Its wall reflections are the linear
root reflections followed by
translations by \(m\alpha^\vee\),
so belong to \(Q^\vee\rtimes W\).
Reflections across adjacent walls act
transitively on its open chambers (alcoves),
by following a generic path across walls.
The alcove given by the positive simple-root
inequalities and highest-root inequality
\(\widetilde\alpha(z)<1\) has vertices
\(0\) and \(\omega_i^\vee/m_i\), where
\(\widetilde\alpha=\sum_i m_i\alpha_i\).
These inequalities give a single alcove,
since all positive roots are bounded above
by the highest root there. In \(F_4\)
the \(m_i\) are \(2,3,4,2\), as one
checks on the roots from NA13-C04.09.
For example the highest root is
\(\varepsilon_1+\varepsilon_2\) for simple
roots \(\varepsilon_2-\varepsilon_3\),
\(\varepsilon_3-\varepsilon_4\),
\(\varepsilon_4\), and
\((\varepsilon_1-\varepsilon_2-\varepsilon_3-\varepsilon_4)/2\).
Only the vertex \(0\) belongs to \(Q^\vee\),
since the other vertex has a fractional
simple-root value \(1/m_i\).
An affine Weyl element preserving
this alcove therefore fixes \(0\)
and preserves the ordinary dominant
chamber, hence is the identity.
The action on all alcoves is thus
simply transitive. For alcoves whose
closures contain \(x\), a small ball
about \(x\) permits a gallery between
any two, crossing only walls through
\(x\). It follows that the stabilizer
of \(x\) is generated by the reflections
in these walls. Their ordinary linear
parts are precisely the reflections of
roots with \(\alpha(t)=1\), which
are the roots of \(C_K(t)^\circ\).
This proves the asserted connectedness.

For \(g\) fixed by the real structure
\(\sigma=\theta\omega\), its polar factors
are also fixed by \(\sigma\): it is an
antiunitary operator and commutes with
\(g^\dagger g\) and its positive powers.
The unitary factor belongs to \(C_K(t)\).
Contracting the positive polar factor
to the identity gives a path in the
real group, so its real points are
connected.

At a noncompact real factor the local
compactness test above gives a split
cocharacter. The subgroup generated
by its contracted radical and all
conjugates is dense in the connected
real group: its closure contains a
nontrivial connected unipotent group,
whose Lie algebra and its conjugates
generate a nonzero ideal in the
simple real Lie algebra. The closure
is therefore open and then full.
For globally isotropic \(F_4\), this
local generation at real factors and
splitting at finite factors complete
the global additive approximation
argument above.

For a globally anisotropic \(F_4\)
over a number field, the adjoint
lattice proof gives a compact adelic
quotient. Apply the lattice projection
argument to the product of the
noncompact real factors outside \(S\).
In projective recurrence use a
nontrivial split cocharacter in
each factor. Its conjugates and the
contracting radicals make the
invariant Lie subspace a sum of
simple factor ideals, by the same
argument as in the inner type \(A\)
proof. Every factor is connected,
and those not in the sum would
give a discrete projection to a
nonempty subproduct, contradicting
the noncompactness at \(S\).
Thus the projection is dense in
the whole noncompact real product.
The compact-quotient isolation
argument kills the contracting
radicals, hence the full product.
Omit these isolated factors and
use the lattice and isolation
arguments at each split finite
factor. Only the compact real
factors remain. The rational torus
parameter argument given for \(G_2\)
gives weak approximation in their
connected product. This completes
strong approximation for all
number-field forms of \(F_4\).

\medskip\noindent
\textit{Weak approximation at finitely many factors in characteristic zero.}
Let \(k\) be a number field and \(T\) a finite set of places.
The closure of \(G(k)\) in \(\prod_{v\in T}G(k_v)\) contains
a neighborhood of the identity. We give the argument for all
the semisimple groups considered here. For each \(v\), choose
a basis of \(\operatorname{Lie}G(k_v)\) made up of regular
semisimple elements. Over a splitting field these elements
contain the conjugates of the part of a Cartan subalgebra
off the finitely many root hyperplanes. The differential of
conjugation in the root-space directions, together with the
Cartan directions, is onto at each such element. Thus the
regular semisimple set contains a nonempty Zariski-open
subset, defined over the base field, and its local points
are dense in the Lie algebra. Consequently we can choose
the basis by successive choices outside the preceding linear
span. Approximate these bases simultaneously by
\(X_1,\ldots,X_d\in\operatorname{Lie}G(k)\), using weak
approximation of the scalar coordinates. Regular
semisimplicity and the basis conditions are open locally,
so both hold for the resulting tuple at every place of \(T\).
For \(T\) empty there is nothing to check.

The connected centralizer \(A_i=Z_G(X_i)^\circ\) is a
\(k\)-torus by the absolute root decomposition, and
\(X_i\in\operatorname{Lie}(A_i)(k)\). Let \(L_i/k\) be a finite
Galois splitting field of \(A_i\). Multiplication of the
conjugate entries gives the norm morphism
\[
 \operatorname{Res}_{L_i/k}(A_{i,L_i})\longrightarrow A_i .
\]
After extending scalars it is the multiplication from a
product of copies of \(A_i\), so its differential is onto.
The product, in the chosen order, of these norm morphisms
is a \(k\)-morphism into \(G\). At the tuple of identities its
differential has image containing all the \(X_i\), and is
surjective at every completion in \(T\). The local inverse
function theorem in smooth coordinates implies that its
image over the local product contains an open neighborhood
of the identity. Each source group is a restriction of
scalars of a split torus and consequently its rational
points are dense in the corresponding finite product of
local points: choose coordinates in \(L_i^\times\), apply
additive weak approximation over \(L_i\), and use openness
of the nonzero condition. This proves the asserted openness of
the diagonal closure. In particular it contains the connected
components of the identities at archimedean places, embedded
in the product.

If \(G\) is globally anisotropic, the compact adelic quotient
proved above implies that this closure is cocompact:
projection to the factors in \(T\) gives a continuous
surjection from the compact quotient onto the quotient by
the closure. The latter quotient is discrete by openness,
hence finite. Let any of its nonarchimedean places be
isotropic. The homomorphism from the local group, embedded
in the product, to the permutation group of that finite
coset space kills every element in a parabolic radical.
Indeed in characteristic zero a unipotent matrix has its
\(n\)-th root given by its finite logarithm and exponential
for every positive integer \(n\). The entire polynomial
one-parameter subgroup stays in the radical: it does so on
all integer powers of the original element, and these
parameters are Zariski dense in the affine line. Take \(n\)
to be the order of the finite permutation group. This
puts each local \(G(k_v)^+\) in the diagonal closure.
In particular full local generation and connectedness of
the remaining real groups suffice for the needed weak
approximation, without a separate arithmetic argument at
each finite place.

\medskip\noindent
\textit{Real groups in types \(E_6,E_7,E_8\).}
Let \(G\) be simply connected of one of these absolute types.
Use the compact conjugation \(\omega\) and write its real structure
after conjugation as \(\theta\omega\), where \(\theta\) is a holomorphic
involution commuting with \(\omega\), by Lie algebra Results J03--J04.
These automorphisms lift to the simply connected group: over the
complex field inner automorphisms lift along the central isogeny
and a diagram automorphism acts on the pinning, as in C08 and F06.
Compact conjugation takes the simple positive root vector to
the negative of the corresponding negative root vector and acts
as conjugate inversion on the split torus, which gives the group
action by the same root words. Thus the lift of the given Lie
algebra real structure is the real structure of \(G\).

Apply the polar decomposition of the preceding \(F_4\) argument
to the adjoint group \(G_{\mathrm{ad}}(\mathbb C)\). Its compact
subgroup \(K_{\mathrm{ad}}\) is connected by the polar retraction.
The compact fixed subgroup \(K=G(\mathbb C)^\omega\) maps onto it:
the differential of the covering restricts to an isomorphism
of the compact real Lie algebras, so the image is an open
subgroup of the connected target. The central kernel is
fixed by compact conjugation, since its coordinates in a
split pinned torus are roots of unity and compact conjugation
is conjugate inversion there. It follows that \(K\) is the full
inverse image of \(K_{\mathrm{ad}}\). The positive part of the
adjoint polar path lifts uniquely from the identity along
the finite central covering, which has invertible differential
everywhere. This gives a retraction of the connected
\(G(\mathbb C)\) to \(K\), so \(K\) is connected. The lifted
path remains fixed by \(\theta\omega\) for an element of
\(G(\mathbb R)\), by uniqueness of lifts and the equivariance
of the adjoint polar decomposition. Thus each real point
is joined to a point of \(K^\theta\).

We verify that \(K^\theta\) is connected. When \(\theta\) is
inner, it is implemented by some \(t\in K\). In fact an
implementing element in the adjoint group is fixed by \(\omega\)
because \(\theta\) commutes with \(\omega\) and the adjoint center
is trivial. It therefore has a lift in \(K\). The connectedness
of \(C_K(t)\) follows from the affine Weyl stabilizer calculation
used for \(F_4\) above. The simply connected torus has exponential
lattice the coroot lattice by C07. To check that the corresponding
affine Weyl group has trivial alcove stabilizers, the highest root
coefficients in the Bourbaki numbering are
\[
  \begin{array}{c|l}
  E_6 &(1,2,2,3,2,1)\\
  E_7 &(2,2,3,4,3,2,1)\\
  E_8 &(2,3,4,6,5,4,3,2).
  \end{array}
\]
Here vertices \(1,3,4,\ldots,r\) are on the chain and vertex
\(2\) attaches to \(4\). The coefficients follow by root
reflection in the Dynkin matrix. For instance each listed
positive vector has square length two and all its simple
pairings nonnegative, so it is the dominant root. In these
simply laced diagrams such a vector is a root: successive
simple reflections with positive pairing reduce its height
until it is simple, since a pairing of two can occur only
for a vector equal to that simple root by Cauchy--Schwarz.
All root lengths-squared are normalized to two.
For a coefficient \(m_i>1\), the corresponding alcove vertex
\(\omega_i^\vee/m_i\) has noninteger simple-root value. At
the coefficient-one vertices, the coweights
\(\omega_1^\vee,\omega_6^\vee\) of \(E_6\) have nonintegral
thirds in their simple-coroot coordinates, and
\(\omega_7^\vee\) of \(E_7\) has nonintegral halves.
Indeed the first column of the \(E_6\) inverse Cartan matrix is
\((4/3,1,5/3,2,4/3,2/3)\), the sixth its diagram reflection,
and the seventh column for \(E_7\) is
\((1,3/2,2,3,5/2,2,3/2)\).
Hence \(0\) is the only vertex of the fundamental alcove in the
coroot lattice. The wall-gallery argument from \(F_4\) now
identifies every point stabilizer with the wall reflections
through that point, and proves that \(C_K(t)\) is connected.

The only outer involution to check is in type \(E_6\).
Let \(S\) be a maximal torus of \(K^{\theta,\circ}\).
In the Lie algebra of the centralizer of \(S\) in \(K\), the
\(\theta\)-fixed Lie subalgebra is exactly \(\operatorname{Lie}S\),
because \(S\) is maximal in the compact connected fixed group.
The derived Lie algebra of the centralizer has no fixed vector:
the compact Lie algebra is the direct sum of its center and
its semisimple derived algebra, and \(\operatorname{Lie}S\)
lies in the center. An involution with zero fixed space acts
by \(-1\), which forces that derived Lie algebra to be abelian
by preservation of brackets. Thus the centralizer has toral
Lie algebra and its connected subgroup is a maximal torus \(T\)
of \(K\). No root for \(T\) vanishes identically on \(S\), as
its root space would centralize \(S\). This gives a regular
vector in \(\operatorname{Lie}S\); its chamber is \(\theta\)-stable.
In this pinning \(\theta=\operatorname{Int}(t)\delta\), with
\(t\in T\) and \(\delta\) the pinned diagram reflection.
The compact torus \(T^\delta\) is connected because \(\delta\)
permutes the simply connected simple-coroot basis. Hence
\(S=T^\delta\). Conjugating \(\theta\) by a torus element changes
\(t\) by \((1-\delta)T\); we may take \(t\in S\) since the two
Lie spaces of eigenvalues \(1,-1\) span \(\operatorname{Lie}T\).

Write \(\mathfrak a\) for the real cocharacter space of \(S\);
its normalization in this calculation is such that a root takes
value \(\exp(2\pi i\alpha(x))\) on \(\exp(2\pi i x)\).
Let \(W\) be the absolute Weyl group. For each orbit \(I\) of
simple roots let \(w_I=\prod_{i\in I}s_i\), with commuting
factors: the nontrivial simple pairs are the orthogonal
pairs \(\{1,6\}\) and \(\{3,5\}\). Their restrictions to
\(\mathfrak a\) are reflections in the projected simple roots
\(\beta_I\). If an absolute positive root is not simple,
some orbit \(I\) has positive scalar product with the sum
of its simple roots, since projection to \(\mathfrak a\)
takes positive roots into the positive cone of these
\(\beta_I\). The corresponding \(w_I\) decreases the absolute
height and preserves positivity, as its only positive inversions
are the simple roots in \(I\). Thus every absolute root is
conjugate by the \(w_I\)'s to a simple root. Its projection
is therefore in the folded root system. The projected simple
roots, in the order \(\{1,6\},\{3,5\},\{4\},\{2\}\), have
squared lengths \(1,1,2,2\), and successive inner products
\(-1/2,-1,-1\), with all other off-diagonal entries zero.
This is the reduced root system of type \(F_4\).
The \(w_I\)'s generate \(W^\delta\): their induced reflections
act transitively on the chambers in \(\mathfrak a\); an
element of \(W^\delta\) stabilizing the fundamental one takes
its regular vectors to the absolute dominant chamber, and
so is the identity in the absolute Weyl group.
Each \(w_I\) has a representative fixed by \(\delta\), by
the pinned compact rank-one Weyl representatives and by
taking commuting products for the pairs. Consequently a
\(\delta\)-fixed absolute root has fixed root space for
\(\delta\), since it is conjugate by these representatives
to a fixed simple root. All other absolute roots occur
in orthogonal pairs; the projected lengths in the two
cases are two and one, respectively. These are exactly the fibres
of the projection: move a projected root to a projected simple root
by the folded Weyl group. Any positive absolute root over that
simple root is supported on its simple orbit. The two-member
orbit is orthogonal, so the only roots there are the simple
ones. This also shows that the long restricted roots have just
one absolute root above them.

Every connected component of \(K^\theta\) has a representative
\(n\) normalizing \(S\): conjugate its image of \(S\) back by
the compact connected identity component. Such \(n\) normalizes
\(T=C_K(S)^\circ\). Its Weyl element \(w\) belongs to \(W^\delta\).
Writing \(n=y n_w\) with a \(\delta\)-fixed representative \(n_w\)
and \(y\in T\), the fixed-point equation is
\(y\delta(y)^{-1}=t w(t)^{-1}\). Thus the possible Weyl elements
are exactly those stabilizing the class of \(t\) in
\(T/(1-\delta)T\), and the kernel of this component calculation is
\(S\), which is connected. In additive coordinates the quotient
torus is \(\mathfrak a/L\), with \(L\) the orthogonal projection
of the absolute coroot lattice. The projections of its simple
coroot generators are precisely the \(\beta_I\), so \(L\) is
the coroot lattice of the dual folded root system
\(\{\beta^\vee\}\) of type \(F_4\).

Write the class of \(t\) as \(x\bmod L\). The Weyl stabilizer
in question is the image of the affine stabilizer of \(x\)
in \(L\rtimes W^\delta\). It is generated by images of the
affine reflections whose walls contain \(x\), by the preceding
\(F_4\) alcove argument applied to the dual root system.
At a long restricted root \(\beta\), the wall condition is
\(\beta(x)\in\mathbb Z\); its unique absolute root is fixed by
\(\delta\), so the corresponding root space is fixed by
\(\theta\). At a short restricted root \(\beta\), the condition
is \(2\beta(x)\in\mathbb Z\); on the two corresponding absolute
root spaces \(\theta\) acts as the swap times the scalar
\(\exp(2\pi i\beta(x))\), so there is a fixed line exactly
under this condition. These are precisely the root spaces
of the complexification of \(K^{\theta,\circ}\) relative
to \(S\); the compact root \(\operatorname{SU}_2\)'s realize
the generating reflections in the connected fixed group.
It follows that every normalizing representative can be
brought into \(S\) by a representative in the connected fixed
group. This proves \(K^\theta\), and thus \(G(\mathbb R)\),
connected in the outer case also.

A noncompact real factor now satisfies
\(\overline{G(\mathbb R)^+}=G(\mathbb R)\). Its nontrivial
split cocharacter from the local compactness test has a
nonzero contracting radical. The closure of the group
generated by the rational parabolic radicals is normal,
and its Lie algebra contains the Lie algebra of that
radical. It is therefore open by real Lie simplicity, and
the connectedness just proved gives the assertion. The
complex factors are split. Hence when all finite
completions of the exceptional group are quasi-split,
strong approximation for a globally isotropic form
follows from the fixed-opposite-radicals argument above.

For a globally anisotropic form over a number field,
assume again that all finite factors are quasi-split.
The compactness of the adelic quotient was proved for all
types above. At the noncompact archimedean factors the lattice
projection proof used for \(F_4\) applies using the just-established
connectedness, real Lie simplicity, and a contracting radical.
That proof isolates their product. Each finite factor is
then isolated by the quasi-split generation and projection
proof. Finite supports are dense in their restricted product.
For the remaining compact real factors the torus
parameter argument given for \(G_2\) uses only the
characteristic-zero Cartan and root decomposition.
It gives dense weak approximation because those real
groups are connected. This finishes the asserted
number-field subcase.

\medskip\noindent
\textit{Local generation in nontrialitarian type \(D\) in
characteristic different from two.}
Here the diagram action fixes one of the vector
representations: for \(D_4\) this means its image
has order at most two, after choosing the fixed
outer vertex as the vector vertex. The split
spin group has the double cover of the orthogonal
group used above. Its inner automorphisms act on
the matrix algebra with orthogonal adjoint
involution through conjugations by orthogonal
similitudes. The diagram involution fixing the
vector vertex also acts there, by conjugation by
an orthogonal reflection. The latter interchanges
the two spin vertices, as is seen by changing
the sign of the last hyperbolic coordinate on
the split torus. Thus the descent cocycle acts on
this algebra with involution. Descent gives a
central simple algebra \(A\) of degree \(2n\)
with an orthogonal involution \(\ast\), and the
group is the simply connected double cover of
\(\{a:a^\ast a=1,\ \operatorname{Nrd}(a)=1\}\).
This description and the central scalar kernel
are checked over the splitting field, where they
are the ordinary vector spin covering.

At a nonarchimedean completion the division
factor of \(A\) is the field or a quaternion
division algebra. Indeed the center-preserving
antiautomorphism of the division corner is
constructed as in type \(C\), and the residue
uniformizer argument there forces degree at
most two. At a real completion the same two
possibilities apply. The field case is the
quadratic spin case above. In the quaternion
case write \(A=\operatorname{End}_B(V)\) with
canonical quaternion conjugation on \(B\).
The Skolem--Noether calculation as in type \(C\)
makes \(\ast\) adjoint to a skew-Hermitian
form \(h\) on \(V\). The sign is the opposite
of the symplectic case: after splitting \(B\)
the matrix \(JH\) in that calculation is
symmetric, so \(\bar H^{\mathsf t}=-H\).
If the spin group is isotropic, a nonzero
weight of a split cocharacter on \(V\) gives
an isotropic vector for \(h\). Complete it to
a hyperbolic pair \(e,f\) with
\(h(e,f)=1,\ h(f,e)=-1\), choosing a partner
of pairing one and correcting its norm by
a multiple of \(e\), using \(2\ne0\).
Let \(W\) be the perpendicular complement.

We compute in the Levi of this pair and the
two opposite parabolic radicals. Let \(w\in W\)
have \(q=h(w,w)\ne0\), so \(\bar q=-q\).
In the ordered coordinates \(e,w,f\), fixing
the remaining perpendicular space, the upper
and lower radical elements include
\[
 U(a,b)=
 \begin{pmatrix}
 1&a&b\\0&1&-q^{-1}\bar a\\0&0&1
 \end{pmatrix},\qquad
 L(a,b)=
 \begin{pmatrix}
 1&0&0\\-q^{-1}\bar a&1&0\\-b&-a&1
 \end{pmatrix},
 \qquad b-\bar b=-a q^{-1}\bar a.
\]
In fact pairing the middle and last columns
gives \(\bar a+q(-q^{-1}\bar a)=0\),
and pairing the last column with itself gives
the displayed condition. The lower formula
is the conjugate by the skew interchange of
the two hyperbolic coordinates. These unipotent
groups lift isomorphically to the spin radicals
under the central cover, as follows in the
split root coordinates and descends.

For \(b\ne0\), direct multiplication gives
\[
 L(-\bar b^{-1}a,\bar b^{-1})\,U(a,b)\,
 L(b^{-1}a,\bar b^{-1})=
 \begin{pmatrix}
 0&0&b\\
 0&1+q^{-1}\bar a b^{-1}a&0\\
 -\bar b^{-1}&0&0
 \end{pmatrix}.
\]
For the condition on either lower factor,
substitute \(a q^{-1}\bar a=\bar b-b\);
its right side becomes \(\bar b^{-1}-b^{-1}\).
In the product the right factor cancels the
first two entries of the top row. The left
factor cancels the last entries of the middle
and bottom rows, leaving the displayed middle
entry. Orthogonality and the pairing of the
first and last columns then give the other
entries, including the minus sign in the
lower-left corner.

Put \(E=F(q)\), a separable quadratic subfield,
where \(F\) is the local base field. For
\(r\in E^\times\) of norm one with \(r\ne1\),
take \(a=1\) and \(b_0=(r-1)^{-1}q^{-1}\).
Then \(\bar b_0=b_0r\), hence
\(b_0-\bar b_0=-q^{-1}\), and the middle
entry in the word is \(r\). For \(r=1\)
use \(a=0,b\in F^\times\).
Taking a ratio with the latter interchange
realizes the rotation \(w\mapsto wr\)
in the Levi, with some scale on the
hyperbolic pair.

These rotations on individual lines generate
the isometry group of \(W\).
If \(u,v\) have the same nonzero norm,
one of \(u-v,u+v\) has nonzero norm because
the two norms add to \(4h(u,u)\).
For \(z=u-v\) of nonzero norm write
\(c=h(z,u)\), so \(h(z,z)=c-\bar c\).
The rotation on \(z\) with right multiplier
\(c^{-1}\bar c\) sends \(u\) to \(v\):
its component on \(z\) is
\(z\,h(z,z)^{-1}c\), whose change is \(-z\).
The multiplier has norm one and commutes with
\(h(z,z)\) because \(c\) and \(\bar c\) commute.
For \(z=u+v\) the same operation sends \(u\)
to \(-v\), and a rotation by \(-1\) on
the line of \(v\) finishes. Iterate on
perpendicular spaces. A nonzero nondegenerate
skew-Hermitian space over \(B\) has a vector
of nonzero norm: otherwise polarization
makes every value \(h(x,y)\) fixed by
quaternion conjugation, hence scalar, and
right scaling by nonscalars forces all
values to be zero. Thus this argument
also supplies the diagonalization and
the choice of \(w\) above. The complement
\(W\) is nonzero in the absolute
rank range \(D_n\), \(n\geq4\).

We check the remaining hyperbolic scales,
including their spin lifts. In the word
with fixed \(r\ne1\) replace \(a=1\)
by arbitrary \(a\in B^\times\) and
\(b_0\) by \(a b_0\bar a\).
The admissibility equation still holds
and the middle entry is the same \(r\).
The ratio of the two words therefore
acts trivially on \(W\) and scales \(e\)
by \(a b_0\bar a b_0^{-1}\).
The scalar scales \(c\in F^\times\)
are generated by the elementary rank-one
words with \(a=0\). Correcting by the
scalar \(\operatorname{Nrd}(a)^{-1}\)
leaves the scale \([a,b_0]\).
As \(r\) varies, \(b_0\) runs through
the elements of \(E\) with
\(b_0-\bar b_0=-q^{-1}\).
This affine line represents every
nonscalar element of \(E^\times\)
modulo a central scalar, since its
pure component is any fixed nonzero
multiple of \(q\). Thus we obtain
\([B^\times,E^\times]\) on these
hyperbolic scales.

The subgroup generated by these
commutators is normal in \(B^\times\):
for fixed \(b\in E^\times\),
\(c[a,b]c^{-1}=[ca,b][c,b]^{-1}\).
Its closure is the norm-one group.
If \(F=\mathbb R\), every norm-one
quaternion is conjugate to an element
of norm one in \(E=\mathbb C\),
and there it is a quotient \(b/\bar b\),
a commutator with a quaternion inducing
complex conjugation. At a finite place,
if \(E/F\) is unramified, use \(E\)
for the unramified quadratic subfield
in the valuation-layer proof above.
All commutators in that proof have
one entry in \(E^\times\), taking
the residue lifts in \(E\).
If \(E/F\) is ramified, take a
uniformizer \(\pi_E\) of value \(1/2\)
in \(B\). Its conjugation is the
nontrivial automorphism of the
quadratic residue field \(\ell\).
The residue norm-one elements thus
come from commutators with \(\pi_E\).
On an even positive layer \(2m\)
the same commutators with
\(1+\varpi^m x\) give the whole
trace-zero space, since conjugation
on its leading residue is that
automorphism. On an odd layer
\(2m+1\), use
\(1+c\varpi^m\pi_E\in E^\times\)
against units of the unramified
quadratic subfield. The leading
coefficients, apart from one fixed
nonzero factor, are
\(\bar c(\eta-1)\) with
\(\eta\in\ell^\times\) of norm one,
by the nondivisible-layer
commutator calculation above.
They span \(\ell\) over the base
residue field: any affine line
meets the unit norm conic in at
most two points, because the norm
of its nonzero direction is
nonzero, while the conic has at
least three points. Products add
these leading coefficients.
Successive correction of the
layers proves density in the
norm-one subgroup in this
ramified case as well.

The spin preimage of the
hyperbolic scale subgroup is
\[
 \{(t,s)\in B^\times\times F^\times:
              \operatorname{Nrd}(t)=s^2\}.
\]
To check the cover, split \(B\).
The scale subgroup then acts on
two isotropic coordinates by
\(\operatorname{GL}_2\), and on
the dual coordinates inversely.
In the Clifford exterior
representation its lift is
\(\bigwedge t\) times a scalar
\(s^{-1}\). The Clifford norm
condition is \(\det(t)=s^2\):
for diagonal \(t\) it is the
product of the two hyperbolic
Clifford norms computed above,
and elementary determinant-one
changes have norm one. The
character \(s\) descends, since
its square is the reduced norm
and the covering subgroup is
geometrically connected. More
explicitly over the splitting
field it is the variety
\(\operatorname{SL}_2\times\mathbb G_m\),
using \(t\operatorname{diag}(s^{-2},1)\)
and \(s\); two square-root
characters with value one at
the identity agree.

The scalar rank-one words lift
\(c\) to \((c,c)\), because
the possible signs in \(s=\pm c\)
are constant on the algebraic
parameter space and the ratio
with the word at \(c=1\)
starts at the identity. The
commutator construction with
fixed \(r\) and varying
\(a\in B^\times\) likewise
lifts the norm-one scale to
\(([a,b_0],1)\). All words,
including the scalar correction,
are morphisms in \(a\) on
the geometrically connected
unit group, and the second
coordinate squares to one
and equals one at \(a=1\).
Taking the closure supplies
all pairs \((b,1)\) with
\(\operatorname{Nrd}(b)=1\).
Together with the scalar
pairs these supply every
\((t,s)\) in the covering
subgroup, by setting \(b=t/s\);
they include its scalar
central kernel.

Now take any spin Levi point
for the chosen hyperbolic
pair. On the isometry group
it acts on the pair by
\(t,\bar t^{-1}\) and on
\(W\) by an isometry.
The rotation words remove
the entire \(W\)-action.
The remaining scale lift
is in the radical-generated
closure by the preceding
argument. The relative
Bruhat argument then
finishes local density:
the centralizer of a
maximal split torus
containing the pair scale
lies in this Levi, and
the relative root rays
supply the unipotent and
Weyl factors. By global
additive approximation
this also proves sufficiency
in globally isotropic
nontrialitarian \(D\) over
global fields of characteristic
different from two.

The local isotropy test for these same
nontrialitarian forms is elementary. Over a
nonarchimedean field the quaternion division
algebra is unique by the cyclic valuation
description above. The unramified quadratic
subfield is unique up to isomorphism, and
matching the uniformizer parameters uses
surjectivity of the norm on its units.
Over \(\mathbb R\) uniqueness is the quaternion
case of C04.12 of the associative algebra
notes. For a pure \(q\ne0\) in this division
algebra write \(E=F(q)\) and present
\(B=(E/F,c)\) with a conjugating generator
\(j\) of square \(c\).
For \(d\in F^\times\) outside the norm group
of \(E\), the quaternion algebra \((E/F,d)\)
is division, by its norm
\(N_E(x)-dN_E(y)\).
An isomorphism to \(B\) can be adjusted
by Skolem--Noether to fix \(E\).
The conjugating generator must go to
\(jz\) with \(z\in E^\times\), since
the centralizer of \(E\) is \(E\).
Hence \(d=cN_E(z)\). Now
\(\bar a q a\) takes the values \(qN_E(a)\)
for \(a\in E^\times\) and \(cqN_E(z)\)
for \(a=jz\). It therefore represents
every scalar multiple of \(q\). This
is the pure-line argument for the
local bound, as in
\href{https://www.math.uni-bielefeld.de/LAG/man/188.pdf}
{Becher--Mahmoudi, Section~6}.
Diagonalize a skew-Hermitian form as
above. Among four pure diagonal
entries there is a nontrivial linear
dependence over \(F\), since the pure
space has dimension three. Realize
each coefficient times its entry
as \(\bar a q a\), taking \(a=0\)
when the coefficient is zero.
The resulting nonzero vector has
norm zero. Thus every quaternionic
factor in this absolute rank range
is isotropic, at a finite place or
over \(\mathbb R\).

Over a number field, a globally
anisotropic nontrialitarian group
has compact adelic quotient by
the adjoint lattice proof.
Every finite completion is
isotropic, by the quadratic or
quaternionic local test.
At a noncompact real place
the local radicals generate densely,
so its group is connected, being
the closure of a subgroup generated
by connected unipotent groups.
At a compact real place only a
definite quadratic model can occur,
and its spin group is connected
by the plane rotations above.
The projective recurrence,
Lie ideal and isolation arguments
of the \(F_4\) case therefore
isolate the product of noncompact
outside real factors.
After omitting those factors,
at each finite place the
arithmetic lattice projection is
open and of finite index by
the characteristic-zero Lie argument,
which uses only the presence of
the local split cocharacter.
The radicals lie in that
projection and generate densely,
so it is the full local
group. The same contraction
isolates the finite factors
from the compact remainder.
Weak approximation at the
connected compact real factors
again follows from the rational
torus parameters. This proves
the anisotropic number-field case.

We also give the compactness
argument for a globally anisotropic
nontrialitarian group over a
function field of odd characteristic.
Its defining orthogonal involution
\(\ast\) on \(A\) has no primitive
idempotent \(e\) with \(e^\ast e=0\):
otherwise \(p=e-ee^\ast/2\) and
\(p^\ast\) are orthogonal idempotents
as in type \(C\), and the inverse
scales give a split torus. These
scales have reduced norm one, since
the two idempotents have equal
rank after splitting. The torus
also lifts up to isogeny to the
spin group. Thus \(e^\ast e\) is
nonzero, and its rank-one image
cannot give a nilpotent,
since then a projector \(e_0\)
onto that line would satisfy
\(e_0^\ast e_0=0\), by the
same calculation as in type \(C\).
The self-adjoint primitive
projector obtained from \(e^\ast e\)
gives a global division corner
with involution. Skolem--Noether
then gives an \(\epsilon\)-Hermitian
form on the underlying module,
as in type \(C\). It is anisotropic
on vectors, since an isotropic
line would again give
\(e^\ast e=0\).

Apply the function-field lattice
argument to its full isometry
group, using the form coordinates
on vectors and the right
division-algebra multiplications
as tensors. The determinant
absolute value is one, since
the determinant on the module
is a power of the reduced norm,
which is \(\pm1\) on isometries.
Some nonzero component of the
form on each rational vector
gives the invariant polynomial
bound. Thus the vector isometry
group has compact adelic quotient.
Its identity component has
compact quotient by correcting
the finitely many rational
components as in the quadratic
case. The finite \(\mu_2\)
cover is handled by the integral
torsor and finite square-class
argument already given there.
Every finite completion is
isotropic by the local test.
For the quaternionic ones
the positive-characteristic
lattice projection argument
applies by the scalar
\(\operatorname{SL}_2\) lift discussed
above. The other factors use
the quadratic spin argument.
All outside factors can
therefore be isolated, proving
strong approximation for the
nontrialitarian forms over
function fields of odd
characteristic as well.

\medskip\noindent
\textit{Nontrialitarian type \(D\) in characteristic two.}
We use the quadratic structure together with its polar structure;
compare the quadratic bounds in \cite[Corollary 6.8]{dolphinquadratic}. For a split
even-dimensional regular quadratic space with form \(q_0\) and alternating
polar form \(h_0\), let \(\sigma_0\) be the adjoint involution and choose
\(\ell_0\) such that \(h_0(v,\ell_0 v)=q_0(v)\). Such a matrix is obtained
by writing a quadratic polynomial as a triangular matrix; the
polar identity gives \(\ell_0+\sigma_0(\ell_0)=1\). On the
self-adjoint matrices put \(f_0(s)=\operatorname{Tr}(\ell_0s)\).
It is independent of the triangular presentation: changing the
presentation changes \(\ell_0\) by \(a+\sigma_0(a)\), which pairs
to zero with every self-adjoint matrix in characteristic two.
The pair \((\sigma_0,f_0)\) is therefore preserved by conjugations
by quadratic similitudes. Conversely, if conjugation by \(g\)
preserves the adjoint involution, then \(\sigma_0(g)g=c\) is
scalar, so \(g\) is a polar similitude of multiplier \(c\).
If it also preserves \(f_0\), the difference
\(g^{-1}\ell_0g-\ell_0\) is orthogonal to the self-adjoint space
under the trace pairing. This orthogonal space consists of
the \(a+\sigma_0(a)\): they are orthogonal by the trace
identity, and rank--nullity gives equality of dimensions.
Such a difference contributes zero to the quadratic value.
It follows that \(q_0(gv)=c q_0(v)\).

In the split spin model the inner automorphisms act by projective
quadratic similitudes. The diagram involution fixing the vector
vertex also acts this way, by swapping the two basis vectors of
one hyperbolic plane: on the diagonal torus it changes the sign of
that one coordinate and interchanges the two spin vertices.
For \(D_4\) with nontransitive diagram image choose a fixed
outer vertex as the vector vertex. Thus the descent in F06 gives
a central simple algebra \(A\) of degree \(2n\), an involution
\(\sigma\) whose split adjoint form is alternating, and a linear
functional \(f\) on \(\operatorname{Sym}(A,\sigma)\) satisfying
\[
 f(a+\sigma(a))=\operatorname{Trd}_A(a).
\]
The group is the spin cover of the identity component of the
isometries of this pair, as is checked after splitting, where
this is the ordinary vector spin cover with kernel \(\mu_2\).
In characteristic two this is a group-scheme kernel. The
root and even-sum torus lattice calculation for the split spin
model above applies to the roots \(\pm\varepsilon_i\pm\varepsilon_j\);
thus the cover is the simply connected form by C07. Conjugation
by a similitude of multiplier \(c\) acts on an even product of
\(2j\) Clifford vectors by applying the similitude to each vector
and multiplying by \(c^{-j}\). The Clifford relations show that
this action respects the even Clifford multiplication and its
reversion, making the cover equivariant for the descent.

Over a global function field \(k\) of characteristic two,
\([k:k^2]=2\), as in the additive approximation argument.
The same is true of \(F=\mathbb F_{2^a}((t))\), by separating even
and odd exponents. Over either field \(L\), the division factor
of an algebra with a first-kind involution has degree at most
two. Indeed the involution identifies the algebra with its
opposite, so its scalar Brauer cocycle has class of order
dividing two. For a finite separable Galois splitting field
\(E/L\), the field \(E^{1/2}=E L^{1/2}\) gives the same
Galois group over \(L^{1/2}\). Under the isomorphism given
by squaring, scalar base change of the cocycle becomes
its square, hence is a coboundary. Thus \(L^{1/2}\)
splits the algebra. The division degree divides two: a
split module of dimension \(d\) over \(L^{1/2}\) for the
division algebra of degree \(d\), viewed over \(L\), has
dimension \(2d\) divisible by \(d^2\).

We need the following elementary quadratic bound in these fields.
Over \(k\), any homogeneous quadratic polynomial in \(m>4\)
variables has a nonzero zero. Take a separating parameter
\(t\) and a basis of \(k/\mathbb F_q(t)\), of dimension \(d\).
Expand each unknown in this basis with polynomial coordinates
of degree at most \(N\). After clearing the fixed denominators
in the quadratic coefficients and in the basis products, each of
the \(d\) equations has polynomial degree at most \(2N+C\) in
\(t\), for a constant \(C\) independent of \(N\). Equating
coefficients gives at most \(d(2N+C+1)\) quadratic equations
in \(md(N+1)\) scalar unknowns over \(\mathbb F_q\).
For large \(N\) their total degrees are less than the number
of unknowns. Their number of common zeros is divisible by the
characteristic: sum the product of \(1-P^{q-1}\), over all
unknowns and all defining polynomials \(P\). Each monomial
has some variable of exponent less than \(q-1\), for which
the sum over \(\mathbb F_q\) is zero (also at exponent zero).
Since the origin is a solution, another exists. The chosen
coordinates give a nonzero vector over \(k\).
For \(F=\mathbb F_q((t))\) approximate the coefficients by
ones in the dense field \(\mathbb F_q(t)\). Normalize
each resulting nonzero solution to maximum coordinate norm one.
These vectors form a closed subset of the compact product of
unit balls; a convergent subsequence gives a nonzero solution
over \(F\). In either field the norm of a quaternion division
algebra \(B\) represents every nonzero scalar: apply the bound to
\(\operatorname{Nrd}_B(x)-a z^2\) in five variables. A nonzero
zero has \(z\ne0\) because \(B\) is division.

When the division factor is the field, \(f\) in the adjoint
description recovers a regular scalar quadratic form and the
bound just proved makes it isotropic in absolute type \(D_n\),
\(n\ge4\); the scalar Clifford argument above applies. Suppose
then that the division factor is a quaternion algebra \(B\)
with conjugation denoted by a bar. As in the type \(C\)
matrix calculation above, \(\sigma\) on \(\operatorname{End}_B(V)\)
is adjoint to a nondegenerate Hermitian form \(h\), alternating
in the quaternion sense that \(h(v,v)\) is a central scalar.
By the nondegenerate matrix reduced trace pairing,
\(f(s)=\operatorname{Trd}_A(\ell s)\) for some \(\ell\in A\).
Its defining identity gives \(\ell+\sigma(\ell)=1\).
Use the generalized quadratic form
\[
 Q(v)=[h(v,\ell v)]\in B/L,
\]
where brackets are classes modulo central scalars. It satisfies
\[
 Q(va)=\bar a Q(v)a,\qquad
 Q(v+w)=Q(v)+Q(w)+[h(v,w)],\qquad
 \operatorname{Trd}_B Q(v)=h(v,v).
\]
Here the trace on \(B/L\) is well-defined, and
\(\bar a L a\subseteq L\) by the norm equation.
These identities follow by applying
\(\ell+\sigma(\ell)=1\), since any quaternion differs from its
conjugate by a central scalar. A change of representing \(\ell\)
by an alternating element \(a+\sigma(a)\) does not change \(Q\).
Conversely, for a self-adjoint change \(\Delta\), the condition
\(h(v,\Delta v)\in L\) for all \(v\) makes the diagonal entries
of the Hermitian matrix of \(h(-,\Delta-)\) central.
Such matrices are exactly \(Y+\bar Y^{\mathsf t}\): off-diagonal
entries are assigned once, and diagonal entries lift along
the surjective quaternion trace. Undoing the matrix of \(h\)
gives \(\Delta=a+\sigma(a)\). Thus among \(h\)-isometries,
preserving \(Q\) is equivalent to preserving the semitrace \(f\)
by conjugation.

The alternating quaternionic Hermitian form \(h\) diagonalizes
by taking complements of nonzero-norm vectors. Such a vector
always exists in a nonzero nondegenerate space: if all norms
were zero, polarization and right scaling would give
\(\operatorname{Trd}_B(h(v,w)a)=0\) for every \(v,w,a\), forcing
\(h=0\). Norm surjectivity makes every diagonal two-dimensional
nondegenerate summand isotropic. An isotropic vector can be
completed to a hyperbolic pair by correcting a partner of
pairing one by a multiple of the vector; the needed scalar is
obtained by the quaternion trace. Hence \(h\), in quaternion
dimension \(n\ge4\), has a totally isotropic subspace of dimension
two. On it \(Q\) has zero polar form, and its values have trace
zero. If a nonzero vector already has value zero we are done.
Otherwise take a basis and representatives \(q,r\) of its two
values. Both are trace-zero noncentral quaternions. Then \(q^2\)
is a nonsquare scalar, since \(q^2=c^2\) would make \(q+c\)
nilpotent. Write \(r^2=c^2+q^2 d^2\) with \(d\ne0\), using
\([L:L^2]=2\). By Skolem--Noether there is \(a\in B^\times\)
with \(a^{-1}qa=(r+c)/d\). Multiplying \(a\) on the left by
an element of \(L(q)^\times\) allows us to prescribe its norm:
the norm on that purely inseparable field is the square map
with image \(L^\times\). Prescribe \(\operatorname{Nrd}(a)=d\).
Then \(\bar a q a=r+c\), so the two basis values cancel on
a nonzero vector. This proves isotropy of \(Q\) over both fields.

Complete a \(Q\)-isotropic vector \(e\) to a \(Q\)-hyperbolic
pair. Take a partner \(f\) with \(h(e,f)=1\); replacing it by
\(f+ea\) adds \([a]\) to its \(Q\)-value, so choose \(a\) to
make that value zero. In \(eB\oplus W\oplus fB\), the formula is
\(Q(x,y,z)=[\bar xz]+Q_W(y)\). Acting by inverse scalar scales on
\(e,f\) gives a split cocharacter up to isogeny of the spin cover,
so \(Q\)-isotropy gives group isotropy as used here.

We now prove local generation for the quaternionic model over \(F\).
For \(w\in W\) and a representative \(b\in B\) of \(Q_W(w)\), the
two opposite radical elements on column coordinates are
\[
\begin{split}
 U(w,b)(x,y,z)&=(x+h(w,y)+bz,\ y+wz,\ z),\\
 L(w,b)(x,y,z)&=(x,\ y+wx,\ z+h(w,y)+bx).
\end{split}
\]
The identities for \(Q\) show that these preserve both \(Q\)
and \(h\): for \(U\) the mixed change in the first coordinate
contributes \([\overline{h(w,y)}z]=[h(y,w)z]\), and the remaining
term is \([\bar z(\bar b+b)z]=0\). Each family contracts along
one direction of the hyperbolic scale cocharacter. Thus it lifts
to the corresponding spin radical, since the split central cover
is an isomorphism on each unipotent radical, which descends.

For \(b\ne0\), multiply the following three radical elements:
\[
 L(wb^{-1},\bar b^{-1})\ U(w,b)\
 L(w\bar b^{-1},\bar b^{-1}).
\]
Both outer parameters are admissible, since
\(Q_W(wa)=[\bar a b a]\). Acting on \(e,f\) in order gives
\(e\mapsto f\bar b^{-1}\) and \(f\mapsto eb\).
On \(W\) the formula is
\[
 R(w,b)y=y+w b^{-1}h(w,y).
\]
Indeed the right lower factor adds
\(b^{-1}h(w,y)\) to the last coefficient. The upper factor
then cancels the first coefficient and gives the displayed
middle one. The left lower factor cancels the last coefficient
by \(h(w,w)=b+\bar b\).
Dividing this word by the one with \(w=0,b=1\) gives a Levi
element with \(W\)-action \(R(w,b)\).

These transformations generate all \(Q_W\)-isometries.
For \(u,v\) with \(Q_W(u)=Q_W(v)\), suppose \(u,v\ne0\).
If \(b=h(u+v,u)\ne0\), it represents \(Q_W(u+v)\), because
\(h(u,u)\) is central and \(h(v,u)\) is conjugate to \(h(u,v)\).
Thus \(R(u+v,b)u=v\). If \(b=0\), choose \(w\) outside the
two kernels \(h(w,u)=0\) and \(h(v,w)=0\) and take any nonzero
representative \(b'\) of \(Q_W(w)\). For \(u'=R(w,b')u\) we have
\[
 h(u'+u,u)\ne0,\qquad h(v+u',u')\ne0.
\]
The first follows since \(u'+u=w\,b'^{-1}h(w,u)\).
For the second, \(h(u',u')=h(u,u)=h(v,u)\) by the original
zero equation and isometry. The remaining term is
\(h(v,w)b'^{-1}h(w,u)\ne0\).
Apply the first case twice. To decompose an isometry, align
one nonzero \(h\)-norm vector with its image and then induct
on its perpendicular complement. Thus the \(W\)-action of every
spin Levi point can be removed by radical words with some
accompanying hyperbolic scales.

For the remaining hyperbolic scales take a fixed \(w\) of
nonzero \(h\)-norm, which exists since \(W\) is nonzero and
nondegenerate. Its representative \(b\) has nonzero trace, so
\(F(b)=E\) is a separable quadratic field. For \(a\in B^\times\)
replace \((w,b)\) by \((wa,\bar a b a)\); the two \(R\)'s are
equal. The two Weyl words therefore have ratio fixing \(W\)
and acting on \(e\) by \(\bar a b a b^{-1}\).
Scalar scales come from the words with \(w=0,b\in F^\times\).
After scalar correction we get scales \([\bar a,b]\).
Replacing \(b\) by \(b+c\), \(c\in F\), gives, up to scalar
multiples, every noncentral element of \(E^\times\). The
commutator closure is all of \(\operatorname{SL}_1(B)(F)\):
the ramified and unramified valuation-layer calculations given
above for \([B^\times,E^\times]\) apply in characteristic two
as well. They use only the separable quadratic \(E\), the
quadratic finite residue extension, and its norm-one conic.

The spin scale lift has the split exterior-algebra description
already computed above, giving pairs \((t,s)\) with
\(\operatorname{Nrd}(t)=s^2\). Descent of \(s\) follows from
its uniqueness as a square-root character on the connected
scale-cover subgroup. On \(F\)-points its square root is
unique and continuous on the set of squares, since
\(v(s_1^2-s_2^2)=2v(s_1-s_2)\). Hence the scalar scale words
and the limits of the norm-one commutator words yield all the
spin scale lifts. Together with the \(W\)-rotations this
gives the whole spin Levi in the radical closure.
The relative root argument above then gives all local points:
the centralizer of a maximal split torus containing the pair
scale lies in this Levi, and the relative-root rays supply the
unipotent and Weyl factors. Every global group in this case
was proved isotropic, so the additive approximation argument
with these local generation statements completes strong
approximation for nontrialitarian type \(D\) in characteristic two.

\medskip\noindent
\textit{Trialitarian type \(D_4\).}
At a finite completion \(F\) with transitive diagram action,
let \(E/F\) be the cubic field fixed by the stabilizer of a
diagram vertex. Over \(E\) the vertex is fixed and gives the
vector model discussed above.

\emph{The Albert realization of the local trialitarian form.}
Let \(F\) be nonarchimedean and let \(E/F\) be its diagram cubic
field. There is a simply connected reference group \(M_0\) of
the given diagram type inside the automorphism group of an
Albert algebra \(J_0\). In fact take the split Albert algebra
with its diagonal frame. Permutations of its diagonal coordinates
act by simultaneous row and column permutations on the Hermitian
matrices. They preserve the unit and the cubic norm in
NA13-C04.07, since the trace of a triple product is cyclic and
reversal conjugates the product and leaves its trace unchanged.
Twist this action on the split algebra by the permutations of
the three embeddings of \(E\). The pointwise stabilizer of the
descended diagonal algebra is \(M_0\): after scalar extension it is
the frame stabilizer of NA13-C04.09 and its Appendix C.
The three Peirce modules are the vector and the two half-spin
modules, so the diagram action, and in particular the action on
the three outer central characters, is the specified permutation
action. Consequently \(M_0\) and the given group \(G\) are inner
forms of one another by Result~\ref{res:F06}. The Peirce module
indexed by the distinguished embedding of \(E\) is a genuine
eight-dimensional vector representation of \(M_{0,E}\),
with its invariant nonsingular norm.

For clarity, one can perform the central lifting calculation
without lifting through a possibly inseparable central cover
on separable-closure points. Write \(C=\operatorname{Res}_{E/F}
\mathbb G_m\). The three scalar characters identify the center
\(Z\) of \(M_0\) with the subgroup of \(C\) given by \(z^2=1\)
and \(N_{E/F}(z)=1\): over a separable closure the three nonzero
characters of \((\mu_2)^2\) have product one. Enlarge \(M_0\)
to the central product \(\widetilde M=(M_0\times C)/
\{(z,z^{-1}):z\in Z\}\). Its map to the adjoint group has kernel
the smooth torus \(C\). The three Peirce representations, extended
by the respective scalar coordinates of \(C\), are representations
of this central product. In particular over \(E\) it acts on the
chosen Peirce block by similarities and \(C\) acts on it by one of
its scalar coordinates.

Lift the inner cocycle for \(G\) to \(\widetilde M\) over a
separable closure. Such lifts exist because their fibers are
torsors under the split torus \(C\) there. Their multiplication
defect is a two-cocycle in \(C\). Evaluation at one embedding
after restriction to its stabilizer identifies the defect class
with the Brauer class of the descended vector endomorphism
algebra over \(E\). The restriction to one coordinate detects whether the defect
is a coboundary. Here is the cochain argument. Label the coordinates
by the embeddings of \(E\), and regard a lift of a Galois element
\(\gamma\) on a coordinate \(i\) as an arrow \(i\to\gamma i\).
Allow a scalar on each arrow and use the degree-two defect for the
multiplication of arrows, with the Galois action on the second
scalar in a composition. The cocycle identity is exactly
associativity of this multiplication. At the distinguished
coordinate \(i_0\), triviality of the restricted defect gives a
multiplicative lift for its stabilizer. Choose once and for all a
lift of a transporting arrow \(r_i:i_0\to i\) to each coordinate.
Lift any arrow \(\gamma:i\to j\) by first taking the inverse of the
fixed lift of \(r_i\), then the multiplicative lift of
\(r_j^{-1}\gamma r_i\), and then the fixed lift of \(r_j\).
These lifts compose, giving one scalar adjustment on every
coordinate for each \(\gamma\). Thus the original defect on \(C\)
is a coboundary. This also shows that a restriction of a class from
\(F\) is represented on all coordinates by the diagonal cochain.
The norm on this induced coefficient group is the corestriction:
it multiplies the scalar defects of the permuted tensor factors.

The square of the scalar defect over \(E\) is a coboundary by
the multiplier of the invariant Peirce norm. Its norm
corestriction vanishes as well: the norm character on \(C\)
is trivial on \(Z\), so extends to \(\widetilde M\) by being
trivial on \(M_0\). Applying this character to the lifts gives a
one-cochain whose coboundary is the norm of their defects.
A \(2\)-torsion local Brauer class is either zero or the unique
quaternion division class. Indeed in the cyclic valuation description
above the division representative of degree \(n\) has an unramified
splitting field of degree \(n\), and the lift of its cyclic generator
has \(n\)-th power a uniformizer. The scalar class is determined by
that power modulo norms: rescaling the lift multiplies the power
by a norm, and with power one the lifts give genuine cyclic descent.
Tensor powers raise the power to the corresponding power of the
uniformizer. Norm valuations from the unramified field are divisible
by \(n\), while the \(n\)-th power of a base-field uniformizer is
a norm. The class therefore has order \(n\), giving \(n=1\) or
\(n=2\); in degree two there is just the single unramified quadratic
field and the same uniformizer class, since the norm is onto
on units by the residue trace argument above. The division class
over \(E\) is the restriction of the class over \(F\), since the
quaternion division algebra over \(F\) stays division over a cubic
extension. The norm on a
restriction multiplies its Brauer class by \(3\): tensor the
three scalar-extended copies, or multiply their scalar
defects after permuting the tensor factors. This contradicts
zero corestriction for a nonzero class. Thus the defect class
vanishes and the lifts may be adjusted to a genuine cocycle
in \(\widetilde M\).

The quotient \(\widetilde M/M_0=C/Z\) is isomorphic to \(C\).
Indeed the homomorphism \(x\mapsto x^2/N_{E/F}(x)\) has kernel
\(Z\): its norm is \(N_{E/F}(x)^{-1}\), and the kernel equations
then give \(x^2=1\). On the split coordinate lattice its matrix
\(2I-\mathbf 1\mathbf 1^t\) has determinant \(-4\), so the map is
surjective as a torus isogeny, also in characteristic two.
The image of the cocycle in this quotient is a coboundary by
Hilbert 90 over \(E\). The fiber over the adjusting point is a
smooth \(M_0\)-torsor, so it has a point over the separable
closure. Use this lift to conjugate the cocycle accordingly. It now takes values in
\(M_0\). Twisting the action of \(M_0\) on \(J_0\) by this cocycle
descends another Albert algebra and the descended pointwise
stabilizer is \(G\), since its adjoint descent datum is the inner
twist we started with. The preceding local Albert calculation
shows that this ambient Albert algebra is split.

\emph{Isotropy inside the split Albert group.}
We prove the required split-torus assertion over every local
field. Compare the embedding criterion in
\cite[Proposition 2.1]{alsaodyf4}. Let \(P=\operatorname{Aut}(J_{\mathrm{split}})\) and let \(M\)
be the stabilizer of the ordered split diagonal frame. Over an algebraic
closure the \(D_4\) subgroup just constructed is conjugate to \(M\).
We will find a nonempty open set of conjugates of a fixed split Borel
\(B_0=T_0U_0\) meeting \(M\) in a smooth neutral component with
positive-dimensional projection to the torus.

Use the integral coordinates \((A,B,D)\) of Appendix C of
\emph{Nonassociative Algebras}, in which the norm is
\(N=\det A+\det B+\det D-\operatorname{tr}(ABD)\) and the unit is
\((I,0,0)\). Take the split torus generated by determinant-one
diagonals acting by \((A,B,D)\mapsto(gAg^{-1},gB,Dg^{-1})\)
and by the diagonal octonion automorphisms multiplying the three
Zorn upper coordinates by \(y_0,y_1,y_2\) with product one.
This is a maximal torus, since its dimension is four.
With zero-based matrix indices, put
\[
 \epsilon_0=(1,0),\qquad \epsilon_1=(0,1),\qquad
 \epsilon_2=(-1,-1).
\]
The coordinate weights of \(A_{ij},B_{ij},D_{ji}\) for its
two diagonal parameters are respectively
\((\epsilon_i-\epsilon_j,0)\), \((\epsilon_i,\epsilon_j)\),
and \((-\epsilon_i,-\epsilon_j)\). Declare the positive roots
to be the ones with lexicographically positive weights.
The line of functionals spanned by \(\ell(A,B,D)=A_{20}\)
is \(B_0\)-stable: this coordinate has smallest weight and
\(U_0\) raises weights. Therefore
\(\ell(v_0)/\ell(v_1)\) is a \(B_0\)-invariant rational function
on the \(P\)-orbit of the ordered frame \((v_0,v_1,v_2)\).
It is nonconstant. The transformation by \(g=I+E^{\mathrm{mat}}_{20}\)
moves the frame to one where \(\ell(v_0)=1,\ell(v_1)=0\);
permuting the two indicated diagonal labels before applying \(g\)
interchanges these two values. Such permutations are realized
by determinant-one signed permutation matrices. Since the frame
orbit has dimension \(52-28=24\), an orbit of \(B_0\) cannot
have dimension 24. Indeed such an orbit would be dense, contradicting
the nonconstant invariant. Thus every \(B_0\)-stabilizer in this
orbit has dimension at least five.

Here is an integer rank certificate for a frame in this orbit.
Write \(g_{ij}(t)\) for the above transformation by the elementary
matrix \(I+tE^{\mathrm{mat}}_{ij}\) and let \(p_i(t)\) act
simultaneously on the three octonion blocks by the Zorn
automorphism
\[
 (a,u,v,b)\longmapsto
 \left(a-tv_i,\,
   u+t(a-b)e_i-t^2v_i e_i,\,
   v+t(e_i\times u),\, b+tv_i\right).
\]
For example, in the Zorn basis \(a_0,b_0,u_j,v_j\) the two
idempotents map to \(a_0+t u_i,b_0-t u_i\). The upper basis
maps to \(u_j+t(e_i\times e_j)_{\mathrm{lower}}\), and the lower
basis maps to \(v_j-t\delta_{ij}a_0+t\delta_{ij}b_0-
t^2\delta_{ij}u_i\). The Zorn multiplication table gives
orthogonal idempotents and the same products of the resulting
upper and lower bases: the opposite mixed products are
\(\delta_{j\ell}\) times the appropriate new idempotent,
and the upper-upper and lower-lower products are respectively
the lower cross product and the negative upper cross product.
The inverse map is \(p_i(-t)\). This verifies the automorphism
and preservation of the three-block norm by the identities in
NA13-C04.07--08. Apply in order
\[
\begin{gathered}
 g_{02}(1),\ g_{10}(1),\ g_{20}(-1),\ g_{01}(-1),\
 g_{10}(1),\ p_0(1),\ g_{10}(-1),\\
 p_1(1),\ p_2(-1),\ g_{20}(1)
\end{gathered}
\]
to the split frame, obtaining \((v_0,v_1,v_2)\).

For reproducibility, use row-major order for each of the three
matrices, and the order \(A,B,D\) on the 27 concatenated coordinates.
The following specifies the columns of the evaluation matrix without
any choices of Chevalley signs. For each weight \(\lambda\), list in
lexicographic order the pairs \((r,s)\) with coordinate weights
\(w_r-w_s=\lambda\). Put an unknown \(x_{rs}\) in each such position of
\(X\), and solve
\[
 X(I,0,0)=0,\qquad \sum_r \frac{\partial N}{\partial z_r}(z)
 (Xz)_r=0
\]
by coefficient comparison. For the second equation use the six
determinant monomials in each matrix and the 27 negative terms
\(A_{ij}B_{jk}D_{ki}\). Row reduction over \(\mathbb{Q}\) gives
24 positive weight kernels of dimension one, with the normalized
free-coordinate matrices integral, and a zero weight kernel of
dimension four, with the basis given by setting each free coordinate
in turn to one and the others to zero. The normalization likewise
has unit free coordinates and integral entries in weight zero.
Reducing the 28 matrices in any characteristic gives independent
solutions, by the separate weights and the free coordinates, and
they span the positive and zero weight parts of the tangent space
described in NA13-C04.09. Order the 24 positive weight kernels
lexicographically, then append the four zero weight vectors in
free-coordinate order. For each such matrix use the column
\((Xv_0,Xv_1,Xv_2)\). In the following lists all indices are
zero-based. They give minors of the first 24 columns and of the
full matrix, respectively:

For \(\operatorname{Lie}U_0\), a minor with determinant \(-3\) has ordered indices
\[
\begin{array}{rl}
\text{rows:}&0,\,1,\,2,\,3,\,4,\,5,\,7,\,9,\,10,\,11,\,12,\,13,\\
&14,\,16,\,19,\,27,\,28,\,29,\,30,\,37,\,43,\\
\text{columns:}&0,\,1,\,2,\,3,\,4,\,5,\,6,\,7,\,8,\,9,\,10,\,11,\\
&12,\,13,\,14,\,15,\,16,\,18,\,19,\,20,\,22.
\end{array}
\]

For \(\operatorname{Lie}U_0\), a minor with determinant \(2\) has ordered indices
\[
\begin{array}{rl}
\text{rows:}&0,\,1,\,2,\,3,\,4,\,5,\,9,\,10,\,11,\,12,\,13,\,14,\\
&16,\,19,\,27,\,28,\,29,\,30,\,32,\,37,\,43,\\
\text{columns:}&0,\,1,\,2,\,3,\,4,\,5,\,6,\,7,\,8,\,9,\,10,\,11,\\
&12,\,13,\,14,\,15,\,16,\,17,\,18,\,19,\,21.
\end{array}
\]

For \(\operatorname{Lie}B_0\), a minor with determinant \(-3\) has ordered indices
\[
\begin{array}{rl}
\text{rows:}&0,\,1,\,2,\,3,\,4,\,5,\,6,\,7,\,9,\,10,\,11,\,12,\\
&13,\,14,\,16,\,19,\,27,\,28,\,29,\,30,\,32,\,37,\,43,\\
\text{columns:}&0,\,1,\,2,\,3,\,4,\,5,\,6,\,7,\,8,\,9,\,10,\,11,\\
&12,\,13,\,14,\,15,\,16,\,18,\,19,\,20,\,22,\,25,\,26.
\end{array}
\]

For \(\operatorname{Lie}B_0\), a minor with determinant \(-2\) has ordered indices
\[
\begin{array}{rl}
\text{rows:}&0,\,1,\,2,\,3,\,4,\,5,\,7,\,9,\,10,\,11,\,12,\,13,\\
&14,\,16,\,19,\,27,\,28,\,29,\,30,\,32,\,33,\,37,\,43,\\
\text{columns:}&0,\,1,\,2,\,3,\,4,\,5,\,6,\,7,\,8,\,9,\,10,\,11,\\
&12,\,13,\,14,\,15,\,16,\,17,\,18,\,19,\,21,\,25,\,26.
\end{array}
\]

One way to check the displayed minors is to replace each \(i,j\)
entry below and right of each pivot \(a\) by
\((a m_{ij}-m_{i k}m_{k j})/a_{\mathrm{prev}}\), with the initial
previous pivot one, swapping rows when needed and recording the
sign. The last surviving entry gives the indicated determinant.
Each pair of determinants is relatively prime. Hence over every
field the ranks are at least 21 for \(U_0\) and 23 for \(B_0\).

The two rank inequalities are open conditions on conjugating
elements. They hold on a nonempty open set for the subgroup \(G\)
already embedded in \(P\), since it becomes conjugate to \(M\)
over a separable closure. The split big cell makes \(P(F)\)
Zariski dense, so choose an \(F\)-rational such conjugating element.
Conjugate \(B_0\) by it (or its inverse as appropriate), and let
\(V=B_0^g\cap G\). Its dimension is at least five by the invariant
above and its tangent dimension at the identity is at most five.
Thus \(V\) is smooth at the identity even geometrically, and its
smooth identity component has dimension five. The corresponding
unipotent intersection has tangent dimension at most three.
It follows that the image of \(V^\circ\) in the split torus
\(B_0^g/U_0^g\) is positive dimensional.

In a faithful triangular representation of this Borel the smooth
group \(V^\circ(F)\) is Zariski dense, by the local implicit
function theorem at the identity. Its image in some diagonal
coordinate is therefore infinite. The roots of unity in a
nonarchimedean local field are finite (in characteristic \(p\)
they inject into the residue field, and in characteristic zero a
sufficiently small unit ball is torsion-free by the logarithm).
Choose an element of \(V^\circ(F)\) with a diagonal entry of
infinite multiplicative order. Its semisimple part is defined
over \(F\) because all its eigenvalues are in \(F\), and belongs
to \(V\) by the matrix Jordan decomposition in closed algebraic
subgroups. Its powers have Zariski closure with a positive-dimensional
identity component inside a diagonal torus over \(F\).
This gives a nontrivial split torus in \(G\). In particular the
cubic-field Levi argument applies at every trialitarian place,
with no prior isotropy assumption.

In characteristic different from two we can now use the
eight-dimensional quadratic form for the group over \(E\). The local dimension
bound proved by the ternary norm argument above shows
its anisotropic kernel has dimension at most four.
Moreover, a four-dimensional anisotropic form over
such a field is a scalar multiple of a quaternion
division norm: write a ternary subform as \(\lambda\)
times the pure norm. This represents every nonzero
class except \(-\lambda\) up to squares, so the
remaining diagonal coefficient must have square class
\(\lambda\). In particular that kernel has square
signed discriminant. The signed discriminant decides
whether Galois exchanges the two half-spin modules:
the even Clifford volume element of an orthogonal
basis has square the signed discriminant, and its
two eigenspaces separate the half-spin modules over
a splitting field. Thus with transposition over \(E\)
the anisotropic kernel has dimension two. The model
with three hyperbolic planes is quasi-split.

With trivial diagram action over \(E\), the cubic
field is cyclic over \(F\). The two half-spin
matrix-algebra classes are Galois conjugates of the
now split vector class, by evaluation of the same
central defect on the conjugate characters. These
classes are the components of the even Clifford
algebra: over a split field its action on even and
odd exterior powers gives the two full matrix
algebras, and conjugation on each is unchanged by a
scalar lift of the spin action. If the anisotropic
kernel were a scaled quaternion division norm
\(\lambda\langle1,-a,-b,ab\rangle\), its even Clifford
components would both have the division class:
the products of the first orthogonal generator
with the second and with the third anticommute
and square to \(\lambda^2a,\lambda^2b\).
They generate a division quaternion subalgebra
in each component, which already has dimension four.
Adding a hyperbolic plane preserves both Brauer
classes. To see this, realize the full Clifford
algebra of the new form as the two-by-two matrix
algebra of the old Clifford algebra, with each old
vector sent to \(\operatorname{diag}(x,-x)\)
and the hyperbolic generators to the two off-diagonal
matrix units (with pairing normalized to one).
The even algebra has even diagonal entries and
odd off-diagonal entries. Each factor has a
corner from the first diagonal entry isomorphic to
the corresponding old even factor, so the Brauer
classes agree. This rules out dimension four.
Square discriminant rules out dimension two,
so the group over \(E\) is split.

In characteristic two the same conclusion about \(G_E\) follows
from the quadratic-pair descent above. The scalar algebra is split,
so the regular form has dimension eight. We spell out the
corresponding Arf and Clifford calculation. In a symplectic
basis of a quadratic space put \(s_i^2=a_i\), \(t_i^2=b_i\)
and \(s_it_i+t_is_i=1\) in the Clifford algebra. The pairs of
different indices commute. Consequently
\[
 z=\sum_i s_it_i,\qquad z^2+z=\sum_i a_i b_i,
 \qquad zv=v(z+1)
\]
for each basis vector \(v\). In the split exterior model \(z\)
separates the even and odd spin modules. Thus the quadratic
étale algebra of their labels is \(F[z]\) for this equation,
and the right-hand side modulo \(x^2+x\) is the Arf invariant.
Here and in the following local calculation the base field
can be \(E\).

For any separable quadratic field \(L=E[i]\), \(i^2+i=d\),
the norm group is the kernel of the map \(c\mapsto[d,c)\),
where the quaternion algebra \([d,c)\) has generators
\(i,j\) with \(j^2=c\) and \(ji=(i+1)j\). The split criterion
follows by representing \(L\) on itself: the operator
inducing its nontrivial automorphism is multiplication by
a scalar followed by that automorphism, and its square is the
scalar's norm. Conversely if the algebra splits, the same
description applies to its two-dimensional module and gives
the norm condition. The classes satisfy
\([d,c)+[e,c)=[d+e,c)\): in their tensor product use \(i_1+i_2,j_1\)
for the indicated subalgebra, and \(i_2,j_1j_2\) for its
centralizer, which is \([e,c^2)\) and splits. Multiplicativity
in \(c\) follows by taking \(i_1,j_1j_2\) in the tensor product
\([d,c_1)\otimes[d,c_2)\); its centralizer has a nontrivial
idempotent from \(i_1+i_2\), satisfying \(x^2+x=0\), and splits.
Over the local field every separable quadratic field embeds
in the unique quaternion division algebra: the unramified
one already occurs in the cyclic valuation construction, and
a ramified one splits it by rescaling the uniformizer, as in
the preceding quadratic spin argument. A degree-two splitting
field embeds by the module argument there. Its norm character
to the two-element Brauer subgroup is therefore nonzero.
Distinct quadratic extensions yield distinct norm characters,
since their sum is the nonzero character for the sum of their
Artin--Schreier classes. Their norm subgroups consequently
multiply to \(E^\times\).

The elementary bound above leaves a nonsingular anisotropic
kernel of dimension at most four. An anisotropic binary
quadratic form is a scalar multiple of a separable quadratic
norm, with nonzero Arf invariant. If a nonsingular
four-dimensional anisotropic form had nonzero Arf invariant,
its two binary summands would be norms from distinct
quadratic fields. Their norm subgroups multiply to the full
group, so the scaled norms would share a nonzero value. This
contradicts anisotropy. Hence a four-dimensional anisotropic
kernel has Arf invariant zero. In a symplectic basis \(u,v,x,y\)
with \(u^2=a,v^2=b,x^2=c,y^2=d\), the two scaled binary norms
then come from the same quadratic field, and the form is
\(a\) times the reduced norm of \([ab,c/a)\). This quaternion
must be division in the anisotropic case.

For this four-dimensional form the even Clifford algebra
contains the quaternion generators \(uv,ux/a\), with relations
\((uv)^2+uv=ab\), \((ux/a)^2=c/a\), and
\((ux/a)(uv)=(uv+1)(ux/a)\). The central element \(uv+xy\)
satisfies the split Artin--Schreier equation for the Arf
invariant. Its two idempotents are nonzero by the Clifford
monomial basis. The quaternion maps injectively to each
component, and each has dimension four, so both components
have the division Brauer class. Adding a hyperbolic pair
gives the two-by-two matrix Clifford construction above,
now with old vectors acting by \(\operatorname{diag}(v,v)\).
The even part has even diagonal entries and odd
off-diagonal entries. The displayed commutation with \(z\)
matches an old \(+\) idempotent in the first corner with
an old \(-\) idempotent in the second, and vice versa.
An invertible old vector gives inverse odd maps between
these opposite corners. Hence the old even factors are
full corners of the respective new factors; adding
hyperbolic planes preserves the two Brauer classes.

If the diagram action over \(E\) is a transposition, the Arf
invariant is nonzero, so the anisotropic kernel is binary and
the group is quasi-split. If it is trivial, the three outer
matrix-algebra classes are conjugate under the cyclic
diagram action over \(F\). They are all split, since the
chosen vector algebra is split. The even Clifford
identification and the preceding calculation rule out a
four-dimensional anisotropic kernel, and zero Arf invariant
rules out dimension two. The eight-dimensional form is
then split.

The Albert embedding just constructed gives isotropy over \(F\).
In a minimal \(F\)-parabolic, the roots vanishing on
a maximal \(F\)-split torus form the root system
of the derived Levi, as in
Results~\ref{res:D08} and~\ref{res:D09}. Its simple
system is the subset of absolute simple roots
that vanish on this torus. The diagram action preserves this
subset: changing the absolute Borel and torus in the
rational parabolic only changes the positive system
there by a Levi conjugation. The derived Levi is
anisotropic, since a split torus in it would enlarge
the original split torus, and the subset is proper by
isotropy. The only possible nonempty such subsets here
are the central vertex or the three outer vertices.
In the first case the kernel is an anisotropic \(A_1\)
over \(F\), hence the norm-one group of a quaternion
division algebra by the rank-one matrix descent above.
It remains anisotropic over the cubic field. In the
second case the kernel is a restriction of scalars
from \(E\) of an anisotropic \(A_1\) over \(E\). Indeed
the derived Levi is simply connected (its simple
coroots are a subset of the basis of the ambient
coroot lattice), and its three factors descend by the
transitive permutation. A split torus in the fixed
factor over \(E\) would give a split cocharacter over
\(F\) by taking its conjugates in the other factors.
After extension to \(E\) this fixed \(A_1\) thus stays
anisotropic. Both cases contradict the quasi-splitness
of the group over \(E\): the centralizer of an
\(E\)-split torus in a quasi-split group is quasi-split.
To see the latter assertion, enlarge the torus to a
maximal split torus, choose an \(E\)-Borel containing
its centralizer (a torus in the quasi-split case),
and intersect it with the smaller torus's centralizer;
the positive root spaces of weight zero give a Borel
of that centralizer over \(E\).

Thus every finite completion with transitive trialitarian
diagram action is quasi-split. At the other completions the
diagram action fixes a vector vertex, so the nontrialitarian
argument gives local generation for isotropic groups. For a
globally isotropic trialitarian group all completions are
isotropic, and these local generation statements complete
the earlier additive approximation argument for it.

\medskip\noindent
\textit{Torsors over the maximal unramified extension.}
Let \(F\) be a nonarchimedean local field with finite residue
field, and put \(K=F^{\rm nr}\) in a fixed separable closure.
For every finite separable extension \(M/K\), one has
\(\operatorname{Br}(M)=0\). Indeed, write \(M=K(a)\) with a
separable primitive element. Its minimal polynomial and the
finitely many structure constants of a central simple algebra
over \(M\) descend to a finite unramified subextension of \(K/F\)
and to the extension \(F'\) of that subfield generated by \(a\).
Centrality and simplicity descend by field extension. We have
\(M=F'K\). The extension \(F'K/F'\) is unramified and its
residue field contains all finite extensions of the residue
field of \(F'\), so \(F'K=(F')^{\rm nr}\). By the local
division-algebra calculation above, the division factor
of a central simple algebra over \(F'\) has an unramified
maximal subfield and splits there. Hence the descended algebra
splits over \(M\).

We deduce \(H^1(K,T)=1\) for every \(K\)-torus \(T\), in
all characteristics. Let \(L/K\) be a finite Galois splitting
extension with group \(\Gamma\), let \(A=L^\times\) be written
as an additive \(\Gamma\)-module, and let \(M_T=X_*(T)\).
Then \(T(L)=A\otimes_{\mathbb Z} M_T\) with diagonal action.
Split Hilbert 90 over \(L\) identifies \(H^1(K,T)\) with
\(H^1(\Gamma,A\otimes M_T)\).

For every \(J\leq\Gamma\), Hilbert 90 gives \(H^1(J,A)=0\).
Also \(H^2(J,A)=0\), since the relative crossed-product
class lies in \(\operatorname{Br}(L^J)=0\). For this
identification, a scalar cocycle gives the crossed-product
algebra with basis elements \(u_\gamma\), multiplication
\(u_\gamma a=\gamma(a)u_\gamma\) and
\(u_\gamma u_\tau=c_{\gamma,\tau}u_{\gamma\tau}\).
If it is split, realize it on a vector space of dimension
\([L:L^J]\) over \(L^J\). This space is one-dimensional
over the included copy of \(L\), and the \(u_\gamma\)'s
then act as scalar multiples of the Galois operators, so
\(c\) is a coboundary.

For an \(\ell\)-group \(J\leq\Gamma\), all positive-degree
cohomology of \(A\) vanishes and \(N_J A=A^J\). For \(J\)
cyclic of order \(\ell\), the periodic resolution alternates
the operators \(\sigma-1\) and \(N_J=1+\cdots+\sigma^{\ell-1}\);
the result follows from degrees one and two. For larger
\(J\), take a central subgroup \(C\) of order \(\ell\).
Since \(H^i(C,A)=0\) for \(i>0\), taking \(C\)-invariants
on an injective \(J\)-resolution produces an injective
\(J/C\)-resolution of \(A^C\). Here restriction preserves
injectives because induction is exact, and \(C\)-invariants
preserve injectives because inflation is exact. Thus the
positive-degree cohomology of \(J\) on \(A\) agrees with
that of \(J/C\) on \(A^C\). This argument applies to
the preimage of each subgroup of \(J/C\), so the degree-one
and degree-two vanishing hypotheses persist. Induction
and the formula \(N_J=N_{J/C}N_C\) give the claim.

Now fix an \(\ell\)-Sylow subgroup \(P\) of \(\Gamma\).
Choose a surjection \(Q\to A\) of \(P\)-modules with \(Q\)
a free \(\mathbb Z[P]\)-module, and write \(B\) for its
kernel. The map \(Q^P\to A^P\) is onto: choose a lift
before applying \(N_P\). Since free group-ring modules
are induced, the long exact sequence gives
\(H^i(P,B)=0\) for \(i>0\). The underlying abelian group
of \(B\) is free. Multiplication by \(\ell\) on \(B\)
therefore gives \(H^1(P,B/\ell B)=0\). The sole simple
\(\mathbb F_\ell[P]\)-module is trivial: for a central
element \(\sigma\) of order \(\ell\), the endomorphism
\(\sigma-1\) is central and nilpotent, hence kills a
simple module, and one inducts on \(P/\langle\sigma\rangle\).
By composition series, vanishing of
\(\operatorname{Ext}^1_{\mathbb F_\ell[P]}(\mathbb F_\ell,B/\ell B)\)
gives the same vanishing with any finite-length first
argument. The left-ideal quotients of this group ring
have finite length, so Baer's criterion makes \(B/\ell B\)
injective. Equivalently every homomorphism from a left
ideal extends to the ring, which extends maps on arbitrary
submodules by adjoining one cyclic generator at a time.

Taking diagonal tensor product with
\(M_T/\ell M_T\) preserves injectives, by the exact tensor
adjunction with its dual. Thus
\(H^i(P,(B\otimes M_T)/\ell)=0\) for \(i>0\).
The abelian group \(B\otimes M_T\) is free. Multiplication
by \(\ell\) is consequently onto on these positive-degree
cohomology groups before reduction, and those groups are
annihilated by \(|P|\), by the index argument below for the
trivial subgroup. They vanish. Tensoring
\(0\to B\to Q\to A\to0\) by \(M_T\) is exact, and
\(Q\otimes M_T\) is still induced, since
\(g\otimes m\mapsto g\otimes g^{-1}m\)
removes the second diagonal action on each regular summand.
Hence \(H^1(P,A\otimes M_T)=0\).
For a \(\Gamma\)-module \(D\), the maps sending \(d\) to
\(\sum_{gP}d\otimes[gP]\) and \(d\otimes[gP]\) to \(d\)
factor multiplication by \([\Gamma:P]\) through
\(D\otimes\mathbb Z[\Gamma/P]\). The latter module is the
coinduced restriction to \(P\); its \(\Gamma\)-cohomology
is the \(P\)-cohomology of \(D\), by the exact
restriction--coinduction adjunction on injective resolutions.
Apply this observation to \(D=A\otimes M_T\) and each
Sylow subgroup. The indices have greatest common divisor
one, so \(H^1(\Gamma,D)=0\), as claimed.

\medskip\noindent
\textit{Quasi-split exceptional reference groups.}
We use the ordered root words of \cite[Section 7]{steinbergregular}
to prove the required conjugacy assertion.
We show \(H^1(K,G_0)=1\) when \(G_0\) is simply connected
and is split of type \(E_6,E_7,E_8\) or quasi-split of
outer type \(E_6\). Work with the pinning of the quasi-split
form from Result~\ref{res:F06}. Order the simple roots so
that each diagram orbit is consecutive. In outer \(E_6\)
the two nontrivial orbits are \(\{1,6\}\) and \(\{3,5\}\)
in the Bourbaki numbering; each is a pair of orthogonal
roots, so the two rank-one subgroups commute.
Writing \(n_i\) for the pinned rank-one Weyl representative
and \(x_{\alpha_i}\) for its positive root map, put
\[
 S(t)=x_{\alpha_1}(t_1)n_1\cdots x_{\alpha_r}(t_r)n_r,\qquad
 c=s_1\cdots s_r.
\]
Moving the representatives to the right gives a root word
on \(\beta_i=s_1\cdots s_{i-1}\alpha_i\), with coordinates
obtained by conjugation and no change of the parameters,
followed by \(n_c=n_1\cdots n_r\). For the fundamental
weights \(\omega_i\) one has \(\beta_i=(1-c)\omega_i\):
the later reflections fix \(\omega_i\), and the earlier
ones fix it after subtracting the \(\alpha_i\) term.
The \(\beta_i\) are triangular with diagonal one in the
simple-root basis, so \(1-c\) is invertible over
\(\mathbb Q\) on the weight space.

Consider the trace on the simple module of highest
weight \(\omega_j\), with highest weight line of
dimension one, as in Results~\ref{res:E02} and~\ref{res:E06}.
A contribution from weight \(\mu\) to the coefficient of
\(\prod_i t_i^{m_i}\) in the trace returns to that weight
after first going to \(c\mu\) and then adding
\(\sum_i m_i\beta_i\). Thus \(\mu=\sum_i m_i\omega_i\),
with \(m_i\geq0\), is dominant. The contribution from the
highest weight is exactly \(d_j t_j\), with \(d_j\ne0\).
Indeed \(n_c\) moves the highest line to weight
\(\omega_j-\beta_j\), and
\(\langle\omega_j,\beta_j^\vee\rangle=1\).
This lower line is the lower weight line in the
rank-one calculation: the two weight spaces have
dimension one by Weyl conjugacy, and the rank-one word
on a vector of highest weight one has a nonzero
linear raising coefficient, including in characteristics
two and three. For every other trace monomial the weight
\(\mu\) is strictly less than \(\omega_j\) in the dominance
order. Pairing with \(2\rho^\vee\), which is positive on
simple roots and on fundamental weights, shows that every
variable in such a monomial has smaller height than
\(\omega_j\). The traces are therefore triangular,
and prescribe a unique tuple of parameters for each
tuple of trace values.

These traces on a maximal torus separate geometric
semisimple classes. Indeed \(\mathbb Z[X]^W\) has basis
the orbit sums of dominant weights. The trace characters
of the fundamental modules have leading orbit-sum
coefficient one, as does the product with any prescribed
highest dominant weight. Subtracting the lower orbit
sums inductively shows, as in Result~\ref{res:E09},
that the fundamental characters generate that invariant
ring. The orbit-sum basis commutes with reduction of
coefficients. Over an algebraic closure, distinct finite
Weyl orbits in the torus are separated by an invariant
Laurent polynomial: interpolate a function with values
zero and one on the two orbits, then multiply its
Weyl translates. It follows that any two elements with
matching fundamental traces have geometrically conjugate
semisimple parts. In fact the trace of an element equals
the trace of its semisimple part, by upper triangularizing
the commuting unipotent part on each eigenspace. If the
semisimple part has torus centralizer, the commuting
unipotent part is trivial.

Now let \(z_\gamma\) be a Galois cocycle with values
in \(G_0(K_s)\), where \(K_s\) is a separable closure.
Twisting by its inner conjugations gives a connected
smooth \(K\)-group with a strongly regular semisimple
\(K\)-point \(y\). To check existence, descend the group
and the strongly regular semisimple open to a finite
local subfield of \(K\); the rational points of the
connected smooth group over that field are Zariski
dense, by smooth local coordinates near the identity.
Over an algebraic closure the open is nonempty:
choose a torus point outside the fixed loci of
nonidentity Weyl elements and the root equations
\(\alpha=1\). The conjugation map from the regular
part of the torus has surjective differential at
such points, because the factors in the root
directions are \(1-\alpha(y)\). The centralizer is
the torus: its tangent space is toral, and any
remaining component would normalize and yield a
nontrivial Weyl element fixing that point.

View \(y\) in \(G_0(K_s)\) via the twisted descent, so
\(z_\gamma\gamma(y)z_\gamma^{-1}=y\). Its fundamental
trace values lie in \(K_s\) and are permuted according
to the diagram action. Solve the triangular equations
for \(S(t)\). All the parameters are in \(K_s\) by
triangularity, and uniqueness makes the tuple
equivariant: Galois permutes the pinned factors within
each consecutive, commuting orbit. The modules' trace
characters are permuted by the same diagram action.
Hence \(y_0=S(t)\) lies in \(G_0(K)\), and by the
preceding paragraph \(y_0\) is strongly regular
semisimple in the geometric class of \(y\). Its
transporter to \(y\) is smooth, since it is geometrically
a torsor under the torus centralizer. Choose
\(a\in G_0(K_s)\) conjugating \(y_0\) to \(y\).
Then \(a^{-1}z_\gamma\gamma(a)\) lies in the
\(K\)-torus \(Z_{G_0}(y_0)\). The torus calculation
consumes this cocycle. This proves \(H^1(K,G_0)=1\).

For an inner form over \(K\) of one of these
quasi-split reference groups it is the adjoint
cocycle that must be trivialized. In the inner
types \(E_6,E_7\), the reference is split and
\(Z=\ker(G_0\to G_{0,\mathrm{ad}})\) is \(\mu_m\)
with \(m=3,2\), respectively, by the weight-to-root
lattice quotient of Result~\ref{res:C07}. Set
\(C=\mathbb G_m\) and consider
\[
 E=(G_0\times C)/\{(z,z^{-1}):z\in Z\}.
\]
The map \(E\to G_{0,\mathrm{ad}}\) has smooth kernel \(C\);
the quotient \(E/G_0\) is \(C/Z\simeq\mathbb G_m\).
For a cocycle in \(G_{0,\mathrm{ad}}(K_s)\) choose lifts
in \(E(K_s)\), possible by smoothness. Their failure
of the cocycle equation lies in the scalar \(H^2(K,C)\),
which is zero by the Brauer calculation; correct
the lifts to an \(E\)-cocycle. Its image in \(C/Z\) is
killed by Hilbert 90, and the correcting point lifts
to \(E(K_s)\) since the quotient has smooth kernel \(G_0\).
The corrected cocycle is thus in \(G_0\), where it is
trivial. This method does not require lifting
through an inseparable central isogeny on points.
In type \(E_8\) the center is trivial, and the previous
simply connected calculation already applies.

In outer type \(E_6\) let \(D/K\) be the separable
quadratic diagram field, provided the diagram action is
still nontrivial over \(K\), and put \(C=\operatorname{Res}_{D/K}
\mathbb G_m\). The \(E_6\) diagram reflection inverts
the order-three character group of the center: the
two end fundamental weights are nonzero modulo roots
and their sum lies in the root lattice, as is checked
from the simple Cartan matrix. Hence we embed the center as
\[
 Z=\{z\in C:z^3=1,\ z\bar z=1\}
\]
with the bar denoting the involution of \(D/K\).
The map \(x\mapsto x^2/\bar x\) on \(C\) has that kernel
and defines \(C/Z\simeq C\): on the two split coordinates
its character matrix has determinant three. Use the
same \(E\) as above. Now \(H^2(K,C)=\operatorname{Br}(D)=0\)
and \(H^1(K,C/Z)=1\), by restriction of scalars and Hilbert 90
(the split product coordinates give the coinduced Galois
module). The argument again kills the adjoint cocycle.
If the diagram action has become trivial, the preceding
split-reference argument applies. Consequently every
group under discussion has its quasi-split reference form
over \(F^{\rm nr}\).

\medskip\noindent
\textit{The affine root calculation for the local exceptional forms.}
We give the needed unramified descent for these forms. For the
general local method compare \cite{prasadunramified}.
In this paragraph assume the simply connected reference group is
of type \(E_6,E_7,E_8\) and is split, or is the quasi-split form of
type \({}^2E_6\) for a ramified quadratic diagram field \(L/F\).
We also use its adjoint form. Fix the pinned Borel and torus, denoted
by \(B\) and \(T\), and use the integral root group coordinates
of C08. For split groups an affine root is \(\alpha+j\),
\(j\in\mathbb Z\); its root group is the set of root parameters
of value at least \(j\), and its positive half-apartment is
\(\alpha(x)+j\geq 0\). In the ramified outer case the two-element
root orbits of the pinned involution are orthogonal. The fixed
roots give groups \(\mathbf G_a\), and the paired roots give
groups \(\operatorname{Res}_{L/F}\mathbf G_a\): in a paired root
coordinate the two entries are \(u,\bar u\). This follows from
the absolute root coordinates by ordering the two commuting
members of each orbit together. The orbit assertion for all
roots follows from the simple orbits \(\{1,6\},\{3,5\}\):
the reflections in the orbit sums reduce each positive root
to a simple root by successive height decrease, and their
pinned representatives are fixed by the involution. Consequently
the involution fixes the coordinate of a fixed root. The
relative root system is \(F_4\); long roots use \(v_F(u)\), and
short roots \(v_L(u)\) normalized by \(v_L(\pi)=1\) for a
uniformizer \(\pi\) of \(F\), so the short values lie in
\(\tfrac12\mathbb Z\). The affine arrangement is that of the dual
\(F_4\), since the short root functionals are doubled. The
rank-one computations are those in \(\mathrm{SL}_2(F)\) and
\(\mathrm{SL}_2(L)\), respectively, using the two commuting
absolute rank-one factors in the latter case.

Write \(I\) for the subgroup reducing into the positive Borel
in the split integral model, intersected with the rational group
in the ramified case. The two ordered big cells about that Borel
give coordinates in \(I\): every positive root parameter has
nonnegative value, every negative root parameter has positive
value, and the torus coordinates are units. These conditions can
also be checked after extension to \(L\) in the ramified case,
and give precisely the indicated integral and half-integral
thresholds. Let \(N\) be the relative torus normalizer. Then
\(I\cap N=T_0\), the subgroup of unit torus coordinates, and
\(I,N\) generate the group. In fact an arbitrary root parameter
can be moved to an integral layer by conjugating by a sufficiently
large power of a uniformizer cocharacter, and rational Bruhat
decomposition uses the root groups and \(N\).
In the split simply connected case \(N/T_0\) is the ordinary affine
Weyl group, with translation lattice \(Q^\vee\). In the split
adjoint case the translation lattice is \(P^\vee\), and the alcove
stabilizer in that extended affine group is
\(\Omega=P^\vee/Q^\vee\). Its representatives normalize \(I\)
by the root thresholds. The type-preserving subgroup is generated
by \(I\) and the ordinary affine Weyl representatives. It is also
the image of the simply connected group times the adjoint unit
torus, by the root generators. Its intersection with \(N\) gives
exactly the ordinary affine group modulo \(T_0\): a rational lift
of a normalizer element under the central map normalizes the
preimage torus, and has coroot translation. Thus the whole
adjoint group has this type-preserving subgroup as a normal
subgroup with quotient \(\Omega\).

In the ramified outer case the torus is described by the
coroot basis permuted by the involution (or the fundamental
coweight basis in the adjoint case). The long single factors give
integral translations and a short paired factor gives the
average of its two cocharacters, by using the uniformizer of
\(L\). Thus the two translation lattices are the projections
of \(Q^\vee,P^\vee\) to the invariant apartment. They coincide:
the order-three quotient \(P^\vee/Q^\vee\) has involution \(-1\),
so both twice and three times the projection of every
coweight belongs to the projected coroot lattice. That lattice
is the coroot lattice for the dual \(F_4\) affine arrangement.
In particular \(\Omega=1\) in the ramified case, also for the
adjoint group.

For completeness we verify the affine double cosets and the
geometry used below. Use \(W\) for the ordinary affine Weyl
group in this paragraph. For an affine simple root \(a\) with
underlying root \(\alpha\), write \(U_a\) for its root layer
at the threshold \(j\), and \(U_{a+}\) for its next, strictly
higher layer. There is a representative \(n_s\) of the simple
affine reflection \(s\) in the corresponding rank-one group.
All positive affine roots except \(a\) stay positive under \(s\).
In a big cell factorization of \(I\), put \(U_a\) first, using
positive roots first for a finite simple root and negative
roots first for the added affine simple root. Every other root
factor and the unit torus then lies in
\(I\cap n_s I n_s^{-1}\). In particular
\[
 I n_w I n_s I=I n_w U_a n_s I .
\]
When \(w(a)\) is positive this is just \(I n_{ws} I\).
If \(w(a)\) is negative, parameters in \(U_{a+}\) give this
same double coset, and parameters \(t\) of exact value \(j\)
give the other possibility \(I n_w I\). The latter assertion
is the following rank-one multiplication, with \(\varpi^m\)
the uniformizer power of value \(j\) in the field of the
root parameter:
\[
 \begin{pmatrix}1&t\\0&1\end{pmatrix}
 \begin{pmatrix}0&-\varpi^m\\\varpi^{-m}&0\end{pmatrix}
 =
 \begin{pmatrix}1&0\\t^{-1}&1\end{pmatrix}
 \begin{pmatrix}t\varpi^{-m}&0\\0&t^{-1}\varpi^m\end{pmatrix}
 \begin{pmatrix}1&-\varpi^{2m}/t\\0&1\end{pmatrix}.
\]
The first factor is sent by \(w\) into \(I\), and the latter
two belong to \(I\). This proves equality with the union of
the two double cosets in the descending-length case.

Here is a way to see the disjointness without tacitly assuming
the local double-coset theorem. For an interior rational point
\(x\) of the standard alcove, \(I\) consists of points which,
on scaling by the cocharacter \(x\) evaluated at \(\pi\) in
a sufficiently ramified scalar extension, are integral with
torus reduction. This is an equivalence by the two big cells:
torus reduction places the element in a big cell, and after
unscaling the unique root coordinates have exactly the
positive-depth conditions for \(x\). In the ramified case use
the absolute torus reduction over that scalar extension.
A paired root parameter has positive scaled depth on its two
coordinates simultaneously. A point in the interior need only
avoid relative affine walls; the equivalence imposing torus
reduction holds whether or not an additional absolute wall
is met there. The same description with \(w x\) gives
\(n_w I n_w^{-1}\), since the difference of the scalings
after conjugation by \(n_w\) is an integral Weyl representative
times a unit torus element. Hence the intersection has the
pointwise maximum of the two thresholds in one and the same
big-cell chart, and unchanged torus parameters. Count in these
coordinates modulo a sufficiently deep common congruence
subgroup in the split integral model. This chart and its inverse
are integral on the reducing big cell. Therefore equality
modulo that congruence subgroup is equality of the coordinates
to that precision. In each root-parameter field, crossing a
level changes the count by the residue-field cardinality \(q\).
That cardinality is the same for each root in the ramified
case. In the two opposite root directions there is one such
factor for each intervening affine wall, so
\[
       |I n_w I/I|=q^{\ell(w)}.
\]
If two double cosets coincide, their lengths therefore agree.
The multiplication just proved detects right descents, since
a right simple multiplication contains the original double
coset exactly when the length decreases. For two equal
double cosets of positive length take a common right descent.
Their shortened double cosets coincide by the equality of the
two unions, since the two lengths differ. Induction reduces to
\(I\cap N=T_0\) and proves disjointness. Generation now gives
the full disjoint decomposition of the type-preserving group.

We describe the affine complex explicitly to record the needed
metric and fixed-face assertions. Its chambers are \(G/I\)
for the type-preserving group, and a face of cotype \(J\) is
a coset of \(P_J=\bigcup_{w\in W_J}I n_w I\), where \(W_J\) is
the finite group generated by the simple walls of that face.
The double-coset multiplication shows \(P_J\) is the subgroup
generated by \(I\) and those simple lifts. Face identifications
are by coset inclusion. The intersection rule for faces
of a chamber follows from \(P_J\cap P_{J'}=P_{J\cap J'}\),
by the disjoint double cosets and
\(W_J\cap W_{J'}=W_{J\cap J'}\) in the Coxeter complex. The chambers \(n_wI\) form one copy
of the Euclidean Coxeter complex, and its translates are
apartments. On each such apartment the face identifications
are exactly the usual Coxeter ones by disjointness and by
\(N\cap P_J\) modulo \(T_0\) being \(W_J\). Any two chambers
lie in a common apartment: reduce the first to \(I\) and
write a representative of the other in \(I n_w I\). For a
fixed chamber \(C\) in an apartment \(A\), the double-coset
word \(w\) from \(C\) sends any chamber to the chamber of \(A\)
at that word from \(C\). For a face of cotype \(J\) this
prescription is independent of the incident chamber: the
double-coset multiplication by a simple generator in \(J\)
keeps the word in the same right \(W_J\)-coset, and these
generators join the choices in the face coset. Thus the map
extends on every simplex by the unique type-preserving
isometry to its image. It fixes \(A\), and it is an isometry
on any apartment containing \(C\). This gives a retraction
to \(A\) that does not lengthen piecewise Euclidean paths.
It follows that the intrinsic distance between two points
is their Euclidean distance in an apartment containing them.

For the square-distance midpoint inequality, take an apartment
containing the endpoints \(y,z\), with midpoint \(m\), and use a
chamber of that apartment containing \(m\) as the retraction
center. If \(x\) is any point, its distance to \(m\) is
preserved by this retraction, by using an apartment containing
its chamber and the center chamber. Its distances to \(y,z\)
do not increase. The midpoint equality in the Euclidean
apartment gives the CAT(0) midpoint inequality. This is a
complete geodesic metric: each panel has finitely many chambers
over a finite residue field by the rank-one root layer
calculation, the links of faces are finite by the finite
spherical Coxeter groups \(W_J\), and all cells have the same
finite set of shapes. To see completeness, the barycentric
coordinate of a vertex is globally Lipschitz, with a uniform
constant, by extending it by zero on simplices omitting that
vertex and checking the finitely many shapes on each simplex.
At a point in a chamber, one vertex has coordinate at least
the reciprocal of the number of vertices in that chamber.
Thus a metric ball of fixed small radius stays in a vertex
star, a finite union of compact cells. A Cauchy sequence
therefore converges. In a complete CAT(0) space a finite
group of isometries fixes the unique center of any finite
orbit: the midpoint inequality makes radius-minimizing
sequences Cauchy, and implies uniqueness of the center.

We apply this to a twist over a finite unramified extension
\(F'/F\) containing the cocycle and, in the split-apartment
case, splitting the reference group. Frobenius with the
cocycle is an action of a finite Galois group by isometries,
so it fixes a point and preserves that point's unique open
face. Its barycenter is fixed. Conjugate the face into the
standard alcove, and write \(x\) for its barycenter.
The subgroup \(P_x=P_J\) of the type-preserving group fixes
\(x\), as its faces are the cosets just described. Its
apartments meeting the face have the finite group \(W_x=W_J\)
generated by the walls through \(x\).
The remaining adjoint stabilizer is \(P_x\) times those
length-zero representatives which fix \(x\). The length-zero
class of a stabilizer element preserves the set of vertex
types of the face. Its representative therefore preserves
that face of the standard alcove and fixes its barycenter.
Removing the representative leaves a type-preserving
stabilizer.

\medskip\noindent
\textit{The split residue and stable alcoves.}
For split reference over \(F'\), choose a further totally
ramified scalar extension with uniformizer root
\(\pi^{1/d}\), where \(d x\) is an integral cocharacter,
and conjugate by \((d x)(\pi^{1/d})\) before reduction.
This scalar extension does not change the residue field.
A root factor in \(I\) becomes integral. For any normalizer
representative fixing \(x\), the new torus valuation is zero,
by the affine fixing equation. Thus \(P_x\) is integral after
scaling and has a reduction map. Its image is the split
reductive group with torus \(T\) and roots
\(\Psi=\{\alpha:\alpha(x)\in\mathbb Z\}\).
Indeed \(I\) reduces to the torus and exactly the positive
half of \(\Psi\) determined by the standard alcove, and the
\(W_x\) representatives reduce to the corresponding finite
Weyl representatives. One can realize this reductive group
as the connected centralizer in the split residue group of
the diagonalizable subgroup \((d x)(\mu_d)\). The big-cell
coordinates give its two opposite root cells by retaining
exactly those coordinates with trivial grading under that
diagonalizable subgroup. Thus the centralizer is smooth
there and its translates give smoothness, also when
\(\mu_d\) is not reduced. The roots are the indicated closed
subsystem, and their opposite rank-one subgroups and torus
give the split root datum of its connected part by the same
root-cell construction as C08. This gives the finite Bruhat
decomposition, so the inverse image of the indicated
Borel under reduction is exactly \(I\), and the chambers
at the face are its flag variety's rational points,
\(P_x/I=M_\Psi(k')/B_\Psi(k')\).

The twisted Galois action descends this residue description
to the original finite residue field. In fact a representative
of an adjoint stabilizer has the preceding factorization
through \(P_x\) and an alcove symmetry; the latter becomes
integral after scaling and normalizes the residue root
data. If unramified Frobenius itself acts on the split
diagram, use its image of the scaling cocharacter at the
image face. The same torus-valuation cancellation is precisely
the equation that the combined action fixes \(x\). Thus the
action on the residue group is semilinear, and equivariant
on its Borel variety. Lang's theorem supplies a rational
Borel, hence an incident alcove fixed by the twisted action.

Conjugate this alcove to the standard one. The part of its
stabilizer preserving all types is \(I\), and its remaining
quotient is \(\Omega\). For the inner \(E\) types Frobenius
acts trivially on \(\Omega\). An element of \(\Omega\) can
be represented by \(t=\lambda(\pi)n_w\) with \(n_w\) a
pinned Weyl lift of finite order up to a finite torus sign.
The finite order of the affine action implies
\(\sum_{j=0}^{m-1}w^j\lambda=0\) for its order \(m\), so
\(t^m=n_w^m\) has finite order. The representative therefore
defines an unramified cocycle at a sufficiently divisible
finite level. The difference between the given cocycle and
this representative lies in \(I\) with its action twisted
by \(t\). The split Iwahori has a smooth model with connected
special fiber: in the open integral big cell containing the
entire special fiber, replace each negative root coordinate
by that coordinate divided by \(\pi\). This is the dilatation
along the special Borel; the group law is induced by the
condition that reduction be in the Borel. A stabilizing
representative preserves all the root depths, so its
conjugation extends over this chart by the ordered root
product, and hence over the whole model by using the generic
fiber on the complement. The same holds for the inverse.
Consequently the difference cocycle descends an unramified
torsor of this smooth model, whose special fiber has a point
by Lang and whose generic fiber has a point by Hensel's lemma.
This conjugates the original cocycle to \(t\). In particular
\(E_8\), where \(\Omega\) is trivial, is split.

For the unramified outer \(E_6\), \(\Omega=C_3\) and the
pinned Frobenius acts on it by inversion. Its generator
value is removed by a coboundary in \(\Omega\), since
multiplication by \(2\) is invertible in \(C_3\). Lift
that conjugation to the alcove stabilizer, then use the
same smooth Iwahori descent to eliminate the resulting
\(I\)-cocycle. The form is then quasi-split.

In the inner cases the torus of the standard apartment under
the representative \(t\) gives a split torus whose rank is
the number of orbits of the alcove vertices minus one. This
rank is maximal. Every fixed point sits in a preserved face,
where its affine coordinates are constant on the vertex
orbits. The fixed complex is thus a countable union of closed
cells of dimension at most this number. Conversely an \(F\)-split
torus of rank \(r\) gives a group \(\mathbb Z^r\) of uniformizer
cocharacters acting by independent translations on an apartment
over \(F'\). These translations commute with the twisted Galois
action. The center of the Galois orbit of a point of that
apartment belongs to every translation's closed convex minimum
set, since the whole orbit belongs to that set. The translations
on the fixed space are therefore semisimple, with escape of
their displacements at this center equal to that of the
translation lattice in that apartment. They act properly,
and the flat torus theorem of \emph{Geometric Group Theory},
Result E11, gives an \(r\)-flat in the fixed space. Baire's
theorem prevents that flat from being covered by countably
many lower-dimensional closed cells. This proves maximality.

The two nontrivial index calculations are as follows, in the
Bourbaki numbering with node \(2\) attached to \(4\). For
\(E_6\) the highest root has coefficients
\((1,2,2,3,2,1)\); the extended diagram has three arms and
\(\Omega=C_3\) rotates them with vertex orbits
\(\{0,1,6\},\{2,3,5\},\{4\}\). On invariant directions the
simple-root values are \((-2a-b,a,a,b,a,-2a-b)\).
Apply simple reflections \(1,3,4,6,5,4,2,3,4\), successively,
to obtain \((0,b,0,a,0,0)\).
For \(E_7\) the highest root has coefficients
\((2,2,3,4,3,2,1)\), and the nontrivial symmetry reverses
\(0,1,3,4,5,6,7\) and fixes \(2\). The invariant values
\((a,d,b,c,b,a,-2a-3b-2c-d)\) become
\((d,0,c,b,0,a,0)\) after reflections
\(7,6,5,4,2,3,1,4,3,5,4,6\).
These calculations use \(r_i\mapsto r_i-A_{ji}r_j\)
for reflection \(j\); the coefficients of the highest roots
are the dominant lattice roots of length squared two.
Take positive free parameters to read the vanishing simple
subsystems in the dominant chamber. Thus the black simple
nodes are \(\{1,3,5,6\}\) for \(E_6\) and \(\{2,5,7\}\)
for \(E_7\).

\medskip\noindent
\textit{The ramified outer form, including residue characteristic two.}
Here the type quotient is trivial, so the original fixed face
is fixed pointwise and we can use its barycenter \(x\) directly.
The twisted cocycle then belongs to \(P_x\) relative to the
quasi-split reference group. Untwisted unramified Frobenius
fixes the reference apartment, because the diagram field is
totally ramified and one may extend Frobenius fixing it.
For this scaling start with the ramified diagram field \(L\)
and choose a uniformizer \(\pi_L\). Take \(d\) such that \(2d x\)
is an integral cocharacter and conjugate by
\((2d x)(\pi_L^{1/d})\). The polynomial for the chosen
root of \(\pi_L\) is Eisenstein over \(L\); after taking its
compositum with the finite unramified extension, the auxiliary
scaling has not changed the residue field. With the base-field
normalization \(v_F(\pi_L)=1/2\), the conjugation shifts root
values by precisely their pairing with \(x\). Then \(I\) and \(P_x\) map into the
split integral group and can be reduced. We use the connected
image of \(P_x\), without equating it to all fixed points of
a folded residue action.

Over the algebraic closure of the finite residue field, the
image of \(I\) is the product of the unit torus image and the
surviving root coordinate images. The former is connected
because the basis torus coordinates have single or paired
Galois orbits. Each latter image is a one-parameter group:
a fixed root has a free leading coefficient in an allowed
integral layer, and a paired root has a free coefficient
in an allowed half-integral layer with the other absolute
coordinate a fixed nonzero scalar multiple of it. These
coefficients can be lifted over finite unramified extensions.
The group \(W_x\) is generated by the simple affine walls of the standard alcove through \(x\).
For such a wall \(\beta+j=0\), its rank-one lift in \(P_x\) with parameter \(t\) of value \(j\) conjugates the Iwahori wall layer \(v(u)\ge j\) onto the opposite layer \(v(u)\ge -j\), by \(u\mapsto -t^{-2}u\) in rank-one coordinates.
Both signed layers thus lie in \(P_x\) and have nonnegative scaled depth. Their reductions are one-parameter groups by the same leading-coefficient calculation.
The lift is a three-term root word in these two layers with all factors of scaled depth zero.
Since \(P_x\) is generated by \(I\) and the marked wall lifts, its geometric reduction image is generated by the unit torus image, the surviving Iwahori root images,
and both signed depth-zero wall images for these walls, all connected algebraic subgroups.
It is also a constructible subgroup of the split residue
group because \(P_x\) is a finite union of \(I n_w I\).
Such a subgroup is closed: it contains a dense open subset
of its closure, so translating that open and taking products
with inverses gives the identity component of the closure,
and then every component in its closure also occurs.
Here the group generated by connected subgroups is connected,
so its closure is connected. With the reduced structure it
is smooth over the perfect residue field. This verifies
the hypotheses for Lang on the reduction image, including
when the ramification is wild.

After this Lang correction we still work over a finite unramified
extension, enlarged once to contain a lift, and the cocycle
has trivial reduction. For clarity we justify the correction
at all further depths. An element of \(P_x\) whose scaled
reduction is the identity has the unique coordinates of the
absolute big cell, and after unscaling they have the paired
relative conditions just described. In particular, a root
parameter has strictly positive scaled depth
\(v(u)+\beta(x)\), and the torus coordinates are units near
\(1\). Conversely all such coordinates give an element of
the kernel: every root factor fixes \(x\) by the depth
condition, and they all reduce to identity. The same
coordinate characterization applies to every higher
congruence subgroup of the scaled kernel. In passing to
one consecutive quotient, multiplication adds the first
nonzero scaled coordinates, since every product of two
strictly positive-depth terms lies in a deeper layer.
A fixed root gives a residue line exactly when that scaled
depth is possible from an integer value, and a paired
root gives one line when it is possible from a half-integer
value, its two leading coefficients related by the fixed
nonzero factor. The torus coordinates give the same additive
lines by writing each unit as \(1+z\) in its single or paired
field factor. All coefficients permitted in such a layer
can be prescribed independently by lifting in the big cell.
It follows that every successive quotient is a vector space
over the residue field, with semilinear unramified Galois action.
This account works in characteristic two: only leading
coefficients in the two paired factors are compared, and
no averaging by the order of the ramified involution is used.

For each of these vector spaces additive Galois descent,
or additive Hilbert 90 after choosing a descended basis,
kills the cocycle in that layer. Lift the correcting vector
to the kernel at its depth. The resulting cocycle is in
the next layer. All corrections take place over the fixed
complete finite unramified extension containing the first
lift; the correcting elements converge to identity with
depth, and their successive product converges in the
closed kernel. Thus the entire cocycle is trivial. This
proves that the ramified outer \(E_6\) is quasi-split as well.

\medskip\noindent
\textit{Local generation in the two nonsplit inner exceptional
indices.}
We give the local generation consequence of the two
nonsplit diagrams, using only the inner type \(A\)
calculations already proved. Let \(F\) again denote the
local field itself. For an inner simply connected group
with one of these diagrams, write \(I\) for the
anisotropic simple roots. Choose a distinguished node
\(i\) perpendicular to the span of \(I\), as specified
below. In the absolute simply laced root system, any
positive relative multiple of the restriction of
\(\alpha_i\) has absolute representatives of the form
\(m\alpha_i+\eta\), with \(\eta\) in the span of \(I\)
and \(m\) a positive integer. Orthogonality and the
common root square length two force \(m=1,\eta=0\).
The same assertion holds on the opposite ray. The
two relative ray groups of Result~\ref{res:D09} are
then precisely the descended absolute root groups for
\(\pm\alpha_i\). They are \(F\)-split and generate
the split rank-one subgroup: its absolute coroot
\(\alpha_i^\vee\) is primitive in the cocharacter lattice,
and the two descended nontrivial root groups give
opposite split unipotent groups in this rank-one form.
The root words put \(\alpha_i^\vee(F^\times)\) in \(G(F)^+\).

Let \(L\) be the Levi of the maximal parabolic omitting
\(i\). Its derived group is simply connected by the
simple-coroot basis. The quotient \(L/L_{\rm der}\)
is a split one-dimensional torus, and
\(\alpha_i^\vee\) gives a section, because its image
generates the quotient of the cocharacter lattice by
the span of the remaining simple coroots. Thus \(L(F)\)
is generated by \(L_{\rm der}(F)\) and
\(\alpha_i^\vee(F^\times)\). Take the ambient maximal
split torus in this Levi. Its relative ray subgroups
in the derived Levi lie in the corresponding ambient
relative ray subgroups, by Result~\ref{res:D09}.
Consequently density of the group generated by the
relative rays in \(L_{\rm der}(F)\) puts \(L(F)\) in
\(\overline{G(F)^+}\). The minimal Levi of \(G\) lies
in \(L\), so the relative Bruhat decomposition then
gives \(G(F)=\overline{G(F)^+}\).

In the index \(^{1}E_{6,2}^{16}\) the black nodes are
\(\{1,3,5,6\}\), and \(i=2\) is perpendicular to them.
Deleting \(2\) leaves the \(A_5\) chain with two black
\(A_2\) blocks and one distinguished midpoint. Its
derived Levi group is \(\operatorname{SL}_2(B)\) for a
degree-three central division algebra \(B\), by the
inner type \(A\) matrix descent and rank calculation
above. In the index \(E_{7,4}^{9}\) the black nodes are
\(\{2,5,7\}\), and \(i=3\) is perpendicular to them.
Deleting \(3\) leaves isolated distinguished node \(1\)
and the chain \(2,4,5,6,7\), with black nodes alternating
with distinguished ones in this chain. The derived
Levi is a product of split \(\operatorname{SL}_2\)
and \(\operatorname{SL}_3(B)\) for a quaternion
division algebra \(B\). The earlier inner type \(A\)
computation gives the required local generation in
both cases, in the topology of \(F\).

\medskip\noindent
\textit{Closed orbits in positive characteristic.}
We first prove the invariant lifting needed here, using the
Steinberg-module construction of \cite[Sections 3.2--3.4]{franjouvdK}. Let \(H\) be a connected semisimple algebraic group
and let \(L\to k\) be a surjection of finite-dimensional rational
\(H\)-modules, with trivial action on \(k\). We will show that
some \(\operatorname{Sym}^dL\) contains an invariant mapping to
\(1\). We may extend to an algebraic closure: for a fixed \(d\),
taking invariants is the kernel of the difference between
the coaction and the trivial coaction, so it commutes with field
extension. We may also pull back along the surjective simply
connected covering, which does not change the invariants. We
can thus work with a split simply connected semisimple group.

For dominant weights \(\lambda,\mu\), put
\(\nabla(\mu)=H^0(H/B^-,\mathcal L_\mu)\) and
\(\Delta(\lambda)=\nabla(-w_0\lambda)^*\). Here \(w_0\) is the
longest Weyl element. We need
\[
 \operatorname{Ext}^1_H(\Delta(\lambda),\nabla(\mu))=0 .
\]
To see this, restriction to \(B^-\) is exact and its right
adjoint is rational induction, so induction preserves
injectives. Its derived functors are the cohomology of the
associated sheaves on \(H/B^-\). Indeed, the quotient torsor is
locally trivial on the big cell and its translates. Cofree
\(B^-\)-comodules are \(k[B^-]\otimes W\), for vector spaces \(W\).
They have associated sheaves given by pushforward of
\(\mathcal O_H\otimes W\) along this affine torsor; these are
acyclic because \(H\) is affine. Every rational comodule
embeds into a cofree one by its coaction, so injectives are
summands of such modules. This establishes the derived
functor identification. By Kempf vanishing in E10, induction
of the character \(\mu\) has no positive derived functors.
Adjunction on an injective resolution therefore computes the
displayed extension group as
\[
 H^1(B^-,\Delta(\lambda)^*\otimes k_\mu).
\]
The split torus invariants are exact by weight decomposition.
The normalized bar cochains for the unipotent radical of \(B^-\)
have one copy per slot of the augmentation ideal of its
coordinate ring. That ideal has only weights in the nonzero
positive root cone: in ordered negative root coordinates
each coordinate function has the corresponding positive root
weight. Thus a positive-degree cochain of weight zero in this
calculation would require \(\lambda-\mu\) to lie in the nonzero
positive root cone, since the coefficient module has weights
in \(\mu-\lambda+\mathbb N\Phi^+\). If \(\lambda-\mu\) is in
that cone, dualize the extension. This replaces
\((\lambda,\mu)\) by \((-w_0\mu,-w_0\lambda)\), whose difference
is not in that cone. Every extension considered is finite
dimensional, so duality applies directly. This proves the
extension vanishing.

Let \(N=\ker(L\to k)\). For \(r\) sufficiently large, all weights
of \(N\) become dominant upon adding
\(\sigma_r=(p^r-1)\rho\). Put \(S_r=\nabla(\sigma_r)\).
The Frobenius pushforward computation in E10, with
\(\lambda=0\), identifies this module with \(L(\sigma_r)\).
The highest weight of its dual is \(-w_0\sigma_r=\sigma_r\),
so \(S_r=S_r^*\) up to isomorphism, and
\(\Delta(\sigma_r)\simeq S_r\). Filter \(N|_{B^-}\) by
one-dimensional characters, ordering its weight spaces
from lower to higher weights. The negative root groups can
only lower a weight, so this is an invariant filtration.
The tensor identity
\[
 S_r\otimes N=\operatorname{Ind}_{B^-}^H(k_{\sigma_r}\otimes N)
\]
follows by sending a section with values in \(N\) to the
section expressed in the ambient \(H\)-representation at each
coset. Kempf vanishing on the character quotients makes
induction exact along the filtration. Its resulting
successive quotients are the costandard modules for the
dominant weights \(\sigma_r+\mu\). The extension vanishing
just proved and the extension exact sequences give
\[
 \operatorname{Ext}^1_H(\operatorname{End}(S_r)^*,N)
 \simeq \operatorname{Ext}^1_H(S_r,S_r\otimes N)=0 .
\]
The displayed tensor adjunction holds for extensions since
tensoring by a finite-dimensional module is exact and has
exact adjoints given by tensoring with its dual. Consequently
evaluation at the identity of \(\operatorname{End}(S_r)\)
lifts from \(\operatorname{End}(S_r)^*\to k\) to an equivariant
map \(\operatorname{End}(S_r)^*\to L\). The determinant is a
homogeneous invariant polynomial on \(\operatorname{End}(S_r)\)
of degree \(\dim S_r\), with value \(1\) at the identity.
Its image under the symmetric power of the lift is the
required invariant in \(\operatorname{Sym}^{\dim S_r}L\).

Now let \(H\subset P=\operatorname{SL}_m\) be a closed
semisimple subgroup. Choose a finite-dimensional right
translation submodule \(V\subset k[P]\) containing generators
of \(I(H)\), and let \(W=V\cap I(H)\). The stabilizer of \(W\)
is \(H\): preservation of \(W\) is equivalent to preservation
of the ideal it generates, and the stabilizer of that ideal
under right translation is \(H\). These assertions also hold
on points over any \(k\)-algebra since all the defining
vector-space intersections commute with scalar extension.
The stabilizer of \(\bigwedge^{\dim W}W\) as a line in
\(\bigwedge^{\dim W}V\) is the same subgroup, by the
exterior description of the subspace \(W\). The line
character of \(H\) is trivial because a semisimple group
has no characters, as is seen over a splitting field by its
root subgroups. Thus \(H\) is the scheme stabilizer of a
vector. Its orbit \(X=P/H\) is quasi-affine as an open
subvariety of the orbit closure \(Y\) in that vector space.
For the orbit description, the geometric orbit is open in
its closure by constructibility and transitivity. Formation
of the closure commutes with field extensions, and the orbit
map to that open subscheme is smooth and surjective:
over an algebraic closure its differential has the dimension
of the orbit as rank, since the scheme stabilizer is smooth.
This can be checked at the identity and then translated.
Descent gives the orbit and its smooth quotient map over
\(k\).

Choose finitely many \(f_i\in k[Y]\) vanishing on \(Y\setminus X\)
whose nonvanishing opens cover \(X\). Their pullbacks to \(P\)
generate the unit ideal of \(k[P]\). The inverse image of the
line \(k\cdot1\) under
\[
 k[P]^{\oplus s}\longrightarrow k[P],\qquad
 (a_i)\longmapsto\sum_i a_i f_i
\]
is a rational right \(H\)-module surjecting onto \(k\cdot1\).
Choose a finite submodule still surjecting. Invariant lifting
applied to it gives an invariant of degree \(d\). Map a
symmetric monomial in the coefficient tuples \((a_i)\) to the
product of the linear expressions \(\sum_i a_iT_i\) in
\(k[P][T_1,\ldots,T_s]\). The resulting polynomial is
\(H\)-invariant, so its coefficients are in \(k[P]^H\).
Substitution \(T_i=f_i\) gives \(1\) by the choice of lift.
The quotient map \(P\to X\) is an \(H\)-torsor, and descent
of regular functions gives \(k[P]^H=\mathcal O(X)\).
Therefore the \(f_i\) generate the unit ideal of \(\mathcal O(X)\).
For each \(i\), \(\mathcal O(X)_{f_i}=k[Y]_{f_i}\): the natural
inclusions in both directions follow by restricting to the
same affine open \(D_Y(f_i)=D_X(f_i)\). The corresponding
principal opens cover \(\operatorname{Spec}\mathcal O(X)\)
and identify it with \(X\). Thus \(X\) is affine.
Finally the finite-dimensional span, under the left action of
\(P\), of a finite generating set of \(k[X]\) gives a
\(P\)-equivariant closed immersion of \(X\) into the dual
of that span. The image of the identity coset is a \(k\)-vector
with stabilizer \(H\) and closed orbit, as required.

\medskip\noindent
\textit{A rational contraction in this orbit.}
Let \(\lambda:\mathbb G_m\to P\) be a \(k\)-cocharacter and
suppose \(x\in X(k)\) has a limit \(x_0=\lim_{t\to0}\lambda(t)x\)
in \(X\). Denote by \(U_\lambda(P)\) the subgroup contracted
to \(1\) by conjugation with \(\lambda(t)\). We claim that
\(x=u x_0\) for some \(u\in U_\lambda(P)(k)\).
To prove it, give regular functions the weights specified
by \(f(\lambda(t)y)=t^a f(y)\). The fiber of the limit map
at \(x_0\) is closed in \(X\); it is defined by all
negative-weight functions and the weight-zero functions
vanishing at \(x_0\). Choose homogeneous lifts \(y_i\) of a
basis of the cotangent space at the smooth \(k\)-point
\(x_0\). The completed local ring is \(k[[y_1,\ldots,y_n]]\).
Every homogeneous function has an expansion with monomials
of that weight, by comparing coefficients under the split
torus. Therefore the ideal of the fiber in this completion
is generated by exactly those variables of nonpositive
weight: every monomial of nonpositive weight with nonzero
degree involves such a variable. The fiber is consequently
smooth at \(x_0\), and its tangent space there is the
positive-weight subspace. The orbit map of \(U_\lambda(P)\)
at \(x_0\) has differential onto this space, by the smooth
orbit map of \(P\) and weight decomposition. Its geometric
image contains an open neighborhood of \(x_0\) in the
fiber. Since the path \(\lambda(t)x\) is in the fiber and
approaches \(x_0\), it meets that neighborhood for some
nonzero \(t\) over an algebraic closure. The group
\(U_\lambda(P)\) is normalized by \(\lambda\), so \(x\)
lies in its geometric orbit of \(x_0\).

The scheme of \(u\)'s carrying \(x_0\) to \(x\) is a
torsor under \(U_\lambda(P)\cap\operatorname{Stab}_P(x_0)\).
As \(x_0\) is fixed by \(\lambda\), this intersection is the
contracted radical for the same cocharacter in the
reductive group \(\operatorname{Stab}_P(x_0)\). The equality
follows directly from the limit definition and closedness of
the stabilizer in \(P\). The relative parabolic calculation
of D08 identifies it with a smooth \(k\)-split unipotent
group. Its torsors over \(k\) are trivial: successive quotients
of a split filtration are additive groups, whose torsors are
trivial by additive descent. It follows that \(u\) can be
chosen over \(k\). In particular, if \(\lambda\ne1\) then
\(\operatorname{Stab}_P(x)\) contains the nontrivial
\(k\)-split torus given by \(u\lambda u^{-1}\).

\medskip\noindent
\textit{Compactness over function fields from a closed orbit.}
We use a finite-dimensional representation of a special linear
group in which the given semisimple group is the stabilizer of
a point of a closed orbit. We will justify this representation
and the contraction property used below.

Let \(k/k_0\) be a finite separable extension with
\(k_0=\mathbb F_q(t)\), obtained by choosing a separating
parameter. Via restriction of scalars we can first work over
\(k_0\). In fact the local coordinates of the restricted group at
a place of \(k_0\) are the product over the places of \(k\) above
it, by the factorization of \(k\otimes_{k_0}(k_0)_v\).
This identifies rational points, and also the restricted products,
since after clearing denominators the integral models agree
outside a finite set. The restricted group is semisimple, as it
becomes a product after separable scalar extension. A nontrivial
\(k_0\)-cocharacter into it would give a nontrivial
\(k\)-cocharacter into the original group by scalar-restriction
adjunction. Thus anisotropy is preserved in this step.

Write \(H\) for the resulting anisotropic semisimple group over
\(k_0\). Embed it as a closed subgroup of \(P=\operatorname{SL}_d\);
the determinant of a faithful matrix representation is trivial,
since an algebraic character vanishes on the geometric root
groups which generate a semisimple group. Assume \(H\) is the
scheme stabilizer of \(b\) in a \(P\)-module \(V\), with \(P b\)
closed, as constructed above. Choose fixed linear coordinates
in \(V\) adapted to the weights of the diagonal torus of \(P\).

Put \(R=\mathbb F_q[t]\) and \(O=\mathbb F_q[[1/t]]\). For
\(g\in H(\mathbb A_{k_0})\), the finite local lattices
\(g_v^{-1}O_v^d\) are recovered from their intersection with
\(k_0^d\), as in the polynomial lattice calculation above.
Choose an \(R\)-basis of that intersection, with matrix
\(\gamma\). Then \(g_v\gamma\) is integral with integral inverse
at every finite place. The determinant of \(\gamma\) is a
nonzero constant, since all finite valuations are zero. We
normalize one basis column by the inverse constant, so that
\(\gamma\in P(k_0)\).

At infinity the basis can be chosen so that
\[
 g_\infty\gamma = k_\infty\operatorname{diag}
          (t^{n_1},\ldots,t^{n_d}),\qquad
 k_\infty\in\operatorname{GL}_d(O),\quad
 \sum_j n_j=0.
\]
For this slightly stronger form of lattice reduction, put a
shortest nonzero vector in the first coordinate by an
\(\operatorname{GL}_d(O)\)-change, and make its first
coordinate a power \(t^{n_1}\) by a unit. The vector is
primitive over \(R\); extend it to an \(R\)-basis and project
along that coordinate. A lift of a nonzero projected vector
can have its first coordinate reduced to absolute value
strictly smaller than \(|t^{n_1}|\), by subtracting the
polynomial part in units of \(t^{n_1}\). Consequently the
projected lattice has no nonzero vector shorter than the first
vector. Induct on the dimension to obtain a reduced basis for
the projection, and lift each projected basis vector with the
first coordinate reduced in this manner. After division of
each column by its power of \(t\), the result is integral
with integral inverse: in the first row the off-diagonal
entries are integral, and the lower diagonal block is the
unimodular matrix from induction. Normalizing the determinant
by a constant only changes \(k_\infty\) by a unit matrix.
The determinant-one property implies \(\sum_j n_j=0\) and \(k_\infty\in\operatorname{SL}_d(O)\).

Suppose that these integer exponent vectors are unbounded as
\(g\) varies. Take a sequence with Euclidean norms tending to
infinity and with normalized vectors converging to a nonzero
direction \(e\) of coordinate sum zero. The vectors
\(b_\gamma=\gamma^{-1}b=(g\gamma)^{-1}b\) are rational. Their
finite coordinates have a uniform denominator in \(R\),
because \(g_v\gamma\) is in the compact integral group and
the representation is integral outside finitely many places.
In each weight space of weight \(\chi\) the infinity norm of
\(b_\gamma\) is at most a common constant times
\(q^{-\langle\chi,n\rangle}\). If
\(\langle\chi,e\rangle>0\), this tends to zero; after
multiplication by the fixed denominator, a nonzero coordinate
polynomial would have infinity norm at least one. Thus all
such weight components are eventually zero. There are only
finitely many weights. Choose a nonzero integral vector
\(m\) of coordinate sum zero with the same strict signs as
\(-e\) for all nonvanishing pairings and with zero pairing
where \(e\) has zero pairing. This is possible since the
linear equations are rational and the remaining conditions
are open within their solution space. For every sufficiently
late \(b_\gamma\), the cocharacter
\(\lambda(a)=\operatorname{diag}(a^{m_j})\) has a limit on
\(b_\gamma\) as \(a\to0\): all its negative-weight components
are zero. By closedness the limit lies in the orbit. The
preceding rational contraction argument then puts a conjugate of
\(\lambda\) in the stabilizer of \(b_\gamma\). This contradicts
anisotropy, since that stabilizer is the rational conjugate of
\(H\).

Thus the integer exponents are uniformly bounded. Hence all
the \(g\gamma\) lie in a fixed compact adelic subset of
\(P(\mathbb A_{k_0})\). The rational vectors \(b_\gamma\) in a
compact adelic set are finite, by discreteness of the
additive diagonal proved above. Choose one \(\gamma_i\) for
each occurring vector. If \(\gamma^{-1}b=\gamma_i^{-1}b\)
then \(\gamma\gamma_i^{-1}\in H(k_0)\). The matrices
\(g\gamma\gamma_i^{-1}\) lie in the intersection of \(H\)
with the union of finitely many compact translates of the
matrix set; that intersection is compact by closedness.
This proves compactness of \(H(\mathbb A_{k_0})/H(k_0)\), and
therefore of \(G(\mathbb A_k)/G(k)\).

\medskip\noindent
\textit{The local mixing step for the remaining isotropic types.}
Assume that the local group \(J\) is simply connected, has
\(\overline{J(F)^+}=J(F)\), and that in a maximal split torus apartment
one longest relative root has a single absolute root over it.
In the exceptional diagrams below and in quasi-split
trialitarian \(D_4\) such a root is supplied by the fixed
simple root orthogonal to the anisotropic roots, or by a
diagram-fixed simple root in the quasi-split case. After
normalizing the length using the simply laced absolute
system, the indicated root lies in the split cocharacter
subspace, so its projection has full absolute length.
No distinct absolute root can have the same projection and
length. The four relative root sets used here can be checked
directly on the simple-root coefficient vectors. Starting with
the simple unit vectors, apply the reflection rule replacing
only entry \(i\) of \(n\) by \(n_i-\sum_j A_{ij}n_j\), until
the set stops changing. In the inner cases keep the white
coordinates read off above; in the outer cases sum within each
diagram orbit. Discard zero and repetitions. This gives
respectively \(12,48,48,12\) roots for the two nonsplit inner
indices, outer \(E_6\), and trialitarian \(D_4\), with no root
alongside twice itself. The corresponding rank-one subgroup is split
\(\operatorname{SL}_2\), by descent of the unique opposite
absolute root groups and the simple connectedness of the
absolute root embedding. Choosing a Weyl conjugate gives this
subgroup at the highest relative root, since the relative
systems in question are reduced and their highest roots are
long. The highest root group lies in the center of the
radical of a maximal standard parabolic: its positive
coefficients contain the omitted simple root, and the sum of
this highest root with any positive radical root is beyond
the maximal height. The relative commutator formula therefore
gives centrality.

In the unitary quotient representation used above, take a
weak limit of positive powers of a regular element of the
split torus. Its range is fixed by the contracted half of
the indicated split rank-one subgroup. The explicit
\(\operatorname{SL}_2\) Mautner identity above then makes
the range fixed by both halves and by the rank-one
diagonal. That diagonal and its inverse contract the radicals
of its pair of proper opposite parabolics. The relative-ray
generation argument therefore makes the range fixed by
the dense subgroup \(J(F)^+\), hence by \(J(F)\). This proves mixing, and the rescaling
argument above puts every conjugate of some noncentral
element in the projected subgroup.

The faithful primitive action on maximal parabolic cosets
proved above applies to this parabolic as well. The abelian
normal subgroup given by the center of its radical has
conjugates generating \(J(F)\) topologically: it contains
the rank-one upper group, whose conjugate gives the lower
group, whence the diagonal and both contracted radicals.
These generate \(J(F)\) just as before. The contraction
commutators also make \(J(F)\) topologically perfect.
The primitive-action argument thus gives the full projected
subgroup. This supplies the same local projection and
isolation argument over function fields once the finite
local structure has been checked.

\medskip\noindent
\textit{Completion of sufficiency.}
Suppose \(G\) is simply connected and \(G_S\) is noncompact.
For the groups not settled by the earlier classical
calculations, the finite local calculations give
\(\overline{G(k_v)^+}=G(k_v)\) and positive split rank at every
nonarchimedean place. The real groups are connected; at each
noncompact real place the closure of its subgroup generated
by the parabolic radicals is the whole group, by normality
and the Lie ideal argument. The complex factors are split.
When \(G\) is globally isotropic, every completion is
isotropic. The earlier approximation by the two fixed
opposite radicals therefore proves density outside \(S\).

When \(G\) is globally anisotropic over a number field, the
adjoint lattice construction gives compactness of the adelic
quotient. First let \(J_\infty\) be the product of the noncompact
archimedean factors outside \(S\), and take \(T\) to include them
and \(S\). Each nonempty subproduct of their lattice projection is
nondiscrete, as shown by compactness and injectivity of the
diagonal in a single completion. The projective recurrence
argument with independent split cocharacters shows that the
Lie algebra of the closed projected subgroup is a sum of
the simple factor ideals. An omitted ideal would give a
discrete projection to the complementary subproduct,
after dividing by the connected factors already in the
closure, contrary to nondiscreteness there. Thus the
projection is full, and isolation by contractions gives the
whole noncompact archimedean subproduct in the outside
adelic closure.

Enlarge the omitted set to \(S'\) by adding these already
isolated places. Since the earlier closure contains
\(J_\infty\), its projection to the adeles outside \(S'\)
is closed and equals the corresponding diagonal closure.
For a finite outside place \(v\), use
\(T=S'\cup\{v\}\). Its projected arithmetic lattice has
nondiscrete closure. Projective recurrence applied to the
space of small logarithms makes this space invariant under
all local conjugates of a split cocharacter. Contraction
then gives local Lie ideals and opens the subgroup. Its index
is finite by cocompactness. Divisibility of the radicals and
local generation give the full factor. The contracting radical
isolation argument puts that factor, embedded with identity
in the other coordinates outside \(S'\), in the closure there.
It follows that the original closure contains each such
factor, since we have already isolated the kernel
\(J_\infty\) of the projection.

The isolated finite factors give their restricted product by
density of finite supports. The remaining factors are compact
archimedean ones. The weak approximation argument above
contains a neighborhood of the identity in the closure of
the diagonal projection to their product. That product is
connected, hence the projection is dense. Quotienting by
the already isolated factors is an open quotient by a
subgroup contained in the closure, so the projected closure
is closed as well and is the full remaining product.

Over a global function field the closed-orbit lattice
argument gives the same compactness for globally anisotropic
\(G\). At each outside place \(v\) take \(T=S\cup\{v\}\).
The projected arithmetic lattice at \(v\) again has a
nondiscrete closure with a compact quotient carrying finite
invariant measure. The mixing and rescaling argument
gives the full local closure. The long-root case of the
local mixing step applies in the remaining exceptional and
trialitarian forms. The nontrialitarian quadratic and
quaternionic spin cases use the hyperbolic pair described
above. In characteristic two the scalar matrices on
the quaternionic pair act by the rank-one scalar
\(\operatorname{SL}_2\) in the polar isometry group: the
two elementary directions are \(U(0,b)\), \(L(0,b)\) with
\(b\) scalar. They lift to their spin root groups. After splitting the quaternion algebra, this upper line
is the ordinary root for \(e_1+e_2\) in the hyperbolic
four-space, and the opposite line is its negative root.
The resulting split root \(\operatorname{SL}_2\) descends:
its lift is unique, because two lifts differ by a
morphism from the connected smooth group to the finite
central \(\mu_2\). Its diagonal contracts the two opposite
parabolic radicals. This supplies the same rank-one
arguments, if they are needed. The compact-product
isolation then puts every local factor in the outside
closure, and finite supports are dense. This proves the
stated strong approximation criterion.
\end{proof}


\backmatter
{\small\begin{thebibliography}{99}
\addcontentsline{toc}{chapter}{Bibliography}

\bibitem{borel}
A. Borel, \emph{Linear Algebraic Groups}, second enlarged edition,
Graduate Texts in Mathematics 126, Springer, 1991.
\href{https://doi.org/10.1007/978-1-4612-0941-6}{DOI}.

\bibitem{conrad}
B. Conrad, \emph{Reductive Group Schemes}, in \emph{Autour des schémas en groupes}, vol.~I,
Panoramas et Synthèses 42--43, 2014.
\href{https://math.stanford.edu/~conrad/papers/luminysga3.pdf}{Author's text}.

\bibitem{conradfeng}
B. Conrad and T. Feng, \emph{Reductive Groups over Fields},
Stanford course notes, version of 5 December 2020.
\href{https://math.stanford.edu/~conrad/249BW16Page/handouts/AMSalggroups2.pdf}{Course text}.

\bibitem{cgp}
B. Conrad, O. Gabber, and G. Prasad, \emph{Pseudo-reductive Groups},
second edition, New Mathematical Monographs 26, Cambridge University Press, 2015.
\href{https://www.cambridge.org/core/books/pseudoreductive-groups/CA287151E422CA75682AE7A460008504}{Publisher}.

\bibitem{etingof}
P. Etingof et al., \emph{Introduction to Representation Theory},
Student Mathematical Library 59, American Mathematical Society, 2011.
\href{https://arxiv.org/abs/0901.0827}{Authors' text}.

\bibitem{gs}
P. Gille and T. Szamuely, \emph{Central Simple Algebras and Galois Cohomology},
second edition, Cambridge Studies in Advanced Mathematics 165,
Cambridge University Press, 2017. \href{https://doi.org/10.1017/9781316661277}{DOI}.

\bibitem{haboush}
W. J. Haboush, A short proof of the Kempf vanishing theorem,
\emph{Inventiones Mathematicae} 56 (1980), 109--112.
\href{https://experts.illinois.edu/en/publications/a-short-proof-of-the-kempf-vanishing-theorem}{Publication record}.

\bibitem{jantzen}
J. C. Jantzen, \emph{Representations of Algebraic Groups}, second edition,
Mathematical Surveys and Monographs 107, American Mathematical Society, 2003.
\href{https://doi.org/10.1090/surv/107}{DOI}.

\bibitem{kalethaprasad}
T. Kaletha and G. Prasad, \emph{Bruhat--Tits Theory: A New Approach},
New Mathematical Monographs 44, Cambridge University Press, 2023.
\href{https://www.cambridge.org/core/books/bruhattits-theory/C4FFE5CC52F0DE23D68CC818146771FC}{Publisher}.

\bibitem{milne}
J. S. Milne, \emph{Algebraic Groups: The Theory of Group Schemes of Finite Type over a Field},
Cambridge Studies in Advanced Mathematics 170, Cambridge University Press, 2017,
corrected 2022 printing. \href{https://www.jmilne.org/math/Books/iag.html}{Author's book page}.

\bibitem{milneerrata}
J. S. Milne, \emph{Errata for the 2022 Printing of Algebraic Groups}.
\href{https://www.jmilne.org/math/Books/iagerrata2022.html}{Author's errata}.

\bibitem{morris}
D. W. Morris, \emph{Introduction to Arithmetic Groups},
Deductive Press, version 1.0, 2015.
\href{https://arxiv.org/abs/math/0106063v6}{Author's text}.

\bibitem{prr}
V. Platonov, A. Rapinchuk, and I. Rapinchuk,
\emph{Algebraic Groups and Number Theory}, volume I, second edition,
Cambridge Studies in Advanced Mathematics 205, Cambridge University Press, 2023.
\href{https://www.cambridge.org/core/books/algebraic-groups-and-number-theory/DCBB6BE68D2326C60EFEE9C9B5CD176B}{Publisher}.

\bibitem{rapinchuk}
A. S. Rapinchuk, On strong approximation for algebraic groups, 2012.
\href{https://arxiv.org/abs/1207.4425v1}{Author's text}.

\bibitem{springer}
T. A. Springer, \emph{Linear Algebraic Groups}, second edition,
Birkhäuser, 1998. \href{https://doi.org/10.1007/978-0-8176-4840-4}{DOI}.

\bibitem{stacks}
The Stacks Project Authors, \emph{The Stacks Project},
Section 35.6. \href{https://stacks.math.columbia.edu/tag/0CDQ}{Galois descent}.

\bibitem{takeuchi}
M. Takeuchi, A correspondence between Hopf ideals and sub-Hopf algebras,
\emph{Manuscripta Mathematica} 7 (1972), 251--270.


\bibitem{dolphinquadratic}
A. Dolphin, \emph{Generalised quadratic forms and the \(u\)-invariant}.
\href{https://arxiv.org/abs/1610.06915v2}{Article preprint}.

\bibitem{alsaodyf4}
S. Alsaody, V. Chernousov, and A. Pianzola,
\emph{\(R\)-triviality of groups of type \(F_4\) arising from the first Tits construction}.
\href{https://www.math.ualberta.ca/~chernous/ChernousovPublications/F4arxiv.pdf}{Author copy}.

\bibitem{steinbergregular}
R. Steinberg, \emph{Regular elements of semi-simple algebraic groups},
Publ. Math. IHES 25 (1965), 49--80.
\href{https://www.numdam.org/item/PMIHES_1965__25__49_0/}{Article}.

\bibitem{prasadunramified}
G. Prasad, \emph{A new approach to unramified descent in Bruhat--Tits theory}.
\href{https://arxiv.org/abs/1611.07430v7}{Article preprint}.

\bibitem{franjouvdK}
V. Franjou and W. van der Kallen, \emph{Power reductivity over an arbitrary base}.
\href{https://arxiv.org/abs/0806.0787v2}{Article preprint}.

\end{thebibliography}
\par}
\end{document}
